\chapter{Bức tranh toàn cảnh: đề cương này nói về điều gì}

\section{Ý chính trong một đoạn}

Đề cương hỏi: khi nhà quản lý quỹ phải chọn một cách chia tiền vào cổ phiếu, chọn sai tốn kém đến mức nào, và có tốn hơn khi thị trường căng thẳng không. Tác giả gọi cái giá đó là rủi ro lựa chọn mô hình, sẽ đo nó trên thị trường chứng khoán Việt Nam theo từng chế độ thị trường, tìm các biến có liên hệ thống kê với nó và thử các cách giảm nó.

Chưa có kết quả nào. Đề tài có thể bác bỏ chính giả thuyết của mình mà vẫn trả lời được câu hỏi ban đầu.

\section{Vấn đề: chọn một cách chia tiền}

Hình dung bạn quản lý một quỹ và cần chia tiền vào rổ cổ phiếu. Bạn có thể chia đều, dùng công thức để tìm cách chia “tối ưu”, hoặc thêm nhận định riêng của mình vào công thức. Mỗi cách như thế là một mô hình danh mục.

Bạn phải chọn một cách trước mùa giải. Đề cương giả định nhà đầu tư chọn, ở đầu mỗi cửa sổ đánh giá sáu tháng, mô hình có lợi suất tương đương chắc chắn ngoài mẫu cao nhất trong 36 tháng liền trước. Hết mùa giải mới thấy một cách khác lẽ ra tốt hơn.

Khoản chênh giữa cách tốt nhất nhìn lại sau thực tế và cách đã chọn gọi là mức hối tiếc (regret). Rủi ro lựa chọn mô hình là mức hối tiếc kỳ vọng.

Vì sao đáng lo? DeMiguel và cộng sự (2009) đánh giá 14 mô hình trên bảy bộ dữ liệu và không thấy mô hình nào vượt nhất quán quy tắc chia đều vốn cho N tài sản (quy tắc 1/N). Khi không mô hình nào thắng ổn định, chọn mô hình là một quyết định mang rủi ro riêng.

\section{Chế độ thị trường: trời quang và mùa bão}

Thị trường có thời tiết: phần lớn thời gian trời quang, có những giai đoạn là mùa bão với biến động cao và lợi suất thấp.

Đề cương gọi hai trạng thái này là chế độ bình thường và chế độ căng thẳng. Chế độ căng thẳng là khái niệm thống kê, do mô hình chuyển chế độ ước lượng từ lợi suất VN-Index: chế độ có phương sai điều kiện cao nhất và lợi suất kỳ vọng thấp hơn. Khủng hoảng là sự kiện lịch sử. Đề tài đối chiếu hai khái niệm, không dùng thay nhau.

Nếu thứ hạng của các cách chia tiền đổi nhiều theo thời tiết, chọn bằng quá khứ gần sẽ dễ sai hơn khi có bão. Đó là giả thuyết H1: rủi ro lựa chọn mô hình ở chế độ căng thẳng cao hơn ở chế độ bình thường.

\section{Khoảng trống: ai đã đo, ai chưa đo}

Đề cương chỉ ra hai chỗ trống. Thước đo rủi ro mô hình hiện có đo chênh lệch giữa các dự báo rủi ro, chưa đo hậu quả của việc chọn sai lên kết quả danh mục. Nghiên cứu danh mục xếp hạng mô hình theo bình quân toàn mẫu, chưa xét thứ hạng đổi ra sao giữa các chế độ. Ở Việt Nam, mỗi công trình xét riêng một mô hình, một phương pháp hoặc một hiện tượng.

Tuyên bố tính mới được rào cẩn thận: “theo hiểu biết của thí sinh”, dựa trên truy vấn đến 30/09/2026, chưa gồm Scopus, Web of Science và tạp chí tiếng Việt. Việc đọc toàn văn nằm trong kế hoạch năm thứ nhất.

Việt Nam là bối cảnh kiểm chứng, không phải lý do của khoảng trống. Thanh khoản hạn chế và nhiều cú sốc trong lịch sử giao dịch ngắn có thể làm chênh lệch giữa các mô hình lớn; đổi lại, mẫu có ít giai đoạn căng thẳng độc lập.

\section{Ba câu hỏi và bốn giả thuyết}

Ba câu hỏi:

\begin{itemize}
\item
  Câu hỏi 1: rủi ro lựa chọn mô hình thay đổi thế nào giữa chế độ bình thường và chế độ căng thẳng?
\item
  Câu hỏi 2: sau khi kiểm soát biến động thực hiện, thanh khoản, tương quan cực trị và (ở mức thăm dò) bầy đàn có liên hệ thống kê với rủi ro lựa chọn mô hình không, và giải thích được bao nhiêu phần chênh lệch giữa hai chế độ?
\item
  Câu hỏi 3: cơ chế nào giảm được rủi ro lựa chọn mô hình sau chi phí giao dịch?
\end{itemize}

H1 trả lời câu hỏi 1; H2 và H4 trả lời câu hỏi 2; H3 trả lời câu hỏi 3. Nói gọn: H1 hỏi bão có làm việc chọn khó hơn không; H2 hỏi thanh khoản cạn và tương quan cao có đi cùng nỗi khó đó không; H4 (thăm dò) hỏi cổ phiếu có “chạy theo đàn” không; H3 hỏi kết hợp mô hình hoặc danh mục vững chắc có giúp ích sau chi phí không.

\section{Cách làm: bảy cách chia tiền thi đấu với nhau}

Đề tài định lượng trên dữ liệu thị trường công khai, không khảo sát.

Dữ liệu là giá đóng cửa điều chỉnh theo ngày của cổ phiếu trên HOSE (kể cả mã đã hủy niêm yết) và VN-Index, giai đoạn dự kiến 2007--2026. Mỗi ngày tái cân bằng, rổ gồm 30 cổ phiếu có giá trị giao dịch bình quân ba tháng và vốn hóa tự do chuyển nhượng lớn nhất, xác định chỉ bằng dữ liệu có tại ngày đó.

Mô hình chuyển chế độ Markov chấm xác suất thị trường đang căng thẳng, ước lượng lại tại mỗi ngày tái cân bằng chỉ bằng dữ liệu đến ngày đó.

Bảy mô hình M1 đến M7 thi đấu: chia đều 1/N; phương sai nhỏ nhất với hiệp phương sai co rút; trung bình-phương sai của Markowitz; Black-Litterman; Black-Litterman có chuyển chế độ; danh mục vững chắc với nhiều niềm tin tiên nghiệm; giảm tỷ trọng rủi ro theo xác suất chế độ căng thẳng.

Mỗi mô hình được chấm bằng lợi suất tương đương chắc chắn sau chi phí giao dịch: mức lợi suất chắc chắn mà nhà đầu tư ngại rủi ro coi là ngang giá với danh mục.

Từ bảy điểm số, đề tài tính mức hối tiếc R, độ phân tán D (độ lệch chuẩn của bảy điểm) và kích thước tập tin cậy mô hình (số ứng viên còn lại sau sơ loại thống kê). Hai thước đo sau chỉ báo hiệu bất định; mức hối tiếc mới đo rủi ro.

Đề tài kiểm định bằng hồi quy có sai số chuẩn bền vững, bootstrap khối, hiệu chỉnh so sánh bội và một phân phối rỗng 1.000 lần tái lấy mẫu để tách rủi ro lựa chọn thật khỏi nhiễu ước lượng.

\section{Điều chưa biết, nói thẳng}

Những điều chưa ai biết, kể cả tác giả:

\begin{itemize}
\item
  Chưa có kết quả; mọi kết quả ở dạng “nếu dữ liệu ủng hộ”.
\item
  Mức hối tiếc chỉ tính được từ khoảng năm thứ năm của mẫu (cần 104 tuần ước lượng và 36 tháng lịch sử ngoài mẫu), tức khoảng 28 cửa sổ đánh giá. Dưới 10 cửa sổ độc lập ở chế độ căng thẳng thì kết quả là không kết luận.
\item
  Tác giả đã dùng dữ liệu 2019--2025 trong bài trước; đăng ký trước giảm thiên lệch này, không loại bỏ nó.
\item
  Liên hệ tìm được, nếu có, là liên hệ có điều kiện, không phải nhân quả.
\item
  Độ dài các cửa sổ và hệ số ngại rủi ro bằng 5 là giá trị dự kiến, sẽ biện minh sau ước lượng thử.
\item
  Kết luận chỉ cho Việt Nam; thủ tục giám sát đề xuất chưa được kiểm nghiệm với người dùng; đề tài không quan sát hành vi chọn mô hình của tổ chức và không đánh giá khung quy định.
\end{itemize}

Nếu giả thuyết bị bác bỏ, câu trả lời vẫn có giá trị: rủi ro không đổi theo chế độ, hoặc 1/N đủ dùng ở Việt Nam.

\section{Vì sao điều này đáng quan tâm}

Giám sát rủi ro mô hình đang thay đổi. Cục Dự trữ Liên bang Mỹ thay thư SR 11-7 (04/04/2011) bằng SR 26-2 (17/04/2026), theo hướng quản lý rủi ro mô hình tương xứng với quy mô và mức rủi ro của từng tổ chức. Các văn bản này chỉ là chuẩn tham chiếu quốc tế.

Ở Việt Nam, Thông tư 83/2025/TT-NHNN (hiệu lực từ 01/07/2026, thay Thông tư 13/2018/TT-NHNN) quy định hệ thống kiểm soát nội bộ của ngân hàng thương mại, có khuôn khổ quản lý rủi ro mô hình. Người dùng kết quả của đề tài là công ty quản lý quỹ, nên tác giả sẽ xác định khung áp dụng cho họ trong năm thứ nhất.

Nếu dữ liệu ủng hộ, đề tài có ba nhóm kết quả: khung khái niệm và thước đo (lý thuyết); bằng chứng Việt Nam về mức độ, thời điểm và các kênh (thực nghiệm); thủ tục giám sát mô hình ba bước: khai báo tập mô hình, đo độ phân tán và tập tin cậy định kỳ, kích hoạt kết hợp mô hình khi căng thẳng (thực tiễn).

\section{Cách dùng bài này để ôn thi}

Bài đọc đi theo thứ tự mục của đề cương nộp ngày 30/09/2026; tiêu đề mỗi phần ghi số mục gốc, ví dụ “Mục 4.1 của đề cương”. Mỗi phần mở bằng kết luận, giải thích từ bước đầu, rồi kết bằng khối Ghi nhớ khi đi thi và Tự kiểm tra.

Phần ghi “Hình dung” hoặc “Ví dụ giả định” chỉ là công cụ giải thích; con số trong đó do người viết đặt, không phải số liệu của đề cương.

Hãy tự trả lời câu Tự kiểm tra trước khi đọc đáp án, và tập nói bằng lời của mình: rủi ro lựa chọn mô hình là gì, vì sao đo bằng mức hối tiếc, chế độ căng thẳng khác khủng hoảng ở đâu, vì sao đây mới là đề cương.

\begin{tcolorbox}[breakable]
\textbf{Ghi nhớ khi đi thi}

• Câu hỏi gốc: chọn sai cách chia tiền tốn bao nhiêu, và có tốn hơn khi thị trường căng thẳng không.

• Rủi ro lựa chọn mô hình đo bằng mức hối tiếc: chênh lệch lợi suất tương đương chắc chắn giữa mô hình tốt nhất sau thực tế và mô hình chọn trước.

• Câu hỏi 1 ứng với H1, câu hỏi 2 với H2 và H4, câu hỏi 3 với H3; H4 chỉ thăm dò.

• Bảy mô hình M1 đến M7 thi đấu trên dữ liệu HOSE; chưa có kết quả nào.

• Liên hệ có điều kiện, không phải nhân quả; dưới 10 cửa sổ độc lập ở chế độ căng thẳng thì gọi là không kết luận.
\end{tcolorbox}

\textbf{Tự kiểm tra}

{\itshape 1. Rủi ro lựa chọn mô hình khác rủi ro mô hình ở điểm nào?\par}

Đáp án: rủi ro mô hình là khả năng có hậu quả bất lợi từ mô hình sai hoặc dùng sai; rủi ro lựa chọn mô hình là tổn thất kỳ vọng, đo bằng mức hối tiếc, do phải chọn trước khi biết mô hình nào tốt nhất.

{\itshape 2. Chế độ căng thẳng và khủng hoảng có dùng thay nhau được không?\par}

Đáp án: không. Khủng hoảng là sự kiện lịch sử; chế độ căng thẳng là trạng thái thống kê do mô hình ước lượng.

{\itshape 3. Đề cương có kết quả thực nghiệm chưa, và nếu giả thuyết bị bác bỏ thì sao?\par}

Đáp án: chưa. Nếu bị bác bỏ, đề tài vẫn trả lời được: rủi ro không đổi theo chế độ, hoặc 1/N đủ dùng.

{\itshape 4. Vì sao Việt Nam được chọn làm bối cảnh?\par}

Đáp án: Việt Nam là bối cảnh kiểm chứng, không phải lý do của khoảng trống; thanh khoản hạn chế và nhiều cú sốc có thể làm chênh lệch giữa các mô hình lớn.

\chapter{Mục 1 và mục 4.4 của đề cương: Tên đề tài, lĩnh vực và phạm vi}

Đề tài “Rủi ro lựa chọn mô hình trong quản trị danh mục đầu tư theo chế độ thị trường: Bằng chứng từ thị trường chứng khoán Việt Nam” thuộc chuyên ngành Kinh tế và Quản lý, nhóm tài chính định lượng và quản trị rủi ro, với hai giao thoa có giới hạn: một phép thử về kinh tế hành vi và phần thiết kế thủ tục giám sát mô hình về quản lý.

Giám khảo có thể hỏi: đề tài thuộc lĩnh vực nào, giao thoa với hành vi ở mức nào, và đề tài không làm gì. Bạn cần trả lời được cả ba.

\section{Tên đề tài, đọc từng cụm}

Tên đề tài có 27 từ, trong giới hạn 30 từ của mẫu đề cương. Bản tiếng Anh dùng khi công bố: Model Selection Risk in Portfolio Management across Market Regimes: Evidence from the Vietnamese Stock Market. Chuyên ngành Kinh tế và Quản lý, mã số 9310116.01QTD, Trường Quốc tế, Đại học Quốc gia Hà Nội.

Để nhớ tên, hãy tách thành bốn cụm, mỗi cụm trả lời một câu hỏi.

{\small\begin{longtable}[]{@{}
  >{\raggedright\arraybackslash}p{(\columnwidth - 4\tabcolsep) * \real{0.3768}}
  >{\raggedright\arraybackslash}p{(\columnwidth - 4\tabcolsep) * \real{0.3411}}
  >{\raggedright\arraybackslash}p{(\columnwidth - 4\tabcolsep) * \real{0.2821}}@{}}
\toprule\noalign{}
\begin{minipage}[b]{\linewidth}\raggedright
Cụm trong tên
\end{minipage} & \begin{minipage}[b]{\linewidth}\raggedright
Thuật ngữ gốc
\end{minipage} & \begin{minipage}[b]{\linewidth}\raggedright
Trả lời câu hỏi
\end{minipage} \\
\midrule\noalign{}
\endhead
\bottomrule\noalign{}
\endlastfoot
Rủi ro lựa chọn mô hình & Model selection risk & Đo cái gì \\
Quản trị danh mục đầu tư & Portfolio management & Trong hoạt động nào \\
Theo chế độ thị trường & Across market regimes & So sánh giữa những hoàn cảnh nào \\
Bằng chứng từ thị trường chứng khoán Việt Nam & Evidence from the Vietnamese stock market & Bằng dữ liệu của đâu \\
\end{longtable}}

Cụm thứ nhất là thứ được đo: mức hối tiếc kỳ vọng khi phải chọn một cách chia tiền trước lúc biết cách nào tốt nhất.

Cụm thứ hai nói đây là việc của người quản lý danh mục: quyết định rổ cổ phiếu gồm những mã nào và mỗi mã chiếm bao nhiêu phần.

Cụm thứ ba nói câu hỏi không dừng ở “rủi ro lớn hay nhỏ” mà hỏi rủi ro đổi thế nào khi thị trường chuyển từ chế độ bình thường (trời quang) sang chế độ căng thẳng (mùa bão).

Cụm thứ tư giới hạn phạm vi: dữ liệu và kết luận thuộc thị trường chứng khoán Việt Nam, cụ thể là cổ phiếu niêm yết trên HOSE. Việt Nam là bối cảnh kiểm chứng, không phải lý do của khoảng trống.

\section{Lĩnh vực: tài chính định lượng và quản trị rủi ro}

“Định lượng” nghĩa là câu trả lời đến từ số liệu và mô hình thống kê. Đề tài không dùng khảo sát hay phương pháp định tính, nên không có thang đo nào cần dịch thuật.

Hình dung lĩnh vực này như phòng kiểm tra chất lượng của một nhà máy làm ra quyết định đầu tư. Phòng kiểm tra không hỏi sản phẩm đẹp không, mà hỏi quy trình làm ra nó đáng tin đến mức nào khi điều kiện thay đổi. Đối tượng đo của đề tài là chất lượng của một quyết định, không phải lợi nhuận của một cổ phiếu.

\section{Giao thoa với kinh tế hành vi: một phép thử giới hạn}

Đề tài giao thoa với kinh tế hành vi ở đúng một phép thử, đề cương gọi là “phép thử giới hạn”. Nắm chắc điểm này để không nói quá.

Đề tài đo độ phân tán chéo của lợi suất và dùng nó làm thước đo bầy đàn ở cấp thị trường. Trong một ngày, nếu mọi cổ phiếu cùng tăng hay cùng giảm gần một nhịp, lợi suất sát nhau và độ phân tán chéo thấp; nếu mỗi mã đi một hướng, độ phân tán cao.

Bầy đàn, theo đề cương, là việc nhà đầu tư làm theo số đông thay vì dựa vào thông tin riêng, như đàn cừu chạy theo con đầu đàn. Khi nhiều người làm theo số đông, các cổ phiếu có xu hướng đi cùng nhau và độ phân tán chéo bị nén lại.

Câu hỏi của phép thử: sau khi kiểm soát biến động, độ phân tán chéo của lợi suất có đi cùng rủi ro lựa chọn mô hình không. Đó là giả thuyết H4, chỉ được báo cáo như kết quả thăm dò.

Đề tài không quan sát hành vi chọn mô hình của tổ chức đầu tư: không phỏng vấn, không khảo sát các quỹ. Nếu giám khảo hỏi “đề tài có chứng minh nhà đầu tư bầy đàn không”, câu trả lời đúng là không, và đề tài không tuyên bố làm điều đó.

Đề tài vẫn đưa bầy đàn vào vì đây là một kênh hành vi khả dĩ, dù bằng chứng trái chiều giữa các thị trường. Christie và Huang (1995) không thấy bầy đàn khi giá biến động mạnh ở Mỹ; Chang và cộng sự (2000) thấy bầy đàn ở Hàn Quốc và Đài Loan, không thấy ở Mỹ và Hồng Kông; Vo và Phan (2017) thấy bầy đàn trên HOSE, rõ hơn khi thị trường giảm. Vì bằng chứng trái chiều, đề cương chỉ xét kênh này ở mức thăm dò.

\section{Giao thoa với quản lý: thiết kế cơ chế giám sát}

Giao thoa thứ hai nằm ở phần thiết kế thủ tục giám sát mô hình: ngoài việc đo, đề tài dự kiến rút ra một quy trình mà tổ chức đầu tư có thể dùng để kiểm soát việc chọn mô hình.

Quy trình gồm ba bước: khai báo tập mô hình; đo độ phân tán và tập tin cậy mô hình định kỳ; kích hoạt kết hợp mô hình khi thị trường ở chế độ căng thẳng. Đây là đề xuất suy ra từ H3, chưa được kiểm nghiệm với người dùng, và chỉ có nếu dữ liệu ủng hộ.

Phần này có hai giới hạn: khung giám sát của Cục Dự trữ Liên bang Mỹ chỉ là tham chiếu quốc tế, và đề tài không đánh giá mức đầy đủ của khung quy định Việt Nam.

\section{Phạm vi thực nghiệm}

Đối tượng thực nghiệm là danh mục cổ phiếu niêm yết trên HOSE. Đề tài không xét phái sinh, trái phiếu doanh nghiệp hay tài sản ngoài cổ phiếu.

Trái phiếu doanh nghiệp xuất hiện trong đề cương chỉ như một cú sốc lịch sử (thị trường đóng băng quý IV/2022), không phải tài sản trong danh mục. Lợi suất trái phiếu chính phủ kỳ hạn một năm cũng chỉ là đầu vào làm lãi suất phi rủi ro.

\begin{tcolorbox}[breakable]
\textbf{Ghi nhớ khi đi thi}

• Lĩnh vực: tài chính định lượng và quản trị rủi ro; chuyên ngành Kinh tế và Quản lý, mã số 9310116.01QTD.

• Giao thoa với kinh tế hành vi chỉ ở một phép thử: độ phân tán chéo của lợi suất có đi cùng rủi ro lựa chọn mô hình sau khi kiểm soát biến động không.

• Đề tài không quan sát hành vi chọn mô hình của tổ chức đầu tư.

• Giao thoa với quản lý nằm ở thủ tục giám sát mô hình; đó là đề xuất chưa kiểm nghiệm.

• Phạm vi: cổ phiếu niêm yết trên HOSE.
\end{tcolorbox}

\textbf{Tự kiểm tra}

{\itshape 1. Đề tài giao thoa với kinh tế hành vi ở điểm nào, và ở mức nào?\par}

Đáp án: ở một phép thử giới hạn: độ phân tán chéo của lợi suất (thước đo bầy đàn cấp thị trường) có đi cùng rủi ro lựa chọn mô hình sau khi kiểm soát biến động không.

{\itshape 2. Đề tài có quan sát cách các tổ chức đầu tư thực sự chọn mô hình không?\par}

Đáp án: không.

{\itshape 3. Giao thoa với quản lý nằm ở đâu?\par}

Đáp án: ở phần thiết kế thủ tục giám sát mô hình.

{\itshape 4. Nêu chuyên ngành, nhóm lĩnh vực và phạm vi thực nghiệm.\par}

Đáp án: chuyên ngành Kinh tế và Quản lý; nhóm tài chính định lượng và quản trị rủi ro; cổ phiếu niêm yết trên HOSE.

\chapter{Mục 4.1 của đề cương: Lý do chọn đề tài (phần một, từ Markowitz đến quy tắc 1/N)}

Mục 4.1 mở đầu bằng một chuỗi lập luận: tổ chức đầu tư chia tiền bằng mô hình thống kê; mô hình nào cũng có sai số ở đầu vào; DeMiguel và cộng sự (2009) so sánh 14 mô hình và không thấy mô hình nào vượt nhất quán quy tắc 1/N. Từ đó, việc chọn mô hình là một quyết định có rủi ro riêng.

Phần này giải thích từng mắt xích của chuỗi đó. Phần hai giải thích các cú sốc ở Việt Nam, khung giám sát, khoảng trống và ba câu hỏi.

\section{Bước 1: tổ chức đầu tư chia tiền bằng mô hình}

Một quỹ đầu tư cầm tiền của nhiều người và phải quyết định mua cổ phiếu nào, mỗi mã bao nhiêu phần. Tỷ lệ của mỗi mã gọi là trọng số danh mục.

Mô hình danh mục là một cách chia tiền: nhận dữ liệu quá khứ, trả ra một dãy tỷ lệ. Cùng dữ liệu, các cách chia khác nhau cho ra những dãy tỷ lệ khác nhau.

Điểm khởi đầu là khung trung bình-phương sai của Markowitz (1952): chọn danh mục có lợi suất kỳ vọng cao nhất ứng với một mức phương sai cho trước. Hãy nhớ nguyên cách diễn đạt này.

\section{Bước 2: trung bình-phương sai nghĩa là gì}

Cần ba khái niệm: lợi suất, lợi suất kỳ vọng và phương sai.

Lợi suất là phần trăm lãi hoặc lỗ của khoản đầu tư trong một khoảng thời gian.

\begin{tcolorbox}[breakable]
Ví dụ giả định (số tự đặt để minh họa, không phải số liệu của đề cương):

Mua cổ phiếu giá 100, cuối kỳ giá 103 thì lợi suất là 3\%.
\end{tcolorbox}

Lợi suất kỳ vọng là mức bạn trông đợi trung bình theo ước tính của mình; nó nhìn về tương lai nên không bao giờ biết chắc.

Phương sai đo mức lợi suất dao động quanh mức trung bình. Hai cổ phiếu có thể cùng lợi suất trung bình, nhưng mã thứ nhất luôn quanh mức đó còn mã thứ hai lúc lãi cao, lúc lỗ nặng; mã thứ hai có phương sai lớn hơn. Trong tài chính, phương sai được dùng làm thước đo rủi ro.

\begin{tcolorbox}[breakable]
Ví dụ giả định (số tự đặt để minh họa, không phải số liệu của đề cương):

Hai cổ phiếu A và B cùng có lợi suất kỳ vọng 1\% mỗi tháng. Lợi suất tháng của A thường từ 0\% đến 2\%; của B từ lỗ 8\% đến lãi 10\%. Nếu muốn lợi suất ổn định, bạn chọn A vì phương sai nhỏ hơn.
\end{tcolorbox}

Khi có nhiều cổ phiếu, cần thêm hiệp phương sai: đại lượng đo hai cổ phiếu có xu hướng đi cùng nhau hay ngược nhau. Hai mã đi cùng nhau giống hai người đạp xe song song, một người tăng tốc thì người kia thường cũng tăng tốc.

Hiệp phương sai quyết định lợi ích của việc dàn trải. Chia tiền vào hai mã hay đi ngược nhau thì lỗ của mã này được bù bởi lãi của mã kia; chia vào hai mã luôn đi cùng nhau thì dàn trải chẳng giúp gì.

Ghép lại: khung trung bình-phương sai tìm cách chia tiền cho lợi suất kỳ vọng cao nhất trong giới hạn dao động mà bạn chịu được.

Công thức cần hai đầu vào: lợi suất kỳ vọng của từng mã và ma trận hiệp phương sai của mọi cặp mã. Cả hai không quan sát được trực tiếp; người dùng phải ước lượng từ dữ liệu quá khứ.

\section{Bước 3: sai số ước lượng truyền sang trọng số}

Ước lượng từ quá khứ luôn có sai số, và sai số khi ước lượng lợi suất kỳ vọng và hiệp phương sai truyền sang trọng số danh mục (Michaud, 1989; Kan và Zhou, 2007).

Hình dung bạn chia bột làm bánh bằng một chiếc cân lệch vài gram. Nếu công thức khuếch đại mọi sai lệch, chiếc bánh sẽ hỏng. Công thức trung bình-phương sai nhạy với sai số đầu vào theo đúng cách đó.

Phần tổng quan (mục 4.2, nhóm thứ nhất) nói rõ: tối ưu hóa dồn trọng số vào tài sản có sai số ước lượng lớn (Michaud, 1989), và thay tham số tổng thể bằng ước lượng mẫu làm hiệu quả ngoài mẫu giảm mạnh (Kan và Zhou, 2007).

“Trong mẫu” là dữ liệu dùng để ước lượng mô hình; “ngoài mẫu” là dữ liệu mô hình chưa thấy, tức tương lai thật. Một cách chia tiền có thể trông hoàn hảo trong mẫu vì được “vẽ” vừa khít quá khứ; bài kiểm tra thật là ngoài mẫu.

\begin{tcolorbox}[breakable]
Ví dụ giả định (số tự đặt để minh họa, không phải số liệu của đề cương):

Một học sinh học thuộc đáp án đề năm ngoái và làm đề đó được 10 điểm (trong mẫu). Đề năm nay khác, em chỉ được 5 điểm (ngoài mẫu). Mô hình ước lượng quá sát quá khứ cũng có thể như vậy.
\end{tcolorbox}

Đề cương không nói trung bình-phương sai là mô hình tồi. Sai số đầu vào là vấn đề chung của mọi mô hình phân bổ, và không mô hình nào thắng ổn định.

\section{Bước 4: các cách khắc phục, và vì sao chúng chưa chấm dứt câu chuyện}

Vì sai số đầu vào là vấn đề thật, tài liệu đề xuất ba hướng khắc phục, và cả ba có mặt trong tập bảy mô hình của đề tài.

\begin{itemize}
\item
  Co rút hiệp phương sai (Ledoit và Wolf, 2004): kéo ước lượng về gần một ma trận mục tiêu có cấu trúc để bớt cực đoan.
\item
  Black-Litterman (Black và Litterman, 1992): kết hợp lợi suất cân bằng của mô hình định giá tài sản vốn với quan điểm của nhà đầu tư. Bessler và cộng sự (2017) thấy cách này vượt cả trung bình-phương sai lẫn 1/N về hiệu quả ngoài mẫu trên dữ liệu đa tài sản toàn cầu.
\item
  Danh mục vững chắc (Garlappi và cộng sự, 2007): mô hình hóa nhà đầu tư né tránh mơ hồ với nhiều niềm tin tiên nghiệm, thu được danh mục ổn định hơn.
\end{itemize}

Kết luận ngay sau đó mới là điểm chính: không hướng nào thắng ổn định (DeMiguel và cộng sự, 2009). Pflug và cộng sự (2012) còn chứng minh quy tắc 1/N xấp xỉ tối ưu theo nghĩa vững chắc khi mức mơ hồ về mô hình đủ cao. Điều này mở ra câu hỏi của đề tài.

\section{Bước 5: nghiên cứu DeMiguel và cộng sự (2009)}

DeMiguel và cộng sự (2009) đánh giá 14 mô hình trên bảy bộ dữ liệu và không thấy mô hình nào vượt nhất quán quy tắc 1/N, xét theo tỷ số Sharpe, lợi suất tương đương chắc chắn hoặc vòng quay danh mục.

Quy tắc 1/N đơn giản đến bất ngờ: có N tài sản thì chia đều, mỗi tài sản 1/N vốn, như bỏ mỗi giỏ một quả trứng. Không cần ước lượng gì.

Đừng nói 1/N “tốt hơn” nói chung. Điều nghiên cứu này nói chỉ là không mô hình nào vượt nhất quán 1/N trong bảy bộ dữ liệu.

Ba tiêu chí của họ như sau.

Tỷ số Sharpe là lợi suất vượt lãi suất phi rủi ro trên mỗi đơn vị độ lệch chuẩn. Lãi suất phi rủi ro là mức lãi gần như chắc chắn; đề tài dùng lợi suất trái phiếu chính phủ kỳ hạn một năm. Phần lợi suất vượt là phần thưởng vì chấp nhận rủi ro; độ lệch chuẩn là căn bậc hai của phương sai.

Sharpe trả lời: mỗi đơn vị rủi ro mang lại bao nhiêu phần thưởng. Hai tay đua cùng về đích sau một giờ, người đi đường bằng đạt cùng thành tích với ít rủi ro hơn người đi đường đèo, nên có Sharpe cao hơn.

\begin{tcolorbox}[breakable]
Ví dụ giả định (số tự đặt để minh họa, không phải số liệu của đề cương):

Danh mục X có lợi suất vượt 6\% mỗi năm, độ lệch chuẩn 12\%: Sharpe bằng 0,5. Danh mục Y có lợi suất vượt 6\%, độ lệch chuẩn 20\%: Sharpe bằng 0,3. Cùng phần thưởng, X rủi ro thấp hơn nên Sharpe cao hơn.
\end{tcolorbox}

Tiêu chí thứ hai là lợi suất tương đương chắc chắn: mức lợi suất chắc chắn mà nhà đầu tư coi là ngang giá với khoản đầu tư có rủi ro; càng cao càng tốt. Khác Sharpe, nó phản ánh mức ngại rủi ro của chính nhà đầu tư.

Ở mục 4.3.6, đề tài cho đại lượng này một công thức gắn với hệ số ngại rủi ro $\rho$ và chọn nó làm thước đo chính. Lý do: khi lợi suất vượt trội trung bình âm, đúng ở chế độ căng thẳng, danh mục biến động hơn lại có Sharpe ít âm hơn, nên Sharpe không còn xếp đúng thứ tự chất lượng. Sharpe chỉ là thước đo phụ.

Tiêu chí thứ ba là vòng quay danh mục, tỷ lệ danh mục được mua bán lại. Vòng quay cao thì tốn nhiều chi phí giao dịch.

Ba tiêu chí hỏi ba điều khác nhau: phần thưởng trên mỗi đơn vị rủi ro, mức tương đương chắc chắn, và chi phí mua bán. Không mô hình nào thắng 1/N nhất quán ở cả ba.

\section{Bước 6: từ “không mô hình nào thắng” đến “chọn mô hình là một quyết định có rủi ro”}

Đây là bước suy luận quan trọng nhất của mục 4.1, gói trong một câu: khi không mô hình nào thắng ổn định, việc chọn mô hình trở thành một quyết định mang rủi ro riêng.

Nếu có một mô hình thắng mọi lúc, bạn chỉ cần tìm nó. Nếu không, mô hình tốt nhất có thể đổi theo thời điểm hoặc tập dữ liệu; bạn phải chọn, và có thể chọn sai.

Chọn sai nghĩa là tiền bị chia theo cách kém hơn một cách khác lẽ ra có sẵn. Rủi ro thị trường là cổ phiếu có thể giảm giá; rủi ro lựa chọn mô hình là bạn có thể đã chọn sai cách chia tiền.

\begin{tcolorbox}[breakable]
Ví dụ giả định (số tự đặt để minh họa, không phải số liệu của đề cương):

Ba cách chia tiền P, Q, R cho kết quả sau mùa giải 4\%, 7\% và 3\%. Nhà quản lý chọn P vì P thắng ba năm gần nhất. Cách tốt nhất sau thực tế là Q, nên mức hối tiếc là 7\% − 4\% = 3\% (ở đây đo thô bằng lợi suất; đề cương đo bằng lợi suất tương đương chắc chắn).
\end{tcolorbox}

Mức hối tiếc không phải là lỗ: nhà quản lý vẫn lãi 4\%. Vì mô hình chọn trước nằm trong tập mô hình, mức hối tiếc không bao giờ âm. Ngoại lệ ở mục 4.3.8: hai cơ chế giảm thiểu nằm ngoài tập mô hình khi kiểm định H3 nên mức hối tiếc của chúng có thể âm.

Ba khái niệm dễ lẫn được đề cương định nghĩa ngay trong mục 4.1:

{\small\begin{longtable}[]{@{}
  >{\raggedright\arraybackslash}p{(\columnwidth - 4\tabcolsep) * \real{0.1713}}
  >{\raggedright\arraybackslash}p{(\columnwidth - 4\tabcolsep) * \real{0.1657}}
  >{\raggedright\arraybackslash}p{(\columnwidth - 4\tabcolsep) * \real{0.6631}}@{}}
\toprule\noalign{}
\begin{minipage}[b]{\linewidth}\raggedright
Khái niệm
\end{minipage} & \begin{minipage}[b]{\linewidth}\raggedright
Thuật ngữ gốc
\end{minipage} & \begin{minipage}[b]{\linewidth}\raggedright
Nghĩa theo đề cương
\end{minipage} \\
\midrule\noalign{}
\endhead
\bottomrule\noalign{}
\endlastfoot
Rủi ro mô hình & Model risk & Khả năng phát sinh hậu quả bất lợi từ quyết định dựa trên mô hình sai hoặc dùng sai (theo SR 11-7) \\
Bất định mô hình & Model uncertainty & Việc nhiều mô hình khả dĩ cho kết quả khác nhau \\
Rủi ro lựa chọn mô hình & Model selection risk & Tổn thất kỳ vọng, đo bằng mức hối tiếc, do phải chọn một mô hình trước khi biết mô hình nào tốt nhất \\
\end{longtable}}

Bất định mô hình là tình trạng: nhiều mô hình và chúng bất đồng. Rủi ro lựa chọn mô hình là hậu quả kỳ vọng của tình trạng đó khi bạn phải chọn một. Độ phân tán mô hình và kích thước tập tin cậy mô hình chỉ báo hiệu bất định, không đo rủi ro.

\section{Một điểm về tính thận trọng của lập luận}

Nếu được hỏi, hãy nói ra giới hạn của lập luận: kết luận dựa trên DeMiguel và cộng sự (2009) với bảy bộ dữ liệu, không gồm Việt Nam, và không chứng minh quy luật đúng ở mọi nơi.

Lập luận cũng không nói 1/N luôn tốt hơn. 1/N vừa là mô hình M1, vừa là một khả năng đối lập của H1: nếu 1/N thắng ổn định, H1 không được ủng hộ.

\begin{tcolorbox}[breakable]
\textbf{Ghi nhớ khi đi thi}

• Markowitz (1952): chọn danh mục có lợi suất kỳ vọng cao nhất với phương sai cho trước.

• Sai số ước lượng lợi suất kỳ vọng và hiệp phương sai truyền sang trọng số danh mục (Michaud 1989; Kan và Zhou 2007).

• DeMiguel và cộng sự (2009): 14 mô hình, bảy bộ dữ liệu, không mô hình nào vượt nhất quán 1/N; Pflug và cộng sự (2012): 1/N xấp xỉ tối ưu khi mơ hồ cao.

• Kết luận: chọn mô hình là một quyết định có rủi ro riêng.

• Ở chế độ căng thẳng Sharpe không xếp đúng thứ tự chất lượng nên chỉ là thước phụ; lợi suất tương đương chắc chắn là thước chính.
\end{tcolorbox}

\textbf{Tự kiểm tra}

{\itshape 1. Khung trung bình-phương sai của Markowitz (1952) chọn danh mục theo tiêu chí nào?\par}

Đáp án: danh mục có lợi suất kỳ vọng cao nhất với phương sai cho trước.

{\itshape 2. DeMiguel và cộng sự (2009) đánh giá bao nhiêu mô hình trên bao nhiêu bộ dữ liệu, và kết luận gì?\par}

Đáp án: 14 mô hình, bảy bộ dữ liệu; không mô hình nào vượt nhất quán 1/N về Sharpe, lợi suất tương đương chắc chắn hoặc vòng quay.

{\itshape 3. Tỷ số Sharpe được định nghĩa thế nào, và vì sao đề tài chỉ dùng nó làm thước phụ?\par}

Đáp án: lợi suất vượt lãi suất phi rủi ro trên mỗi đơn vị độ lệch chuẩn. Khi lợi suất vượt trội trung bình âm (chế độ căng thẳng), Sharpe không còn xếp đúng thứ tự chất lượng, nên chỉ báo cáo ở cửa sổ có lợi suất vượt trội trung bình dương.

{\itshape 4. Từ kết quả của DeMiguel và cộng sự (2009), đề cương suy ra điều gì?\par}

Đáp án: việc chọn mô hình là một quyết định có rủi ro riêng.

{\itshape 5. Độ phân tán mô hình và kích thước tập tin cậy mô hình là chỉ báo của rủi ro hay của bất định?\par}

Đáp án: của bất định.

\chapter{Mục 4.1 của đề cương: Lý do chọn đề tài (phần hai, cú sốc, giám sát, khoảng trống và ba câu hỏi)}

Nửa sau của mục 4.1 làm bốn việc: nêu các cú sốc lịch sử để thấy “mùa bão” ở Việt Nam trông thế nào; nói khung giám sát rủi ro mô hình đang thay đổi; chỉ ra hai khoảng trống; rồi đặt ba câu hỏi.

\section{Khủng hoảng lặp lại, và ở Việt Nam có ba loại cú sốc}

Khủng hoảng tài chính lặp lại qua nhiều thế kỷ (Reinhart và Rogoff, 2009); Laeven và Valencia (2020) ghi nhận 151 cuộc khủng hoảng ngân hàng hệ thống trong giai đoạn 1970--2017. Con số này nói chung về thế giới, không gắn riêng với Việt Nam.

Thị trường chứng khoán Việt Nam trải qua ít nhất ba loại cú sốc trong hai thập niên gần đây. Đề cương dẫn bốn công trình:

{\small\begin{longtable}[]{@{}
  >{\raggedright\arraybackslash}p{(\columnwidth - 2\tabcolsep) * \real{0.2063}}
  >{\raggedright\arraybackslash}p{(\columnwidth - 2\tabcolsep) * \real{0.7937}}@{}}
\toprule\noalign{}
\begin{minipage}[b]{\linewidth}\raggedright
Công trình
\end{minipage} & \begin{minipage}[b]{\linewidth}\raggedright
Điều đề cương nói về công trình đó
\end{minipage} \\
\midrule\noalign{}
\endhead
\bottomrule\noalign{}
\endlastfoot
Vo và Phan (2017) & Giai đoạn 2005--2015, xét tác động của khủng hoảng tài chính toàn cầu lên hành vi nhà đầu tư \\
Dao và cộng sự (2021) & Ghi nhận số ca mắc COVID-19 hằng ngày tác động tiêu cực đến lợi suất cổ phiếu \\
Quỹ Tiền tệ Quốc tế (2023) & Ghi nhận thị trường trái phiếu doanh nghiệp đóng băng trong quý IV/2022; bất động sản chiếm tỷ trọng lớn nhất trong thị trường này \\
Vu và cộng sự (2024) & Các đợt sụt giảm giá cổ phiếu giai đoạn 2007--2022 lan truyền giữa các ngành \\
\end{longtable}}

Ba loại cú sốc là khủng hoảng tài chính toàn cầu, đại dịch COVID-19 năm 2020 và căng thẳng trái phiếu doanh nghiệp từ quý IV/2022; Vu và cộng sự (2024) bổ sung cách các đợt sụt giảm lan giữa các ngành. Ba cú sốc này cũng là các giai đoạn nằm trong mẫu dữ liệu ở mục 4.3.3.

Đề tài không đo trực tiếp các sự kiện này. Chúng là các giai đoạn tham chiếu để kiểm tra xem chế độ căng thẳng do mô hình ước lượng có trùng với biến cố thực tế không.

\section{Khủng hoảng và chế độ căng thẳng: hai khái niệm khác nhau}

Đây là chỗ dễ lẫn nhất. Khủng hoảng là sự kiện lịch sử, có tên và ngày tháng. Chế độ căng thẳng là trạng thái thống kê do mô hình ước lượng từ lợi suất VN-Index. Hai khái niệm không dùng thay nhau.

Khủng hoảng giống một cơn bão có tên mà báo chí đưa tin; chế độ căng thẳng giống việc trạm khí tượng ghi nhận “đang là mùa bão” dựa trên số đo gió và áp suất. Hai thứ có thể trùng nhau nhưng không nhất thiết: trạm đo có thể báo mùa bão trước khi cơn bão có tên xuất hiện, hoặc có bão có tên mà số đo chưa đủ ngưỡng.

Đề tài kiểm định rủi ro lựa chọn mô hình theo chế độ căng thẳng, rồi đối chiếu chế độ đó với các cuộc khủng hoảng; chương 3 của luận án làm việc đối chiếu này (Bảng 4 của đề cương).

Tham số của mô hình chuyển chế độ được ước lượng lại tại mỗi ngày tái cân bằng chỉ bằng dữ liệu đến ngày đó, để tránh thiên lệch nhìn trước. Chế độ căng thẳng vì vậy do dữ liệu xác định, không do tên gọi của khủng hoảng.

\section{Giám sát rủi ro mô hình đang thay đổi}

Cục Dự trữ Liên bang Mỹ ban hành thư giám sát SR 11-7 ngày 04/04/2011 và thay bằng SR 26-2 ngày 17/04/2026, theo hướng quản lý rủi ro mô hình tương xứng với quy mô, mức độ phức tạp và mức rủi ro mô hình của từng tổ chức.

Hình dung một bộ quy tắc an toàn phải áp cho mọi cửa hàng, từ quán nước đầu ngõ đến siêu thị. Hướng “tương xứng với rủi ro” nghĩa là mức kiểm soát tùy mức rủi ro của từng nơi. Đề cương chỉ nêu hướng thay đổi, không mô tả chi tiết văn bản.

Định nghĩa rủi ro mô hình mà đề tài dùng cũng lấy từ SR 11-7: khả năng phát sinh hậu quả bất lợi từ quyết định dựa trên mô hình sai hoặc dùng sai. Văn bản này là nguồn của định nghĩa nền tảng của đề tài.

Các văn bản này không ràng buộc tổ chức Việt Nam; đề tài dùng chúng làm chuẩn tham chiếu quốc tế. Nếu được hỏi “SR 11-7 có áp dụng cho quỹ ở Việt Nam không”, câu trả lời là không.

\section{Thông tư 83/2025/TT-NHNN và khoảng cách với đối tượng của đề tài}

Ở Việt Nam, Thông tư 83/2025/TT-NHNN của Ngân hàng Nhà nước, ban hành ngày 31/12/2025, có hiệu lực từ 01/07/2026, thay Thông tư 13/2018/TT-NHNN, quy định hệ thống kiểm soát nội bộ của ngân hàng thương mại và lần đầu thiết lập khuôn khổ quản lý rủi ro mô hình cho các ngân hàng.

Có một điểm lệch: khung này áp dụng cho ngân hàng, còn người dùng kết quả của đề tài là công ty quản lý quỹ và tổ chức đầu tư cổ phiếu.

Đề cương xử lý độ lệch này bằng một cam kết: trong quý 1 năm thứ nhất, tác giả xác định khung quy định áp dụng cho các tổ chức đó và mức họ dùng mô hình định lượng để phân bổ danh mục, từ báo cáo của cơ quan quản lý và tài liệu công bố của công ty quản lý quỹ.

Đề tài không đánh giá mức đầy đủ của khung quy định, tức không kết luận khung hiện hành “đủ” hay “thiếu”. Đề cương chưa nói khung nào áp dụng cho công ty quản lý quỹ; đó là việc của năm thứ nhất, người đọc không nên điền vào chỗ trống này.

\section{Khoảng trống cấu trúc, ở hai chỗ}

Phần then chốt: tài liệu hiện có để ngỏ hai vấn đề.

\textbf{Chỗ thứ nhất: thước đo rủi ro mô hình chưa đo hậu quả của việc chọn sai.} Các thước đo hiện có (Boucher và cộng sự, 2014; Danielsson và cộng sự, 2016) đo chênh lệch giữa các dự báo rủi ro, chưa đo hậu quả của việc chọn sai mô hình lên trọng số và kết quả danh mục.

Hai trạm khí tượng đưa ra hai dự báo khác nhau. Đo độ lệch giữa hai dự báo chưa cho biết buổi dã ngoại thiệt hại thế nào nếu bạn nghe nhầm trạm. Thước đo hiện có đo độ lệch giữa dự báo; đề tài muốn đo thiệt hại của buổi dã ngoại.

Danielsson và cộng sự (2016) thấy rủi ro mô hình tăng theo mức bất định của thị trường và nhỏ khi thị trường yên ắng, nhưng chỉ cho các mô hình dự báo rủi ro thị trường và rủi ro hệ thống. Boucher và cộng sự (2014) đề xuất hiệu chỉnh độ đo rủi ro theo kết quả kiểm định ngược. Cả hai nói về dự báo rủi ro, không về quyết định phân bổ.

\textbf{Chỗ thứ hai: nghiên cứu danh mục xếp hạng mô hình theo bình quân toàn mẫu.} DeMiguel và cộng sự (2009) và Garlappi và cộng sự (2007) so sánh mô hình bằng thứ hạng bình quân toàn mẫu, chưa xét thứ hạng đổi thế nào giữa các chế độ.

Một học sinh đứng đầu bình quân cả năm có thể đứng đầu nhờ học kỳ hai, còn học kỳ một chỉ đứng giữa bảng. Nếu cuộc thi diễn ra đúng học kỳ một, thứ hạng cả năm có thể làm bạn chọn sai. Thứ hạng toàn mẫu cũng có thể giấu khác biệt giữa trời quang và mùa bão.

\begin{tcolorbox}[breakable]
Ví dụ giả định (số tự đặt để minh họa, không phải số liệu của đề cương):

Hai mô hình A và B được chấm trên 10 giai đoạn. Ở 8 giai đoạn bình thường A đứng trên B; ở 2 giai đoạn căng thẳng B đứng trên A. Tính bình quân toàn mẫu thì A thắng, nhưng chọn ngay trước một giai đoạn căng thẳng thì A là lựa chọn sai.
\end{tcolorbox}

Câu hỏi của đề tài nằm ở giao điểm của hai chỗ trống: đo hậu quả của việc chọn mô hình, và đo theo chế độ.

\section{Bảng tóm tắt bốn khoảng trống}

Mục 4.2 của đề cương tóm tắt bốn khoảng trống trong Bảng 1; ba câu hỏi và phần đóng góp ở mục 4.4 đều dựa vào bảng này.

{\small\begin{longtable}[]{@{}
  >{\raggedright\arraybackslash}p{(\columnwidth - 4\tabcolsep) * \real{0.0797}}
  >{\raggedright\arraybackslash}p{(\columnwidth - 4\tabcolsep) * \real{0.4250}}
  >{\raggedright\arraybackslash}p{(\columnwidth - 4\tabcolsep) * \real{0.4953}}@{}}
\toprule\noalign{}
\begin{minipage}[b]{\linewidth}\raggedright
Mã
\end{minipage} & \begin{minipage}[b]{\linewidth}\raggedright
Nhóm nghiên cứu
\end{minipage} & \begin{minipage}[b]{\linewidth}\raggedright
Khoảng trống
\end{minipage} \\
\midrule\noalign{}
\endhead
\bottomrule\noalign{}
\endlastfoot
(1) & Rủi ro mô hình (Boucher và cộng sự, 2014; Danielsson và cộng sự, 2016) & Thước đo chưa gắn với hậu quả của quyết định phân bổ danh mục theo chế độ thị trường \\
(2) & Danh mục ở Việt Nam (Ta và Vuong, 2020; Nguyen và cộng sự, 2020; Goutte và cộng sự, 2025; Nguyen và Nguyen, 2026a) & Đánh giá từng mô hình, chưa đặt các mô hình vào một tập cạnh tranh để đo rủi ro lựa chọn theo chế độ \\
(3) & Chế độ thị trường ở Việt Nam (Hoang và Luu, 2024) & Mô hình hóa biến động VN-Index, chưa nối chế độ với chất lượng quyết định danh mục \\
(4) & Bầy đàn ở Việt Nam (Vo và Phan, 2017) & Bằng chứng bầy đàn trên HOSE, chưa gắn với rủi ro của việc chọn mô hình \\
\end{longtable}}

Điểm chung: các công trình đo từng mô hình, từng chế độ hoặc từng hiện tượng riêng lẻ; chưa công trình nào đo rủi ro của việc chọn giữa các mô hình theo chế độ thị trường.

\section{Việt Nam là bối cảnh kiểm chứng, không phải lý do của khoảng trống}

Câu dễ bị hiểu sai: Việt Nam là bối cảnh kiểm chứng, không phải lý do của khoảng trống.

Câu này nói hai điều. Hai khoảng trống là chuyện của cách đo và cách so sánh mô hình, không bắt nguồn từ Việt Nam. Còn Việt Nam phù hợp để kiểm chứng vì thanh khoản hạn chế và nhiều cú sốc trong lịch sử giao dịch ngắn, nên chênh lệch giữa các mô hình có thể lớn. Lịch sử ngắn cũng làm mẫu có ít giai đoạn căng thẳng độc lập, và đề cương thừa nhận điều này làm công suất kiểm định thấp (mục 4.3.10).

Nếu được hỏi “vì sao chọn Việt Nam”, đừng trả lời “vì Việt Nam chưa có nghiên cứu”.

\section{Ba câu hỏi và cú chuyển từ một mô hình sang nhiều mô hình}

Đề cương hỏi ba câu:

\begin{itemize}
\item
  Câu hỏi 1: rủi ro lựa chọn mô hình thay đổi thế nào giữa chế độ bình thường và chế độ căng thẳng.
\item
  Câu hỏi 2: sau khi kiểm soát biến động thực hiện, các biến kênh (thanh khoản, tương quan cực trị và, ở mức thăm dò, bầy đàn) có liên hệ thống kê với rủi ro lựa chọn mô hình không, và giải thích được bao nhiêu phần chênh lệch giữa hai chế độ.
\item
  Câu hỏi 3: cơ chế nào giảm được rủi ro lựa chọn mô hình sau chi phí giao dịch.
\end{itemize}

Câu hỏi 1 hỏi “bão có làm việc chọn khó hơn không”; câu hỏi 2 hỏi “nếu khó hơn, những dấu hiệu nào đi cùng”; câu hỏi 3 hỏi “làm gì để bớt thiệt”.

“Liên hệ thống kê” không có nghĩa là nguyên nhân. Đề tài gọi kết quả là liên hệ có điều kiện. Nếu dữ liệu cho thấy thanh khoản cạn đi cùng rủi ro lựa chọn cao (điều H2 dự kiến), đề tài cũng không khẳng định thanh khoản cạn gây ra rủi ro đó.

Mục 4.1 kết bằng câu: kết quả có thể làm tư liệu tham khảo cho thủ tục giám sát mô hình của công ty quản lý quỹ. “Tư liệu tham khảo” là mức khiêm tốn; đề cương không nói kết quả sẽ thay đổi quy định.

Đề tài chuyển từ đánh giá một mô hình sang câu hỏi chọn giữa nhiều mô hình. Bài trước của tác giả (Nguyen và Nguyen, 2026a) đánh giá một mô hình, Black-Litterman có chuyển chế độ; đề cương này hỏi “chọn giữa nhiều mô hình thì rủi ro đến đâu”. Mục 8.2 của đề cương trình bày sự chuẩn bị của tác giả cho bước chuyển này.

Bước chuyển giống việc đi từ chấm điểm một vận động viên sang đánh giá huấn luyện viên phải chọn ai ra sân.

\section{Những điều mục 4.1 không nói}

Để tránh nói quá, hãy nhớ những gì mục 4.1 không khẳng định:

\begin{itemize}
\item
  Không đưa ra kết quả thực nghiệm nào.
\item
  Không nói khung quy định của Việt Nam là thiếu.
\item
  Không nói SR 11-7 hay SR 26-2 ràng buộc tổ chức Việt Nam.
\item
  Không nói khoảng trống bắt nguồn từ đặc điểm của Việt Nam.
\item
  Không khẳng định chưa ai làm; tuyên bố tính mới ở mục 4.2 chỉ là “theo hiểu biết của thí sinh”, với phạm vi tìm kiếm có giới hạn.
\end{itemize}

\begin{tcolorbox}[breakable]
\textbf{Ghi nhớ khi đi thi}

• Laeven và Valencia (2020): 151 cuộc khủng hoảng ngân hàng hệ thống, 1970--2017.

• Ba loại cú sốc ở Việt Nam: khủng hoảng tài chính toàn cầu (Vo và Phan, 2017), COVID-19 (Dao và cộng sự, 2021), trái phiếu doanh nghiệp đóng băng quý IV/2022 (Quỹ Tiền tệ Quốc tế, 2023); sụt giảm lan giữa các ngành 2007--2022 (Vu và cộng sự, 2024).

• SR 11-7 (04/04/2011) được thay bằng SR 26-2 (17/04/2026); chỉ là chuẩn tham chiếu quốc tế.

• Thông tư 83/2025/TT-NHNN (hiệu lực 01/07/2026) áp dụng cho ngân hàng thương mại; người dùng kết quả là công ty quản lý quỹ; đề tài không đánh giá khung quy định.

• Hai khoảng trống: thước đo chưa đo hậu quả lên kết quả danh mục; xếp hạng theo bình quân toàn mẫu, chưa theo chế độ.

• Việt Nam là bối cảnh kiểm chứng, không phải lý do của khoảng trống.
\end{tcolorbox}

\textbf{Tự kiểm tra}

{\itshape 1. Đề cương phân biệt “khủng hoảng” và “chế độ căng thẳng” thế nào?\par}

Đáp án: khủng hoảng là sự kiện lịch sử; chế độ căng thẳng là trạng thái thống kê do mô hình ước lượng. Đề tài đối chiếu, không dùng thay nhau.

{\itshape 2. SR 11-7 và SR 26-2 là gì, và có ràng buộc tổ chức Việt Nam không?\par}

Đáp án: thư giám sát của Cục Dự trữ Liên bang Mỹ (04/04/2011 và 17/04/2026), theo hướng quản lý rủi ro mô hình tương xứng với từng tổ chức. Không ràng buộc tổ chức Việt Nam; chỉ là chuẩn tham chiếu.

{\itshape 3. Vì sao Thông tư 83/2025/TT-NHNN không khớp hoàn toàn với đối tượng của đề tài?\par}

Đáp án: thông tư áp dụng cho ngân hàng thương mại, còn người dùng kết quả là công ty quản lý quỹ và tổ chức đầu tư cổ phiếu. Năm thứ nhất, tác giả xác định khung áp dụng cho nhóm này.

{\itshape 4. Nêu hai chỗ của khoảng trống cấu trúc.\par}

Đáp án: thước đo rủi ro mô hình hiện có đo chênh lệch giữa các dự báo rủi ro, chưa đo hậu quả của chọn sai lên kết quả danh mục; nghiên cứu danh mục xếp hạng theo bình quân toàn mẫu, chưa xét thứ hạng theo chế độ.

{\itshape 5. Vì sao Việt Nam được chọn, và có phải lý do của khoảng trống không?\par}

Đáp án: Việt Nam là bối cảnh kiểm chứng, không phải lý do của khoảng trống. Thanh khoản hạn chế và nhiều cú sốc có thể làm chênh lệch giữa các mô hình lớn; đổi lại, mẫu có ít giai đoạn căng thẳng độc lập.

\chapter{Mục 4.1 của đề cương (tiếp): Mục tiêu và câu hỏi nghiên cứu}

Đề tài có một mục tiêu tổng quát, bốn mục tiêu cụ thể và ba câu hỏi nghiên cứu, ứng với bốn giả thuyết H1 đến H4. Đây là xương sống của đề cương: gần như mọi câu hỏi của giám khảo về phương pháp, kết quả hay đóng góp đều quay về đây.

\section{Mục tiêu chung: ba động từ}

Nguyên văn: đo lường và giải thích rủi ro lựa chọn mô hình trong quản trị danh mục đầu tư trên thị trường chứng khoán Việt Nam qua các chế độ thị trường, rồi đánh giá các cơ chế giảm rủi ro đó.

Ba động từ là ba bước: đo lường (có một con số cho rủi ro), giải thích (con số đổi theo chế độ thế nào, biến nào đi cùng), đánh giá (cách giảm rủi ro nào có hiệu quả). Giống bác sĩ đo nhiệt độ, tìm triệu chứng đi kèm, rồi thử thuốc hạ sốt.

Hai chi tiết cần để ý. “Qua các chế độ thị trường” nghĩa là đề tài so sánh giữa các chế độ, không đo một con số chung. “Giải thích” chỉ có nghĩa là chỉ ra các biến có liên hệ thống kê, vì đề cương gọi kết quả là liên hệ có điều kiện, không phải nhân quả.

\section{Bốn mục tiêu cụ thể}

Bốn mục tiêu cụ thể, nguyên văn:

{\itshape (i) Xây dựng thước đo rủi ro lựa chọn mô hình cho danh mục và xác định các chế độ thị trường ở Việt Nam.

(ii) So sánh rủi ro lựa chọn mô hình giữa chế độ căng thẳng và chế độ bình thường.

(iii) Xác định các biến kênh có liên hệ thống kê với rủi ro lựa chọn mô hình, gồm biến động, thanh khoản, tương quan cực trị và bầy đàn.

(iv) Đánh giá các cơ chế giảm thiểu và rút ra hàm ý quản trị.\par}

\subsection{Mục tiêu (i): dựng thước đo và xác định chế độ}

Không có thước đo thì không so sánh được giữa các chế độ. Mục 4.3.6 dự kiến ba thước đo: mức hối tiếc (kết quả chính), độ phân tán mô hình và kích thước tập tin cậy mô hình (kết quả phụ).

Việc thứ hai là xác định chế độ bằng mô hình chuyển chế độ Markov: ước lượng xác suất thị trường đang ở chế độ căng thẳng, tức chế độ có phương sai điều kiện cao nhất và lợi suất kỳ vọng thấp hơn chế độ bình thường. Muốn nói “rủi ro ở vùng bão cao hơn”, bạn cần cả thước đo lẫn bản đồ phân vùng.

\subsection{Mục tiêu (ii): so sánh chế độ căng thẳng với chế độ bình thường}

Mục tiêu này kiểm định H1 và trả lời câu hỏi 1. Đề tài không giả định đáp án là “có”: khả năng đối lập là rủi ro không đổi theo chế độ, hoặc 1/N thắng ổn định. Nếu dữ liệu nói vậy, đề tài báo cáo vậy.

\subsection{Mục tiêu (iii): xác định biến kênh}

Mục tiêu này tìm các biến có liên hệ thống kê với rủi ro lựa chọn mô hình: biến động, thanh khoản, tương quan cực trị và bầy đàn.

Một chỗ tinh tế có thể bị hỏi: trong câu hỏi 2, biến động thực hiện là thứ được kiểm soát, không phải biến kênh cần kiểm định. Ba biến kênh được kiểm định là thanh khoản, tương quan cực trị và bầy đàn (thăm dò).

Muốn biết món ăn ngon hơn nhờ gia vị nào, bạn giữ nhiệt độ nấu cố định rồi mới so từng gia vị. “Kiểm soát biến động thực hiện” có tinh thần đó.

Lý do: theo Forbes và Rigobon (2002), hệ số tương quan phụ thuộc vào biến động. Khả năng đối lập của H2 chính là tương quan tăng chỉ vì biến động tăng; không kiểm soát biến động thì không tách được hai hiệu ứng.

\subsection{Mục tiêu (iv): đánh giá cơ chế giảm thiểu và rút ra hàm ý quản trị}

Mục tiêu cuối ứng với H3. Mục 4.3.8 nêu ba cơ chế (kết hợp mô hình, danh mục vững chắc, quy tắc chuyển chế độ), đánh giá bằng lợi suất tương đương chắc chắn, mức hối tiếc, độ sụt giảm tối đa, giá trị thiệt hại kỳ vọng mức 95\%, vòng quay và hiệu quả sau chi phí giao dịch.

Hàm ý quản trị có giới hạn: thủ tục giám sát ba bước là đề xuất suy ra từ H3, chưa được kiểm nghiệm với người dùng, và chỉ có nếu dữ liệu ủng hộ.

\section{Ba câu hỏi nghiên cứu}

Ba câu hỏi, nguyên văn:

\begin{itemize}
\item
  Câu hỏi 1: Rủi ro lựa chọn mô hình thay đổi thế nào giữa chế độ bình thường và chế độ căng thẳng?
\item
  Câu hỏi 2: Sau khi kiểm soát biến động thực hiện, các biến kênh (thanh khoản, tương quan cực trị và, ở mức thăm dò, hành vi bầy đàn) có liên hệ thống kê với rủi ro lựa chọn mô hình không, và chúng giải thích được bao nhiêu phần chênh lệch giữa hai chế độ?
\item
  Câu hỏi 3: Cơ chế nào giảm được rủi ro lựa chọn mô hình sau chi phí giao dịch?
\end{itemize}

\subsection{Câu hỏi 1: so sánh hai chế độ}

Câu hỏi 1 so sánh hai trạng thái. Giả thuyết H1 là rủi ro ở chế độ căng thẳng cao hơn; khả năng đối lập là không đổi theo chế độ. Kết quả chính đo bằng mức hối tiếc R, kết quả phụ bằng độ phân tán D và kích thước tập tin cậy mô hình.

Mức hối tiếc là chênh lệch lợi suất tương đương chắc chắn giữa mô hình tốt nhất sau thực tế và mô hình chọn trước. Mô hình chọn trước là mô hình có lợi suất tương đương chắc chắn ngoài mẫu cao nhất trong 36 tháng liền trước đầu cửa sổ đánh giá (mục 4.3.6).

\subsection{Câu hỏi 2: biến kênh và phần chênh lệch giải thích được}

Câu hỏi 2 có hai nửa: các biến kênh có liên hệ thống kê với rủi ro lựa chọn mô hình không, sau khi kiểm soát biến động; và chúng giải thích được bao nhiêu phần chênh lệch giữa hai chế độ.

Chênh lệch giữa hai chế độ là hệ số \texorpdfstring{$\beta$}{beta} trong phương trình hồi quy của mục 4.3.7. Hình dung \texorpdfstring{$\beta$}{beta} là khoảng cách giữa hai bậc thang, trời quang và mùa bão; nửa thứ hai hỏi khoảng cách đó ngắn lại bao nhiêu khi đưa các biến kênh vào phương trình.

Đề tài đo phần đó bằng hiệu giữa \texorpdfstring{$\beta$}{beta} khi phương trình không có và khi có biến kênh.

Câu hỏi 2 ứng với H2 và H4. H2 dự kiến hệ số của độ phi thanh khoản dương và hệ số của tương quan đuôi âm; H4 là giả thuyết thăm dò về độ phân tán chéo của lợi suất.

\subsection{Câu hỏi 3: cơ chế giảm thiểu}

Câu hỏi 3 ứng với H3: ở chế độ căng thẳng, sau chi phí giao dịch, kết hợp mô hình và danh mục vững chắc cho mức hối tiếc và độ sụt giảm tối đa thấp hơn, lợi suất tương đương chắc chắn cao hơn so với chọn một mô hình. Cụm “sau chi phí giao dịch” quan trọng vì một cách giảm rủi ro có thể mất lợi thế khi tính phí.

\section{Câu hỏi và giả thuyết}

H1 trả lời câu hỏi 1; H2 và H4 trả lời câu hỏi 2; H3 trả lời câu hỏi 3 (Bảng 2 của đề cương).

{\small\begin{longtable}[]{@{}
  >{\raggedright\arraybackslash}p{(\columnwidth - 6\tabcolsep) * \real{0.0797}}
  >{\raggedright\arraybackslash}p{(\columnwidth - 6\tabcolsep) * \real{0.0797}}
  >{\raggedright\arraybackslash}p{(\columnwidth - 6\tabcolsep) * \real{0.4717}}
  >{\raggedright\arraybackslash}p{(\columnwidth - 6\tabcolsep) * \real{0.3690}}@{}}
\toprule\noalign{}
\begin{minipage}[b]{\linewidth}\raggedright
Câu hỏi
\end{minipage} & \begin{minipage}[b]{\linewidth}\raggedright
Giả thuyết
\end{minipage} & \begin{minipage}[b]{\linewidth}\raggedright
Nội dung giả thuyết
\end{minipage} & \begin{minipage}[b]{\linewidth}\raggedright
Thước đo chính
\end{minipage} \\
\midrule\noalign{}
\endhead
\bottomrule\noalign{}
\endlastfoot
Câu hỏi 1 & H1 & Rủi ro lựa chọn mô hình ở chế độ căng thẳng cao hơn ở chế độ bình thường & Mức hối tiếc R (chính); độ phân tán D, kích thước tập tin cậy mô hình (phụ) \\
Câu hỏi 2 & H2 & Sau khi kiểm soát biến động thực hiện, hệ số độ phi thanh khoản dương và hệ số tương quan đuôi âm & Dấu và khoảng tin cậy của hai hệ số; phần chênh lệch giữa hai chế độ do biến kênh giải thích \\
Câu hỏi 2 & H4 (thăm dò) & Sau khi kiểm soát biến động thực hiện, độ phân tán chéo cao đi cùng rủi ro lựa chọn mô hình cao & Độ phân tán chéo của lợi suất; thước đo bầy đàn \\
Câu hỏi 3 & H3 & Ở chế độ căng thẳng, sau chi phí giao dịch, kết hợp mô hình và danh mục vững chắc cho mức hối tiếc, độ sụt giảm tối đa thấp hơn và lợi suất tương đương chắc chắn cao hơn so với chọn một mô hình & Lợi suất tương đương chắc chắn; mức hối tiếc; độ sụt giảm tối đa; thiệt hại kỳ vọng 95\%; vòng quay \\
\end{longtable}}

Về H2: hệ số của độ phi thanh khoản được kỳ vọng dương vì thanh khoản cạn làm các mô hình có vòng quay khác nhau cách xa nhau; hệ số của tương quan đuôi được kỳ vọng âm vì tương quan cao làm các danh mục chỉ mua hội tụ. Hai dấu ngược nhau nên dễ nhầm. Cơ sở: Amihud (2002), Longin và Solnik (2001), Forbes và Rigobon (2002).

Độ phi thanh khoản Amihud (2002) càng cao thì thanh khoản càng cạn. Tương quan đuôi là mức đi cùng nhau của lợi suất ở những ngày xấu nhất.

Tương quan đuôi cao giống cảnh mưa to, mọi người cùng chạy vào một mái hiên: khi mọi cổ phiếu cùng rơi, các danh mục chỉ mua trông giống nhau, khoảng cách giữa các mô hình thu hẹp. Vì vậy H2 dự kiến dấu âm.

\section{Hình 1: khung phân tích và mũi tên không có giả thuyết}

Ánh xạ câu hỏi và giả thuyết được vẽ thành Hình 1, khung phân tích ở mục 4.3.2: chế độ tác động lên rủi ro lựa chọn mô hình (H1); biến kênh tác động lên quan hệ đó (H2, H4); cơ chế giảm thiểu tác động lên quan hệ giữa rủi ro và kết quả danh mục (H3).

Chi tiết thường bị hỏi: mũi tên nét đứt từ rủi ro đến kết quả danh mục không có giả thuyết riêng, vì mức hối tiếc được định nghĩa bằng chính chênh lệch kết quả danh mục.

\section{Vì sao có ba câu hỏi mà bốn giả thuyết}

Đề cương không viết mỗi câu hỏi một giả thuyết. Cách hiểu của bản đọc: câu hỏi 2 chứa hai nhóm biến kênh có trạng thái bằng chứng khác nhau. Thanh khoản và tương quan đuôi vào H2; bầy đàn, có bằng chứng trái chiều, vào H4 với nhãn thăm dò.

Cơ sở của H4 là Vo và Phan (2017), Chang và cộng sự (2000) và Christie và Huang (1995). Khả năng đối lập: bầy đàn không xuất hiện ở tần suất dữ liệu của đề tài, như trên dữ liệu Mỹ.

\begin{tcolorbox}[breakable]
\textbf{Ghi nhớ khi đi thi}

• Mục tiêu tổng quát: đo lường, giải thích rủi ro lựa chọn mô hình qua các chế độ thị trường ở Việt Nam, rồi đánh giá cơ chế giảm rủi ro.

• Bốn mục tiêu cụ thể: (i) thước đo và chế độ; (ii) so sánh hai chế độ; (iii) biến kênh; (iv) cơ chế giảm thiểu và hàm ý quản trị.

• Câu hỏi 1 -- H1; câu hỏi 2 -- H2 và H4; câu hỏi 3 -- H3.

• Câu hỏi 2 kiểm soát biến động thực hiện, hỏi liên hệ của biến kênh và phần chênh lệch giữa hai chế độ mà chúng giải thích.

• H4 chỉ thăm dò; liên hệ chỉ là liên hệ có điều kiện, không phải nhân quả.

• Mũi tên từ rủi ro đến kết quả danh mục ở Hình 1 không có giả thuyết riêng.
\end{tcolorbox}

\textbf{Tự kiểm tra}

{\itshape 1. Nêu mục tiêu chung của đề tài.\par}

Đáp án: đo lường và giải thích rủi ro lựa chọn mô hình trong quản trị danh mục trên thị trường chứng khoán Việt Nam qua các chế độ thị trường, rồi đánh giá các cơ chế giảm rủi ro đó.

{\itshape 2. Nêu bốn mục tiêu cụ thể.\par}

Đáp án: (i) xây dựng thước đo và xác định chế độ thị trường; (ii) so sánh rủi ro giữa chế độ căng thẳng và chế độ bình thường; (iii) xác định biến kênh: biến động, thanh khoản, tương quan cực trị, bầy đàn; (iv) đánh giá cơ chế giảm thiểu và rút ra hàm ý quản trị.

{\itshape 3. Câu hỏi nào ứng với giả thuyết nào?\par}

Đáp án: câu hỏi 1 -- H1; câu hỏi 2 -- H2 và H4; câu hỏi 3 -- H3.

{\itshape 4. Câu hỏi 2 có hai nửa nào, và biến động đóng vai trò gì?\par}

Đáp án: biến kênh có liên hệ thống kê với rủi ro lựa chọn mô hình không; và chúng giải thích bao nhiêu phần chênh lệch giữa hai chế độ (hệ số \texorpdfstring{$\beta$}{beta}). Biến động thực hiện là biến kiểm soát.

{\itshape 5. Mũi tên từ rủi ro đến kết quả danh mục trong Hình 1 có giả thuyết riêng không?\par}

Đáp án: không; nó suy ra từ định nghĩa mức hối tiếc.

\chapter{Mục 4.3 và 4.4 của đề cương: Phương pháp nhìn tổng thể, đối tượng và phạm vi}

Đề tài dùng phương pháp định lượng trên dữ liệu thị trường công khai, theo năm bước: thu thập và xử lý dữ liệu; xác định chế độ thị trường; ước lượng tập mô hình cạnh tranh; đo rủi ro lựa chọn mô hình và kiểm định giả thuyết; đánh giá cơ chế giảm thiểu. Phạm vi là danh mục cổ phiếu niêm yết trên HOSE, cấu hình cơ sở là chỉ mua.

Phần này là bản đồ; chi tiết ở các mục 4.3.3 đến 4.3.10 và các phần sau của bài đọc.

\section{Nguyên văn đoạn phương pháp}

Tóm lại bằng một câu: đề tài xác định chế độ bằng mô hình chuyển chế độ Markov, ước lượng bảy mô hình danh mục trên cửa sổ trượt, đo rủi ro lựa chọn bằng mức hối tiếc, độ phân tán mô hình và kích thước tập tin cậy mô hình, rồi kiểm định bằng hồi quy có sai số chuẩn bền vững, bootstrap khối và hiệu chỉnh so sánh bội.

\section{Định lượng, dữ liệu công khai, và không có khảo sát}

Định lượng nghĩa là câu trả lời đến từ số liệu và mô hình thống kê. Dữ liệu gồm giá đóng cửa điều chỉnh theo ngày và khối lượng giao dịch của cổ phiếu trên HOSE, kể cả mã đã hủy niêm yết; chỉ số VN-Index; lợi suất trái phiếu chính phủ kỳ hạn một năm làm lãi suất phi rủi ro. Nguồn là HOSE và một nhà cung cấp dữ liệu thương mại.

Không có khảo sát hay phương pháp định tính, nên không có thang đo cần dịch thuật, và đề tài không phỏng vấn nhà quản lý quỹ về cách họ chọn mô hình. Dữ liệu không chứa thông tin cá nhân.

\section{Năm bước của quy trình}

Hình 3 của đề cương vẽ năm bước này thành sơ đồ quy trình.

{\small\begin{longtable}[]{@{}
  >{\raggedright\arraybackslash}p{(\columnwidth - 4\tabcolsep) * \real{0.1254}}
  >{\raggedright\arraybackslash}p{(\columnwidth - 4\tabcolsep) * \real{0.7282}}
  >{\raggedright\arraybackslash}p{(\columnwidth - 4\tabcolsep) * \real{0.1464}}@{}}
\toprule\noalign{}
\begin{minipage}[b]{\linewidth}\raggedright
Bước
\end{minipage} & \begin{minipage}[b]{\linewidth}\raggedright
Việc làm
\end{minipage} & \begin{minipage}[b]{\linewidth}\raggedright
Mục của đề cương
\end{minipage} \\
\midrule\noalign{}
\endhead
\bottomrule\noalign{}
\endlastfoot
1 & Thu thập và xử lý dữ liệu thị trường & 4.3.3 \\
2 & Xác định chế độ thị trường & 4.3.4 \\
3 & Ước lượng tập mô hình danh mục cạnh tranh trên cửa sổ trượt & 4.3.5 \\
4 & Đo rủi ro lựa chọn mô hình và kiểm định bốn giả thuyết & 4.3.6 và 4.3.7 \\
5 & Đánh giá cơ chế giảm thiểu & 4.3.8 \\
\end{longtable}}

Mục 4.3.10 không phải bước mới; nó nêu độ tin cậy, giới hạn và khả năng thực hiện.

\subsection{Bước 1: thu thập và xử lý dữ liệu}

Điểm đáng nhớ là cách chọn rổ: tại mỗi ngày tái cân bằng, rổ gồm 30 cổ phiếu có giá trị giao dịch bình quân ba tháng và vốn hóa tự do chuyển nhượng lớn nhất, xác định chỉ bằng dữ liệu có tại ngày đó, để tránh thiên lệch sống sót và thiên lệch nhìn trước.

Nếu chỉ chọn cổ phiếu còn tồn tại đến hôm nay, bạn bỏ qua các công ty đã biến mất và kết quả trông đẹp hơn thực tế: đó là thiên lệch sống sót. Nếu chọn rổ bằng thông tin của tương lai, bạn “nhìn trước”.

Tái cân bằng là điều chỉnh lại danh mục theo kế hoạch, ở đây là hằng tháng. Chế độ ước lượng trên dữ liệu ngày, hiệp phương sai trên dữ liệu tuần; kiểm tra độ bền dùng dữ liệu tháng và tái cân bằng hằng quý.

Giai đoạn mẫu dự kiến 2007--2026, gồm khủng hoảng tài chính toàn cầu, đại dịch COVID-19 và căng thẳng trái phiếu doanh nghiệp từ quý IV/2022. Rổ VN30 và VN100 chỉ dùng để đối chiếu và kiểm tra độ bền.

\subsection{Bước 2: xác định chế độ thị trường}

Mô hình chuyển chế độ Markov (Hamilton, 1989) giả định thị trường ở một trong các trạng thái ẩn và chuyển giữa chúng theo xác suất. Ta không thấy “thời tiết”, chỉ thấy lợi suất hằng ngày; mô hình đoán thời tiết nào khớp với chuỗi lợi suất và cho xác suất mỗi ngày là ngày bão.

Ba điểm kỹ thuật:

\begin{itemize}
\item
  Tham số được ước lượng lại theo cửa sổ mở rộng tại mỗi ngày tái cân bằng, chỉ bằng dữ liệu đến ngày đó; xác suất lọc tại ngày t được tính từ các tham số này.
\item
  Số chế độ (hai hoặc ba) do tiêu chí thông tin quyết định trên giai đoạn ước lượng ban đầu và cố định khi đăng ký trước.
\item
  Chế độ căng thẳng là chế độ có phương sai điều kiện cao nhất và lợi suất kỳ vọng thấp hơn chế độ bình thường.
\end{itemize}

Kiểm định điểm gãy Bai-Perron trên toàn mẫu chỉ để mô tả và đối chiếu với các cú sốc lịch sử, không làm biến chế độ trong kiểm định giả thuyết. Định nghĩa thay thế theo ngưỡng sụt giảm và biến động thực hiện nằm trong kiểm tra độ bền.

\subsection{Bước 3: bảy mô hình trên cửa sổ trượt}

Bảy cách chia tiền cùng được giao một khoản vốn ảo và xem cùng dữ liệu (Bảng 3 của đề cương).

{\small\begin{longtable}[]{@{}
  >{\raggedright\arraybackslash}p{(\columnwidth - 4\tabcolsep) * \real{0.0991}}
  >{\raggedright\arraybackslash}p{(\columnwidth - 4\tabcolsep) * \real{0.5289}}
  >{\raggedright\arraybackslash}p{(\columnwidth - 4\tabcolsep) * \real{0.3720}}@{}}
\toprule\noalign{}
\begin{minipage}[b]{\linewidth}\raggedright
Mã
\end{minipage} & \begin{minipage}[b]{\linewidth}\raggedright
Mô hình
\end{minipage} & \begin{minipage}[b]{\linewidth}\raggedright
Nguồn
\end{minipage} \\
\midrule\noalign{}
\endhead
\bottomrule\noalign{}
\endlastfoot
M1 & Phân bổ đều 1/N & DeMiguel và cộng sự (2009) \\
M2 & Phương sai nhỏ nhất với hiệp phương sai co rút & Ledoit và Wolf (2004) \\
M3 & Trung bình-phương sai với ước lượng mẫu & Markowitz (1952) \\
M4 & Black-Litterman với quan điểm từ mô hình dự báo & Black và Litterman (1992) \\
M5 & Black-Litterman có chuyển chế độ & Mở rộng từ Nguyen và Nguyen (2026a) \\
M6 & Danh mục vững chắc với nhiều niềm tin tiên nghiệm & Garlappi và cộng sự (2007) \\
M7 & Giảm tỷ trọng rủi ro theo xác suất chế độ căng thẳng & Kritzman và cộng sự (2012) \\
\end{longtable}}

Cửa sổ trượt: mỗi mô hình được ước lượng trên một đoạn quá khứ, đánh giá trên đoạn kế tiếp, rồi cả hai dịch về phía trước, như người thợ may đo vải trên đoạn đã có, cắt đoạn tiếp theo, rồi dịch thước. Giá trị dự kiến: ước lượng 104 tuần, đánh giá 6 tháng không chồng lấn; cửa sổ trượt hằng tháng chỉ là phân tích phụ.

Các mô hình được đối xử như nhau: siêu tham số của M2, M4, M5, M6 và M7 chọn từ một lưới cho trước bằng cùng quy tắc và cùng ngân sách tinh chỉnh. M5 do tác giả phát triển nên kết quả được báo cáo cả khi có lẫn khi không có M5.

\subsection{Bước 4: đo rủi ro và kiểm định}

Mỗi cửa sổ đánh giá, mỗi mô hình có một điểm lợi suất tương đương chắc chắn sau chi phí giao dịch. Từ bảy điểm đó, đề tài tính ba thước đo:

\begin{itemize}
\item
  Mức hối tiếc R: chênh lệch giữa điểm của mô hình tốt nhất sau thực tế và điểm của mô hình chọn trước; thước đo rủi ro, kết quả chính của H1.
\item
  Độ phân tán mô hình D: độ lệch chuẩn của bảy điểm; chỉ báo bất định, kết quả phụ.
\item
  Kích thước tập tin cậy mô hình (Hansen và cộng sự, 2011) ở mức ý nghĩa 10\%: số ứng viên chưa thể loại sau vòng sơ loại thống kê. Danh sách dài nghĩa là chưa phân biệt nổi các mô hình.
\end{itemize}

Tập tin cậy lớn hơn ở chế độ căng thẳng có thể chỉ do dữ liệu kém thông tin hơn, nên được báo cáo như thước đo mô tả, kèm số quan sát và công suất.

\section{Hồi quy với sai số chuẩn bền}

Hồi quy là đường xấp xỉ mối liên hệ giữa các đại lượng. Đề tài hồi quy thước đo rủi ro (D hoặc R) theo biến chế độ và các biến kênh; hệ số \texorpdfstring{$\beta$}{beta} là chênh lệch của thước đo giữa chế độ căng thẳng và chế độ bình thường.

Sai số chuẩn bền vững trước phương sai thay đổi và tự tương quan vẫn đáng tin khi mức dao động đổi theo thời gian và các quan sát dính nhau, như phần lớn dữ liệu tài chính. Giống đo cây bằng một cây thước có thể co giãn và báo kèm sai số có tính đến độ co giãn đó; nhờ vậy không báo sai số quá nhỏ và kết luận quá lạc quan.

\section{Bootstrap khối}

Bootstrap là bốc lại nhiều lần từ chính túi dữ liệu để xem kết quả dao động bao nhiêu. Bootstrap khối bốc từng nắm ngày liền nhau để giữ tính liên tục của chuỗi thời gian.

Hai biến thể: bootstrap lấy mẫu lặp theo khối liên tiếp (Ledoit và Wolf, 2008) để so sánh lợi suất tương đương chắc chắn và Sharpe; bootstrap khối dừng (Politis và Romano, 1994) cho hệ số hồi quy.

Bootstrap còn dùng để dựng phân phối rỗng: kết quả sẽ trông thế nào nếu thực ra không có gì xảy ra. Đề cương dự kiến 1.000 lần tái lấy mẫu chuỗi lợi suất của bảy mô hình, sau khi loại chênh lệch trung bình giữa các mô hình trong từng chế độ, để tách rủi ro lựa chọn thật khỏi nhiễu ước lượng.

\section{Hiệu chỉnh so sánh bội}

Khi kiểm định nhiều giả thuyết cùng lúc, khả năng gặp kết quả “có ý nghĩa” chỉ do may tăng lên, như tung đồng xu nhiều lần thì dễ gặp một chuỗi mặt sấp. Đề cương dùng hiệu chỉnh Holm (1979) ở mức 10\% cho bốn giả thuyết.

Đề tài còn dùng tỷ số Sharpe hiệu chỉnh của Bailey và López de Prado (2014) cho thiên lệch chọn lọc: thử nhiều mô hình rồi chỉ khoe mô hình tốt nhất thì kết quả phải bị hạ xuống.

\section{Đăng ký trước}

Đăng ký trước là bỏ dự đoán và cách kiểm tra vào phong bì có dấu thời gian trước khi xem dữ liệu mới. Kế hoạch phân tích được đăng ký trên OSF Registries vào quý 4 năm thứ nhất, ở chế độ khóa có dấu thời gian.

Giới hạn nói thẳng: tác giả đã dùng dữ liệu 2019--2025 (hai giai đoạn căng thẳng) trong Nguyen và Nguyen (2026a), nên tệp đăng ký khai rõ những kết quả đã xem. Đăng ký trước giảm thiên lệch, không loại được nó.

Tệp đăng ký cố định độ dài cửa sổ, siêu tham số, số chế độ và ngưỡng xác suất lọc, mức ý nghĩa, nhóm kết quả chính, quy tắc hiệu chỉnh so sánh bội và quy tắc gọi kết quả là không kết luận. Dữ liệu sau ngày đăng ký là mẫu kiểm tra riêng; mọi thay đổi ghi là lệch khỏi kế hoạch.

\section{Đối tượng và phạm vi}

Mục 4.4 của đề cương nêu đối tượng và phạm vi:

Đối tượng nghiên cứu là rủi ro lựa chọn mô hình trong quản trị danh mục cổ phiếu, xét qua các chế độ thị trường, cùng các biến kênh và cơ chế giảm thiểu. Phạm vi không gian là cổ phiếu niêm yết trên HOSE (Sở Giao dịch Chứng khoán Thành phố Hồ Chí Minh); phạm vi thời gian là 2007--2026, tùy dữ liệu sẵn có.

Danh mục chỉ mua là cấu hình cơ sở; các ràng buộc khác nằm trong kiểm tra độ bền. Chỉ mua nghĩa là chỉ nắm giữ cổ phiếu, không bán khống thứ mình chưa sở hữu, nên trọng số không âm.

Đề tài không xét phái sinh, trái phiếu doanh nghiệp và tài sản ngoài cổ phiếu. Trái phiếu doanh nghiệp chỉ xuất hiện như một cú sốc lịch sử.

\section{Vì sao tác giả chọn cách làm này}

Đề cương không có mục riêng giải thích “vì sao chọn cách này”; lý do nằm rải rác ở mục 4.3.

So sánh nhiều mô hình cùng lúc vì câu hỏi là chọn giữa các mô hình (khoảng trống thứ hai). Xác định chế độ chỉ bằng dữ liệu đến từng ngày để tránh thiên lệch nhìn trước. Dùng phân phối rỗng vì kỳ vọng của mức hối tiếc vẫn dương khi các mô hình ngang chất lượng: giá trị lớn nhất của các đại lượng có nhiễu tăng theo phương sai của nhiễu. Mức hối tiếc dương tự nó chưa chứng minh có rủi ro lựa chọn thật.

Bảy học sinh học ngang nhau làm cùng một bài kiểm tra. Điểm cao nhất vẫn cao hơn điểm của người bạn chọn trước, chỉ vì may rủi. Phân phối rỗng đo khoảng cách do may rủi đó để biết phần thật còn lại bao nhiêu.

Vì vậy đề cương chia D và R cho sai số chuẩn bootstrap của chênh lệch giữa các mô hình trong cùng cửa sổ. H1 được ủng hộ khi thỏa đồng thời ba điều kiện: \texorpdfstring{$\beta$}{beta} của R dương và vượt phân vị 95\% của phân phối rỗng; khoảng tin cậy bootstrap khối của \texorpdfstring{$\beta$}{beta} loại trừ 0; \texorpdfstring{$\beta$}{beta} của D cùng dấu.

\section{Giới hạn của phương pháp, theo đề cương}

Các giới hạn ở mục 4.3.10, cần thuộc:

\begin{itemize}
\item
  Ít giai đoạn căng thẳng độc lập nên công suất kiểm định thấp; dưới 10 cửa sổ độc lập ở chế độ căng thẳng thì kết quả là không kết luận.
\item
  Kết luận giới hạn trong phạm vi Việt Nam.
\item
  Nhận diện chế độ tự nó là một dạng rủi ro mô hình, nên có định nghĩa thay thế trong kiểm tra độ bền.
\item
  Khả năng thực thi trên HOSE phụ thuộc thanh khoản, nên đề tài dùng rổ thanh khoản cao và kiểm tra độ nhạy theo chi phí.
\end{itemize}

Nếu giả thuyết bị bác bỏ, đề tài vẫn trả lời được câu hỏi: rủi ro không đổi theo chế độ, hoặc 1/N đủ dùng ở Việt Nam.

\begin{tcolorbox}[breakable]
\textbf{Ghi nhớ khi đi thi}

• Định lượng trên dữ liệu công khai; không khảo sát.

• Năm bước: dữ liệu, xác định chế độ, ước lượng tập mô hình, đo rủi ro và kiểm định, đánh giá cơ chế giảm thiểu.

• Chế độ: mô hình chuyển chế độ Markov; bảy mô hình M1 đến M7 trên cửa sổ trượt; ba thước đo: mức hối tiếc, độ phân tán, tập tin cậy mô hình.

• Kiểm định: hồi quy với sai số chuẩn bền vững, bootstrap khối, hiệu chỉnh Holm, phân phối rỗng 1.000 lần.

• Phạm vi: cổ phiếu trên HOSE, chỉ mua; không phái sinh, trái phiếu doanh nghiệp hay tài sản khác.

• Giới hạn: ít giai đoạn căng thẳng; dưới 10 cửa sổ độc lập thì không kết luận; kết luận chỉ cho Việt Nam.
\end{tcolorbox}

\textbf{Tự kiểm tra}

{\itshape 1. Phương pháp của đề tài là định lượng hay định tính, và có khảo sát không?\par}

Đáp án: định lượng trên dữ liệu công khai; không khảo sát hay định tính.

{\itshape 2. Liệt kê năm bước của quy trình.\par}

Đáp án: dữ liệu; xác định chế độ; ước lượng tập mô hình trên cửa sổ trượt; đo rủi ro và kiểm định; đánh giá cơ chế giảm thiểu.

{\itshape 3. Nêu ba thước đo rủi ro lựa chọn mô hình.\par}

Đáp án: mức hối tiếc (kết quả chính của H1), độ phân tán mô hình và kích thước tập tin cậy mô hình (kết quả phụ).

{\itshape 4. Phạm vi của đề tài gồm và không gồm những gì?\par}

Đáp án: gồm danh mục cổ phiếu trên HOSE, chỉ mua; không gồm phái sinh, trái phiếu doanh nghiệp và tài sản ngoài cổ phiếu.

{\itshape 5. Vì sao đề tài cần phân phối rỗng?\par}

Đáp án: vì kỳ vọng của mức hối tiếc vẫn dương khi các mô hình ngang chất lượng; phân phối rỗng 1.000 lần tái lấy mẫu tách rủi ro thật khỏi nhiễu ước lượng.

\chapter{Mục 4.4 của đề cương (tiếp): Đóng góp dự kiến}

Đề cương xác định đóng góp bằng cách so với các công trình gần nhất và nêu bốn điểm khác biệt: một mô hình so với bảy mô hình; so sánh phương pháp so với đo bất định theo chế độ; chế độ của biến động so với chất lượng quyết định; và khác biệt với bài trước của chính tác giả. Mọi đóng góp đều là dự kiến.

Giám khảo có thể hỏi “đóng góp của bạn là gì, khác ai”. Bạn cần biết các công trình được so, khác biệt với bài trước của tác giả, và những gì đề cương không tuyên bố.

\section{Đóng góp được xác định bằng cách so sánh}

Đóng góp luôn đi kèm một công trình cụ thể để đối chiếu, giống khi xin vào đội bóng bạn không chỉ nói “tôi đá tốt” mà nói “đội đã có những tiền đạo giỏi việc này, còn tôi làm việc khác”.

Đề cương viết “dựa trên phần tóm tắt của các công trình gần nhất”: tác giả mới đọc tóm tắt và thông tin thư mục, việc đọc toàn văn nằm trong kế hoạch năm thứ nhất.

\section{Điểm khác biệt thứ nhất: một mô hình so với bảy mô hình}

Ta và Vuong (2020) áp dụng một mô hình Black-Litterman cho Việt Nam, với quan điểm dự báo từ mô hình tự hồi quy trung bình trượt tích hợp. Đề tài đặt Black-Litterman cùng sáu mô hình khác vào một tập cạnh tranh (Bảng 3, M1 đến M7).

Một bên theo dõi một vận động viên chạy thế nào; bên kia tổ chức giải chạy bảy người để xem huấn luyện viên chọn nhầm thì thiệt bao nhiêu.

Khác biệt này không nói Black-Litterman tốt hay kém; nó nói câu hỏi khác nhau. Đây cũng là khoảng trống thứ hai: các công trình về danh mục ở Việt Nam đánh giá từng mô hình, chưa đặt vào một tập cạnh tranh.

\section{Điểm khác biệt thứ hai: so sánh phương pháp so với đo bất định theo chế độ}

Goutte và cộng sự (2025) so sánh năm phương pháp ra quyết định đa tiêu chí với bốn cách gán trọng số trên rổ VN100, giai đoạn 01/2015--11/2023. Đề tài đo mức bất định khi chọn mô hình theo từng chế độ thị trường.

Một bên đặt các phương pháp cạnh nhau để so; bên kia hỏi kết quả so sánh dao động ra sao khi thời tiết đổi, và buộc phải chọn một phương án thì trả giá bao nhiêu. Đừng nói thêm về kết quả của Goutte và cộng sự ngoài những gì đề cương nêu.

\section{Điểm khác biệt thứ ba: biến động của chỉ số so với chất lượng quyết định}

Hoang và Luu (2024), một bản thảo công bố trực tuyến, phân tích VN-Index giai đoạn 2012--2023 và thấy mô hình một chế độ dự báo biến động tốt hơn ở kỳ hạn ngắn, mô hình chuyển chế độ tốt hơn ở kỳ hạn trung và dài.

Đề tài nối chế độ với chất lượng quyết định phân bổ: không dùng chế độ để dự báo biến động, mà hỏi quyết định chia tiền tốt hay xấu tùy chế độ ra sao. Đây là khoảng trống thứ ba.

Một bên dự báo thời tiết ngày mai; bên kia hỏi khi trời xấu, quyết định chọn lộ trình có sai nhiều hơn không.

\section{Bảng tóm tắt ba điểm khác biệt}

{\small\begin{longtable}[]{@{}
  >{\raggedright\arraybackslash}p{(\columnwidth - 4\tabcolsep) * \real{0.1577}}
  >{\raggedright\arraybackslash}p{(\columnwidth - 4\tabcolsep) * \real{0.3711}}
  >{\raggedright\arraybackslash}p{(\columnwidth - 4\tabcolsep) * \real{0.4712}}@{}}
\toprule\noalign{}
\begin{minipage}[b]{\linewidth}\raggedright
Công trình
\end{minipage} & \begin{minipage}[b]{\linewidth}\raggedright
Công trình đó làm gì (theo tóm tắt)
\end{minipage} & \begin{minipage}[b]{\linewidth}\raggedright
Đề tài này làm gì khác
\end{minipage} \\
\midrule\noalign{}
\endhead
\bottomrule\noalign{}
\endlastfoot
Ta và Vuong (2020) & Áp dụng một mô hình Black-Litterman & Đặt Black-Litterman cùng sáu mô hình khác vào một tập cạnh tranh (Bảng 3) \\
Goutte và cộng sự (2025) & So sánh các phương pháp ra quyết định đa tiêu chí trên VN100 & Đo mức bất định khi chọn mô hình theo từng chế độ thị trường \\
Hoang và Luu (2024) & Phân tích chế độ của biến động chỉ số & Nối chế độ với chất lượng của quyết định phân bổ danh mục \\
\end{longtable}}

Ba công trình này cùng Nguyen và cộng sự (2020) và bài của chính tác giả là các công trình “gần nhất” trong Bảng 1 của đề cương.

\section{Điểm khác biệt thứ tư: so với bài trước của chính tác giả}

Bài trước, Nguyen và Nguyen (2026a), áp dụng Black-Litterman với mô hình định giá tài sản vốn có chuyển chế độ cho Việt Nam giai đoạn 2019--2025 và chỉ đánh giá một mô hình. Đề cương nêu ba khác biệt:

{\itshape 1. Đặt M5 (Black-Litterman có chuyển chế độ, mở rộng từ bài trước) cạnh sáu mô hình khác; M5 thành một ứng viên, không còn đứng một mình.

2. Mở mẫu về trước năm 2019 nếu dữ liệu cho phép (mẫu dự kiến 2007--2026).

3. Đo chênh lệch giữa các mô hình thay vì đánh giá riêng M5.\par}

Kết quả của bài trước chỉ làm đối chứng, không được báo cáo lại như kết quả mới của luận án, để giữ ranh giới giữa công bố đã có và đóng góp mới.

Vì đã dùng dữ liệu 2019--2025 (hai giai đoạn căng thẳng), tác giả không thể khẳng định chưa xem kết quả ở chế độ căng thẳng. Đăng ký trước, khai rõ những kết quả đã xem, giảm chứ không loại thiên lệch này.

M5 được xử lý như mọi mô hình khác, kết quả báo cáo cả khi có lẫn khi không có M5; kiểm định H1 còn được lặp lại trên tập chỉ gồm mô hình tĩnh M1, M2, M3, M4, M6.

\section{Tuyên bố tính mới: giới hạn cần nói ra}

Bốn điểm khác biệt khác với tuyên bố tính mới. Tuyên bố tính mới ở mục 4.2: theo hiểu biết của thí sinh, chưa có nghiên cứu nào đo rủi ro lựa chọn mô hình theo chế độ thị trường cho quyết định phân bổ danh mục trên thị trường chứng khoán Việt Nam.

Hãy giữ nguyên các từ rào. Căn cứ là truy vấn trên SSRN, RePEc/IDEAS, ScienceDirect, Oxford Academic, Wiley, Springer, Taylor \& Francis, MDPI và arXiv đến ngày 30/09/2026.

Phạm vi chưa gồm Scopus, Web of Science và tạp chí tiếng Việt, nên nhận định không loại trừ nghiên cứu nằm ngoài phạm vi đó.

Tra mục lục vài giá sách trong thư viện, bạn có thể nói “tôi chưa thấy”, không thể nói “không có”.

Vì vậy đừng nói “đề tài chưa ai làm”; hãy nói “theo hiểu biết của tác giả, trong phạm vi tìm kiếm đã nêu, chưa thấy”.

\section{Đóng góp nằm ở đâu: ba nhóm kết quả dự kiến}

Đừng nhầm “đóng góp” với “kết quả đã có”. Nếu dữ liệu ủng hộ, đề tài dự kiến ba nhóm kết quả:

\begin{itemize}
\item
  Lý thuyết: khung khái niệm và thước đo rủi ro lựa chọn mô hình theo chế độ.
\item
  Thực nghiệm: bằng chứng Việt Nam về mức độ, thời điểm và các kênh.
\item
  Thực tiễn: thủ tục giám sát mô hình ba bước (khai báo tập mô hình; đo độ phân tán và tập tin cậy định kỳ; kích hoạt kết hợp mô hình khi căng thẳng), suy ra từ H3, chưa kiểm nghiệm với người dùng.
\end{itemize}

Sản phẩm dự kiến (mục 7.3): hai đến ba bài báo, một đến hai báo cáo hội thảo và luận án; số bài là mục tiêu, phụ thuộc dữ liệu và phản biện.

\section{Một cách tự hỏi trước khi đi thi}

Nếu được hỏi “đóng góp lý thuyết cụ thể là gì”, hãy trả lời bằng thước đo, không bằng khẩu hiệu: một thước đo rủi ro lựa chọn mô hình theo chế độ, dựa trên mức hối tiếc, lấp khoảng trống thứ nhất (thước đo hiện có chưa đo hậu quả của chọn sai lên kết quả danh mục).

Nếu được hỏi “kết quả là 1/N đủ dùng thì có đóng góp không”: có. Đề tài vẫn trả lời được câu hỏi: rủi ro không đổi theo chế độ, hoặc 1/N đủ dùng ở Việt Nam.

\begin{tcolorbox}[breakable]
\textbf{Ghi nhớ khi đi thi}

• Ta và Vuong (2020): một mô hình, đề tài bảy mô hình; Goutte và cộng sự (2025): so sánh phương pháp, đề tài đo bất định theo chế độ; Hoang và Luu (2024): chế độ của biến động, đề tài chế độ và chất lượng quyết định.

• So với Nguyen và Nguyen (2026a): M5 cạnh sáu mô hình khác, mẫu mở về trước 2019, đo chênh lệch giữa các mô hình.

• Kết quả của bài trước chỉ làm đối chứng, không báo cáo lại như kết quả mới.

• Tính mới chỉ “theo hiểu biết của thí sinh”, phạm vi tìm kiếm có giới hạn.

• Ba nhóm kết quả dự kiến chỉ có nếu dữ liệu ủng hộ.
\end{tcolorbox}

\textbf{Tự kiểm tra}

{\itshape 1. Nêu các điểm khác biệt giữa đề tài và các công trình gần nhất.\par}

Đáp án: Ta và Vuong (2020) áp dụng một mô hình Black-Litterman, đề tài đặt nó cùng sáu mô hình khác vào tập cạnh tranh; Goutte và cộng sự (2025) so sánh phương pháp đa tiêu chí trên VN100, đề tài đo bất định khi chọn mô hình theo chế độ; Hoang và Luu (2024) phân tích chế độ của biến động, đề tài nối chế độ với chất lượng quyết định; so với bài trước của tác giả, đề tài đo chênh lệch giữa bảy mô hình thay vì đánh giá riêng M5.

{\itshape 2. Đề tài khác bài trước của chính tác giả (Nguyen và Nguyen, 2026a) thế nào?\par}

Đáp án: M5 cạnh sáu mô hình khác; mẫu mở về trước 2019 nếu được; đo chênh lệch giữa các mô hình. Kết quả bài trước chỉ làm đối chứng.

{\itshape 3. Tuyên bố tính mới của đề cương có những rào nào?\par}

Đáp án: “theo hiểu biết của thí sinh”; truy vấn trên một số cơ sở dữ liệu đến 30/09/2026, chưa gồm Scopus, Web of Science và tạp chí tiếng Việt; mới đọc tóm tắt và thông tin thư mục.

{\itshape 4. Đóng góp về thực tiễn dự kiến là gì, và đã được kiểm nghiệm chưa?\par}

Đáp án: thủ tục giám sát ba bước (khai báo tập mô hình; đo độ phân tán và tập tin cậy định kỳ; kích hoạt kết hợp mô hình khi căng thẳng); suy ra từ H3, chưa kiểm nghiệm với người dùng.

\chapter{Mục 4.5 của đề cương: Bố cục của đề cương và luận án}

Đề cương theo mẫu của Trường Quốc tế gồm 11 mục; phần nghiên cứu nằm ở mục 4 (từ 4.1 đến 4.5). Luận án dự kiến có sáu chương và phần Kết luận, đi từ giới thiệu và khung khái niệm, đến dữ liệu và chế độ, rủi ro lựa chọn theo chế độ, biến kênh, rồi cơ chế giảm thiểu và hàm ý quản trị.

Mục này như mục lục của một cuốn sách chưa viết: nó cho biết mỗi phần nằm đâu và trả lời câu hỏi nào. Nhớ được ánh xạ này, bạn trả lời được nhiều câu hỏi “luận án sẽ viết gì” mà không phải học thuộc.

\section{Bố cục của đề cương}

Đề cương gồm 11 mục theo mẫu:

{\small\begin{longtable}[]{@{}
  >{\raggedright\arraybackslash}p{(\columnwidth - 4\tabcolsep) * \real{0.0797}}
  >{\raggedright\arraybackslash}p{(\columnwidth - 4\tabcolsep) * \real{0.3389}}
  >{\raggedright\arraybackslash}p{(\columnwidth - 4\tabcolsep) * \real{0.5815}}@{}}
\toprule\noalign{}
\begin{minipage}[b]{\linewidth}\raggedright
Phần
\end{minipage} & \begin{minipage}[b]{\linewidth}\raggedright
Tiêu đề
\end{minipage} & \begin{minipage}[b]{\linewidth}\raggedright
Việc chính
\end{minipage} \\
\midrule\noalign{}
\endhead
\bottomrule\noalign{}
\endlastfoot
Mục 1--3 & Tên đề cương, người thực hiện, đơn vị công tác & Tên đề tài (không quá 30 từ) và thông tin thí sinh \\
Mục 4 & Đề cương sơ bộ dự kiến & 4.1 Lý do, câu hỏi, mục tiêu; 4.2 Tổng quan và khoảng trống; 4.3 Phương pháp và giả thuyết; 4.4 Đối tượng, phạm vi, đóng góp; 4.5 Bố cục luận án \\
Mục 5 & Lý do chọn cơ sở nghiên cứu & Ba lý do chọn Trường Quốc tế \\
Mục 6 & Mục tiêu và mong muốn khi học nghiên cứu sinh & Năng lực nghiên cứu và người dùng kết quả \\
Mục 7 & Dự định và kế hoạch & Học phần, chuyên đề, nghiên cứu khoa học, kế hoạch bốn năm (Bảng 5) \\
Mục 8 & Kiến thức chuyên môn của người dự tuyển & Kinh nghiệm nghiên cứu, công trình (Bảng 6, 7), sự chuẩn bị \\
Mục 9 & Dự kiến việc làm và hướng nghiên cứu sau tốt nghiệp & Định hướng nghề nghiệp và ba hướng mở rộng \\
Mục 10 & Đề xuất người hướng dẫn & Người hướng dẫn chính và thứ hai (dự kiến) \\
Mục 11 & Danh mục tài liệu tham khảo & Tiếng Việt và tiếng Anh \\
\end{longtable}}

Cột “Việc chính” do bản đọc tóm tắt từ nội dung từng mục.

Phần lớn bài đọc này giải thích mục 4. Các mục 5 đến 10 là phần mang tính tuyển sinh, nói về cơ sở đào tạo, kế hoạch, sự chuẩn bị và người hướng dẫn.

Đề cương giống bản vẽ thiết kế, luận án là ngôi nhà. Bản vẽ có phần nền (4.1, 4.2), phần cách xây (4.3), phần hình dung ngôi nhà (4.4, 4.5), và phần về nơi xây, kế hoạch thi công, năng lực người thiết kế (mục 5 đến 10).

\section{Luận án dự kiến: sáu chương và Kết luận}

Bảng 4 của đề cương liệt kê sáu chương và phần Kết luận, với tiểu mục đến ba chữ số. Bản đọc gom chúng thành năm nội dung:

{\small\begin{longtable}[]{@{}
  >{\raggedright\arraybackslash}p{(\columnwidth - 4\tabcolsep) * \real{0.1100}}
  >{\raggedright\arraybackslash}p{(\columnwidth - 4\tabcolsep) * \real{0.2201}}
  >{\raggedright\arraybackslash}p{(\columnwidth - 4\tabcolsep) * \real{0.6699}}@{}}
\toprule\noalign{}
\begin{minipage}[b]{\linewidth}\raggedright
Nội dung
\end{minipage} & \begin{minipage}[b]{\linewidth}\raggedright
Chương của luận án
\end{minipage} & \begin{minipage}[b]{\linewidth}\raggedright
Chủ đề
\end{minipage} \\
\midrule\noalign{}
\endhead
\bottomrule\noalign{}
\endlastfoot
1 & Chương 1-2 & Giới thiệu, tổng quan và khung khái niệm \\
2 & Chương 3 & Dữ liệu, chế độ thị trường và phương pháp \\
3 & Chương 4 & Rủi ro lựa chọn mô hình theo chế độ (kiểm định H1) \\
4 & Chương 5 & Liên hệ với các biến kênh (kiểm định H2 và H4) \\
5 & Chương 6 và Kết luận & Cơ chế giảm thiểu và hàm ý quản trị (kiểm định H3) \\
\end{longtable}}

Lưu ý: “Chương” ở đây là chương của luận án. Đề cương không chia chương; nó chia mục 1 đến 11. Khi giám khảo hỏi “chương 4”, đó là chương 4 của luận án.

\section{Nội dung 1 (Chương 1--2): giới thiệu, tổng quan và khung khái niệm}

Chương 1 giới thiệu bối cảnh, vấn đề, mục tiêu, câu hỏi, phạm vi và đóng góp. Chương 2 tổng quan tài liệu và dựng khung khái niệm, trong đó tiểu mục 2.2.1 phân biệt ba loại rủi ro dựa trên Derman (1996), Cont (2006), Kerkhof và cộng sự (2010) và Boucher và cộng sự (2014).

Ba loại rủi ro: rủi ro đặc tả (mô hình sai dạng), rủi ro ước lượng (sai số tham số) và rủi ro lựa chọn (chọn sai trong tập mô hình cạnh tranh). Hình dung bạn may một bộ đồ:

\begin{itemize}
\item
  Rủi ro đặc tả: chọn sai khuôn, ví dụ khuôn áo cho người cao gầy trong khi bạn thấp đậm; dạng mô hình sai từ đầu.
\item
  Rủi ro ước lượng: khuôn đúng nhưng thợ đo sai vài centimet; tham số có sai số.
\item
  Rủi ro lựa chọn: có nhiều khuôn đều hợp lý, nhưng phải chọn một khuôn trước khi biết khuôn nào vừa nhất.
\end{itemize}

Ba tầng này không loại trừ nhau trong thực tế; đề tài tập trung vào tầng thứ ba. Sai số ước lượng là chuyện đã biết, đề tài chuyển sang câu hỏi chọn giữa nhiều mô hình.

Derman (1996) là ghi chú thực hành về rủi ro mô hình trong định giá phái sinh; Cont (2006) đo bất định mô hình bằng độ đo rủi ro nhất quán; Kerkhof và cộng sự (2010) tách rủi ro ước lượng và rủi ro đặc tả sai; Boucher và cộng sự (2014) hiệu chỉnh độ đo rủi ro theo kiểm định ngược. Đề cương không ghép từng công trình với từng loại rủi ro, nên bạn cũng đừng ghép.

Chương 2 tổng quan theo đúng sáu nhóm ở mục 4.2: sai số ước lượng và tối ưu hóa danh mục; rủi ro và bất định mô hình; kết hợp mô hình, tập tin cậy và thiên lệch chọn lọc; chế độ thị trường và khủng hoảng; thanh khoản, tương quan cực trị và bầy đàn; nghiên cứu về Việt Nam. Tiểu mục 2.7 khép lại bằng khoảng trống, khung phân tích và giả thuyết.

\section{Nội dung 2 (Chương 3): dữ liệu, chế độ thị trường và phương pháp}

Chương 3 mô tả dữ liệu và cách loại thiên lệch sống sót, ước lượng chế độ (mục 4.3.4), đối chiếu chế độ với các cú sốc lịch sử, rồi trình bày tập bảy mô hình, các thước đo, phương pháp kiểm định và các lệch khỏi kế hoạch đăng ký trước.

Chương này gắn với mục tiêu (i), và là nơi chế độ căng thẳng (trạng thái thống kê) gặp khủng hoảng (sự kiện lịch sử).

Kiểm định điểm gãy Bai-Perron trên toàn mẫu xuất hiện ở đây chỉ để mô tả và đối chiếu, không làm biến chế độ ở chương 4 đến 6. Chế độ dùng cho kiểm định do mô hình Markov ước lượng lại tại mỗi ngày tái cân bằng.

Trước khi ước lượng các mô hình danh mục, tác giả báo cáo số giai đoạn căng thẳng rơi vào phần mẫu tính được mức hối tiếc. Khủng hoảng tài chính toàn cầu chủ yếu phục vụ ước lượng chế độ.

\section{Nội dung 3 (Chương 4): rủi ro lựa chọn mô hình theo chế độ}

Nội dung 3 là lõi thực nghiệm: chương 4 ước lượng bảy mô hình, tính các thước đo cho từng chế độ và kiểm định H1, trả lời câu hỏi 1.

Đây là chương bảy cách chia tiền thi đấu thật: mỗi cửa sổ có bảng điểm lợi suất tương đương chắc chắn, từ đó tính mức hối tiếc, độ phân tán và tập tin cậy, rồi so hai chế độ.

H1 được ủng hộ khi thỏa ba điều kiện: \texorpdfstring{$\beta$}{beta} của R dương và vượt phân vị 95\% của phân phối rỗng; khoảng tin cậy bootstrap khối loại trừ 0; \texorpdfstring{$\beta$}{beta} của D cùng dấu. Kiểm định được lặp lại trên tập mô hình tĩnh M1, M2, M3, M4 và M6.

Dưới 10 cửa sổ độc lập ở chế độ căng thẳng thì kết quả là không kết luận; vì vậy đừng nói trước về kết quả chương 4.

\section{Nội dung 4 (Chương 5): liên hệ với các biến kênh}

Chương 5 hồi quy các thước đo rủi ro theo biến chế độ và biến kênh (thanh khoản, tương quan cực trị, bầy đàn), có kiểm soát biến động, để kiểm định H2 và H4, trả lời câu hỏi 2.

Phần chênh lệch giữa hai chế độ do biến kênh giải thích là hiệu \texorpdfstring{$\beta$}{beta}0 − \texorpdfstring{$\beta$}{beta}1, giữa hệ số \texorpdfstring{$\beta$}{beta} khi phương trình không có và khi có biến kênh, kèm khoảng tin cậy bootstrap khối.

Biến chế độ phụ thuộc biến động và lợi suất, còn các biến kênh cũng phụ thuộc biến động, nên chúng dễ trùng thông tin. Đề tài báo cáo hệ số phóng đại phương sai với ngưỡng 5; vượt ngưỡng thì báo cáo hai phương trình tách riêng và không diễn giải riêng từng hệ số.

Giống kiểm tra xem hai nhân chứng có cùng nghe một nguồn tin không: nếu có, lời họ không phải hai bằng chứng độc lập.

Kết quả chương 5 là liên hệ có điều kiện, không phải nhân quả; H4 chỉ là thăm dò.

\section{Nội dung 5 (Chương 6 và Kết luận): cơ chế giảm thiểu và hàm ý quản trị}

Chương 6 đánh giá ba cơ chế giảm thiểu để kiểm định H3, đề xuất thủ tục giám sát mô hình và hàm ý cho tổ chức đầu tư và cơ quan giám sát. Phần Kết luận trả lời ba câu hỏi, nêu hạn chế (số giai đoạn căng thẳng, phạm vi Việt Nam) và hướng nghiên cứu tiếp.

Ba cơ chế (mục 4.3.8): kết hợp mô hình theo trọng số Bayes hoặc trong tập tin cậy, đối chứng là kết hợp 1/N với trung bình-phương sai; danh mục vững chắc tối ưu cho tình huống xấu nhất trong nhiều niềm tin tiên nghiệm; quy tắc chuyển chế độ theo Kritzman và cộng sự (2012), đối chứng là giá trị của chuyển chế độ theo Tu (2010).

Tiêu chí: lợi suất tương đương chắc chắn, mức hối tiếc, độ sụt giảm tối đa, giá trị thiệt hại kỳ vọng mức 95\%, vòng quay và hiệu quả sau chi phí giao dịch; chi phí được tham số hóa và kiểm tra độ nhạy.

Hai cơ chế (ii) và (iii) nằm ngoài tập mô hình khi kiểm định H3 nên mức hối tiếc của chúng có thể âm; tính không âm chỉ đúng khi mô hình chọn trước thuộc tập M.

\section{Nhìn lại: năm nội dung trả lời ba câu hỏi}

Gom năm nội dung thành một sơ đồ:

{\small\begin{longtable}[]{@{}
  >{\raggedright\arraybackslash}p{(\columnwidth - 4\tabcolsep) * \real{0.5661}}
  >{\raggedright\arraybackslash}p{(\columnwidth - 4\tabcolsep) * \real{0.0931}}
  >{\raggedright\arraybackslash}p{(\columnwidth - 4\tabcolsep) * \real{0.3408}}@{}}
\toprule\noalign{}
\begin{minipage}[b]{\linewidth}\raggedright
Câu hỏi hoặc mục tiêu
\end{minipage} & \begin{minipage}[b]{\linewidth}\raggedright
Giả thuyết
\end{minipage} & \begin{minipage}[b]{\linewidth}\raggedright
Nội dung của luận án
\end{minipage} \\
\midrule\noalign{}
\endhead
\bottomrule\noalign{}
\endlastfoot
Dựng khái niệm và thước đo (mục tiêu i) & không có & Nội dung 1 (Chương 1-2) \\
Dữ liệu và xác định chế độ (mục tiêu i) & không có & Nội dung 2 (Chương 3) \\
Câu hỏi 1: rủi ro thay đổi thế nào giữa hai chế độ & H1 & Nội dung 3 (Chương 4) \\
Câu hỏi 2: biến kênh có liên hệ không, giải thích bao nhiêu phần chênh lệch & H2 và H4 & Nội dung 4 (Chương 5) \\
Câu hỏi 3: cơ chế nào giảm được rủi ro & H3 & Nội dung 5 (Chương 6 và Kết luận) \\
\end{longtable}}

Bảng này là cách sắp xếp của bản đọc; đề cương không có bảng như vậy.

\section{Liên hệ với kế hoạch bốn năm}

Bảng 5 của đề cương nêu kế hoạch bốn năm. Đề cương không ghi năm nào ứng với chương nào, nhưng thứ tự công việc cho thấy:

\begin{itemize}
\item
  Năm 1: xác nhận dữ liệu và giấy phép; xác định khung quy định cho công ty quản lý quỹ; học bổ sung 9 tín chỉ; ước lượng thử trên 2019--2025; làm sạch dữ liệu, đọc toàn văn; đăng ký trước kế hoạch phân tích.
\item
  Năm 2: xác định chế độ, ước lượng bảy mô hình; đo rủi ro lựa chọn, kiểm định H1, viết bài báo 1.
\item
  Năm 3: biến kênh (H2, H4); cơ chế giảm thiểu (H3); viết bài báo 2 và 3.
\item
  Năm 4: hoàn thiện luận án, kết luận và hàm ý quản trị; chỉnh sửa và bảo vệ.
\end{itemize}

Thời gian đào tạo là bốn năm. Dự phòng: dữ liệu không đủ thì thu hẹp rổ hoặc giai đoạn mẫu ngay quý 1 năm thứ nhất; quá ít giai đoạn căng thẳng thì áp dụng quy tắc không kết luận; bài bị từ chối thì chuyển tạp chí khác trong năm.

\section{Điều cần tránh khi nói về bố cục}

Ba lỗi hay gặp:

Nói như thể luận án đã viết xong. Các chương là kế hoạch.

Nhầm mục của đề cương với chương của luận án.

Nói trước kết quả. Chương 4 hay 5 có thể kết thúc bằng “không kết luận” nếu quá ít giai đoạn căng thẳng độc lập.

\begin{tcolorbox}[breakable]
\textbf{Ghi nhớ khi đi thi}

• Đề cương: 11 mục theo mẫu; phần nghiên cứu ở mục 4 (4.1 đến 4.5).

• Luận án dự kiến: sáu chương và Kết luận (Bảng 4 của đề cương).

• Chương 1--2: giới thiệu, tổng quan, khung khái niệm; ba loại rủi ro: đặc tả, ước lượng, lựa chọn.

• Chương 3: dữ liệu, chế độ, đối chiếu với cú sốc lịch sử, phương pháp.

• Chương 4: bảy mô hình, thước đo theo chế độ, H1.

• Chương 5: biến kênh, H2 và H4, kiểm soát biến động.

• Chương 6 và Kết luận: cơ chế giảm thiểu, H3, hàm ý quản trị, hạn chế, hướng tiếp.
\end{tcolorbox}

\textbf{Tự kiểm tra}

{\itshape 1. Luận án dự kiến gồm mấy chương và có mấy nội dung?\par}

Đáp án: sáu chương và phần Kết luận, chia thành năm nội dung.

{\itshape 2. Luận án phân biệt ba loại rủi ro nào?\par}

Đáp án: rủi ro đặc tả (mô hình sai dạng), rủi ro ước lượng (sai số tham số), rủi ro lựa chọn (chọn sai trong tập mô hình cạnh tranh), ở tiểu mục 2.2.1 của luận án.

{\itshape 3. Chương nào của luận án kiểm định H1, chương nào kiểm định H2 và H4, chương nào kiểm định H3?\par}

Đáp án: Chương 4 kiểm định H1; Chương 5 kiểm định H2 và H4; Chương 6 kiểm định H3.

{\itshape 4. Phần Kết luận của luận án gồm những gì?\par}

Đáp án: trả lời ba câu hỏi nghiên cứu, hạn chế (số giai đoạn căng thẳng, phạm vi Việt Nam) và hướng nghiên cứu tiếp.

{\itshape 5. Chương 3 của luận án đối chiếu chế độ với điều gì?\par}

Đáp án: với các cú sốc lịch sử: khủng hoảng tài chính toàn cầu, đại dịch COVID-19 và căng thẳng trái phiếu doanh nghiệp.

\chapter{Mục 4.1 của đề cương (tiếp): Ba loại “rủi ro” dễ nhầm và hai chữ “căng thẳng” dễ lẫn}

Ngay trong mục 4.1, đề cương định nghĩa ba thuật ngữ gần nhau (rủi ro mô hình, bất định mô hình, rủi ro lựa chọn mô hình) và tách hai khái niệm hay bị dùng thay nhau (khủng hoảng, chế độ căng thẳng). Nắm chắc năm khái niệm này thì phần còn lại dễ theo.

Giám khảo có thể hỏi ngay: rủi ro lựa chọn mô hình khác gì rủi ro mô hình nói chung, và độ phân tán mô hình đo cái gì.

\section{Mô hình là gì, và vì sao có thể “rủi ro”}

Mô hình ở đây là một công thức hay quy trình tính toán biến dữ liệu thành quyết định: giá cổ phiếu mấy năm qua đi vào, tỷ lệ tiền cho từng mã đi ra.

Mô hình giống công thức nấu ăn: cùng giỏ nguyên liệu, mỗi công thức cho ra một món. Công thức viết sai, hoặc công thức đúng mà dùng nhầm bếp, thì bữa ăn hỏng.

“Bữa ăn hỏng” là rủi ro mô hình: khả năng phát sinh hậu quả bất lợi từ quyết định dựa trên mô hình sai hoặc dùng sai. Định nghĩa lấy từ thư giám sát SR 11-7 của Cục Dự trữ Liên bang Mỹ (Board of Governors of the Federal Reserve System, 2011).

“Sai” là công thức không phản ánh đúng thực tế; “dùng sai” là công thức tạm ổn nhưng áp vào chỗ không phù hợp. Cả hai đều dẫn tới quyết định kém.

\section{Bất định mô hình: nhiều công thức, nhiều đáp án}

Bất định mô hình là tình trạng nhiều mô hình khả dĩ cho kết quả khác nhau.

Bạn hỏi năm nhà tư vấn, mỗi người một cách tính, nhận về năm khuyến nghị khác nhau. Không ai chắc chắn sai; điều khó là không biết nên tin ai.

Bất định mô hình mô tả tình trạng của thế giới, chưa phải khoản thua lỗ. Thua lỗ chỉ xuất hiện khi bạn phải chọn và chọn trật.

\section{Rủi ro lựa chọn mô hình: thứ mà đề tài đo}

Rủi ro lựa chọn mô hình là tổn thất kỳ vọng, đo bằng mức hối tiếc, phát sinh vì phải chọn một mô hình trước khi biết mô hình nào tốt nhất.

Trước mùa giải bạn nộp một cách chia tiền duy nhất và theo nó suốt mùa. Đề cương mô phỏng quy tắc nộp cách có kết quả ngoài mẫu tốt nhất trong 36 tháng gần nhất; đây là quy tắc mô phỏng, không phải mô tả cách các tổ chức thật chọn.

Mức hối tiếc là hiệu giữa kết quả của cách tốt nhất nhìn lại và kết quả của cách bạn đã nộp: bằng không nếu chọn đúng, dương nếu chọn trật.

\begin{tcolorbox}[breakable]
Ví dụ giả định (số tự đặt để minh họa): ba cách chia tiền A, B, C có điểm sau mùa giải là 5, 8 và 3.

Chọn A vì A thắng giai đoạn trước thì mức hối tiếc là 8 − 5 = 3; chọn B thì bằng 0.
\end{tcolorbox}

Chữ “kỳ vọng” quan trọng: một lần chọn trật có thể chỉ là xui; rủi ro lựa chọn mô hình hỏi trung bình qua nhiều lần chọn bạn mất bao nhiêu vì không biết trước đáp án.

Công thức cụ thể của mức hối tiếc R\textsubscript{t} ở mục 4.3.6. Ở đây chỉ cần ý tưởng: khoảng cách giữa cái bạn có và cái bạn có thể có nếu biết trước.

\section{Độ phân tán mô hình và kích thước tập tin cậy mô hình: chỉ báo của bất định}

Đề cương nói rõ: độ phân tán mô hình và kích thước tập tin cậy mô hình chỉ báo hiệu bất định, không đo rủi ro.

Độ phân tán đo mức các mô hình bất đồng: năm khuyến nghị gần giống nhau thì độ phân tán nhỏ, mỗi người một đằng thì lớn. Nó cho biết tình trạng bất định, chưa cho biết bạn thiệt bao nhiêu.

Tập tin cậy mô hình là danh sách ứng viên chưa thể loại sau vòng sơ loại thống kê. Danh sách dài nghĩa là dữ liệu chưa đủ để loại ai, tức bất định cao.

Vì sao phải tách? Cách hiểu của bản đọc: các công thức có thể bất đồng nhiều nhưng dẫn đến kết quả tương tự, khi đó chọn trật cũng không thiệt mấy. Mức hối tiếc đo hậu quả của việc chọn nên mới là thước đo rủi ro. Đề cương cũng lưu ý tập tin cậy lớn hơn ở chế độ căng thẳng có thể chỉ do dữ liệu kém thông tin.

Bảng tóm tắt:

{\small\begin{longtable}[]{@{}
  >{\raggedright\arraybackslash}p{(\columnwidth - 4\tabcolsep) * \real{0.1760}}
  >{\raggedright\arraybackslash}p{(\columnwidth - 4\tabcolsep) * \real{0.4900}}
  >{\raggedright\arraybackslash}p{(\columnwidth - 4\tabcolsep) * \real{0.3340}}@{}}
\toprule\noalign{}
\begin{minipage}[b]{\linewidth}\raggedright
Thuật ngữ
\end{minipage} & \begin{minipage}[b]{\linewidth}\raggedright
Nghĩa trong đề cương
\end{minipage} & \begin{minipage}[b]{\linewidth}\raggedright
Cách hình dung
\end{minipage} \\
\midrule\noalign{}
\endhead
\bottomrule\noalign{}
\endlastfoot
Rủi ro mô hình & Hậu quả bất lợi từ quyết định dựa trên mô hình sai hoặc dùng sai & Bữa ăn hỏng vì công thức sai hoặc dùng sai \\
Bất định mô hình & Nhiều mô hình khả dĩ cho kết quả khác nhau & Năm nhà tư vấn nói năm đằng \\
Rủi ro lựa chọn mô hình & Tổn thất kỳ vọng, đo bằng mức hối tiếc, do phải chọn trước khi biết mô hình nào tốt nhất & Nộp một cách chia tiền trước mùa giải \\
Độ phân tán mô hình, kích thước tập tin cậy mô hình & Chỉ báo của bất định, không phải của rủi ro & Mức bất đồng giữa các nhà tư vấn; độ dài danh sách chưa thể loại \\
\end{longtable}}

\section{Danielsson và cộng sự (2016): rủi ro mô hình lớn lên theo bất định của thị trường}

Mục 4.2 (nhóm thứ hai) dẫn hai công trình đo rủi ro mô hình. Danielsson và cộng sự (2016) thấy rủi ro mô hình tăng theo mức bất định của thị trường và nhỏ khi thị trường yên ắng.

Giới hạn: phát hiện đó chỉ cho các mô hình dự báo rủi ro thị trường và rủi ro hệ thống, không cho mô hình phân bổ danh mục. Đây là chỗ đề tài bước tiếp.

Hình dung của bản đọc: trời quang thì các công thức dự báo cho đáp án gần nhau, chọn công thức nào cũng ít hệ trọng; trời giông thì chúng tách nhau và chọn sai tốn kém hơn. Công trình này là một cơ sở của H1.

Hãy giới hạn mình ở ba điều đề cương nói: rủi ro mô hình tăng theo bất định, nhỏ khi thị trường yên, và chỉ cho mô hình rủi ro thị trường và rủi ro hệ thống.

\section{Boucher và cộng sự (2014): hiệu chỉnh theo kiểm định ngược}

Boucher và cộng sự (2014) đề xuất hiệu chỉnh độ đo rủi ro theo kết quả kiểm định ngược.

Kiểm định ngược là thử mô hình trên quá khứ: giả vờ đã dùng nó từ nhiều năm trước rồi xem kết quả với dữ liệu đã biết. Nếu mô hình báo rủi ro thấp mà lịch sử cho thấy thua lỗ lớn hơn, có lý do để điều chỉnh con số rủi ro. Đây là giải thích chung; chi tiết cách làm của Boucher và cộng sự đề cương không nêu.

Hai công trình này là các thước đo rủi ro mô hình hiện có, và đều đo chênh lệch giữa các dự báo rủi ro, chưa đo hậu quả của việc chọn sai lên kết quả danh mục. Đó là khoảng trống thứ nhất.

\section{Khủng hoảng khác chế độ căng thẳng}

Đề cương quy ước: khủng hoảng là sự kiện lịch sử; chế độ căng thẳng là trạng thái thống kê do mô hình chuyển chế độ ước lượng.

Khủng hoảng là bản tin về một cơn bão có tên và ngày tháng. Chế độ căng thẳng là kết luận của trạm khí tượng sau khi nhìn số liệu: thời tiết đang ở trạng thái bão, bất kể báo chí gọi là gì. Trạm có thể báo bão vào những ngày không ai đặt tên, và một cơn bão nổi tiếng chưa chắc làm trạm báo bão đúng ngày người ta nhớ.

Lý do giữ ranh giới, theo bản đọc: hai khái niệm thuộc hai loại khác nhau (sự kiện được kể lại và trạng thái do mô hình suy ra), và muốn đối chiếu chúng thì phải coi là hai thứ riêng.

Đề tài kiểm định rủi ro lựa chọn mô hình theo chế độ căng thẳng, rồi đối chiếu chế độ đó với các cuộc khủng hoảng.

Chế độ bình thường là trời quang, chế độ căng thẳng là mùa bão, còn khủng hoảng là tên một số cơn bão trong sử sách.

\section{Vì sao mục này quan trọng với cả đề cương}

Mục toàn định nghĩa này quan trọng vì: mọi giả thuyết về sau dùng các thuật ngữ này (H1 lấy mức hối tiếc làm kết quả chính); đề cương không dùng chúng thay nhau; và các chỉ báo bất định được xếp vào kết quả phụ.

Đề tài đo hệ quả bằng số trên dữ liệu thị trường, không quan sát người ra quyết định. “Phải chọn một mô hình” là mô phỏng một quy tắc chọn, không mô tả cách các tổ chức thật chọn.

\begin{tcolorbox}[breakable]
\textbf{Ghi nhớ khi đi thi}

• Rủi ro mô hình: hậu quả bất lợi từ quyết định dựa trên mô hình sai hoặc dùng sai (SR 11-7, 2011).

• Bất định mô hình: nhiều mô hình khả dĩ cho kết quả khác nhau.

• Rủi ro lựa chọn mô hình: tổn thất kỳ vọng, đo bằng mức hối tiếc, do phải chọn trước khi biết mô hình nào tốt nhất.

• Độ phân tán và kích thước tập tin cậy mô hình chỉ báo bất định, không đo rủi ro.

• Danielsson và cộng sự (2016): rủi ro mô hình tăng theo bất định, chỉ cho mô hình rủi ro thị trường và hệ thống; Boucher và cộng sự (2014): hiệu chỉnh theo kiểm định ngược.

• Khủng hoảng là sự kiện lịch sử; chế độ căng thẳng là trạng thái thống kê; không dùng thay nhau.
\end{tcolorbox}

\section{Tự kiểm tra}

{\itshape 1. Rủi ro mô hình khác bất định mô hình ở chỗ nào?\par}

Đáp án: rủi ro mô hình là khả năng có hậu quả bất lợi từ mô hình sai hoặc dùng sai; bất định mô hình là việc nhiều mô hình khả dĩ cho kết quả khác nhau.

{\itshape 2. Rủi ro lựa chọn mô hình được đo bằng đại lượng nào, và đại lượng đó nghĩa là gì?\par}

Đáp án: mức hối tiếc, tức khoản thiệt vì không chọn mô hình tốt nhất sau thực tế.

{\itshape 3. Độ phân tán mô hình và kích thước tập tin cậy mô hình đo rủi ro hay bất định?\par}

Đáp án: bất định.

{\itshape 4. Danielsson và cộng sự (2016) thấy điều gì, và giới hạn nào đi kèm?\par}

Đáp án: rủi ro mô hình tăng theo bất định của thị trường, nhỏ khi thị trường yên; chỉ cho mô hình dự báo rủi ro thị trường và rủi ro hệ thống.

{\itshape 5. Khủng hoảng và chế độ căng thẳng khác nhau thế nào?\par}

Đáp án: khủng hoảng là sự kiện lịch sử; chế độ căng thẳng là trạng thái thống kê do mô hình ước lượng. Đề tài đối chiếu, không dùng thay nhau.

\chapter{Mục 4.2 của đề cương, nhóm (1): Sai số ước lượng và tối ưu hóa danh mục}

Mọi mô hình phân bổ đều dùng dữ liệu có sai số, và không mô hình nào thắng ổn định. Từ đó đề cương coi việc chọn mô hình là một quyết định có rủi ro riêng.

Phần này đi qua: sai số lọt vào đâu, ba hướng khắc phục, và vì sao chúng vẫn không chốt được câu chuyện. Khái niệm trung bình-phương sai và hiệp phương sai đã giải thích ở phần một của mục 4.1.

\section{Tối ưu hóa trung bình-phương sai: cách chia tiền “khoa học” đầu tiên}

Nhắc lại: khung trung bình-phương sai của Markowitz (1952) chọn danh mục có lợi suất kỳ vọng cao nhất ứng với một mức phương sai cho trước.

Giống tìm tuyến đường nhanh nhất mà độ xóc không vượt một ngưỡng: đặt mức dao động chịu được rồi tìm cách chia tiền lời cao nhất.

Gọi là “tối ưu hóa” vì máy tính tìm ra bộ trọng số tốt nhất theo công thức. Bộ trọng số đó phụ thuộc hai nhóm số đầu vào: lợi suất kỳ vọng của từng mã và hiệp phương sai giữa các mã.

\section{Sai số đi vào từ đâu và truyền đi thế nào}

Các số đầu vào phải ước lượng từ quá khứ, và sai số khi ước lượng chúng truyền sang trọng số danh mục: số đầu vào lệch thì cách chia tiền lệch theo.

Giống dẫn một chiếc xe chạy rất nhanh bằng bản đồ vẽ tay có sai số: xe càng bám tuyệt đối theo bản đồ thì càng xa đường thật.

Michaud (1989): tối ưu hóa dồn trọng số lớn vào tài sản có sai số ước lượng lớn. Công thức không biết con số nào đáng tin; một mã có lợi suất ước lượng cao chỉ do sai số vẫn được dồn nhiều tiền.

\begin{tcolorbox}[breakable]
Ví dụ giả định (số tự đặt để minh họa): hai cổ phiếu X và Y thực ra có lợi suất kỳ vọng như nhau, nhưng do ngẫu nhiên, dữ liệu quá khứ cho X trông cao hơn một chút.

Công thức chỉ nhìn số ước lượng nên dồn phần lớn tiền vào X; khi tương lai đến, X không hơn Y và khoản dồn đó không đem lại gì.
\end{tcolorbox}

Kan và Zhou (2007): thay tham số tổng thể (giá trị thật mà ta không thấy) bằng ước lượng mẫu (con số tính từ dữ liệu hữu hạn) làm hiệu quả ngoài mẫu giảm mạnh.

Trong mẫu thì kết quả luôn đẹp vì công thức đã “nhìn đáp án”; ngoài mẫu mới là bài thi thật, và bài thi đó có thể đạt điểm kém.

Đề cương không nêu chi tiết cách hai nhóm tác giả chứng minh; hãy nhớ các kết luận trên, nền của nhận định “mọi mô hình phân bổ đều có sai số ở đầu vào”.

\section{Ba hướng khắc phục}

Tài liệu đề xuất ba hướng khắc phục; hướng nào cũng muốn công thức bớt tin tuyệt đối vào số đầu vào.

\subsection{Hướng thứ nhất: co rút hiệp phương sai (Ledoit và Wolf, 2004)}

Ledoit và Wolf (2004) co rút ma trận hiệp phương sai về một ma trận mục tiêu có cấu trúc.

Co rút là kéo một ước lượng quá nhiễu về phía một giá trị đơn giản, ổn định hơn. Một học sinh thi thử một lần được điểm rất cao, bạn chưa nên tin em luôn đạt như vậy mà kéo nhận định về phía mức trung bình của lớp.

Nhiều cổ phiếu thì có rất nhiều cặp hiệp phương sai, mỗi cặp là một ước lượng nhiễu. Co rút làm các con số cực đoan bớt cực đoan, nên trọng số ít bị sai số làm lệch. Chi tiết kỹ thuật đề cương không nêu.

Đây là cơ sở của M2 (phương sai nhỏ nhất với hiệp phương sai co rút), và cũng là phương pháp mà Nguyen và cộng sự (2020) dùng cho Việt Nam.

\subsection{Hướng thứ hai: Black-Litterman (Black và Litterman, 1992)}

Black và Litterman (1992) kết hợp lợi suất cân bằng của mô hình định giá tài sản vốn với quan điểm riêng của nhà đầu tư.

Mô hình định giá tài sản vốn nói rằng khi thị trường cân bằng, lợi suất kỳ vọng của mỗi tài sản gắn với mức rủi ro của nó đối với thị trường. Ý chính: nó cho một điểm xuất phát có logic, do thị trường hình thành, thay vì một ước lượng mẫu nhiễu.

Quan điểm của nhà đầu tư là dự báo riêng, chẳng hạn “cổ phiếu ngân hàng sẽ tốt hơn thị trường”. Theo cách hiểu của bản đọc, Black-Litterman lấy cân bằng thị trường làm điểm tựa rồi dịch về phía quan điểm đó.

Giống mua xe: bạn có giá thị trường làm mốc, nghe bạn bè nói chiếc này đáng giá hơn, và bạn chỉnh mốc một chút về phía lời bạn chứ không bỏ hẳn mốc.

Bessler và cộng sự (2017) thấy Black-Litterman vượt cả danh mục trung bình-phương sai lẫn 1/N về hiệu quả ngoài mẫu trên dữ liệu đa tài sản toàn cầu.

M4 và M5 thuộc họ Black-Litterman, và tác giả đã làm đề tài sinh viên về Black-Litterman cho ngành ngân hàng Việt Nam.

\subsection{Hướng thứ ba: né tránh mơ hồ (Garlappi và cộng sự, 2007)}

Garlappi và cộng sự (2007) mô hình hóa nhà đầu tư né tránh mơ hồ với nhiều niềm tin tiên nghiệm và thu được danh mục ổn định hơn.

Né tránh mơ hồ là ngại những tình huống mà chính xác suất cũng chưa rõ, khác với ngại rủi ro (biết xác suất nhưng không thích dao động). Một túi bi biết rõ nửa đỏ nửa xanh khác với một túi không biết tỷ lệ; nhiều người thích đặt cược vào túi thứ nhất.

“Nhiều niềm tin tiên nghiệm” nghĩa là nhà đầu tư giữ nhiều bộ ước lượng cùng lúc thay vì tin một bộ, rồi tối ưu cho tình huống xấu nhất trong số đó (mục 4.3.8). Tiêu chí của Gilboa và Schmeidler (1989) ở nhóm (2) cũng dựa trên nhiều niềm tin tiên nghiệm.

Đây là cơ sở của M6 và của cơ chế giảm thiểu thứ hai.

\section{Không hướng nào thắng ổn định: DeMiguel và cộng sự (2009)}

Sau ba hướng khắc phục là bằng chứng làm mọi thứ lung lay: DeMiguel và cộng sự (2009) đánh giá 14 mô hình trên bảy bộ dữ liệu và không thấy mô hình nào vượt nhất quán 1/N, xét theo Sharpe, lợi suất tương đương chắc chắn hoặc vòng quay.

Quy tắc 1/N chia đều vốn cho N tài sản: không công thức, không ước lượng, nên không có sai số ước lượng để lọt vào.

Mười bốn mô hình phức tạp không mô hình nào thắng nhất quán cách chia trứng đều. Điều này không có nghĩa 1/N luôn tốt nhất, chỉ là trong bảy bộ dữ liệu đó không mô hình nào vượt nổi nó một cách nhất quán.

Pflug và cộng sự (2012) đi xa hơn về lý thuyết: khi mức mơ hồ về mô hình đủ cao, quy tắc 1/N xấp xỉ tối ưu theo nghĩa vững chắc. Kết luận của đề cương: chọn mô hình là một quyết định có rủi ro riêng, vì nếu mô hình thắng cuộc thay đổi theo dữ liệu thì người phải chọn luôn có thể chọn trật.

\section{Vì sao đề cương dựa vào 1/N để đặt lập luận}

1/N có hai vai trong đề cương: là M1 trong tập bảy mô hình, và là một khả năng đối lập của H1 (rủi ro không đổi theo chế độ, hoặc 1/N thắng ổn định; DeMiguel và cộng sự, 2009; Pflug và cộng sự, 2012). Nếu giả thuyết bị bác bỏ, đề tài vẫn trả lời được: 1/N đủ dùng ở Việt Nam.

Đề cương đặt các mô hình phức tạp cạnh một mô hình sơ đẳng và để dữ liệu quyết định; nó không nhận xét trước mô hình nào tốt.

\section{Điều đề cương không nói ở mục này}

Phần này không báo cáo kết quả nào của đề tài, cũng không nói ba hướng khắc phục là sai; nó chỉ nói không hướng nào thắng ổn định.

Vì tác giả mới đọc tóm tắt của phần lớn nguồn, đề cương mô tả các công trình ở mức kết luận chính. Bạn nên làm tương tự: nêu đúng điều đề cương nói, không thêm chi tiết kỹ thuật.

\begin{tcolorbox}[breakable]
\textbf{Ghi nhớ khi đi thi}

• Markowitz (1952): lợi suất kỳ vọng cao nhất với phương sai cho trước.

• Sai số đầu vào truyền sang trọng số: Michaud (1989) dồn trọng số vào tài sản sai số lớn; Kan và Zhou (2007) ước lượng mẫu làm hiệu quả ngoài mẫu giảm mạnh.

• Ba hướng khắc phục: co rút (Ledoit và Wolf, 2004); Black-Litterman (1992; Bessler và cộng sự, 2017); né tránh mơ hồ với nhiều niềm tin tiên nghiệm (Garlappi và cộng sự, 2007).

• DeMiguel và cộng sự (2009): không mô hình nào vượt nhất quán 1/N; Pflug và cộng sự (2012): 1/N xấp xỉ tối ưu khi mơ hồ cao.

• Kết luận của đề cương: chọn mô hình là một quyết định có rủi ro riêng.
\end{tcolorbox}

\section{Tự kiểm tra}

{\itshape 1. Theo Michaud (1989), tối ưu hóa trung bình-phương sai làm gì với tài sản có sai số ước lượng lớn?\par}

Đáp án: dồn trọng số lớn vào các tài sản đó.

{\itshape 2. Black-Litterman kết hợp hai thứ nào?\par}

Đáp án: lợi suất cân bằng của mô hình định giá tài sản vốn và quan điểm của nhà đầu tư.

{\itshape 3. Bessler và cộng sự (2017) thấy gì về Black-Litterman?\par}

Đáp án: vượt trung bình-phương sai và 1/N về hiệu quả ngoài mẫu trên dữ liệu đa tài sản toàn cầu.

{\itshape 4. Garlappi và cộng sự (2007) mô hình hóa nhà đầu tư thế nào và thu được gì?\par}

Đáp án: nhà đầu tư né tránh mơ hồ với nhiều niềm tin tiên nghiệm; thu được danh mục ổn định hơn.

{\itshape 5. Vì sao DeMiguel và cộng sự (2009) quan trọng với lập luận của đề cương?\par}

Đáp án: không mô hình nào vượt nhất quán 1/N trên bảy bộ dữ liệu, nên chọn mô hình là một quyết định có rủi ro riêng.

\chapter{Mục 4.2 của đề cương, nhóm (2): Rủi ro và bất định mô hình đã được đo ở đâu, chưa được đo ở đâu}

Rủi ro mô hình đã được đo cho dự báo và định giá, chưa được đo cho quyết định phân bổ danh mục. Đó là nhận xét trung tâm của nhóm (2), và đề tài nhắm đúng chỗ trống này.

Nhóm này điểm tám công trình: Derman (1996), Cont (2006), Kerkhof và cộng sự (2010), Boucher và cộng sự (2014), Danielsson và cộng sự (2016), Lo và Mueller (2010), Gilboa và Schmeidler (1989), Hansen và Sargent (2008). Với mỗi nguồn, ta hỏi: họ làm gì, nghĩa là gì, vì sao đề cương nhắc.

\section{Bức tranh chung: ba nhóm nguồn}

Cách chia của bản đọc, không phải của đề cương: tám công trình thuộc ba nhóm. Nhóm thứ nhất đo rủi ro mô hình trong định giá và dự báo (Derman, Cont, Kerkhof, Boucher, Danielsson). Nhóm thứ hai cảnh báo về độ chính xác giả (Lo và Mueller). Nhóm thứ ba là lý thuyết ra quyết định khi không chắc mô hình nào đúng (Gilboa và Schmeidler, Hansen và Sargent).

Nhóm một kiểm tra cân sai bao nhiêu; nhóm hai nhắc rằng cân có ba chữ số thập phân vẫn có thể sai; nhóm ba trả lời nên mua hàng thế nào khi không biết cân nào đúng.

\section{Derman (1996): rủi ro mô hình trong định giá phái sinh}

Đề cương mô tả Derman (1996) bàn về rủi ro mô hình trong định giá phái sinh từ góc độ thực hành, không nêu thêm nội dung bên trong.

Phái sinh là hợp đồng tài chính có giá trị phụ thuộc giá của một tài sản khác; định giá phái sinh là dùng mô hình tính xem hợp đồng đáng giá bao nhiêu. Mô hình sai thì giá sai, và Derman viết về rủi ro đó từ góc nhìn người làm nghề.

Đề cương nhắc Derman để cho thấy rủi ro mô hình đã được bàn trong định giá. Ở bố cục luận án (Bảng 4), Derman (1996), Cont (2006) và Boucher và cộng sự (2014) là nền để phân biệt ba loại rủi ro ở chương 1: rủi ro đặc tả, rủi ro ước lượng và rủi ro lựa chọn.

\section{Cont (2006): đo bất định mô hình bằng độ đo rủi ro nhất quán}

Theo đề cương, Cont (2006) đo bất định mô hình bằng độ đo rủi ro nhất quán trên một tập mô hình định giá.

Độ đo rủi ro là một con số tóm tắt mức nguy hiểm của khoản đầu tư, ví dụ mức lỗ có thể gặp trong tình huống xấu. “Nhất quán” nghĩa là con số tuân theo một số quy tắc hợp lý để không ra kết luận vô lý. Đây là nghĩa chung của thuật ngữ; cách làm cụ thể của Cont đề cương không viết.

Chữ bất định mô hình ở đây đúng nghĩa trong mục 4.1: nhiều mô hình khả dĩ cho kết quả khác nhau. Cont biến bất định đó thành một con số; theo bản đọc, ông được nhắc vì đây là một cách đo bất định mô hình đã có sẵn.

\section{Kerkhof và cộng sự (2010), Boucher và cộng sự (2014), Danielsson và cộng sự (2016)}

Kerkhof và cộng sự (2010) tách rủi ro mô hình thành rủi ro ước lượng (tham số tính sai) và rủi ro đặc tả sai (dạng mô hình sai), rồi đưa cả hai vào vốn dự phòng, tức khoản tiền tổ chức phải giữ để bù cho khả năng mô hình sai. Hai công trình sau chuyển câu hỏi sang các mô hình rủi ro thị trường và rủi ro hệ thống.

Rủi ro thị trường là rủi ro thua lỗ do giá biến động; rủi ro hệ thống là rủi ro cả hệ thống tài chính bị ảnh hưởng chứ không chỉ một cá nhân. Đó là nơi người ta dùng mô hình để dự báo tổn thất tiềm tàng.

Như đã nói ở mục 4.1: Boucher và cộng sự hiệu chỉnh độ đo rủi ro theo kết quả kiểm định ngược; Danielsson và cộng sự thấy rủi ro mô hình tăng theo bất định của thị trường, nhỏ khi thị trường yên ắng, nhưng chỉ cho mô hình dự báo rủi ro thị trường và rủi ro hệ thống.

Nhận xét chung của đề cương về năm nguồn đầu: thước đo của họ gắn với dự báo rủi ro và định giá, tức trả lời “mô hình dự báo lệch bao nhiêu”. Đề tài hỏi câu khác: chọn sai mô hình để phân bổ vốn thì danh mục thiệt bao nhiêu.

Các nguồn kia đo độ sai của chiếc cân; đề tài đo người mua chọn nhầm cân thì thiệt bao nhiêu tiền hàng. Hai việc liên quan nhưng không phải một.

\section{Lo và Mueller (2010): cảm giác chính xác giả tạo}

Lo và Mueller (2010) cảnh báo việc mô hình hóa tài chính theo chuẩn mực vật lý tạo ra cảm giác chính xác giả tạo, và phân loại các mức bất định.

Vật lý có định luật rất chính xác: ném một vật, công thức cho quỹ đạo đúng tới nhiều chữ số. Theo bản đọc, mô hình tài chính mượn công cụ toán ấy nên người ta dễ tưởng con số tài chính cũng chính xác như vậy.

Đo chiều cao một đám mây bằng thước cặp chính xác đến phần mười milimét: con số trông chính xác, nhưng đám mây không có mép rõ ràng.

Về phân loại các mức bất định, đề cương chỉ nói có đề xuất phân loại, không nêu từng mức hay mục đích.

Đề cương xếp Lo và Mueller vào nhóm rủi ro và bất định mô hình, không nêu lý do riêng. Đề cương có coi việc nhận diện chế độ là một dạng rủi ro mô hình (mục 4.3.4, nên dùng định nghĩa thay thế để kiểm tra độ bền), nhưng không nối điều này với Lo và Mueller; đừng nói rằng đề cương làm vậy vì họ.

\section{Gilboa và Schmeidler (1989): chọn phương án tốt nhất trong tình huống xấu nhất}

Gilboa và Schmeidler (1989) xây dựng tiêu chí chọn phương án tốt nhất trong tình huống xấu nhất khi nhà đầu tư có nhiều niềm tin tiên nghiệm. Trong tài liệu, tiêu chí này gọi là maxmin.

Với mỗi phương án, lấy kết quả nhỏ nhất (min) trên các cách nhìn; rồi chọn phương án có kết quả nhỏ nhất đó lớn nhất (max). Tức là chọn phương án mà trường hợp tệ nhất vẫn khá nhất.

\begin{tcolorbox}[breakable]
Ví dụ giả định (số tự đặt để minh họa): hai phương án P và Q, ba cách nhìn về thị trường. P cho 9, 4, 1; Q cho 6, 5, 4.

Tình huống xấu nhất của P là 1, của Q là 4. Theo tiêu chí này bạn chọn Q, dù P có kịch bản cao nhất.
\end{tcolorbox}

Niềm tin tiên nghiệm là điều bạn tin về thế giới trước khi nhìn dữ liệu. Chỉ có một niềm tin thì không có “xấu nhất trong nhiều cách nhìn”; giữ nhiều niềm tin vì không biết cái nào đúng thì tiêu chí này cho một quy tắc rõ: chuẩn bị cho cách nhìn bất lợi nhất.

Ở khung lý thuyết (mục 4.3.1), Gilboa và Schmeidler cùng Hansen và Sargent thuộc khối lý thuyết quyết định dưới bất định mô hình, cơ sở của các cơ chế vững chắc. Ý tưởng này gần với Garlappi và cộng sự (2007) ở nhóm (1), và với cơ chế giảm thiểu thứ hai (mục 4.3.8): danh mục vững chắc tối ưu cho tình huống xấu nhất trong nhiều niềm tin tiên nghiệm.

\section{Hansen và Sargent (2008): kiểm soát vững chắc trước sai đặc tả}

Hansen và Sargent (2008) phát triển lý thuyết kiểm soát vững chắc trước sai đặc tả.

Sai đặc tả là mô hình sai ngay ở dạng của nó. Kiểm soát vững chắc là ra quyết định sao cho vẫn ổn dù mô hình sai: người lái tàu biết la bàn có thể lệch nên chọn hướng không dẫn tàu vào đá ngay cả khi la bàn lệch.

Đề cương không nói mối nối từng bước giữa hai nguồn lý thuyết quyết định với M6; bạn chỉ nên nói ý tưởng vững chắc xuất hiện ở M6 và ở cơ chế giảm thiểu thứ hai.

\section{Đề tài chuyển thước đo sang quyết định phân bổ}

Đề cương khép nhóm (2) bằng câu: các thước đo trong nhóm này gắn với dự báo rủi ro và định giá; chưa thước đo nào đo hậu quả của việc chọn mô hình lên kết quả danh mục.

Đây là khoảng trống (1) trong Bảng 1: thước đo hiện có chưa gắn với hậu quả của quyết định phân bổ theo chế độ thị trường.

Câu kết nhóm (2) không kèm cụm “theo hiểu biết”, nhưng tuyên bố tính mới ở cuối mục 4.2 có cụm đó, với phạm vi tìm kiếm giới hạn và việc mới đọc tóm tắt. Theo bản đọc, nên trình bày nhận xét này thận trọng, như nhận định của tác giả.

Việt Nam không phải lý do của khoảng trống. Việt Nam là nơi kiểm chứng (mục 4.1): thanh khoản hạn chế và nhiều cú sốc trong lịch sử giao dịch ngắn nên khoảng cách giữa các mô hình có thể lớn; khoảng trống nằm ở thước đo, không ở địa lý.

\section{Trả lời một câu hỏi thi điển hình}

Hội đồng hỏi: “rủi ro mô hình đã được đo nhiều, đề tài có gì mới?” Bốn ý: rủi ro mô hình đã được đo cho dự báo và định giá; các thước đo đó đo chênh lệch dự báo rủi ro; chưa đo hậu quả của chọn sai lên kết quả danh mục; đề tài chuyển thước đo sang quyết định phân bổ. Không cần khẳng định chưa ai làm, vì tính mới chỉ là theo hiểu biết của tác giả.

\begin{tcolorbox}[breakable]
\textbf{Ghi nhớ khi đi thi}

• Rủi ro mô hình đã được đo cho dự báo và định giá, chưa cho quyết định phân bổ.

• Derman (1996): rủi ro mô hình trong định giá phái sinh, góc thực hành. Cont (2006): đo bất định mô hình bằng độ đo rủi ro nhất quán. Kerkhof và cộng sự (2010): rủi ro ước lượng và rủi ro đặc tả sai, đưa vào vốn dự phòng.

• Boucher và cộng sự (2014), Danielsson và cộng sự (2016): mô hình rủi ro thị trường và rủi ro hệ thống.

• Lo và Mueller (2010): mô hình hóa tài chính theo chuẩn vật lý tạo cảm giác chính xác giả tạo; phân loại các mức bất định.

• Gilboa và Schmeidler (1989): tốt nhất trong tình huống xấu nhất, nhiều niềm tin tiên nghiệm; Hansen và Sargent (2008): kiểm soát vững chắc trước sai đặc tả.

• Đề tài chuyển thước đo sang quyết định phân bổ danh mục (khoảng trống (1)).
\end{tcolorbox}

\section{Tự kiểm tra}

{\itshape 1. Rủi ro mô hình đã được đo cho những việc gì, và chưa được đo cho việc gì?\par}

Đáp án: Đã đo cho dự báo và định giá; chưa đo cho quyết định phân bổ.

{\itshape 2. Tiêu chí của Gilboa và Schmeidler (1989) là gì?\par}

Đáp án: chọn phương án tốt nhất trong tình huống xấu nhất, khi có nhiều niềm tin tiên nghiệm.

{\itshape 3. Lo và Mueller (2010) cảnh báo điều gì?\par}

Đáp án: mô hình hóa tài chính theo chuẩn vật lý tạo cảm giác chính xác giả tạo; họ phân loại các mức bất định.

{\itshape 4. Kerkhof và cộng sự (2010) và Hansen và Sargent (2008) làm gì?\par}

Đáp án: Kerkhof và cộng sự tách rủi ro ước lượng và rủi ro đặc tả sai, đưa vào vốn dự phòng; Hansen và Sargent phát triển kiểm soát vững chắc trước sai đặc tả.

{\itshape 5. Đề tài khác các thước đo trong nhóm (2) ở đâu?\par}

Đáp án: các thước đo đó gắn với dự báo rủi ro và định giá; đề tài đo hậu quả của việc chọn mô hình lên kết quả danh mục.

\chapter{Mục 4.2 của đề cương, nhóm (3): Kết hợp mô hình, chọn mô hình và kiểm định}

Khi nhiều mô hình bất đồng, đề cương nêu ba cách xử lý: kết hợp mô hình, thu hẹp tập mô hình, và hiệu chỉnh thiên lệch chọn lọc. Nhóm (3) điểm các công trình đứng sau từng cách.

Cách nhớ nhanh: cách thứ nhất là không chọn mà trộn; cách thứ hai là chọn một nhóm thay vì một ứng viên; cách thứ ba là tự kiểm tra xem mình có đang tự lừa vì đã thử quá nhiều lần.

\section{Cách thứ nhất: kết hợp mô hình}

Thay vì nghe một nhà tư vấn, bạn nghe cả năm người rồi lấy trung bình có trọng số; người có thành tích tốt hơn nặng ký hơn một chút. Bạn không cần đoán ai đúng, chỉ cần không đặt cược tất cả vào một người.

Theo bản đọc, nếu chọn một mô hình có thể chọn trật thì không chọn một mô hình duy nhất là cách tránh phần nào nguy cơ đó. Đề cương dựa trên tinh thần này ở H3 và cơ chế giảm thiểu (i).

\subsection{Hoeting và cộng sự (1999): trung bình hóa mô hình Bayes}

Hoeting và cộng sự (1999) trình bày phương pháp trung bình hóa mô hình Bayes.

Giải thích chung, không phải chi tiết bài gốc: “trung bình hóa” là lấy trung bình dự báo của nhiều mô hình; “Bayes” là dùng xác suất để biểu diễn mức tin. Bạn bắt đầu với niềm tin tiên nghiệm về mô hình nào có khả năng đúng, cập nhật sau khi nhìn dữ liệu; mô hình giải thích dữ liệu tốt hơn nhận trọng số lớn hơn.

\begin{tcolorbox}[breakable]
Ví dụ giả định (số tự đặt để minh họa): ba mô hình dự báo lợi suất tháng sau là 2\%, 0\% và 1\%; sau khi nhìn dữ liệu, mức tin vào ba mô hình là 50\%, 30\% và 20\%.

Dự báo kết hợp = 0,5 × 2\% + 0,3 × 0\% + 0,2 × 1\% = 1,2\%. Một mô hình hóa ra tệ không kéo cả dự báo đi quá xa vì nó chỉ chiếm một phần.
\end{tcolorbox}

Đề cương không nêu kết quả thực nghiệm của Hoeting và cộng sự. Họ được nhắc vì trọng số Bayes là một cách trộn mô hình trong cơ chế giảm thiểu (i).

\subsection{Avramov (2002): áp dụng vào dự báo lợi suất cổ phiếu}

Avramov (2002) áp dụng phương pháp này vào dự báo lợi suất cổ phiếu và thu được hiệu quả ngoài mẫu tốt hơn việc dùng tiêu chí để chọn một mô hình.

Tiêu chí chọn mô hình là quy tắc chọn một mô hình duy nhất, ví dụ mô hình khớp dữ liệu tốt nhất. Ở Avramov, trộn thắng chọn trên dữ liệu chưa dùng để tính.

Theo bản đọc, đây là bằng chứng cho hướng kết hợp trong dự báo lợi suất, chưa phải bằng chứng cho phân bổ danh mục ở Việt Nam.

\subsection{Timmermann (2006): tổng hợp bằng chứng về kết hợp dự báo}

Timmermann (2006) tổng hợp bằng chứng rằng dự báo kết hợp thường tốt hơn dự báo riêng lẻ.

Kết hợp dự báo là trộn nhiều dự báo của cùng một đại lượng thành một. Chữ “thường” quan trọng: câu này không nói kết hợp luôn thắng mọi dự báo riêng lẻ.

Có một tuyển thủ giỏi nhất nhưng bạn không biết trước là ai. Rút ngẫu nhiên một người thì thành tích kỳ vọng là mức trung bình của đội; cả đội cùng thi thường tốt hơn mức đó. Bạn mất khả năng bằng người giỏi nhất, được việc tránh chọn nhầm người kém.

Đề cương không nêu tên nghiên cứu nào trong tổng hợp đó.

\subsection{Rapach và cộng sự (2010): kết hợp bù phần mất do mô hình bất ổn định}

Rapach và cộng sự (2010) cho thấy kết hợp dự báo bù được phần chất lượng dự báo mất đi do mô hình bất ổn định.

Mô hình bất ổn định là mô hình tốt ở giai đoạn này chưa chắc tốt ở giai đoạn sau, vì quan hệ trong dữ liệu đổi theo thời gian: con đường tốt mùa khô có thể lầy mùa mưa, và bản đồ vẽ mùa khô không còn dùng được.

Rapach và cộng sự là một trong ba cơ sở của H3 (Bảng 2): ở chế độ căng thẳng, sau chi phí giao dịch, kết hợp mô hình và danh mục vững chắc cho mức hối tiếc và độ sụt giảm tối đa thấp hơn, lợi suất tương đương chắc chắn cao hơn so với chọn một mô hình. Đây là giả thuyết cần kiểm định.

\subsection{Tu và Zhou (2011): trộn hai quy tắc đối nghịch}

Tu và Zhou (2011) kết hợp quy tắc 1/N với quy tắc trung bình-phương sai và vượt 1/N trong hầu hết kịch bản mô phỏng.

Hai quy tắc ở hai thái cực: 1/N không dùng ước lượng nên không có sai số ước lượng nhưng cũng không dùng thông tin; trung bình-phương sai khai thác thông tin nhưng nhạy với sai số (nhóm (1)). Theo bản đọc, trộn hai quy tắc là thỏa hiệp giữa an toàn và thông tin.

Kịch bản mô phỏng là dữ liệu máy tính tạo theo quy luật đã biết để thử phương pháp. “Hầu hết” nghĩa là không phải mọi kịch bản. Trong đề cương, đây là đối chứng của cơ chế (i).

\section{Cách thứ hai: thu hẹp tập mô hình bằng tập tin cậy mô hình}

Tập tin cậy mô hình của Hansen và cộng sự (2011) giữ lại những mô hình chưa thể bị loại là kém hơn ở một mức tin cậy cho trước.

Vòng sơ loại một cuộc thi hát: ban giám khảo chỉ loại thí sinh khi chắc chắn kém hơn; ai chưa chắc thì đi tiếp. Danh sách đi tiếp là tập tin cậy mô hình.

Tính chất then chốt: dữ liệu kém thông tin cho tập lớn, dữ liệu nhiều thông tin cho tập nhỏ. Kích thước tập phản ánh phần nào chất lượng dữ liệu.

Điều này khớp với việc đề cương xếp kích thước tập tin cậy vào chỉ báo bất định (mục 4.1), và với cảnh báo ở mục 4.3.6: tập lớn hơn ở chế độ căng thẳng có thể chỉ do dữ liệu kém thông tin, nên kích thước tập được báo cáo như thước đo mô tả, kèm số quan sát và công suất kiểm định.

\begin{tcolorbox}[breakable]
Ví dụ giả định (số tự đặt để minh họa): bảy mô hình; với một năm dữ liệu ít thông tin, không mô hình nào bị loại, tập gồm cả bảy. Với mười năm dữ liệu rõ hơn, bốn mô hình bị loại, tập còn ba.

Tập co từ bảy xuống ba không có nghĩa thị trường bớt rủi ro; nó chỉ cho biết dữ liệu đủ thông tin để phân biệt các mô hình.
\end{tcolorbox}

Đề cương dùng tập tin cậy theo hai cách: kích thước tập ở mức ý nghĩa 10\% ấn định trước là kết quả phụ của H1; và ở cơ chế (i), có thể chỉ kết hợp các mô hình trong tập.

\section{Cách thứ ba: hiệu chỉnh thiên lệch chọn lọc}

\subsection{Thiên lệch chọn lọc và quá khớp khi kiểm định ngược}

Thử rất nhiều chiến lược trên dữ liệu quá khứ rồi báo cáo chiến lược đẹp nhất: nó đẹp một phần vì tốt thật, một phần vì thử nhiều lần thì thế nào cũng có chiến lược gặp may. Phần may đó không lặp lại.

Thiên lệch chọn lọc là sự méo mó khi chỉ nhìn ứng viên thắng sau nhiều lần thử. Quá khớp là khi mô hình bám quá sát quá khứ đến mức nắm cả nhiễu, nên đẹp trên dữ liệu đã dùng và kém trên dữ liệu mới.

\subsection{Bailey và López de Prado (2014): tỷ số Sharpe hiệu chỉnh}

Bailey và López de Prado (2014) định lượng thiên lệch chọn lọc và quá khớp khi kiểm định ngược bằng tỷ số Sharpe hiệu chỉnh.

Tỷ số Sharpe là mức lời vượt lãi suất an toàn trên mỗi đơn vị dao động; cao nghĩa là mỗi đơn vị rủi ro đem lại nhiều lời hơn.

Tỷ số Sharpe hiệu chỉnh trừ bớt phần có thể do may mắn nhờ thử nhiều lần: học sinh thi nhiều đề và chỉ khoe điểm cao nhất thì điểm khoe đó cần “giảm giá”. Công thức đề cương không nêu.

Đề tài dùng công cụ này ở mục 4.3.7, và khả năng đối lập của H3 nói: kết hợp mô hình có thể không vượt việc chọn một mô hình sau chi phí giao dịch và sau hiệu chỉnh thiên lệch chọn lọc.

\subsection{Ledoit và Wolf (2008): khoảng tin cậy bootstrap cho chênh lệch Sharpe}

Ledoit và Wolf (2008) đề xuất khoảng tin cậy bootstrap chuỗi thời gian cho chênh lệch tỷ số Sharpe.

Mô hình A có Sharpe cao hơn B; chênh lệch đó có thật hay do ngẫu nhiên? Khoảng tin cậy là dải giá trị hợp lý cho độ chênh thật; nếu dải chứa số 0 thì chưa thể nói A hơn B.

\begin{tcolorbox}[breakable]
Ví dụ giả định (số tự đặt để minh họa): A có Sharpe 0,6, B có 0,5, chênh 0,1. Khoảng tin cậy từ −0,05 đến 0,25 chứa số 0, nên chưa kết luận được A hơn B.
\end{tcolorbox}

Bootstrap là bốc lại nhiều lần từ chính túi dữ liệu để xem kết quả dao động bao nhiêu. Với chuỗi thời gian, hôm nay liên hệ với hôm qua, nên bootstrap bốc từng nắm ngày liền nhau thay vì từng ngày rời; cách này gọi là bootstrap khối.

Mục 4.3.7 dùng đúng ý này để so sánh lợi suất tương đương chắc chắn và tỷ số Sharpe giữa các mô hình.

\section{Ba cách đặt cạnh nhau}

{\small\begin{longtable}[]{@{}
  >{\raggedright\arraybackslash}p{(\columnwidth - 4\tabcolsep) * \real{0.1655}}
  >{\raggedright\arraybackslash}p{(\columnwidth - 4\tabcolsep) * \real{0.4569}}
  >{\raggedright\arraybackslash}p{(\columnwidth - 4\tabcolsep) * \real{0.3777}}@{}}
\toprule\noalign{}
\begin{minipage}[b]{\linewidth}\raggedright
Cách xử lý
\end{minipage} & \begin{minipage}[b]{\linewidth}\raggedright
Công trình
\end{minipage} & \begin{minipage}[b]{\linewidth}\raggedright
Ý chính
\end{minipage} \\
\midrule\noalign{}
\endhead
\bottomrule\noalign{}
\endlastfoot
Kết hợp mô hình & Hoeting và cộng sự (1999), Avramov (2002), Timmermann (2006), Rapach và cộng sự (2010), Tu và Zhou (2011) & Trộn nhiều mô hình thay vì chọn một \\
Giới hạn tập mô hình & Hansen và cộng sự (2011) & Giữ các mô hình chưa thể loại ở mức tin cậy cho trước \\
Hiệu chỉnh thiên lệch chọn lọc & Bailey và López de Prado (2014), Ledoit và Wolf (2008) & Trừ phần may mắn do thử nhiều lần; khoảng tin cậy cho chênh lệch Sharpe \\
\end{longtable}}

\section{Nhóm này nối vào phần còn lại của đề cương thế nào}

Ba mối nối: trộn mô hình là nền của H3 và cơ chế (i); tập tin cậy mô hình là thước đo thứ ba ở mục 4.3.6; tỷ số Sharpe hiệu chỉnh và bootstrap là công cụ kiểm định ở mục 4.3.7.

Đề cương không khẳng định trộn mô hình sẽ thắng ở Việt Nam; đó là giả thuyết, có khả năng đối lập đi kèm. Nếu số cửa sổ độc lập ở chế độ căng thẳng dưới 10, kết quả gọi là không kết luận (mục 4.3.10). Câu chốt của nhóm (3): các công cụ này có sẵn, song chưa công trình nào dùng chúng để đo tổn thất của việc chọn mô hình theo từng trạng thái thị trường.

\begin{tcolorbox}[breakable]
\textbf{Ghi nhớ khi đi thi}

• Ba cách xử lý bất định mô hình: kết hợp mô hình, giới hạn tập mô hình, hiệu chỉnh thiên lệch chọn lọc.

• Hoeting và cộng sự (1999): trung bình hóa mô hình Bayes; Avramov (2002): áp dụng cho lợi suất cổ phiếu, tốt hơn chọn một mô hình; Timmermann (2006): dự báo kết hợp thường tốt hơn; Rapach và cộng sự (2010): kết hợp bù phần mất do bất ổn định.

• Hansen và cộng sự (2011): tập tin cậy mô hình; dữ liệu kém thông tin cho tập lớn, nhiều thông tin cho tập nhỏ.

• Bailey và López de Prado (2014): tỷ số Sharpe hiệu chỉnh; Ledoit và Wolf (2008): khoảng tin cậy bootstrap cho chênh lệch Sharpe.

• Tu và Zhou (2011): kết hợp 1/N với trung bình-phương sai, vượt 1/N trong hầu hết kịch bản mô phỏng.
\end{tcolorbox}

\section{Tự kiểm tra}

{\itshape 1. Ba cách xử lý bất định mô hình trong nhóm (3) là gì?\par}

Đáp án: kết hợp mô hình, thu hẹp tập mô hình, hiệu chỉnh thiên lệch chọn lọc.

{\itshape 2. Tập tin cậy mô hình là gì, và co giãn thế nào theo thông tin của dữ liệu?\par}

Đáp án: tập các mô hình chưa thể loại là kém hơn ở mức tin cậy cho trước; dữ liệu kém thông tin cho tập lớn, nhiều thông tin cho tập nhỏ.

{\itshape 3. Avramov (2002) thấy gì?\par}

Đáp án: trung bình hóa mô hình Bayes cho dự báo lợi suất cổ phiếu hiệu quả ngoài mẫu tốt hơn dùng tiêu chí chọn một mô hình.

{\itshape 4. Tỷ số Sharpe hiệu chỉnh dùng để làm gì?\par}

Đáp án: định lượng thiên lệch chọn lọc và quá khớp khi kiểm định ngược (Bailey và López de Prado, 2014).

{\itshape 5. Tu và Zhou (2011) kết hợp những quy tắc nào và thấy gì?\par}

Đáp án: 1/N với trung bình-phương sai; vượt 1/N trong hầu hết kịch bản mô phỏng.

\chapter{Mục 4.2 của đề cương, nhóm (4): Chế độ thị trường và khủng hoảng}

Thị trường không hành xử giống nhau mọi lúc; hành vi đổi theo chế độ thì rủi ro chọn mô hình cũng có thể đổi theo. Đó là cơ sở để đề tài so sánh chế độ bình thường với chế độ căng thẳng. Nhóm (4) điểm các công trình về chuyển chế độ (Hamilton, Ang, Guidolin, Kritzman, Tu, Bai và Perron) rồi về khủng hoảng (Longin và Solnik, Forbes và Rigobon).

Hai tầng ý: thị trường có những trạng thái khác nhau và có công cụ để ước lượng trạng thái; lúc khủng hoảng, các con số quen thuộc như tương quan có thể đổi, và phải giải thích sự đổi đó cẩn thận.

\section{Chế độ thị trường là gì}

Chế độ thị trường là thời tiết của thị trường: bình thường là trời quang, giá dao động vừa phải; căng thẳng là mùa bão, dao động mạnh, lợi suất thấp. Bạn không thấy trực tiếp “tên” thời tiết, chỉ thấy dấu hiệu như mưa gió.

Đề cương quy ước (mục 4.3.4): chế độ căng thẳng có phương sai điều kiện cao nhất và lợi suất kỳ vọng thấp hơn chế độ bình thường. Phương sai điều kiện là mức dao động trong từng chế độ, không phải dao động chung cả mẫu.

\section{Hamilton (1989): mô hình chuyển chế độ Markov}

Hamilton (1989) đề xuất mô hình chuyển chế độ Markov: thị trường ở một trong các trạng thái ẩn và chuyển giữa chúng theo xác suất. Ang và Timmermann (2012) tổng quan hướng này.

Trạng thái ẩn nghĩa là bạn không thấy trực tiếp thị trường đang ở đâu mà phải suy ra từ giá, như bác sĩ không thấy bệnh, chỉ thấy triệu chứng. Xác suất chuyển là xác suất mỗi ngày thị trường ở nguyên chế độ hay nhảy sang chế độ khác.

“Markov” là giả định về trí nhớ: xác suất chuyển ngày mai chỉ phụ thuộc chế độ hôm nay, không phụ thuộc cả lịch sử trước đó (nghĩa chung của thuật ngữ).

\begin{tcolorbox}[breakable]
Ví dụ giả định (số tự đặt để minh họa): hai chế độ quang và bão. Hôm nay quang thì ngày mai vẫn quang với xác suất 98\%, sang bão 2\%. Hôm nay bão thì ngày mai vẫn bão 90\%, trở lại quang 10\%.

Vì xác suất nghiêng về giữ nguyên, bão thường kéo dài nhiều ngày. Mô hình ước lượng các xác suất này và mức dao động trong từng chế độ từ dữ liệu giá.
\end{tcolorbox}

Đề tài dùng đúng khung này (mục 4.3.4): ước lượng mô hình chuyển chế độ Markov cho lợi suất VN-Index theo cửa sổ mở rộng; tại mỗi ngày tái cân bằng, tham số được ước lượng lại chỉ bằng dữ liệu đến ngày đó.

Từ đó tính xác suất lọc tại ngày t, tức xác suất thị trường đang ở chế độ căng thẳng; cách làm này loại thiên lệch nhìn trước, tức vô tình dùng thông tin tương lai để quyết định hôm nay. Không được chấm bài bằng đáp án chưa công bố.

\section{Ang và Bekaert (2002): tương quan và biến động tăng ở chế độ xấu}

Ang và Bekaert (2002) thấy tương quan và biến động cùng tăng ở chế độ xấu và giải bài toán phân bổ tài sản có chuyển chế độ.

Ở chế độ xấu, giá các tài sản đi cùng chiều hơn và dao động mạnh hơn. Trời quang thì mỗi người một hướng; bão đến thì mọi người cùng chạy vào một chỗ trú, và đa dạng hóa giúp ít hơn.

Họ không chỉ mô tả hiện tượng mà còn hỏi nên chia tiền thế nào khi thị trường có thể đổi chế độ; đề cương không nêu kết quả của bài toán đó.

Theo bản đọc, tương quan tăng ở chế độ xấu cùng chủ đề với biến kênh tương quan đuôi trong H2. Nhưng cơ sở của H2 trong Bảng 2 là Amihud (2002), Longin và Solnik (2001), Forbes và Rigobon (2002), không gồm Ang và Bekaert.

\section{Guidolin và Timmermann (2007): bốn chế độ}

Guidolin và Timmermann (2007) nhận diện bốn chế độ: sụp đổ, tăng trưởng chậm, tăng giá và phục hồi.

Thị trường không nhất thiết chỉ có hai trạng thái. Đề cương chỉ nêu tên bốn chế độ, không nêu đặc điểm thống kê; đừng tự mô tả thêm.

Ở mục 4.3.4, số chế độ (hai hoặc ba) do tiêu chí thông tin quyết định trên giai đoạn ước lượng ban đầu và cố định khi đăng ký trước; việc nhận diện chế độ cũng được coi là một dạng rủi ro mô hình (mục 4.3.10).

\section{Kritzman và cộng sự (2012): dự báo chế độ và chiến lược động}

Kritzman và cộng sự (2012) dự báo chế độ bằng mô hình Markov, và chiến lược động của họ vượt chiến lược tĩnh khi kiểm định ngược.

Chiến lược tĩnh chia tiền theo quy tắc cố định; chiến lược động đổi cách chia khi dự báo chế độ đổi, ví dụ giảm rủi ro khi dự báo sắp vào bão.

Chú ý chữ “khi kiểm định ngược”: kết quả loại này nói chung có thể chịu nguy cơ quá khớp (nhóm (3)). Đề cương không nói điều đó về Kritzman và cộng sự.

M7 trong Bảng 3 là giảm tỷ trọng rủi ro theo xác suất chế độ căng thẳng (Kritzman và cộng sự, 2012), và cơ chế giảm thiểu (iii) cũng theo quy tắc này.

\section{Tu (2010): bỏ qua chuyển chế độ thì thiệt}

Tu (2010) ước tính nhà đầu tư bỏ qua chuyển chế độ chịu tổn thất lợi suất tương đương chắc chắn thường trên 2\% mỗi năm, có thể tới 10\%.

Lợi suất tương đương chắc chắn là mức lợi suất chắc chắn mà nhà đầu tư coi là ngang giá với khoản đầu tư có rủi ro; càng cao càng tốt. Số 2\% và 10\% là của nghiên cứu của Tu, không phải kết quả của đề tài, không dùng để suy ra gì về Việt Nam; đề cương không nêu dữ liệu đứng sau chúng.

Tu (2010) là bằng chứng rằng tính đến chuyển chế độ có giá trị, và là đối chứng của cơ chế (iii).

\textbf{Bai và Perron (1998): nhiều điểm gãy cấu trúc}

Bai và Perron (1998) cung cấp phương pháp kiểm định nhiều điểm gãy cấu trúc. Điểm gãy là thời điểm quan hệ trong dữ liệu đổi lâu dài, như dòng sông đổi lòng sau một trận lũ. Theo cách phân biệt thông thường: với chuyển chế độ, thị trường nhảy qua lại và có thể quay về; với điểm gãy, quan hệ chuyển sang dạng mới.

Vai trò trong đề tài hạn chế (mục 4.3.4): kiểm định điểm gãy trên toàn mẫu chỉ để mô tả và đối chiếu với các cú sốc lịch sử, không dùng làm biến chế độ trong kiểm định giả thuyết, vì biến chế độ chỉ được dùng dữ liệu đến ngày đó.

Tác giả đã dùng phương pháp Bai-Perron cùng độ phi thanh khoản Amihud trong bản thảo về điểm gãy cấu trúc của thanh khoản (bản thảo số 5, Bảng 7).

\section{Về khủng hoảng: Longin và Solnik (2001)}

Longin và Solnik (2001) bác bỏ phân phối chuẩn đa biến ở đuôi âm của lợi suất.

Phân phối chuẩn là đường hình chuông: phần lớn giá trị gần trung bình, giá trị rất xa hiếm đến mức gần như không xảy ra. “Đa biến” là mô tả nhiều tài sản cùng lúc, kể cả mức chúng đi cùng nhau. “Đuôi âm” là những ngày rất tệ.

Theo bản đọc: mô hình hình chuông cho nhiều tài sản không mô tả đúng những ngày rất tệ. Đề cương không nói họ bác bỏ bằng cách nào.

Tương quan đuôi là mức các tài sản cùng giảm mạnh vào những ngày cực xấu: tương quan thường đo mức cùng nhịp ngày thường, tương quan đuôi chỉ hỏi chuyện gì xảy ra khi ai cũng rơi. Đây là một biến kênh của H2.

\section{Forbes và Rigobon (2002): tương quan phụ thuộc vào biến động, và lây nhiễm}

Forbes và Rigobon (2002) chỉ ra hệ số tương quan phụ thuộc vào biến động và, sau hiệu chỉnh, gần như không còn dấu hiệu lây nhiễm trong ba cuộc khủng hoảng họ khảo sát.

Lây nhiễm là khi cú sốc ở thị trường này kéo thị trường khác theo nhiều hơn mức liên hệ bình thường. Hai người bạn thường gọi nhau mỗi ngày; một người gặp nạn thì người kia gọi nhiều hơn. Nếu số cuộc gọi tăng chỉ vì ai cũng đang hoảng, đó chưa chắc là quan hệ đã đổi.

Khi biến động tăng, hệ số tương quan tính được cũng có thể tăng dù quan hệ cơ bản chưa đổi; vì vậy phải hiệu chỉnh theo biến động trước khi kết luận. Đây chính là khả năng đối lập của H2.

Câu này không nói lây nhiễm không bao giờ tồn tại; đề cương cũng không nêu tên ba cuộc khủng hoảng.

\begin{tcolorbox}[breakable]
Ví dụ giả định (số tự đặt để minh họa): hai thị trường có quan hệ cố định. Tuần yên ắng, tương quan tính được là 0,3; tuần hỗn loạn, giá dao động mạnh gấp ba và tương quan tính được lên 0,6.

Chỉ nhìn 0,3 lên 0,6 thì dễ kết luận quan hệ đã mạnh lên; hiệu chỉnh theo biến động cho phép hỏi mức tăng đó có chỉ do giá dao động mạnh hơn không.
\end{tcolorbox}

Forbes và Rigobon (2002) là cơ sở của H2, và tương quan hiệu chỉnh theo biến động là một biến kênh. Bản thảo số 6 (Bảng 7) của tác giả, về tương quan các chỉ số lồng nhau, chuẩn bị cho kênh này.

\section{Hai nửa ghép lại}

Khủng hoảng là sự kiện lịch sử, chế độ căng thẳng là trạng thái thống kê (mục 4.1). Theo bản đọc, nhóm chuyển chế độ cho công cụ nhận diện trạng thái; nhóm khủng hoảng cho biết lúc căng thẳng, các thước đo phụ thuộc như tương quan phải đọc cẩn thận. Câu chốt của nhóm (4): hành vi thị trường đổi theo chế độ, nhưng chưa nghiên cứu nào kiểm định thứ hạng của các mô hình danh mục có đổi theo hay không.

Đề cương cũng tính đến việc chế độ căng thẳng hiếm: mẫu có ít giai đoạn căng thẳng độc lập nên kiểm định có công suất thấp, và khi số cửa sổ độc lập ở chế độ căng thẳng dưới 10 thì kết quả gọi là không kết luận (mục 4.3.10). Ít cơn bão thì khó rút ra quy luật chắc chắn.

\begin{tcolorbox}[breakable]
\textbf{Ghi nhớ khi đi thi}

• Hành vi thị trường đổi theo chế độ nên rủi ro chọn mô hình có thể đổi theo.

• Hamilton (1989): chuyển chế độ Markov, trạng thái ẩn, xác suất chuyển; Ang và Timmermann (2012) tổng quan.

• Ang và Bekaert (2002): tương quan và biến động cùng tăng ở chế độ xấu; Guidolin và Timmermann (2007): sụp đổ, tăng trưởng chậm, tăng giá, phục hồi.

• Kritzman và cộng sự (2012): chiến lược động vượt tĩnh khi kiểm định ngược; Tu (2010): bỏ qua chuyển chế độ mất lợi suất tương đương chắc chắn thường trên 2\% mỗi năm, có thể tới 10\%; Bai và Perron (1998): nhiều điểm gãy cấu trúc.

• Longin và Solnik (2001): bác bỏ phân phối chuẩn đa biến ở đuôi âm; Forbes và Rigobon (2002): tương quan phụ thuộc biến động, sau hiệu chỉnh gần như không còn lây nhiễm trong ba cuộc khủng hoảng.
\end{tcolorbox}

\section{Tự kiểm tra}

{\itshape 1. Trong mô hình Hamilton (1989), thị trường được mô tả thế nào?\par}

Đáp án: thị trường ở một trong các trạng thái ẩn và chuyển giữa chúng theo xác suất.

{\itshape 2. Guidolin và Timmermann (2007) nhận diện những chế độ nào?\par}

Đáp án: sụp đổ, tăng trưởng chậm, tăng giá và phục hồi.

{\itshape 3. Tu (2010) thấy tổn thất bao nhiêu khi bỏ qua chuyển chế độ?\par}

Đáp án: tổn thất lợi suất tương đương chắc chắn thường trên 2\% mỗi năm, có thể tới 10\%.

{\itshape 4. Forbes và Rigobon (2002) thấy gì về hệ số tương quan và lây nhiễm?\par}

Đáp án: tương quan phụ thuộc biến động; sau hiệu chỉnh, gần như không còn lây nhiễm trong ba cuộc khủng hoảng họ khảo sát.

{\itshape 5. Longin và Solnik (2001) bác bỏ điều gì?\par}

Đáp án: phân phối chuẩn đa biến ở đuôi âm của lợi suất.

\chapter{Mục 4.2 của đề cương, nhóm (5): Tài chính hành vi và bầy đàn}

Bầy đàn là một kênh hành vi có thể làm các mô hình cho kết quả khác nhau, nhưng bằng chứng trái chiều giữa các thị trường. Nhóm (5) điểm bốn nguồn: Kahneman và Tversky (1979), Barberis và Thaler (2003), Christie và Huang (1995), Chang và cộng sự (2000). Cơ sở của giả thuyết thăm dò H4 là Vo và Phan (2017), Chang và cộng sự (2000), Christie và Huang (1995).

Giới hạn cần nhớ ngay (mục 4.4): đề tài không quan sát hành vi chọn mô hình của tổ chức đầu tư; nó chỉ kiểm tra độ phân tán chéo của lợi suất, một thước đo bầy đàn ở cấp thị trường, có đi cùng rủi ro lựa chọn mô hình sau khi kiểm soát biến động hay không.

\section{Bầy đàn là gì}

Bầy đàn là việc nhà đầu tư làm theo số đông thay vì dựa vào thông tin riêng.

Đàn cừu chạy theo con đầu đàn thay vì tự nhìn đường; mỗi con có thể thấy hố nước trước mặt nhưng vẫn chạy theo. Trên thị trường, nhà đầu tư thấy nhiều người mua một cổ phiếu thì mua theo, dù thông tin riêng nói khác.

Đề cương nói bầy đàn “có thể” làm các mô hình cho kết quả khác nhau, không khẳng định, và không giải thích cơ chế từng bước. H4 chỉ dự kiến độ phân tán chéo cao đi cùng rủi ro lựa chọn mô hình cao; chiều liên hệ này cũng cần kiểm định.

\section{Kahneman và Tversky (1979): nền móng của lý thuyết triển vọng}

Kahneman và Tversky (1979) đặt nền móng cho lý thuyết triển vọng.

Lý thuyết triển vọng mô tả cách con người thật sự quyết định khi đối mặt rủi ro, trái với hình dung nhà đầu tư là người tính toán lạnh lùng và nhất quán (nghĩa chung; đề cương không mô tả nội dung bên trong).

Được tặng một khoản tiền rồi mất đúng khoản đó: về số học là hòa, nhưng cảm giác không đối xứng. Lý thuyết triển vọng nghiên cứu những khác biệt kiểu đó giữa tính toán và cảm nhận.

Đề cương không nêu lý do riêng cho việc nhắc công trình này. Đề tài chỉ giao thoa kinh tế hành vi ở một phép thử giới hạn (mục 4.4). Việc chương trình Tiến sĩ Kinh tế và Quản lý tập trung vào kinh tế hành vi là lý do chọn cơ sở đào tạo (mục 5), chuyện khác; đừng nối nó với lý thuyết triển vọng.

\section{Barberis và Thaler (2003): hai trụ cột}

Barberis và Thaler (2003) tổng hợp hai trụ cột của tài chính hành vi: giới hạn của kinh doanh chênh lệch giá và tâm lý nhà đầu tư.

Tài chính hành vi là hướng nghiên cứu cho rằng cảm xúc, thói quen suy nghĩ và những giới hạn thực tế ảnh hưởng đến giá.

Kinh doanh chênh lệch giá là mua rẻ chỗ này, bán đắt chỗ kia; về lý thuyết, nó kéo giá về mức hợp lý. “Giới hạn” nghĩa là trong thực tế người làm việc này không phải lúc nào cũng làm được, nên giá có thể lệch lâu hơn.

Tâm lý, tức sai lệch trong tư duy, là nguồn làm giá lệch. Giá cổ phiếu như chiếc thuyền trên sóng: tâm lý là sóng đẩy thuyền lệch, kinh doanh chênh lệch giá là người chèo kéo thuyền về; người chèo bị giới hạn thì thuyền lệch lâu hơn.

Đề cương dùng công trình này như một tổng hợp về lĩnh vực, không nêu thêm lý do.

\section{Christie và Huang (1995): đo bầy đàn bằng độ phân tán chéo}

Christie và Huang (1995) đo bầy đàn bằng độ phân tán chéo của lợi suất và không thấy bầy đàn khi giá biến động mạnh trên dữ liệu Mỹ.

Công trình này cho cách đo mà đề tài dùng. Độ phân tán chéo của lợi suất là mức các cổ phiếu khác nhau về lợi suất trong cùng một ngày: mã nào cũng tăng gần như nhau thì nhỏ; có mã tăng mạnh, mã giảm mạnh thì lớn.

Vì sao nó đo bầy đàn? Giải thích chung của bản đọc, không phải câu của đề cương: nếu nhà đầu tư chạy theo đám đông, họ mua bán giống nhau, các cổ phiếu đi gần nhau hơn và độ phân tán chéo nhỏ lại. Đề cương không giải thích H4 đọc theo chiều bầy đàn ra sao.

\begin{tcolorbox}[breakable]
Ví dụ giả định (số tự đặt để minh họa): năm cổ phiếu. Ngày thứ nhất lợi suất là 1\%; 1,2\%; 0,9\%; 1,1\%; 1\%. Ngày thứ hai là 4\%; −3\%; 2\%; −5\%; 0\%.

Ngày thứ nhất các mã gần nhau, độ phân tán chéo nhỏ; ngày thứ hai chúng tản ra, độ phân tán lớn.
\end{tcolorbox}

Đề tài dùng độ phân tán chéo làm biến kênh với cơ sở Christie và Huang (1995); kết quả của họ trên dữ liệu Mỹ thành khả năng đối lập của H4: bầy đàn có thể không xuất hiện ở tần suất dữ liệu của đề tài.

\section{Chang và cộng sự (2000): bằng chứng khác nhau theo nước}

Chang và cộng sự (2000) tìm thấy bầy đàn ở Hàn Quốc và Đài Loan nhưng không thấy ở Mỹ và Hồng Kông.

Bốn thị trường, hai kết quả; từ đó đề cương nói bằng chứng trái chiều giữa các thị trường.

Kiểm tra học sinh có chép bài nhau không: hai lớp có dấu hiệu rõ, hai lớp không. Bạn chỉ nói được bằng chứng khác nhau theo lớp.

Đề cương không nêu nguyên nhân của khác biệt hay phương pháp chi tiết; đừng tự giải thích.

Câu chốt của nhóm (5) nêu lý do H4 chỉ ở mức thăm dò: vì bằng chứng trái chiều giữa các thị trường, đề tài chỉ xét kênh này ở mức thăm dò. Không thể mặc định Việt Nam có bầy đàn ở tần suất dữ liệu đã chọn.

\section{H4 trong bức tranh chung}

H4 (Bảng 2): sau khi kiểm soát biến động thực hiện, độ phân tán chéo của lợi suất cao đi cùng rủi ro lựa chọn mô hình cao.

Phân biệt: bầy đàn trong đề tài là thước đo cấp thị trường, không phải quan sát hành vi từng tổ chức; và giao thoa kinh tế hành vi chỉ ở một phép thử giới hạn.

Nếu được hỏi vì sao H4 chỉ thăm dò: bằng chứng trái chiều giữa các thị trường; đề tài không quan sát hành vi chọn mô hình của tổ chức; khả năng đối lập là bầy đàn không xuất hiện ở tần suất dữ liệu của đề tài.

Mọi kết quả là liên hệ có điều kiện, không phải nhân quả (mục 4.3.7). Ngay cả khi H4 có kết quả, không được nói bầy đàn gây ra rủi ro lựa chọn mô hình.

\section{Điều đề cương không nói}

Đề cương không nói Việt Nam có bầy đàn ở tần suất dữ liệu đã chọn; chỉ nói Vo và Phan (2017) thấy bằng chứng bầy đàn trên một mẫu HOSE (nhóm (6)). Cũng không nói tổ chức đầu tư thực tế bầy đàn khi chọn mô hình.

\begin{tcolorbox}[breakable]
\textbf{Ghi nhớ khi đi thi}

• Bầy đàn: làm theo số đông thay vì thông tin riêng; bằng chứng trái chiều giữa các thị trường.

• Kahneman và Tversky (1979): nền móng lý thuyết triển vọng; Barberis và Thaler (2003): giới hạn của kinh doanh chênh lệch giá và tâm lý nhà đầu tư.

• Christie và Huang (1995): đo bầy đàn bằng độ phân tán chéo của lợi suất, không thấy bầy đàn khi giá biến động mạnh trên dữ liệu Mỹ.

• Chang và cộng sự (2000): có ở Hàn Quốc và Đài Loan, không ở Mỹ và Hồng Kông.

• H4 chỉ là thăm dò; đề tài không quan sát hành vi chọn mô hình của tổ chức đầu tư.
\end{tcolorbox}

\section{Tự kiểm tra}

{\itshape 1. Bầy đàn được định nghĩa thế nào trong đề cương?\par}

Đáp án: nhà đầu tư làm theo số đông thay vì dựa vào thông tin riêng.

{\itshape 2. Christie và Huang (1995) đo bầy đàn bằng gì, và thấy gì?\par}

Đáp án: độ phân tán chéo của lợi suất; không thấy bầy đàn khi giá biến động mạnh trên dữ liệu Mỹ.

{\itshape 3. Chang và cộng sự (2000) thấy bầy đàn ở đâu và không thấy ở đâu?\par}

Đáp án: có ở Hàn Quốc và Đài Loan, không ở Mỹ và Hồng Kông.

{\itshape 4. Barberis và Thaler (2003) tổng hợp hai trụ cột nào?\par}

Đáp án: giới hạn của kinh doanh chênh lệch giá và tâm lý nhà đầu tư.

{\itshape 5. Vì sao H4 chỉ là thăm dò?\par}

Đáp án: bằng chứng bầy đàn trái chiều giữa các thị trường; đề tài không quan sát hành vi chọn mô hình của tổ chức; khả năng đối lập là bầy đàn không xuất hiện ở tần suất dữ liệu của đề tài.

\chapter{Mục 4.2 của đề cương, nhóm (6): Nghiên cứu về Việt Nam}

Các công trình về Việt Nam xét riêng một mô hình, một phương pháp hoặc một hiện tượng. Đó là câu mở đầu nhóm (6) và là điểm tựa của khoảng trống. Nhóm này nêu sáu công trình: Vo và Phan (2017), Ta và Vuong (2020), Nguyen và cộng sự (2020), Goutte và cộng sự (2025), Hoang và Luu (2024), và bài của chính tác giả, Nguyen và Nguyen (2026a).

Nhóm này chỉ mô tả điều từng công trình làm. Tác giả mới đọc tóm tắt và thông tin thư mục của phần lớn nguồn, việc đọc toàn văn nằm trong kế hoạch năm thứ nhất; khi nói về từng công trình, hãy giữ đúng mức mô tả của đề cương.

\section{Vo và Phan (2017): bầy đàn trên HOSE}

Vo và Phan (2017) dùng 299 công ty niêm yết trên Sở Giao dịch Chứng khoán Thành phố Hồ Chí Minh (HOSE) giai đoạn 2005--2015 và tìm thấy bầy đàn, rõ hơn khi thị trường giảm.

Leo núi khi trời đẹp, mỗi người tự chọn lối; trời đổ mưa, mọi người dồn theo người dẫn đầu. Ví von chỉ để nhớ: đề cương không nói cơ chế hay nguyên nhân.

Đề cương không nói họ dùng phương pháp đo nào; đừng khẳng định là độ phân tán chéo hay cách khác.

Vo và Phan được nhắc ba lần: là một trong ba cơ sở của H4; là nhóm “bầy đàn ở Việt Nam” trong Bảng 1 (khoảng trống (4): chưa gắn bầy đàn với rủi ro của việc chọn mô hình); và ở mục 4.1, là công trình xét tác động của khủng hoảng tài chính toàn cầu lên hành vi nhà đầu tư trên dữ liệu 2005--2015.

\section{Ta và Vuong (2020): Black-Litterman với quan điểm từ mô hình dự báo}

Ta và Vuong (2020) áp dụng mô hình Black-Litterman với quan điểm dự báo từ mô hình tự hồi quy trung bình trượt tích hợp.

Black-Litterman kết hợp cân bằng thị trường với quan điểm của nhà đầu tư (nhóm (1)). Câu hỏi ở đây là quan điểm lấy từ đâu: Ta và Vuong lấy từ một mô hình dự báo chuỗi thời gian.

Giải thích chung: “tự hồi quy” là dự báo dựa vào giá trị quá khứ của chính chuỗi; “trung bình trượt” là dựa vào các sai số dự báo gần đây; “tích hợp” là xử lý để chuỗi bớt trôi dạt. Giống dự báo thời tiết ngày mai bằng cách nhìn mấy ngày qua và xem các lần đoán trước lệch bao nhiêu.

\begin{tcolorbox}[breakable]
Ví dụ giả định (số tự đặt để minh họa): mô hình dự báo nói cổ phiếu ngân hàng tháng tới có lợi suất cao hơn mức cân bằng một chút; Black-Litterman không bỏ cân bằng thị trường mà dịch lợi suất kỳ vọng của cổ phiếu ngân hàng lên một ít.

Quan điểm ở đây là sản phẩm của mô hình dự báo, không phải ý kiến chủ quan.
\end{tcolorbox}

Đề cương không nêu kết quả của Ta và Vuong. Công trình là một nguồn của khoảng trống (2): áp dụng một mô hình, còn đề tài đặt Black-Litterman cùng sáu mô hình khác vào tập cạnh tranh. Ở M4, quan điểm lấy từ “một mô hình dự báo đơn giản” (mục 4.3.5); đề cương không nói đó là mô hình nào.

\section{Nguyen và cộng sự (2020): ma trận hiệp phương sai co rút}

Nguyen và cộng sự (2020) dùng ma trận hiệp phương sai co rút trên 468 tuần giai đoạn 2011--2019 và thu được danh mục có lợi nhuận cao hơn, rủi ro thấp hơn các phương pháp truyền thống.

Co rút là kéo các ước lượng hiệp phương sai quá nhiễu về mức ổn định hơn (nhóm (1)); công trình này áp dụng ý đó cho Việt Nam.

Đề cương không nêu “phương pháp truyền thống” cụ thể nào; đừng tự điền.

Công trình thuộc nhóm “danh mục ở Việt Nam” trong Bảng 1: đánh giá từng mô hình, chưa đặt vào tập cạnh tranh. M2 dùng hiệp phương sai co rút theo Ledoit và Wolf (2004).

\section{Goutte và cộng sự (2025): ra quyết định đa tiêu chí trên VN100}

Goutte và cộng sự (2025) so sánh năm phương pháp ra quyết định đa tiêu chí với bốn cách gán trọng số trên rổ VN100, giai đoạn 01/2015--11/2023.

Ra quyết định đa tiêu chí là chấm điểm các lựa chọn theo nhiều tiêu chí cùng lúc; cách gán trọng số quyết định tiêu chí nào quan trọng hơn. Tuyển người bằng cách chấm kinh nghiệm, học vấn, thái độ rồi quy ước kinh nghiệm nặng gấp đôi: mỗi quy ước cho một bảng xếp hạng khác. Đề cương không nêu họ dùng tiêu chí nào.

VN100 là rổ 100 cổ phiếu do HOSE công bố. Đề cương không nói năm phương pháp và bốn cách gán trọng số được kết hợp ra sao; đừng tự suy ra số trường hợp.

Đề cương không nêu kết quả. Khác biệt của đề tài: Goutte và cộng sự so sánh phương pháp, đề tài đo mức bất định khi chọn mô hình theo từng chế độ.

\section{Hoang và Luu (2024): một chế độ hay chuyển chế độ khi dự báo biến động}

Hoang và Luu (2024), trong một bản thảo công bố trực tuyến, phân tích VN-Index giai đoạn 2012--2023 và thấy mô hình một chế độ dự báo biến động tốt hơn ở kỳ hạn ngắn, mô hình chuyển chế độ tốt hơn ở kỳ hạn trung và dài.

VN-Index là chỉ số giá chung của HOSE. Mô hình một chế độ giả định thị trường luôn hành xử một kiểu; mô hình chuyển chế độ cho phép đổi trạng thái (nhóm (4)).

Dự báo ngày mai thì nhìn hôm nay là đủ; dự báo cả mùa thì cần biết hệ thời tiết có thể đổi pha. Ví von chỉ để nhớ, đề cương không giải thích kết quả theo hướng đó.

Hai lưu ý: đây là bản thảo, khác bài đã đăng tạp chí, nên gọi đúng tên; và đề cương không định nghĩa kỳ hạn ngắn, trung, dài.

Công trình thuộc nhóm “chế độ thị trường ở Việt Nam” trong Bảng 1: mô hình hóa biến động VN-Index, chưa nối chế độ với chất lượng quyết định danh mục (khoảng trống (3)).

\section{Nguyen và Nguyen (2026a): bài của chính tác giả}

Nguyen và Nguyen (2026a) áp dụng mô hình Black-Litterman với mô hình định giá tài sản vốn có chuyển chế độ cho thị trường Việt Nam giai đoạn 2019--2025 và chỉ đánh giá một mô hình.

Bài này là công trình số 2 ở Bảng 6: Nguyễn Thanh Bình và Nguyễn Văn Trung, Tạp chí Nghiên cứu Chính sách và Phát triển, đã công bố. Bài dùng phương pháp tính toán Bayes xấp xỉ.

Ba vai trong đề cương: M5 mở rộng từ bài này; tác giả đã dùng dữ liệu 2019--2025, gồm hai giai đoạn căng thẳng, nên không thể khẳng định chưa xem kết quả ở chế độ căng thẳng (mục 4.3.5); kết quả bài trước chỉ làm đối chứng.

Khác biệt (mục 4.4): đặt M5 cạnh sáu mô hình khác, mở mẫu về trước 2019 nếu dữ liệu cho phép, đo chênh lệch giữa các mô hình thay vì đánh giá riêng M5.

\section{Bảng khoảng trống nhìn từ nhóm (6)}

Bảng dưới gom các công trình Việt Nam theo ba hàng cuối của Bảng 1.

{\small\begin{longtable}[]{@{}
  >{\raggedright\arraybackslash}p{(\columnwidth - 4\tabcolsep) * \real{0.1520}}
  >{\raggedright\arraybackslash}p{(\columnwidth - 4\tabcolsep) * \real{0.3479}}
  >{\raggedright\arraybackslash}p{(\columnwidth - 4\tabcolsep) * \real{0.5001}}@{}}
\toprule\noalign{}
\begin{minipage}[b]{\linewidth}\raggedright
Nhóm
\end{minipage} & \begin{minipage}[b]{\linewidth}\raggedright
Công trình Việt Nam
\end{minipage} & \begin{minipage}[b]{\linewidth}\raggedright
Khoảng trống theo đề cương
\end{minipage} \\
\midrule\noalign{}
\endhead
\bottomrule\noalign{}
\endlastfoot
Danh mục ở Việt Nam, khoảng trống (2) & Ta và Vuong (2020), Nguyen và cộng sự (2020), Goutte và cộng sự (2025), Nguyen và Nguyen (2026a) & Đánh giá từng mô hình hoặc từng phương pháp; chưa đặt vào tập cạnh tranh để đo rủi ro lựa chọn theo chế độ \\
Chế độ ở Việt Nam, khoảng trống (3) & Hoang và Luu (2024) & Mô hình hóa biến động VN-Index; chưa nối với chất lượng quyết định danh mục \\
Bầy đàn ở Việt Nam, khoảng trống (4) & Vo và Phan (2017) & Bằng chứng bầy đàn trên HOSE; chưa gắn với rủi ro của việc chọn mô hình \\
\end{longtable}}

Điểm chung: các công trình đo từng mô hình, từng chế độ hoặc từng hiện tượng riêng lẻ; chưa công trình nào đo rủi ro của việc chọn giữa các mô hình theo chế độ thị trường.

\section{Dùng tuyên bố tính mới cho cẩn thận}

Đề cương viết “theo hiểu biết của thí sinh”. Căn cứ là truy vấn trên SSRN, RePEc/IDEAS, ScienceDirect, Oxford Academic, Wiley, Springer, Taylor \& Francis, MDPI và arXiv đến ngày 30/09/2026; phạm vi chưa gồm Scopus, Web of Science và tạp chí tiếng Việt, nên không loại trừ nghiên cứu ngoài phạm vi.

Đừng nói “không ai từng làm”; hãy nói “theo hiểu biết của tác giả, trong phạm vi tìm kiếm đã nêu, chưa thấy nghiên cứu nào đo rủi ro lựa chọn mô hình theo chế độ cho quyết định phân bổ danh mục ở Việt Nam”.

Việt Nam là nơi kiểm chứng vì thanh khoản hạn chế và nhiều cú sốc trong lịch sử giao dịch ngắn; đó không phải lý do của khoảng trống.

\section{Các cú sốc ở Việt Nam, như bối cảnh}

Mục 4.1 nêu ba loại cú sốc ở Việt Nam: khủng hoảng tài chính toàn cầu (Vo và Phan, 2017); COVID-19, số ca mắc hằng ngày tác động tiêu cực đến lợi suất cổ phiếu (Dao và cộng sự, 2021); thị trường trái phiếu doanh nghiệp đóng băng quý IV/2022, bất động sản chiếm tỷ trọng lớn nhất (Quỹ Tiền tệ Quốc tế; International Monetary Fund, 2023). Vu và cộng sự (2024) cho thấy các đợt sụt giảm giá 2007--2022 lan truyền giữa các ngành.

Các cú sốc này là giai đoạn tham chiếu để kiểm tra chế độ căng thẳng do mô hình ước lượng có trùng với biến cố thực tế không; đề cương không dùng chúng làm bằng chứng cho giả thuyết nào.

\begin{tcolorbox}[breakable]
\textbf{Ghi nhớ khi đi thi}

• Ở Việt Nam, mỗi công trình xét riêng một mô hình, một phương pháp hoặc một hiện tượng.

• Vo và Phan (2017): 299 công ty trên HOSE, 2005--2015, có bầy đàn, rõ hơn khi thị trường giảm.

• Ta và Vuong (2020): Black-Litterman, quan điểm từ mô hình tự hồi quy trung bình trượt tích hợp; Nguyen và cộng sự (2020): hiệp phương sai co rút, 468 tuần 2011--2019, lợi nhuận cao hơn, rủi ro thấp hơn.

• Goutte và cộng sự (2025): năm phương pháp đa tiêu chí, bốn cách gán trọng số, VN100, 01/2015--11/2023.

• Hoang và Luu (2024), bản thảo trực tuyến, VN-Index 2012--2023: một chế độ tốt hơn ở kỳ hạn ngắn, chuyển chế độ tốt hơn ở kỳ hạn trung và dài.

• Nguyen và Nguyen (2026a): bài của tác giả, Black-Litterman với mô hình định giá tài sản vốn có chuyển chế độ, 2019--2025; chỉ làm đối chứng.
\end{tcolorbox}

\section{Tự kiểm tra}

{\itshape 1. Vo và Phan (2017) dùng mẫu nào và thấy gì?\par}

Đáp án: 299 công ty trên HOSE, 2005--2015; có bầy đàn, rõ hơn khi thị trường giảm.

{\itshape 2. Ta và Vuong (2020) lấy quan điểm cho Black-Litterman từ đâu?\par}

Đáp án: từ mô hình tự hồi quy trung bình trượt tích hợp.

{\itshape 3. Nguyen và cộng sự (2020) dùng dữ liệu nào và thu được gì?\par}

Đáp án: hiệp phương sai co rút, 468 tuần 2011--2019; lợi nhuận cao hơn, rủi ro thấp hơn các phương pháp truyền thống.

{\itshape 4. Hoang và Luu (2024) thấy mô hình nào dự báo biến động tốt hơn ở từng kỳ hạn?\par}

Đáp án: một chế độ tốt hơn ở kỳ hạn ngắn; chuyển chế độ tốt hơn ở kỳ hạn trung và dài (bản thảo trực tuyến, VN-Index 2012--2023).

{\itshape 5. Vì sao tuyên bố tính mới của đề cương phải đi kèm cụm theo hiểu biết của thí sinh?\par}

Đáp án: vì truy vấn chỉ trên một số cơ sở dữ liệu đến 30/09/2026, chưa gồm Scopus, Web of Science và tạp chí tiếng Việt, nên không loại trừ nghiên cứu ngoài phạm vi.

\chapter{Mục 4.2 của đề cương (tiếp): Khoảng trống nghiên cứu, hay chỗ chưa ai đo}

Phần khoảng trống ở cuối mục 4.2 trả lời câu hội đồng nào cũng hỏi: đề tài làm gì mà người khác chưa làm. Câu trả lời gói trong bốn khoảng trống, đánh số (1) đến (4) ở Bảng 1, và một điểm chung: các công trình đo từng mô hình, từng chế độ hoặc từng hiện tượng riêng lẻ, chưa công trình nào đo rủi ro của việc chọn giữa các mô hình theo chế độ thị trường.

Tuyên bố này có hai lớp rào: “theo hiểu biết của thí sinh”, và danh sách nơi đã tìm, nơi chưa tìm. Hãy để ý hai lớp rào này không kém bốn khoảng trống.

\section{Khoảng trống nghiên cứu là gì}

Mỗi công trình đã vẽ kín một khu phố trên bản đồ thị trấn; khoảng trống là mảng trắng chưa ai vẽ. Đề tài tiến sĩ phải chỉ ra mảng trắng ở đâu và vì sao vẽ nó có ích.

Có những chỗ trống vì không ai cần vẽ, nên phần khoảng trống tốt làm ba việc: nêu điều đã làm, chỉ điều còn thiếu, giải thích vì sao điều thiếu đáng quan tâm. Bảng 1 làm hai việc đầu; việc thứ ba nằm ở mục 4.1 (lý do chọn đề tài).

Nhắc lại ba khái niệm mà cả bốn khoảng trống xoay quanh: \textbf{mô hình danh mục} là một cách quyết định chia tiền vào các cổ phiếu; \textbf{chế độ thị trường} là trạng thái thống kê của thị trường, như thời tiết (bình thường là trời quang, căng thẳng là mùa bão); \textbf{rủi ro lựa chọn mô hình} là tổn thất kỳ vọng do phải chọn một cách chia tiền trước khi biết cách nào tốt nhất, đo bằng mức hối tiếc.

\section{Hai vấn đề để ngỏ, nhìn từ trên cao}

Trước khi tách thành bốn hàng, mục 4.1 đã nêu gọn rằng tài liệu hiện có để ngỏ hai vấn đề.

Thứ nhất, các thước đo rủi ro mô hình (Boucher và cộng sự, 2014; Danielsson và cộng sự, 2016) đo chênh lệch giữa các dự báo rủi ro, chưa đo hậu quả của việc chọn sai mô hình lên trọng số và kết quả danh mục.

Ba dự báo viên cho ba xác suất mưa khác nhau. Đo độ chênh giữa ba con số cho biết họ bất đồng bao nhiêu, chưa cho biết buổi dã ngoại hỏng bao nhiêu nếu bạn nghe nhầm người. Thước đo hiện có đo độ bất đồng; đề tài đo buổi dã ngoại hỏng, tức chênh lệch kết quả danh mục.

Thứ hai, các nghiên cứu so sánh danh mục (DeMiguel và cộng sự, 2009; Garlappi và cộng sự, 2007) xếp hạng mô hình theo bình quân toàn mẫu, chưa xét thứ hạng thay đổi ra sao khi thị trường chuyển chế độ.

Đội xếp cao trên bảng điểm cả mùa chưa chắc mạnh hơn khi sân ướt mưa. Xếp hạng bình quân toàn mẫu gộp trời quang và mùa bão vào một con số; đề tài tách hai trạng thái và hỏi thứ hạng đổi thế nào.

\section{Bảng 1: bốn nhóm công trình và bốn khoảng trống}

Bảng dưới giữ nội dung Bảng 1 của đề cương.

{\small\begin{longtable}[]{@{}
  >{\raggedright\arraybackslash}p{(\columnwidth - 4\tabcolsep) * \real{0.3609}}
  >{\raggedright\arraybackslash}p{(\columnwidth - 4\tabcolsep) * \real{0.2437}}
  >{\raggedright\arraybackslash}p{(\columnwidth - 4\tabcolsep) * \real{0.3953}}@{}}
\toprule\noalign{}
\begin{minipage}[b]{\linewidth}\raggedright
Nhóm nghiên cứu
\end{minipage} & \begin{minipage}[b]{\linewidth}\raggedright
Kết quả đã có
\end{minipage} & \begin{minipage}[b]{\linewidth}\raggedright
Khoảng trống
\end{minipage} \\
\midrule\noalign{}
\endhead
\bottomrule\noalign{}
\endlastfoot
Rủi ro mô hình (Boucher và cộng sự, 2014; Danielsson và cộng sự, 2016) & Đo rủi ro mô hình của các mô hình dự báo rủi ro thị trường và rủi ro hệ thống; rủi ro mô hình tăng theo bất định & (1) Thước đo chưa gắn với hậu quả của quyết định phân bổ danh mục theo chế độ thị trường \\
Danh mục ở Việt Nam (Ta và Vuong, 2020; Nguyen và cộng sự, 2020; Goutte và cộng sự, 2025; Nguyen và Nguyen, 2026a) & Đánh giá từng mô hình hoặc từng phương pháp & (2) Chưa đặt các mô hình vào một tập cạnh tranh để đo rủi ro lựa chọn mô hình theo chế độ \\
Chế độ thị trường ở Việt Nam (Hoang và Luu, 2024) & Mô hình hóa biến động của VN-Index & (3) Chưa nối chế độ thị trường với chất lượng quyết định danh mục \\
Bầy đàn ở Việt Nam (Vo và Phan, 2017) & Bằng chứng bầy đàn trên HOSE & (4) Chưa gắn bầy đàn với rủi ro của việc chọn mô hình \\
\end{longtable}}

Mỗi hàng là một nhóm tài liệu; cột giữa là điều nhóm đã làm, cột phải là điều chưa làm, theo tác giả.

\subsection{Khoảng trống (1): thước đo rủi ro mô hình chưa gắn với hậu quả của quyết định phân bổ}

Boucher và cộng sự (2014) hiệu chỉnh độ đo rủi ro theo kết quả kiểm định ngược; Danielsson và cộng sự (2016) thấy rủi ro mô hình tăng theo bất định của thị trường, nhỏ khi thị trường yên ắng.

Kết luận của Danielsson và cộng sự chỉ cho mô hình dự báo rủi ro thị trường và rủi ro hệ thống, tức mô hình dự báo rủi ro, chưa phải mô hình quyết định chia tiền.

Nói bằng lời thường: người ta đã đo “các mô hình bất đồng bao nhiêu khi dự báo rủi ro”, chưa đo “chọn nhầm cách chia tiền thì danh mục thiệt bao nhiêu, và thiệt nhiều hơn hay ít hơn tùy trời quang hay mùa bão”.

Theo bản đọc, mức hối tiếc ở mục 4.3.6 (chênh lệch lợi suất tương đương chắc chắn giữa mô hình tốt nhất sau thực tế và mô hình chọn trước) tính trên kết quả danh mục chứ không dừng ở chênh lệch dự báo, nên có thể đọc như câu trả lời cho khoảng trống (1).

\subsection{Khoảng trống (2): danh mục ở Việt Nam được đánh giá từng mô hình một}

Bốn công trình, theo nhóm (6):

\begin{itemize}
\item
  Ta và Vuong (2020): Black-Litterman với quan điểm từ mô hình tự hồi quy trung bình trượt tích hợp.
\item
  Nguyen và cộng sự (2020): hiệp phương sai co rút trên 468 tuần 2011--2019, lợi nhuận cao hơn, rủi ro thấp hơn phương pháp truyền thống.
\item
  Goutte và cộng sự (2025): năm phương pháp đa tiêu chí, bốn cách gán trọng số, VN100, 01/2015--11/2023.
\item
  Nguyen và Nguyen (2026a), bài của tác giả: Black-Litterman với mô hình định giá tài sản vốn có chuyển chế độ, 2019--2025.
\end{itemize}

Cả bốn “đánh giá từng mô hình hoặc từng phương pháp”; khoảng trống (2) là chưa đặt các mô hình vào một tập cạnh tranh.

Muốn biết thợ nào sửa xe giỏi: đưa xe cho một thợ rồi chấm, hoặc đưa cùng chiếc xe cho bảy thợ, cùng đề bài, rồi so. Chỉ cách thứ hai cho biết chọn nhầm thợ thiệt bao nhiêu. “Bảy thợ” là M1 đến M7 ở Bảng 3; “chọn nhầm thợ” là rủi ro lựa chọn mô hình.

Goutte và cộng sự đã so sánh nhiều phương pháp, nên đề cương không nói “chưa ai so sánh nhiều phương pháp ở Việt Nam”. Đề cương nói họ so sánh phương pháp, còn đề tài đo mức bất định khi chọn mô hình theo từng chế độ.

\subsection{Khoảng trống (3): chế độ ở Việt Nam chưa nối với quyết định danh mục}

Hoang và Luu (2024), bản thảo trực tuyến, phân tích VN-Index 2012--2023: mô hình một chế độ dự báo biến động tốt hơn ở kỳ hạn ngắn, chuyển chế độ tốt hơn ở kỳ hạn trung và dài.

Công trình mô tả thời tiết của thị trường, chưa nối thời tiết ấy với chất lượng quyết định danh mục.

\subsection{Khoảng trống (4): bầy đàn ở Việt Nam chưa gắn với việc chọn mô hình}

Vo và Phan (2017) dùng 299 công ty trên HOSE 2005--2015 và tìm thấy bầy đàn, rõ hơn khi thị trường giảm.

Bầy đàn là làm theo số đông thay vì dựa vào thông tin riêng, như đàn cừu ngoái nhìn con phía trước chứ không nhìn đường.

Đề tài không quan sát hành vi chọn mô hình của tổ chức đầu tư (mục 4.4); phép thử bầy đàn chỉ dùng độ phân tán chéo của lợi suất ở cấp thị trường, sau khi kiểm soát biến động. Vì H4 dựa trên Vo và Phan (2017) và ghi là thăm dò, đề tài chạm vào khoảng trống (4) ở mức thăm dò.

\section{Điểm chung của bốn khoảng trống}

Sau bảng, đề cương viết: bốn khoảng trống có chung một điểm, các công trình đo từng mô hình, từng chế độ hoặc từng hiện tượng riêng lẻ, chưa công trình nào đo rủi ro của việc chọn giữa các mô hình theo chế độ thị trường.

Tóm bằng chữ “riêng lẻ”: mô hình riêng lẻ (khoảng trống (2)), chế độ riêng lẻ (3), hiện tượng riêng lẻ (4); cách gán này là của bản đọc. Đề tài ghép lại: nhiều mô hình cạnh tranh, xét theo chế độ, kèm các biến kênh.

Mục 8.2 diễn đạt bước chuyển này: luận án chuyển trọng tâm từ đánh giá một mô hình sang câu hỏi chọn giữa nhiều mô hình.

\section{Vì sao là Việt Nam: nơi kiểm chứng, không phải lý do của khoảng trống}

Đừng nghĩ khoảng trống là “chưa ai làm ở Việt Nam”. Khoảng trống là một câu hỏi khoa học chung; Việt Nam là nơi thử.

Mục 4.1 nêu lý do: thanh khoản hạn chế và nhiều cú sốc trong một lịch sử giao dịch ngắn, nên chênh lệch giữa các mô hình có thể lớn hơn ở thị trường phát triển. Đề cương dùng chữ “có thể”, không dùng chữ “độc nhất”.

Lịch sử ngắn cũng là giới hạn: mẫu có ít giai đoạn căng thẳng độc lập, nên kiểm định có công suất thấp; khi số cửa sổ độc lập ở chế độ căng thẳng dưới 10, kết quả gọi là không kết luận (mục 4.3.10). Khoảng trống có thật không bảo đảm dữ liệu Việt Nam đủ để lấp nó.

\section{Tuyên bố tính mới và phạm vi tìm kiếm}

Tuyên bố tính mới: theo hiểu biết của thí sinh, chưa có nghiên cứu nào đo rủi ro lựa chọn mô hình theo chế độ thị trường cho quyết định phân bổ danh mục trên thị trường chứng khoán Việt Nam.

Cụm “theo hiểu biết của” biến câu tuyệt đối (“chưa ai làm”) thành câu có điều kiện (“tôi chưa thấy ai làm”). Không ai chứng minh được rằng một công trình không tồn tại; chỉ nói được đã tìm ở đâu và không thấy gì.

\subsection{Tác giả đã tìm ở đâu}

Truy vấn thực hiện đến 30/09/2026 trên chín cơ sở: SSRN, RePEc/IDEAS, ScienceDirect, Oxford Academic, Wiley, Springer, Taylor \& Francis, MDPI và arXiv.

\subsection{Tác giả không tìm ở đâu}

Phạm vi chưa gồm Scopus, Web of Science và các tạp chí tiếng Việt trong nước, nên nhận định không loại trừ nghiên cứu ngoài phạm vi đó.

Tìm chìa khóa ở phòng khách, phòng ngủ và bếp rồi không thấy: kết luận đúng là “không thấy ở ba phòng này”, chưa phải “không có trong nhà”. Tác giả khai rõ phòng nào chưa lục, không nói trong đó có hay không.

\subsection{Mức độ đã đọc}

Tác giả mới đọc tóm tắt và thông tin thư mục của phần lớn nguồn; đọc toàn văn và mở rộng tìm kiếm nằm trong kế hoạch năm thứ nhất.

Vì vậy mọi mô tả về công trình trước trong bản đọc này chỉ đúng như đề cương nói về chúng, không thêm chi tiết.

\subsection{Công trình gần nhất}

Mục 4.4 nhắc các công trình gần nhất: Ta và Vuong (2020), Goutte và cộng sự (2025), Hoang và Luu (2024) và bài của tác giả (Nguyen và Nguyen, 2026a); Nguyen và cộng sự (2020) cùng nằm ở hàng “danh mục ở Việt Nam” của Bảng 1.

Tác giả tự đưa bài của mình vào danh sách gần nhất, nên người đọc thấy ngay đề tài khác bài cũ ở đâu: M5 cạnh sáu mô hình khác, mẫu mở về trước 2019 nếu được, đo chênh lệch giữa các mô hình thay vì đánh giá riêng M5; kết quả bài trước chỉ làm đối chứng.

\section{Bốn điểm khác biệt, đối chiếu với bốn khoảng trống}

Mục 4.4 nêu bốn điểm khác biệt. Bản đọc đối chiếu chúng với Bảng 1 (cách ghép là của bản đọc):

{\itshape 1. Ta và Vuong (2020) áp dụng một mô hình Black-Litterman; đề tài đặt nó cùng sáu mô hình khác vào tập cạnh tranh. Ứng với khoảng trống (2).

2. Goutte và cộng sự (2025) so sánh phương pháp đa tiêu chí trên VN100; đề tài đo bất định khi chọn mô hình theo chế độ. Cũng ứng với (2), ở khía cạnh theo chế độ.

3. Hoang và Luu (2024) phân tích chế độ của biến động; đề tài nối chế độ với chất lượng quyết định phân bổ. Ứng với (3). Điểm thứ tư là khác biệt với bài trước của tác giả.\par}

Khoảng trống (1) gắn với thước đo mức hối tiếc (mục 4.3.6), khoảng trống (4) với giả thuyết thăm dò H4 (Bảng 2).

\section{Điều mục 4.2 không nói}

Để tránh hiểu quá mức, phần này không nói:

\begin{itemize}
\item
  rằng chưa ai trên thế giới nghiên cứu rủi ro lựa chọn mô hình; nó chỉ nói về việc đo theo chế độ cho phân bổ danh mục ở Việt Nam, theo hiểu biết của tác giả;
\item
  rằng nghiên cứu ngoài phạm vi tìm kiếm không tồn tại;
\item
  rằng Việt Nam là nơi duy nhất hay là lý do của khoảng trống;
\item
  rằng các công trình trước sai; cột “Kết quả đã có” mô tả trung tính, cột “Khoảng trống” chỉ nói chúng chưa trả lời một câu hỏi khác;
\item
  và đề tài sẽ tìm thấy gì. Mục 4.3.2 biến khoảng trống thành các giả thuyết có thể bị bác bỏ.
\end{itemize}

\section{Ghi nhớ khi đi thi}

\begin{tcolorbox}[breakable]
Bốn khoảng trống: (1) thước đo rủi ro mô hình chưa gắn với hậu quả quyết định phân bổ theo chế độ; (2) danh mục ở Việt Nam chưa đặt vào tập cạnh tranh; (3) chế độ ở Việt Nam chưa nối với chất lượng quyết định; (4) bầy đàn chưa gắn với rủi ro của việc chọn mô hình.

Điểm chung: các công trình đo từng mô hình, từng chế độ, từng hiện tượng riêng lẻ; chưa công trình nào đo rủi ro của việc chọn giữa các mô hình theo chế độ.

Hai vấn đề để ngỏ: thước đo rủi ro mô hình chỉ đo chênh lệch dự báo rủi ro; nghiên cứu danh mục xếp hạng theo bình quân toàn mẫu, chưa xét thứ hạng đổi theo chế độ.

Việt Nam là nơi kiểm chứng (thanh khoản hạn chế, nhiều cú sốc, lịch sử ngắn), không phải lý do của khoảng trống.

Tính mới chỉ “theo hiểu biết của thí sinh”: chín cơ sở đến 30/09/2026; chưa gồm Scopus, Web of Science, tạp chí tiếng Việt; mới đọc tóm tắt và thông tin thư mục.
\end{tcolorbox}

\section{Tự kiểm tra}

{\itshape 1. Bốn khoảng trống là gì, và điểm chung của chúng?\par}

Đáp án: (1) thước đo rủi ro mô hình chưa gắn với hậu quả quyết định phân bổ theo chế độ; (2) danh mục ở Việt Nam chưa đặt vào tập cạnh tranh; (3) chế độ ở Việt Nam chưa nối với chất lượng quyết định danh mục; (4) bầy đàn chưa gắn với rủi ro của việc chọn mô hình. Điểm chung: chưa công trình nào đo rủi ro của việc chọn giữa các mô hình theo chế độ thị trường.

{\itshape 2. Vì sao chọn Việt Nam, và đó có phải lý do của khoảng trống không?\par}

Đáp án: thanh khoản hạn chế và nhiều cú sốc trong lịch sử giao dịch ngắn, nên chênh lệch giữa các mô hình có thể lớn. Đó là nơi kiểm chứng, không phải lý do của khoảng trống.

{\itshape 3. Tuyên bố tính mới được rào bằng những điều kiện nào?\par}

Đáp án: “theo hiểu biết của thí sinh”; truy vấn chín cơ sở đến 30/09/2026; chưa gồm Scopus, Web of Science và tạp chí tiếng Việt; mới đọc tóm tắt và thông tin thư mục, đọc toàn văn trong năm thứ nhất.

{\itshape 4. Hai vấn đề để ngỏ ở mục 4.1 là gì?\par}

Đáp án: thước đo rủi ro mô hình (Boucher và cộng sự, 2014; Danielsson và cộng sự, 2016) đo chênh lệch dự báo rủi ro, chưa đo hậu quả lên trọng số và kết quả danh mục; nghiên cứu danh mục (DeMiguel và cộng sự, 2009; Garlappi và cộng sự, 2007) xếp hạng theo bình quân toàn mẫu, chưa xét thứ hạng đổi theo chế độ.

{\itshape 5. Đề tài khác bài Nguyen và Nguyen (2026a) ở những điểm nào?\par}

Đáp án: M5 cạnh sáu mô hình khác, mẫu mở về trước 2019 nếu được, đo chênh lệch giữa các mô hình thay vì đánh giá riêng M5; kết quả bài trước chỉ làm đối chứng.

\chapter{Mục 4.3.2 của đề cương: Bốn giả thuyết và khung phân tích}

Từ bốn khoảng trống ở Bảng 1, đề cương đề xuất khung phân tích (Hình 1) và bốn giả thuyết (Bảng 2). H1 hỏi rủi ro lựa chọn mô hình ở chế độ căng thẳng có cao hơn không; H2 và H4 hỏi vì sao, qua các biến kênh; H3 hỏi cách nào giảm được rủi ro. Mỗi giả thuyết đi kèm một khả năng đối lập, tức một lời giải thích khác có thể thắng.

Cả bốn là dự đoán cần kiểm định; chữ “dự kiến” và “nếu dữ liệu ủng hộ” là cách viết đúng.

\section{Giả thuyết phải viết sao cho có thể sai}

Giả thuyết khoa học là dự đoán đủ cụ thể để dữ liệu có thể phản bác. “Ngày mai trời sẽ thay đổi” không thể sai nên vô giá trị; “ngày mai mưa trên 10 mm” thì kiểm tra được. Mỗi giả thuyết của đề cương nêu thước đo nào, dấu hệ số ra sao, và điều gì sẽ bác bỏ nó.

\textbf{Thước đo chính} là đại lượng quyết định phán quyết; \textbf{thước đo phụ} chỉ bổ sung. \textbf{Cơ sở} là công trình làm chỗ dựa lập luận, không phải bằng chứng rằng giả thuyết đúng ở Việt Nam.

\section{Bảng 2: giả thuyết, thước đo và cơ sở lý thuyết}

Bảng dưới giữ nội dung Bảng 2 của đề cương.

{\small\begin{longtable}[]{@{}
  >{\raggedright\arraybackslash}p{(\columnwidth - 4\tabcolsep) * \real{0.3448}}
  >{\raggedright\arraybackslash}p{(\columnwidth - 4\tabcolsep) * \real{0.3448}}
  >{\raggedright\arraybackslash}p{(\columnwidth - 4\tabcolsep) * \real{0.3104}}@{}}
\toprule\noalign{}
\begin{minipage}[b]{\linewidth}\raggedright
Giả thuyết
\end{minipage} & \begin{minipage}[b]{\linewidth}\raggedright
Thước đo chính
\end{minipage} & \begin{minipage}[b]{\linewidth}\raggedright
Cơ sở
\end{minipage} \\
\midrule\noalign{}
\endhead
\bottomrule\noalign{}
\endlastfoot
H1: Rủi ro lựa chọn mô hình ở chế độ căng thẳng cao hơn ở chế độ bình thường & Kết quả chính: mức hối tiếc \(R_{t}\) (chênh lệch lợi suất tương đương chắc chắn giữa mô hình tốt nhất sau thực tế và mô hình chọn trước). Kết quả phụ: độ phân tán mô hình \(D_{t}\), kích thước tập tin cậy mô hình & Danielsson và cộng sự (2016); Hansen và cộng sự (2011) \\
H2: Sau khi kiểm soát biến động thực hiện, hệ số của độ phi thanh khoản dương (thanh khoản cạn làm các mô hình có vòng quay khác nhau cách xa nhau) và hệ số của tương quan đuôi âm (tương quan cao làm các danh mục chỉ mua hội tụ) & Dấu và khoảng tin cậy của hai hệ số; phần chênh lệch giữa hai chế độ (\(\beta\)) do biến kênh giải thích & Amihud (2002); Longin và Solnik (2001); Forbes và Rigobon (2002) \\
H3: Ở chế độ căng thẳng, sau chi phí giao dịch, kết hợp mô hình và danh mục vững chắc cho mức hối tiếc và độ sụt giảm tối đa thấp hơn, lợi suất tương đương chắc chắn cao hơn so với chọn một mô hình & Lợi suất tương đương chắc chắn; mức hối tiếc; độ sụt giảm tối đa; giá trị thiệt hại kỳ vọng mức 95\%; vòng quay & Hoeting và cộng sự (1999); Garlappi và cộng sự (2007); Rapach và cộng sự (2010) \\
H4 (thăm dò): Sau khi kiểm soát biến động thực hiện, độ phân tán chéo của lợi suất cao đi cùng rủi ro lựa chọn mô hình cao & Độ phân tán chéo của lợi suất; thước đo bầy đàn & Vo và Phan (2017); Chang và cộng sự (2000); Christie và Huang (1995) \\
\end{longtable}}

H1 trả lời Câu hỏi 1; H2 và H4 trả lời Câu hỏi 2; H3 trả lời Câu hỏi 3. Câu hỏi 2 hỏi các biến kênh có liên hệ thống kê với rủi ro không và giải thích bao nhiêu phần chênh lệch giữa hai chế độ, tức bao nhiêu phần của hệ số \(\beta\) ở phương trình (5).

\section{H1: mùa bão làm việc chọn cách chia tiền khó hơn}

H1: rủi ro lựa chọn mô hình ở chế độ căng thẳng cao hơn ở chế độ bình thường. Bằng ngôn ngữ thời tiết: chọn cách chia tiền trước mùa giải dễ sai hơn, và sai đắt hơn, khi thị trường vào mùa bão.

Kết quả chính là mức hối tiếc \(R_{t}\). Đầu mỗi cửa sổ đánh giá, nhà đầu tư chọn trước mô hình có lợi suất tương đương chắc chắn ngoài mẫu cao nhất trong 36 tháng liền trước; hết cửa sổ mới biết mô hình nào tốt nhất. Chênh lệch giữa hai kết quả là mức hối tiếc.

Lợi suất tương đương chắc chắn là mức lợi suất chắc chắn mà nhà đầu tư ngại rủi ro coi là ngang giá với danh mục; càng cao càng tốt. Một cách chia tiền lãi to lỗ to, một cách lãi vừa mà đều: thước đo này so hai cách trên cùng một thước vì nó trừ một khoản “phí cho sự bất an” tùy mức ngại rủi ro.

Hai kết quả phụ là độ phân tán mô hình \(D_{t}\) (bảy mô hình bất đồng nhiều hay ít) và kích thước tập tin cậy mô hình (danh sách ứng viên chưa thể loại). Chúng chỉ báo bất định, không đo rủi ro, nên là kết quả phụ.

Cơ sở: Danielsson và cộng sự (2016) thấy rủi ro mô hình tăng theo bất định của thị trường, nhưng chỉ cho mô hình dự báo rủi ro thị trường và rủi ro hệ thống; đề tài thử ý tương tự cho quyết định phân bổ.

Hansen và cộng sự (2011) là nguồn của tập tin cậy mô hình. Dữ liệu kém thông tin cho tập lớn, nên tập lớn hơn ở chế độ căng thẳng chưa chắc là rủi ro lớn hơn; vì vậy kích thước tập chỉ báo cáo dạng mô tả, kèm số quan sát và công suất kiểm định.

Khả năng đối lập: rủi ro không đổi theo chế độ, hoặc một quy tắc đơn giản như 1/N thắng ổn định (DeMiguel và cộng sự, 2009; Pflug và cộng sự, 2012). Nếu vậy, kết luận là 1/N đủ dùng ở Việt Nam (mục 4.3.10).

Tiêu chí ở mục 4.3.7: H1 được ủng hộ khi thỏa đồng thời ba điều kiện: hệ số \(\beta\) của \(R_{t}\) dương và vượt phân vị 95\% của phân phối rỗng; khoảng tin cậy bootstrap khối của nó loại trừ 0; và \(\beta\) của \(D_{t}\) cùng dấu. Phân phối rỗng là “kết quả trông thế nào nếu thực ra không có gì xảy ra”.

Điểm tinh tế: ngay cả khi các mô hình ngang chất lượng, mức hối tiếc vẫn có kỳ vọng dương, vì giá trị lớn nhất của các đại lượng có nhiễu tăng theo phương sai của nhiễu. Chế độ nhiễu hơn tự nhiên cho mức hối tiếc lớn hơn. Vì vậy cần phân phối rỗng, và cần chia \(D_{t}\), \(R_{t}\) cho sai số chuẩn bootstrap để loại tác động cơ học này.

\begin{tcolorbox}[breakable]
Ví dụ giả định (số tự đặt để minh họa): bảy cầu thủ kỹ năng y hệt cùng sút phạt. Ngày quang, kết quả lệch nhau ít; ngày gió lớn, lệch nhiều. Không ai giỏi hơn ai, nhưng người tốt nhất ngày gió vẫn thường ghi điểm cao hơn người tốt nhất ngày quang.

Chọn trước một người rồi so với người tốt nhất trong ngày thì khoảng cách ngày gió lớn hơn, chỉ vì may rủi lớn hơn. Cần thước đối chứng để tách “may rủi lớn hơn” khỏi “chọn khó hơn thật”.
\end{tcolorbox}

\section{H2: hai biến kênh, hai dấu ngược nhau}

Biến kênh là đại lượng có thể giải thích vì sao rủi ro lựa chọn mô hình đổi theo chế độ. Nếu H1 cho thấy mùa bão làm việc chọn khó hơn, H2 hỏi khó hơn qua “đường dẫn” nào: thanh khoản cạn là một đường, tương quan giữa các cổ phiếu là đường khác.

“Sau khi kiểm soát biến động thực hiện”: biến động thực hiện là mức dao động đã quan sát trong cửa sổ; kiểm soát nó là đưa vào hồi quy để tách ảnh hưởng riêng của biến kênh. Ăn kem và đuối nước cùng tăng vào mùa hè; muốn xem riêng kem phải giữ nhiệt độ cố định. Ở đây biến động đóng vai nhiệt độ.

Mục 4.3.7 lưu ý: biến chế độ phụ thuộc biến động và lợi suất, các biến kênh cũng phụ thuộc biến động, nên chúng dễ đi cùng nhau. Đề cương báo cáo hệ số phóng đại phương sai với ngưỡng 5 ấn định trước; vượt ngưỡng thì báo cáo hai phương trình tách riêng và không diễn giải riêng từng hệ số.

H2 có hai nửa với hai dấu: hệ số của độ phi thanh khoản dương; hệ số của tương quan đuôi âm.

Độ phi thanh khoản Amihud (Amihud, 2002) càng cao thì thị trường càng “cạn”: giá dễ bị đẩy mạnh khi có lệnh mua bán. Đề cương không nêu công thức.

Vì sao dương? Thanh khoản cạn làm các mô hình có vòng quay khác nhau cách xa nhau. Vòng quay là tỷ lệ danh mục được mua bán lại. Một cách chia tiền đổi vị trí liên tục, một cách gần như đứng yên; chợ vắng thì mỗi lần đổi tốn hơn, nên hai cách cách xa nhau và chọn nhầm đắt hơn.

Vì sao âm? Tương quan cao làm các danh mục chỉ mua hội tụ. Tương quan đuôi là mức các cổ phiếu cùng rơi khi thị trường xấu nhất; danh mục chỉ mua là danh mục không bán khống. Diễn giải của bản đọc: mọi cổ phiếu cùng đi một hướng thì mô hình nào cũng mua những thứ gần giống nhau, danh mục giống nhau thì chọn nhầm ít thiệt.

Hai dấu ngược nhau không mâu thuẫn: một biến làm các mô hình khác nhau hơn, biến kia làm chúng giống nhau hơn. Cả hai là dự đoán cần kiểm tra.

Thước đo của H2 gồm dấu và khoảng tin cậy của hai hệ số, cộng phần của \(\beta\) do biến kênh giải thích: hiệu giữa hai hệ số \(\beta\) khi chưa có và khi có biến kênh trong phương trình (5), kèm khoảng tin cậy bootstrap khối. Tổng thiệt hại do mùa bão là \(\beta\); thêm biến kênh mà \(\beta\) nhỏ đi thì phần giảm là phần biến kênh “gánh” thay chế độ.

Kết quả chỉ là liên hệ có điều kiện, không phải nhân quả: biến kênh “gánh” một phần của \(\beta\) nghĩa là đi cùng phần đó về thống kê, không chứng minh nó gây ra.

\subsection{Cơ sở, khả năng đối lập và tiêu chí của H2}

Cơ sở: Amihud (2002) cho độ phi thanh khoản; Longin và Solnik (2001) bác bỏ phân phối chuẩn đa biến ở đuôi âm; Forbes và Rigobon (2002) chỉ ra tương quan phụ thuộc biến động.

Khả năng đối lập: tương quan tăng chỉ vì biến động tăng (Forbes và Rigobon, 2002). Theo bản đọc, nếu sau khi kiểm soát biến động mà tương quan đuôi không còn tác động riêng, “kênh tương quan” chỉ là biến động khoác áo khác.

H2 được ủng hộ khi hệ số độ phi thanh khoản dương, hệ số tương quan đuôi âm, và khoảng tin cậy bootstrap khối của cả hai loại trừ 0.

\section{H3: cách giảm rủi ro khi trời bão}

H3: ở chế độ căng thẳng, sau chi phí giao dịch, kết hợp mô hình và danh mục vững chắc cho mức hối tiếc và độ sụt giảm tối đa thấp hơn, lợi suất tương đương chắc chắn cao hơn so với chọn một mô hình. Thay vì đặt cả mùa giải vào một cách chia tiền, nhà đầu tư pha trộn nhiều cách, hoặc chọn cách chịu được tình huống xấu nhất.

H3 trả lời Câu hỏi 3; chi tiết ở mục 4.3.8.

Kết hợp mô hình: trung bình hóa mô hình Bayes (Hoeting và cộng sự, 1999) trộn nhiều mô hình theo trọng số, như nghe cả bảy dự báo viên và cho người đúng nhiều trong quá khứ tiếng nói lớn hơn. Danh mục vững chắc: nhà đầu tư né tránh mơ hồ tối ưu cho tình huống xấu nhất trong nhiều niềm tin tiên nghiệm (Garlappi và cộng sự, 2007).

Thước đo: lợi suất tương đương chắc chắn; mức hối tiếc; độ sụt giảm tối đa (mức giảm lớn nhất từ đỉnh); giá trị thiệt hại kỳ vọng mức 95\% (tổn thất trung bình trong 5\% số ngày xấu nhất); vòng quay. So sánh tính sau chi phí giao dịch, vì một cách pha trộn đẹp trên giấy có thể mất hết lợi thế nếu phải mua bán liên tục.

Chi tiết kỹ thuật (mục 4.3.8): khi kiểm định H3, cơ chế (ii) danh mục vững chắc và (iii) quy tắc chuyển chế độ được loại khỏi tập bảy mô hình, nên mức hối tiếc của chúng có thể âm (mức hối tiếc chỉ chắc chắn không âm khi mô hình chọn trước thuộc tập).

\subsection{Cơ sở, khả năng đối lập và tiêu chí của H3}

Cơ sở: Hoeting và cộng sự (1999), Garlappi và cộng sự (2007), Rapach và cộng sự (2010); nhóm sau cho thấy kết hợp dự báo bù được phần chất lượng mất đi do mô hình bất ổn định.

Khả năng đối lập: kết hợp mô hình không vượt việc chọn một mô hình sau chi phí giao dịch và sau hiệu chỉnh thiên lệch chọn lọc (Bailey và López de Prado, 2014). Nếu lợi thế biến mất sau hai bước đó, H3 không được ủng hộ.

H3 được ủng hộ khi, ở chế độ căng thẳng, cơ chế cho lợi suất tương đương chắc chắn sau chi phí cao hơn, mức hối tiếc thấp hơn và độ sụt giảm tối đa không cao hơn so với chọn một mô hình, với khoảng tin cậy bootstrap khối loại trừ 0.

\section{H4: thăm dò bầy đàn}

H4 ghi rõ là thăm dò: sau khi kiểm soát biến động thực hiện, độ phân tán chéo của lợi suất cao đi cùng rủi ro lựa chọn mô hình cao.

Độ phân tán chéo là mức lợi suất các cổ phiếu khác nhau trong cùng một ngày, thước đo bầy đàn ở cấp thị trường theo Christie và Huang (1995): cừu chạy theo đầu đàn thì sát nhau, độ phân tán nhỏ; mỗi con tự chọn đường thì tản ra.

\subsection{Ba giới hạn của H4}

\begin{itemize}
\item
  Đề tài không quan sát hành vi chọn mô hình của tổ chức đầu tư; đây là phép thử giới hạn ở cấp thị trường.
\item
  H4 chỉ báo cáo như kết quả thăm dò.
\item
  Bằng chứng bầy đàn trái chiều giữa các thị trường (Christie và Huang, 1995; Chang và cộng sự, 2000).
\end{itemize}

Cơ sở: Vo và Phan (2017), Chang và cộng sự (2000), Christie và Huang (1995). Khả năng đối lập: bầy đàn không xuất hiện ở tần suất dữ liệu của đề tài, như Christie và Huang (1995) quan sát trên dữ liệu Mỹ.

\section{Hình 1: khung phân tích, mô tả bằng lời}

Hình 1 là sơ đồ khung phân tích. Đề cương mô tả bốn quan hệ:

{\itshape 1. Chế độ thị trường tác động lên rủi ro lựa chọn mô hình (H1).

2. Các biến kênh tác động lên chính quan hệ đó, tức giải thích một phần chênh lệch giữa hai chế độ (H2 và H4). Mục tiêu (iii) ở mục 4.1 liệt kê biến động, thanh khoản, tương quan cực trị và bầy đàn; biến động thực hiện luôn là biến kiểm soát.

3. Cơ chế giảm thiểu tác động lên quan hệ giữa rủi ro lựa chọn mô hình và kết quả danh mục (H3).

4. Mũi tên nét đứt từ rủi ro lựa chọn mô hình sang kết quả danh mục không mang giả thuyết riêng, vì định nghĩa của mức hối tiếc đã là chênh lệch kết quả danh mục.\par}

Điểm thứ tư: mức hối tiếc là chênh lệch lợi suất tương đương chắc chắn, mà đó chính là thước đo kết quả danh mục, nên rủi ro cao hơn đồng nghĩa kết quả của mô hình đã chọn kém hơn mô hình tốt nhất. Định nghĩa “khoảng cách đến đích” bằng hiệu hai quãng đường thì khoảng cách lớn đồng nghĩa còn xa đích; không cần thí nghiệm để chứng minh.

Gom lại: chế độ thị trường, rồi rủi ro lựa chọn mô hình, rồi kết quả danh mục; biến kênh tác động lên mũi tên thứ nhất, cơ chế giảm thiểu tác động lên mũi tên thứ hai.

\section{Kiểm định cùng lúc một khả năng đối lập cho mỗi giả thuyết}

Đề cương kiểm định cùng lúc một khả năng đối lập cho mỗi giả thuyết:

{\small\begin{longtable}[]{@{}
  >{\raggedright\arraybackslash}p{(\columnwidth - 4\tabcolsep) * \real{0.0797}}
  >{\raggedright\arraybackslash}p{(\columnwidth - 4\tabcolsep) * \real{0.6551}}
  >{\raggedright\arraybackslash}p{(\columnwidth - 4\tabcolsep) * \real{0.2652}}@{}}
\toprule\noalign{}
\begin{minipage}[b]{\linewidth}\raggedright
Giả thuyết
\end{minipage} & \begin{minipage}[b]{\linewidth}\raggedright
Khả năng đối lập
\end{minipage} & \begin{minipage}[b]{\linewidth}\raggedright
Nguồn
\end{minipage} \\
\midrule\noalign{}
\endhead
\bottomrule\noalign{}
\endlastfoot
H1 & Rủi ro không đổi theo chế độ, hoặc một quy tắc đơn giản như 1/N thắng ổn định & DeMiguel và cộng sự (2009); Pflug và cộng sự (2012) \\
H2 & Tương quan có thể tăng chỉ vì biến động tăng & Forbes và Rigobon (2002) \\
H3 & Kết hợp mô hình có thể không vượt việc chọn một mô hình sau chi phí giao dịch và sau hiệu chỉnh thiên lệch chọn lọc & Bailey và López de Prado (2014) \\
H4 & Bầy đàn có thể không xuất hiện ở tần suất dữ liệu của đề tài & Christie và Huang (1995) \\
\end{longtable}}

Thám tử chỉ có một nghi phạm dễ kết tội oan; thám tử giữ sẵn giả thuyết khác buộc mình hỏi bằng chứng nào phân biệt hai bên. Nêu sẵn khả năng đối lập cho người đọc biết đề tài chấp nhận kết quả nào là câu trả lời, kể cả kết quả bất lợi.

Mục 4.3.10 xác nhận: nếu dữ liệu bác bỏ giả thuyết, đề tài vẫn trả lời được câu hỏi, chẳng hạn rủi ro không đổi theo chế độ hoặc 1/N đủ dùng ở Việt Nam.

Hai lớp bảo vệ chung: hiệu chỉnh so sánh bội theo Holm (1979) ở mức 10\% cho bốn giả thuyết, để tránh thử nhiều rồi nhặt cái “có ý nghĩa”; và quy tắc không kết luận khi số cửa sổ độc lập ở chế độ căng thẳng dưới 10.

\section{Ghi nhớ khi đi thi}

\begin{tcolorbox}[breakable]
H1: căng thẳng cao hơn bình thường; kết quả chính là mức hối tiếc \(R_{t}\), phụ là \(D_{t}\) và kích thước tập tin cậy; đối lập: không đổi theo chế độ hoặc 1/N thắng ổn định.

H2: sau khi kiểm soát biến động thực hiện, hệ số độ phi thanh khoản dương, hệ số tương quan đuôi âm; đối lập: tương quan tăng chỉ vì biến động tăng.

H3: ở chế độ căng thẳng, sau chi phí, kết hợp mô hình và danh mục vững chắc cho mức hối tiếc, độ sụt giảm tối đa thấp hơn, lợi suất tương đương chắc chắn cao hơn; đối lập: không vượt sau chi phí và sau hiệu chỉnh thiên lệch chọn lọc.

H4 chỉ thăm dò: độ phân tán chéo cao đi cùng rủi ro lựa chọn cao; đối lập: bầy đàn không xuất hiện ở tần suất dữ liệu.

Hình 1: chế độ tác động lên rủi ro (H1); biến kênh tác động lên quan hệ đó (H2, H4); cơ chế giảm thiểu tác động lên quan hệ rủi ro và kết quả danh mục (H3); mũi tên nét đứt không mang giả thuyết.

Quan hệ chỉ là liên hệ có điều kiện, không phải nhân quả.
\end{tcolorbox}

\section{Tự kiểm tra}

{\itshape 1. Nêu H2 và hai dấu của nó. Vì sao ngược nhau?\par}

Đáp án: sau khi kiểm soát biến động thực hiện, hệ số độ phi thanh khoản dương và hệ số tương quan đuôi âm. Thanh khoản cạn làm các mô hình có vòng quay khác nhau cách xa nhau; tương quan cao làm các danh mục chỉ mua hội tụ.

{\itshape 2. Khả năng đối lập của từng giả thuyết là gì?\par}

Đáp án: H1: không đổi theo chế độ hoặc 1/N thắng ổn định (DeMiguel và cộng sự, 2009; Pflug và cộng sự, 2012). H2: tương quan tăng chỉ vì biến động tăng (Forbes và Rigobon, 2002). H3: kết hợp không vượt chọn một mô hình sau chi phí và hiệu chỉnh thiên lệch chọn lọc (Bailey và López de Prado, 2014). H4: bầy đàn không xuất hiện ở tần suất dữ liệu (Christie và Huang, 1995).

{\itshape 3. Vì sao mũi tên nét đứt trong Hình 1 không mang giả thuyết riêng?\par}

Đáp án: vì định nghĩa của mức hối tiếc đã là chênh lệch kết quả danh mục.

{\itshape 4. Ba câu hỏi nghiên cứu ứng với bốn giả thuyết thế nào?\par}

Đáp án: Câu hỏi 1 với H1; Câu hỏi 2 với H2 và H4; Câu hỏi 3 với H3.

{\itshape 5. H4 có địa vị gì trong đề tài, và đề tài có quan sát hành vi chọn mô hình của tổ chức đầu tư không?\par}

Đáp án: H4 chỉ là thăm dò. Đề tài không quan sát hành vi chọn mô hình của tổ chức; độ phân tán chéo chỉ là phép thử giới hạn ở cấp thị trường.

\chapter{Mục 4.3.9 và 4.3.3 của đề cương: Quy trình năm bước, dữ liệu và phạm vi}

Hình 3 của đề cương tóm tắt năm bước: thu thập và xử lý dữ liệu; xác định chế độ thị trường; ước lượng tập bảy mô hình; đo lường và kiểm định H1, H2, H4; đánh giá cơ chế giảm thiểu và kiểm định H3. Phần này đi qua bản đồ năm bước rồi vào bước 1 (mục 4.3.3): dùng dữ liệu gì, lấy từ đâu, giới hạn ở cổ phiếu nào, và thiết kế ra sao để dữ liệu không tự “nói dối”.

Điều cần nhớ nhất của mục 4.3.3 là hai thiên lệch mà thiết kế rổ cổ phiếu nhắm tránh: thiên lệch sống sót và thiên lệch nhìn trước. Cả hai làm quá khứ trông đẹp hơn thực tế.

\section{Quy trình năm bước, đọc như một công thức nấu ăn}

Thí sinh đăng ký trước kế hoạch phân tích trước khi ước lượng chính thức tập mô hình ở bước 3. Giới hạn và khả năng thực hiện nằm ở mục 4.3.10.

Muốn biết chọn nhầm công thức nấu ăn thiệt bao nhiêu khi trời đổi, bạn mua nguyên liệu (bước 1), xem thời tiết từng ngày (bước 2), nấu cùng lúc bảy món theo bảy công thức (bước 3), so món được chọn với món ngon nhất rồi kiểm tra chênh lệch có thật không (bước 4), cuối cùng thử vài mẹo pha trộn để chữa (bước 5). Thứ tự không đảo được.

Bảng dưới ghép từng bước với mục chi tiết.

{\small\begin{longtable}[]{@{}
  >{\raggedright\arraybackslash}p{(\columnwidth - 4\tabcolsep) * \real{0.1237}}
  >{\raggedright\arraybackslash}p{(\columnwidth - 4\tabcolsep) * \real{0.7251}}
  >{\raggedright\arraybackslash}p{(\columnwidth - 4\tabcolsep) * \real{0.1512}}@{}}
\toprule\noalign{}
\begin{minipage}[b]{\linewidth}\raggedright
Bước
\end{minipage} & \begin{minipage}[b]{\linewidth}\raggedright
Việc chính
\end{minipage} & \begin{minipage}[b]{\linewidth}\raggedright
Mục trong đề cương
\end{minipage} \\
\midrule\noalign{}
\endhead
\bottomrule\noalign{}
\endlastfoot
1 & Thu thập và xử lý dữ liệu & 4.3.3 \\
2 & Xác định chế độ thị trường & 4.3.4 \\
3 & Ước lượng tập bảy mô hình (đã đăng ký trước) & 4.3.5 \\
4 & Đo lường và kiểm định H1, H2, H4 & 4.3.6 và 4.3.7 \\
5 & Đánh giá cơ chế giảm thiểu, kiểm định H3 & 4.3.8 \\
\end{longtable}}

\textbf{Cửa sổ ước lượng} là đoạn dữ liệu có độ dài cố định mà mô hình “thấy” khi tính trọng số, dời dần theo thời gian như khung ảnh trượt dọc cuộn phim. \textbf{Tập mô hình cạnh tranh} là bảy mô hình M1 đến M7 cùng đua trên một dữ liệu. \textbf{Cơ chế giảm thiểu} là cách làm giảm rủi ro lựa chọn mô hình, như kết hợp mô hình hay danh mục vững chắc.

Đề tài hoàn toàn định lượng trên dữ liệu thị trường công khai, không dùng khảo sát hay phương pháp định tính. Vì vậy không phỏng vấn hay bảng hỏi nào cứu được một bộ dữ liệu sai.

\section{Mục 4.3.3: dữ liệu gồm những gì}

Đề tài dùng bốn loại dữ liệu.

\subsection{Giá đóng cửa điều chỉnh theo ngày}

Loại thứ nhất là giá đóng cửa điều chỉnh theo ngày của cổ phiếu niêm yết trên HOSE, kể cả những mã đã hủy niêm yết, và của VN-Index.

Giá đóng cửa là giá cuối phiên. “Điều chỉnh” là sửa giá quá khứ theo các sự kiện của doanh nghiệp như chia cổ tức hay chia tách; đề cương không liệt kê sự kiện nào.

Cổ phiếu giá 100 hôm qua, hôm nay tách thành hai cổ mỗi cổ 50: không điều chỉnh thì bảng giá báo “rớt 50\% trong một đêm”, dù tài sản người nắm giữ không đổi. Giá điều chỉnh sửa những chỗ giả này.

Lợi suất là phần trăm tăng giảm của giá; mọi mô hình làm việc trên lợi suất.

\subsection{HOSE và VN-Index}

HOSE là Sở Giao dịch Chứng khoán Thành phố Hồ Chí Minh; phạm vi đề tài là cổ phiếu niêm yết trên HOSE (mục 4.4).

VN-Index là chỉ số chung của HOSE, như nhiệt kế đo “thân nhiệt” của sàn: không cho biết từng cổ phiếu, nhưng cho biết thị trường đang nóng hay lạnh. Đề tài dùng VN-Index để xác định chế độ (mục 4.3.4).

\subsection{Lãi suất phi rủi ro}

Loại thứ ba là lợi suất trái phiếu chính phủ kỳ hạn một năm, dùng làm lãi suất phi rủi ro.

Lãi suất phi rủi ro là mốc so sánh: khoản đầu tư rủi ro chỉ đáng giá khi kỳ vọng lời hơn mốc đó; phần vượt là lợi suất vượt trội. Mục 4.3.6 cho thấy lợi suất vượt trội có thể âm ở chế độ căng thẳng, và đó là lý do tỷ số Sharpe không làm thước đo chính.

\subsection{Khối lượng giao dịch}

Loại thứ tư là khối lượng giao dịch. Đề cương không nói cụ thể dùng vào đâu; rổ chọn theo giá trị giao dịch và độ phi thanh khoản Amihud có thể cần đến nó, nhưng đó là liên tưởng của bản đọc.

\section{Rổ 30 cổ phiếu và hai thiên lệch}

Tại mỗi ngày tái cân bằng, rổ gồm 30 mã có giá trị giao dịch bình quân ba tháng và vốn hóa tự do chuyển nhượng lớn nhất, xác định chỉ bằng dữ liệu có tại ngày đó để tránh thiên lệch sống sót và thiên lệch nhìn trước.

\textbf{Tái cân bằng} là định kỳ kéo tỷ lệ thực tế, đã lệch do giá đổi, về lại công thức, đồng thời có thể đổi cổ phiếu trong rổ. Danh mục tái cân bằng hằng tháng; kiểm tra độ bền dùng dữ liệu tháng và tái cân bằng hằng quý.

\textbf{Vốn hóa tự do chuyển nhượng} là giá trị thị trường của phần cổ phiếu thực sự giao dịch tự do, loại phần nằm lâu dài trong két của cổ đông lớn (định nghĩa thông thường; đề cương chỉ nêu tên).

Hai tiêu chí cho rổ vừa lớn vừa giao dịch sôi động. Mục 4.3.10 nói khả năng thực thi trên HOSE phụ thuộc thanh khoản, nên đề tài dùng rổ thanh khoản cao và kiểm tra độ nhạy theo chi phí.

\subsection{Thiên lệch sống sót}

Thiên lệch sống sót xảy ra khi dữ liệu chỉ gồm mã còn tồn tại đến hôm nay. Chỉ phỏng vấn những nhà đầu tư còn đứng vững thì quên những người làm y hệt đã phá sản; kết luận tất nhiên quá lạc quan.

\begin{tcolorbox}[breakable]
Ví dụ giả định (số tự đặt để minh họa): 10 cổ phiếu ở năm đầu mẫu, 3 mã sau này bị hủy niêm yết sau khi giá rơi mạnh.

Chỉ lấy 7 mã còn tồn tại ở cuối mẫu rồi quay ngược về quá khứ thì danh mục “lịch sử” không bao giờ chứa 3 mã tệ, và kết quả đẹp hơn điều nhà đầu tư thật đã trải qua.
\end{tcolorbox}

Đề cương tránh bằng hai cách: dữ liệu gồm cả những mã đã hủy niêm yết, và rổ dựng từ các mã đang giao dịch tại ngày tái cân bằng, không từ danh sách còn lại về sau.

\subsection{Thiên lệch nhìn trước}

Thiên lệch nhìn trước xảy ra khi phân tích dùng thông tin mà lúc quyết định chưa ai có: “kiểm tra” chiến lược mua hôm qua nhưng chọn mã bằng cách xem hôm nay mã nào đã tăng.

Cụm “chỉ bằng dữ liệu có tại ngày đó” chặn điều này. Mục 4.3.4 áp dụng cùng nguyên tắc cho chế độ: mô hình chuyển chế độ được ước lượng lại chỉ bằng dữ liệu đến ngày đó. “Ngày nào dùng dữ liệu ngày ấy” là sợi chỉ chạy suốt phần phương pháp.

\subsection{VN30 và VN100}

VN30 và VN100 là hai rổ chỉ số của HOSE. Đề tài chỉ dùng chúng để đối chiếu và kiểm tra độ bền, tính từ thời điểm HOSE công bố các rổ này; tác giả xác nhận thời điểm đó trong quý 1 năm thứ nhất.

Vì sao không dùng thẳng VN30 làm rổ chính? Đề cương không nói. Suy luận của bản đọc: các rổ này chỉ có từ khi HOSE công bố, nên dùng cho giai đoạn trước đó là biết trước thành phần.

\section{Giai đoạn mẫu và ba loại tần suất}

Giai đoạn mẫu dự kiến 2007--2026, tùy dữ liệu sẵn có, bao trùm khủng hoảng tài chính toàn cầu, đại dịch COVID-19 năm 2020 và căng thẳng trái phiếu doanh nghiệp từ quý IV/2022. Tác giả xác nhận nguồn và độ dài dữ liệu trong quý 1 năm thứ nhất, và thu hẹp rổ hoặc giai đoạn nếu không đủ (mục 4.3.10).

Ba sự kiện ứng với ba loại cú sốc ở mục 4.1. Chúng là giai đoạn tham chiếu để đối chiếu với chế độ căng thẳng do mô hình ước lượng; khủng hoảng là sự kiện lịch sử, chế độ căng thẳng là trạng thái thống kê.

\subsection{Phần mẫu tính được mức hối tiếc}

Mục 4.3.6 nói rõ: cần 104 tuần để ước lượng và 36 tháng kết quả ngoài mẫu để chọn mô hình, nên mức hối tiếc chỉ tính được từ khoảng năm thứ năm của mẫu. Với mẫu 2007--2026, đề tài dự kiến khoảng 28 cửa sổ đánh giá độc lập.

Ngoài mẫu là phần dữ liệu không dùng để “dạy” mô hình, chỉ để kiểm tra: học sinh ôn bằng đề cương rồi thi bằng đề mới. Trước khi ước lượng các mô hình danh mục, đề tài báo cáo số giai đoạn căng thẳng rơi vào phần mẫu tính được mức hối tiếc, tức cho biết có đủ “đề thi mùa bão” hay không.

Liên hệ với mục 4.3.10: ít giai đoạn căng thẳng độc lập thì công suất kiểm định thấp, và dưới 10 cửa sổ độc lập ở chế độ căng thẳng thì kết quả gọi là không kết luận.

Vì phần đầu mẫu chỉ dùng để ước lượng và chọn mô hình, khủng hoảng tài chính toàn cầu chủ yếu phục vụ ước lượng chế độ, không phải một “kỳ thi” cho bảy mô hình.

Các độ dài cửa sổ (104 tuần ước lượng, 6 tháng đánh giá, 36 tháng chọn mô hình) là giá trị dự kiến; tác giả biện minh sau khi ước lượng thử ở năm thứ nhất.

\subsection{Ba tần suất dữ liệu}

Ba tần suất cho ba mục đích:

{\small\begin{longtable}[]{@{}
  >{\raggedright\arraybackslash}p{(\columnwidth - 2\tabcolsep) * \real{0.5618}}
  >{\raggedright\arraybackslash}p{(\columnwidth - 2\tabcolsep) * \real{0.4382}}@{}}
\toprule\noalign{}
\begin{minipage}[b]{\linewidth}\raggedright
Việc cần làm
\end{minipage} & \begin{minipage}[b]{\linewidth}\raggedright
Tần suất
\end{minipage} \\
\midrule\noalign{}
\endhead
\bottomrule\noalign{}
\endlastfoot
Ước lượng chế độ thị trường & Dữ liệu ngày \\
Ước lượng hiệp phương sai & Dữ liệu tuần \\
Tái cân bằng danh mục & Hằng tháng \\
Kiểm tra độ bền & Dữ liệu tháng, tái cân bằng hằng quý \\
\end{longtable}}

Nhìn mỗi ngày để biết trời có đổi không; nhìn mỗi tuần để biết hai cổ phiếu có đi cùng nhau không mà đỡ nhiễu từng phiên; đổi cách chia tiền mỗi tháng vì mua bán thường xuyên tốn chi phí. Đề cương không nêu lý do từng lựa chọn; đây chỉ là cách hình dung.

Hiệp phương sai đo mức hai cổ phiếu cùng lên cùng xuống; ma trận hiệp phương sai gom mức đó cho mọi cặp trong rổ, để mô hình biết rủi ro có cộng dồn không.

\section{Những thứ chỉ là biến kiểm soát}

Chỉ số giá tiêu dùng, tỷ giá đô la Mỹ so với đồng Việt Nam và lãi suất điều hành chỉ là biến kiểm soát bổ sung trong kiểm tra độ bền, không nằm trong phân tích chính. Thêm chúng mà kết luận không đổi thì tự tin hơn; đổi thì biết kết luận nhạy với gì.

\section{Nguồn, giấy phép và điều chưa quy định}

Nguồn dữ liệu là HOSE và một nhà cung cấp dữ liệu thương mại (không nêu tên).

Luận án sẽ ghi kèm nguồn của biên độ dao động giá, mức thuế, phí giao dịch theo từng thời kỳ và điều khoản giấy phép dữ liệu; đề cương không liệt kê chúng.

Dữ liệu là dữ liệu thị trường công khai, không chứa thông tin cá nhân. Đề cương không nói thêm về thủ tục đạo đức.

Quý 1 năm thứ nhất (Bảng 5): xác nhận nguồn, độ dài và giấy phép dữ liệu, ngày HOSE công bố rổ VN30 và VN100.

\section{Phạm vi: những gì có trong đề tài và những gì không}

Phạm vi (mục 4.4): cổ phiếu niêm yết trên HOSE; cấu hình cơ sở là danh mục chỉ mua, không bán khống (vay cổ phiếu bán trước, mua lại sau); các ràng buộc khác nằm trong kiểm tra độ bền.

Đề tài không xét phái sinh, trái phiếu doanh nghiệp và tài sản ngoài cổ phiếu. Căng thẳng trái phiếu doanh nghiệp quý IV/2022 nằm trong mẫu như một sự kiện bối cảnh, nhưng trái phiếu không vào danh mục.

\section{Ghi nhớ khi đi thi}

\begin{tcolorbox}[breakable]
Năm bước: (1) dữ liệu, (2) chế độ, (3) ước lượng bảy mô hình, (4) đo lường và kiểm định H1, H2, H4, (5) cơ chế giảm thiểu và H3.

Dữ liệu: giá đóng cửa điều chỉnh theo ngày của cổ phiếu HOSE (kể cả mã đã hủy niêm yết) và VN-Index; lợi suất trái phiếu chính phủ một năm; khối lượng giao dịch.

Rổ: 30 mã lớn nhất theo giá trị giao dịch ba tháng và vốn hóa tự do chuyển nhượng, chỉ dùng dữ liệu có tại ngày đó, tránh thiên lệch sống sót và nhìn trước.

VN30 và VN100 chỉ để đối chiếu và kiểm tra độ bền, từ thời điểm HOSE công bố.

Mẫu 2007--2026; chế độ trên dữ liệu ngày, hiệp phương sai dữ liệu tuần, tái cân bằng tháng; mức hối tiếc tính từ khoảng năm thứ năm, khoảng 28 cửa sổ đánh giá.

Phạm vi: cổ phiếu HOSE, chỉ mua; không phái sinh, trái phiếu doanh nghiệp, tài sản ngoài cổ phiếu.
\end{tcolorbox}

\section{Tự kiểm tra}

{\itshape 1. Hai thiên lệch nào được rổ cổ phiếu nhắm tránh, và đề cương tránh bằng cách nào?\par}

Đáp án: thiên lệch sống sót và thiên lệch nhìn trước; dữ liệu gồm cả mã đã hủy niêm yết, và rổ xác định chỉ bằng dữ liệu có tại ngày tái cân bằng.

{\itshape 2. Rổ cổ phiếu gồm bao nhiêu cổ phiếu, chọn theo tiêu chí nào?\par}

Đáp án: 30 mã có giá trị giao dịch bình quân ba tháng và vốn hóa tự do chuyển nhượng lớn nhất tại mỗi ngày tái cân bằng.

{\itshape 3. VN30 và VN100 được dùng thế nào trong đề tài?\par}

Đáp án: chỉ để đối chiếu và kiểm tra độ bền, từ thời điểm HOSE công bố; thời điểm này xác nhận trong quý 1 năm thứ nhất.

{\itshape 4. Chế độ, hiệp phương sai và danh mục lần lượt dùng dữ liệu tần suất nào?\par}

Đáp án: chế độ trên dữ liệu ngày, hiệp phương sai trên dữ liệu tuần, tái cân bằng hằng tháng; kiểm tra độ bền dùng dữ liệu tháng và tái cân bằng hằng quý.

{\itshape 5. Vì sao mức hối tiếc chỉ tính được từ khoảng năm thứ năm của mẫu?\par}

Đáp án: vì cần 104 tuần để ước lượng và 36 tháng kết quả ngoài mẫu để chọn mô hình. Với mẫu 2007--2026 có khoảng 28 cửa sổ đánh giá; khủng hoảng tài chính toàn cầu chủ yếu phục vụ ước lượng chế độ.

\chapter{Mục 4.3.4 của đề cương: Xác định chế độ thị trường}

Đề tài cần biết mỗi ngày thị trường ở chế độ bình thường hay căng thẳng, mà không “ăn gian” bằng dữ liệu tương lai. Cách làm: mô hình chuyển chế độ Markov (Hamilton, 1989) cho lợi suất VN-Index, ước lượng lại tại mỗi ngày tái cân bằng chỉ bằng dữ liệu đến ngày đó, rồi đọc ra xác suất lọc, tức xác suất thị trường đang ở chế độ căng thẳng.

Chế độ căng thẳng là chế độ có phương sai điều kiện cao nhất và lợi suất kỳ vọng thấp hơn chế độ bình thường: một trạng thái thống kê, không phải tên một cuộc khủng hoảng.

\section{Vì sao phải xác định “thời tiết” của thị trường}

Hành vi thị trường đổi theo chế độ nên rủi ro lựa chọn mô hình có thể đổi theo (nhóm (4) của mục 4.2). Muốn so rủi ro ở hai chế độ, phải có cách đáng tin để biết ngày nào thuộc chế độ nào.

Nhà khí tượng chia ngày thành “quang” và “bão” rồi so thiệt hại của nông dân. Nếu cách chia đã lẫn kết quả vụ mùa (gọi một ngày là bão chỉ vì vụ mùa hôm đó tệ) thì phép so vô nghĩa. Cách xác định chế độ phải độc lập với kết quả danh mục.

Các nguồn nền (đã giới thiệu ở nhóm (4)): Hamilton (1989), Ang và Timmermann (2012), Ang và Bekaert (2002), Guidolin và Timmermann (2007), Kritzman và cộng sự (2012), Tu (2010).

\section{Mô hình chuyển chế độ Markov, giải thích từ đầu}

Mô hình giả định thị trường ở một trong vài trạng thái ẩn và chuyển giữa chúng theo xác suất. Bạn không thấy trạng thái, chỉ thấy lợi suất hằng ngày, và suy ra trạng thái từ cách lợi suất xử sự.

Ngồi trong phòng kín không cửa sổ, chỉ nghe tiếng động bên ngoài: trời quang thì nhỏ và đều, bão thì ồn và thất thường. Bạn đoán mỗi ngày nhiều khả năng là trạng thái nào; mô hình làm đúng việc đó với lợi suất VN-Index.

Mô hình có hai thành phần: mỗi trạng thái có lợi suất kỳ vọng và phương sai riêng; và luật chuyển, tức xác suất ở lại hoặc sang trạng thái khác. Đề cương không viết công thức.

\begin{tcolorbox}[breakable]
Ví dụ giả định (số tự đặt để minh họa): xác suất trời quang ở lại quang ngày kế tiếp là 95\%, bão ở lại bão là 80\%.

Khi đó các đợt bão kéo dài nhiều ngày nhưng ít xảy ra, đợt quang kéo dài hơn. Mô hình ước lượng các xác suất này từ dữ liệu; đề cương không đặt trước con số nào.
\end{tcolorbox}

\subsection{Số chế độ: hai hay ba}

Số chế độ (hai hoặc ba) do tiêu chí thông tin quyết định trên giai đoạn ước lượng ban đầu và cố định khi đăng ký trước.

Tiêu chí thông tin là thước điểm so các mô hình có độ phức tạp khác nhau. Thêm chế độ thì luôn khớp dữ liệu hơn, như áo nhiều nếp gấp luôn vừa người hơn; tiêu chí cộng một khoản “phạt” cho độ phức tạp để mô hình nhiều chế độ chỉ thắng khi cái lợi lớn hơn cái giá. Đề cương không nêu tên tiêu chí cụ thể.

Số chế độ cố định khi đăng ký trước, nên mọi thay đổi sau đó là lệch khỏi kế hoạch (mục 4.3.5). Theo bản đọc, nhờ vậy tác giả khó quay lại đổi số chế độ cho “đẹp” sau khi thấy kết quả.

\section{Tránh thiên lệch nhìn trước}

Nếu ước lượng trên toàn mẫu rồi gán nhãn hôm nay bằng tham số đó, nhãn hôm nay đã chịu ảnh hưởng của dữ liệu năm sau, điều nhà đầu tư lúc đó không thể biết.

Đề cương ước lượng lại tham số theo \textbf{cửa sổ mở rộng} tại mỗi ngày tái cân bằng, chỉ bằng dữ liệu đến ngày đó.

Cửa sổ mở rộng giữ điểm bắt đầu và kéo dài dần, như cuốn nhật ký dày thêm mỗi ngày; cửa sổ trượt giữ độ dài cố định và bỏ dữ liệu cũ, như tờ ghi chú chỉ có những ngày gần nhất.

Ngày tái cân bằng đầu tiên mô hình có ít dữ liệu, càng về sau càng nhiều; mỗi ngày mô hình chỉ dùng dữ liệu nhà đầu tư lúc đó có thể có.

Xác suất lọc tại ngày t là xác suất thị trường đang ở chế độ căng thẳng, tính chỉ bằng thông tin đến ngày t.

Sáng nay nghe tiếng động rồi nói “khả năng cao hôm nay là bão”: đó là xác suất lọc. Đợi cuối tuần nhìn lại cả tuần rồi nói “thứ Hai chắc chắn là bão” thì chính xác hơn nhưng không dùng được để quyết định hôm thứ Hai.

Xác suất lọc là một số từ 0 đến 1 mỗi ngày. Đề tài dùng nó để xây biến S\textsubscript{t} ở phương trình (5): tỷ lệ số ngày trong cửa sổ có xác suất lọc của chế độ căng thẳng vượt ngưỡng đăng ký trước. M5 và M7 cũng đổi hành vi khi xác suất căng thẳng tăng (mục 4.3.7).

\begin{tcolorbox}[breakable]
Ví dụ giả định (số tự đặt để minh họa): ngưỡng là 0,5; cửa sổ có 120 ngày, trong đó 30 ngày có xác suất lọc vượt 0,5.

Khi đó S\textsubscript{t} của cửa sổ bằng 30/120, tức một phần tư. Đề cương chưa nêu giá trị ngưỡng; ngưỡng được đăng ký trước.
\end{tcolorbox}

Phương sai điều kiện là mức dao động của lợi suất khi biết thị trường đang ở chế độ nào: độ ồn riêng của ngày bão khác độ ồn ngày quang, còn độ ồn trung bình cả năm là con số gộp.

Định nghĩa chế độ căng thẳng có hai vế, biến động cao nhất và lợi suất kỳ vọng thấp hơn, khớp hình ảnh mùa bão. Nếu mô hình có ba chế độ, chế độ căng thẳng vẫn là chế độ có phương sai điều kiện cao nhất; đề cương không nói cách gọi chế độ còn lại.

Khủng hoảng là sự kiện lịch sử (các cú sốc ở mục 4.1); chế độ căng thẳng là trạng thái thống kê. Đề tài đối chiếu hai thứ. Theo bản đọc, mô hình gọi một giai đoạn là căng thẳng mà không có sự kiện nào, hoặc ngược lại, là thông tin có giá trị chứ không phải lỗi cần che. Trong luận án, việc đối chiếu này nằm ở chương 3 (tiểu mục 3.2.2, Bảng 4).

\section{Bai-Perron: chỉ để mô tả}

Kiểm định của Bai và Perron (1998) xác định nhiều điểm gãy cấu trúc, như đánh dấu những chỗ dòng sông đổi dòng.

Trong đề tài, kiểm định này chạy trên toàn mẫu, chỉ để mô tả và đối chiếu với các cú sốc lịch sử, không dùng làm biến chế độ trong kiểm định giả thuyết.

Suy luận của bản đọc: chạy trên toàn mẫu là dùng cả dữ liệu tương lai, nên làm biến chế độ sẽ gây thiên lệch nhìn trước.

Tác giả đã dùng phương pháp Bai-Perron trong bản thảo số 5 (Bảng 7) về điểm gãy cấu trúc của thanh khoản thị trường chứng khoán Việt Nam.

\section{Định nghĩa thay thế: chính việc nhận diện chế độ cũng có rủi ro}

Kiểm tra độ bền dùng các định nghĩa chế độ thay thế theo ngưỡng sụt giảm và theo biến động thực hiện: thay vì để mô hình Markov gọi bão, gọi bão khi giá đã giảm quá một mức hoặc dao động thực tế vượt một mức. Kết luận không đổi thì vững hơn.

Mục 4.3.10 thừa nhận việc nhận diện chế độ tự nó là một dạng rủi ro mô hình. Đề tài nói về rủi ro mô hình nên không thể coi bước xác định chế độ là miễn nhiễm.

\section{Những gì mục 4.3.4 không nói}

\begin{itemize}
\item
  Không nói chế độ căng thẳng trùng với cuộc khủng hoảng nào.
\item
  Không nêu tên tiêu chí thông tin hay giá trị ngưỡng xác suất lọc; ngưỡng và số chế độ cố định khi đăng ký trước.
\item
  Không nói đã tìm thấy chế độ nào ở Việt Nam.
\item
  Không nói số giai đoạn căng thẳng có đủ không; con số này được báo cáo trước khi ước lượng các mô hình danh mục.
\end{itemize}

\section{Ghi nhớ khi đi thi}

\begin{tcolorbox}[breakable]
Mô hình chuyển chế độ Markov (Hamilton, 1989) cho lợi suất VN-Index, ước lượng lại theo cửa sổ mở rộng tại mỗi ngày tái cân bằng, chỉ bằng dữ liệu đến ngày đó.

Xác suất lọc tại ngày t: xác suất đang ở chế độ căng thẳng, chỉ dùng thông tin đến ngày t, nên không có thiên lệch nhìn trước.

Số chế độ (hai hoặc ba) được chọn bằng tiêu chí thông tin trên giai đoạn ước lượng ban đầu và cố định khi đăng ký trước.

Chế độ căng thẳng: phương sai điều kiện cao nhất, lợi suất kỳ vọng thấp hơn chế độ bình thường.

Bai-Perron (1998) trên toàn mẫu chỉ để mô tả và đối chiếu, không làm biến chế độ trong kiểm định giả thuyết.

Định nghĩa thay thế theo ngưỡng sụt giảm và biến động thực hiện nằm trong kiểm tra độ bền; khủng hoảng là sự kiện lịch sử, chế độ căng thẳng là trạng thái thống kê.
\end{tcolorbox}

\section{Tự kiểm tra}

{\itshape 1. Chế độ căng thẳng được định nghĩa thế nào?\par}

Đáp án: chế độ có phương sai điều kiện cao nhất và lợi suất kỳ vọng thấp hơn chế độ bình thường.

{\itshape 2. Đề cương tránh thiên lệch nhìn trước khi xác định chế độ bằng cách nào?\par}

Đáp án: ước lượng lại theo cửa sổ mở rộng tại mỗi ngày tái cân bằng, chỉ bằng dữ liệu đến ngày đó; xác suất lọc tại ngày t dùng các tham số này.

{\itshape 3. Số chế độ được chọn thế nào, và khi nào được cố định?\par}

Đáp án: hai hoặc ba, do tiêu chí thông tin quyết định trên giai đoạn ước lượng ban đầu, cố định khi đăng ký trước.

{\itshape 4. Kiểm định điểm gãy Bai-Perron đóng vai trò gì?\par}

Đáp án: chạy trên toàn mẫu, chỉ để mô tả và đối chiếu với các cú sốc lịch sử, không làm biến chế độ trong kiểm định giả thuyết.

{\itshape 5. Khủng hoảng và chế độ căng thẳng khác nhau ở đâu?\par}

Đáp án: khủng hoảng là sự kiện lịch sử; chế độ căng thẳng là trạng thái thống kê do mô hình ước lượng. Không dùng thay nhau; đề tài đối chiếu chúng.

\chapter{Mục 4.3.5 của đề cương: Bảy cách chia tiền cạnh tranh và đăng ký trước}

Mục 4.3.5 làm hai việc: chọn bảy mô hình danh mục (M1 đến M7) cạnh tranh trên cùng dữ liệu; và nói thẳng rằng tác giả không thể khẳng định chưa xem kết quả ở chế độ căng thẳng, nên cam kết đăng ký trước kế hoạch phân tích trên nền tảng OSF Registries vào quý 4 năm thứ nhất, ở chế độ khóa có dấu thời gian.

Đây là chỗ dễ bị hỏi nhất: vì sao bảy mô hình này, và làm sao tin tác giả không chọn phương pháp theo kết quả.

\section{Vì sao cần một “tập” mô hình cạnh tranh}

Muốn đo rủi ro của việc chọn giữa các mô hình thì phải có nhiều mô hình để chọn; một mô hình đơn lẻ không có “rủi ro lựa chọn”.

Cả lớp chỉ có một đội bóng thì không ai hỏi chọn đội nào. Bảy mô hình ở Bảng 3 là bảy đội cùng ra sân.

Đề cương nói tập này bao phủ các hướng chính của tài liệu: quy tắc đơn giản, co rút hiệp phương sai, tối ưu hóa cổ điển, Black-Litterman, danh mục vững chắc và chiến lược theo chế độ. Đề cương không giải thích vì sao đúng bảy chứ không sáu hay tám.

\section{Mô hình danh mục là gì: một “cách chia tiền”}

Mô hình danh mục là quy tắc cho biết, ở mỗi thời điểm, bỏ bao nhiêu phần trăm tiền vào từng cổ phiếu: một “cách chia tiền”.

Nhắc lại vì sao khó: khung trung bình-phương sai của Markowitz (1952) cần ước lượng lợi suất kỳ vọng và hiệp phương sai, sai số ước lượng truyền sang trọng số (Michaud, 1989; Kan và Zhou, 2007), và không hướng khắc phục nào thắng ổn định (DeMiguel và cộng sự, 2009). Bảng 3 gom đại diện của các hướng ấy.

\section{Bảng 3: tập mô hình danh mục cạnh tranh}

{\small\begin{longtable}[]{@{}
  >{\raggedright\arraybackslash}p{(\columnwidth - 4\tabcolsep) * \real{0.0991}}
  >{\raggedright\arraybackslash}p{(\columnwidth - 4\tabcolsep) * \real{0.5289}}
  >{\raggedright\arraybackslash}p{(\columnwidth - 4\tabcolsep) * \real{0.3720}}@{}}
\toprule\noalign{}
\begin{minipage}[b]{\linewidth}\raggedright
Mã
\end{minipage} & \begin{minipage}[b]{\linewidth}\raggedright
Mô hình
\end{minipage} & \begin{minipage}[b]{\linewidth}\raggedright
Nguồn
\end{minipage} \\
\midrule\noalign{}
\endhead
\bottomrule\noalign{}
\endlastfoot
M1 & Phân bổ đều 1/N & DeMiguel và cộng sự (2009) \\
M2 & Phương sai nhỏ nhất với hiệp phương sai co rút & Ledoit và Wolf (2004) \\
M3 & Trung bình-phương sai với ước lượng mẫu & Markowitz (1952) \\
M4 & Black-Litterman với quan điểm từ mô hình dự báo & Black và Litterman (1992) \\
M5 & Black-Litterman có chuyển chế độ & Mở rộng từ Nguyen và Nguyen (2026a) \\
M6 & Danh mục vững chắc với nhiều niềm tin tiên nghiệm & Garlappi và cộng sự (2007) \\
M7 & Giảm tỷ trọng rủi ro theo xác suất chế độ căng thẳng & Kritzman và cộng sự (2012) \\
\end{longtable}}

Mô tả dưới đây dựa vào tên, nguồn trong Bảng 3 và những gì mục 4.2 nói về từng nguồn; chỗ nào chỉ là giải thích khái niệm chung sẽ ghi rõ.

\subsection{M1: chia đều trứng vào N giỏ}

M1 là phân bổ đều 1/N: rổ 30 cổ phiếu thì mỗi mã được 1/30 số tiền. Không dự báo, không tính toán; chia đều trứng vào N giỏ vì không biết giỏ nào sẽ rơi.

Một quy tắc “ngây thơ” có mặt vì DeMiguel và cộng sự (2009) không thấy mô hình nào vượt nó nhất quán, và Pflug và cộng sự (2012) cho thấy nó xấp xỉ tối ưu khi mơ hồ cao. M1 là chuẩn so sánh tự nhiên và gắn với khả năng đối lập của H1.

\subsection{M2: dao động nhỏ nhất, với số liệu được co rút}

M2 là phương sai nhỏ nhất với ma trận hiệp phương sai co rút (Ledoit và Wolf, 2004). Phương sai nhỏ nhất là chia tiền sao cho danh mục dao động ít nhất; không cần ước lượng lợi suất kỳ vọng, chỉ cần biết các cổ phiếu dao động và đi cùng nhau ra sao.

Co rút kéo ước lượng hiệp phương sai từ mẫu về một mốc ổn định hơn, như không tin ngay học sinh luôn giỏi tuyệt đối sau một bài thi thử mà kéo nhận định về phía trung bình lớp (nhóm (1)). Chi tiết kỹ thuật đề cương không nêu.

Ở Việt Nam, Nguyen và cộng sự (2020) dùng hiệp phương sai co rút và thu được kết quả tốt, nhưng đánh giá riêng một phương pháp; đề tài đặt nó vào tập cạnh tranh. Đề cương không nói họ dùng đúng phương pháp của Ledoit và Wolf.

\subsection{M3: bản gốc của Markowitz, dùng số liệu mẫu nguyên xi}

M3 là trung bình-phương sai với ước lượng mẫu (Markowitz, 1952): dùng chính trung bình và hiệp phương sai tính từ mẫu, không làm mượt.

Đây là cách chịu ảnh hưởng nặng nhất của sai số ước lượng (Michaud, 1989; Kan và Zhou, 2007): chia tiền theo nguyên văn điểm thi thử, mã nào điểm cao nhất nhận nhiều nhất, kể cả khi điểm cao chủ yếu do may.

Quy định riêng: chỉ ước lượng trực tiếp M3 khi tỷ lệ số cổ phiếu trên số quan sát tuần không vượt 0,5; vượt ngưỡng thì dùng nghịch đảo giả Moore-Penrose và ghi rõ trong kết quả. Hình dung: quá ít điểm số so với số học sinh cần xếp hạng thì phép tính chuẩn không chạy được, phải dùng công cụ thay thế.

\subsection{M4: Black-Litterman, mốc thị trường cộng quan điểm}

M4 là Black-Litterman với quan điểm từ mô hình dự báo (Black và Litterman, 1992): kết hợp lợi suất cân bằng của mô hình định giá tài sản vốn với quan điểm của nhà đầu tư. Bessler và cộng sự (2017) thấy cách này vượt trung bình-phương sai và 1/N ngoài mẫu trên dữ liệu đa tài sản toàn cầu.

Bắt đầu từ “mức giá hợp lý chung của chợ” rồi pha ý kiến riêng, mạnh yếu tùy độ tự tin.

Quan điểm của M4 lấy từ một mô hình dự báo đơn giản, cố định khi đăng ký trước. Ta và Vuong (2020) dùng mô hình tự hồi quy trung bình trượt tích hợp; đề cương không nói M4 dùng mô hình nào.

\subsection{M5: Black-Litterman có chuyển chế độ, do tác giả phát triển}

M5 là Black-Litterman có chuyển chế độ, mở rộng từ Nguyen và Nguyen (2026a), bài áp dụng Black-Litterman với mô hình định giá tài sản vốn có chuyển chế độ cho Việt Nam 2019--2025.

M4 giữ cách trộn như nhau suốt năm; M5 đổi cách trộn theo thời tiết, nên M5 đổi hành vi khi xác suất căng thẳng tăng. Đề cương không mô tả chi tiết cách đổi.

Vì M5 do tác giả phát triển, đề cương xử lý nó như mọi mô hình khác và báo cáo kết quả cả khi có lẫn khi không có M5. Nếu M5 làm nên kết quả, người đọc sẽ thấy.

Tác giả đã dùng mô hình định giá tài sản vốn có chuyển chế độ và phương pháp tính toán Bayes xấp xỉ kết hợp Monte Carlo chuỗi Markov trong bài trước. Ước lượng thử ở quý 2 năm thứ nhất quyết định giữ phương pháp này hay thay bằng một phương pháp nhanh hơn cho M5 (mục 4.3.10).

Kết quả bài trước chỉ làm đối chứng, không báo cáo lại như kết quả mới.

\subsection{M6: chịu được tình huống xấu nhất trong nhiều niềm tin}

M6 là danh mục vững chắc với nhiều niềm tin tiên nghiệm (Garlappi và cộng sự, 2007): nhà đầu tư né tránh mơ hồ, tối ưu cho tình huống xấu nhất trong nhiều niềm tin, và được danh mục ổn định hơn.

Không chắc nên tin dự báo nào, bạn chia tiền sao cho ngay cả khi dự báo xấu nhất thành thật, thiệt hại vẫn nhỏ nhất: tinh thần của tiêu chí Gilboa và Schmeidler (1989) ở nhóm (2); cách nối này là của bản đọc.

\subsection{M7: giảm tỷ trọng rủi ro theo xác suất chế độ căng thẳng}

M7 theo Kritzman và cộng sự (2012), những người dự báo chế độ bằng mô hình Markov và thấy chiến lược động vượt chiến lược tĩnh khi kiểm định ngược.

M7 như người lái xe giảm ga khi thấy trời sắp bão: xác suất căng thẳng tăng thì danh mục giảm tỷ trọng rủi ro. Đề cương không mô tả chi tiết cách giảm.

\subsection{Mô hình tĩnh và mô hình đổi theo chế độ}

Năm mô hình tĩnh M1, M2, M3, M4, M6 không đổi hành vi theo xác suất căng thẳng; M5 và M7 chỉ đổi hành vi khi xác suất căng thẳng tăng. Đề cương lặp lại kiểm định H1 trên tập chỉ gồm năm mô hình tĩnh để loại tác động của M5 và M7 (mục 4.3.7).

Nếu ở mùa bão mô hình động khác hẳn mô hình tĩnh chỉ vì nó được thiết kế để phản ứng với bão, thì “rủi ro lựa chọn lớn hơn ở chế độ căng thẳng” có thể là hệ quả của thiết kế chứ không phải của thị trường. Bỏ M5 và M7 ra là tách điều đó khỏi kết quả.

\section{Cách các mô hình được ước lượng và so sánh công bằng}

Mỗi mô hình được ước lượng trên cửa sổ 104 tuần và đánh giá trên cửa sổ 6 tháng không chồng lấn; cửa sổ trượt hằng tháng chỉ dùng cho phân tích phụ (mục 4.3.6, Hình 2). Các độ dài là dự kiến, biện minh sau ước lượng thử ở năm thứ nhất.

Ước lượng là “dạy” mô hình bằng 104 tuần quá khứ; đánh giá là cho mô hình “thi” trên 6 tháng tiếp theo. Không chồng lấn để các kỳ thi độc lập, nên số cửa sổ độc lập đếm được không phồng lên vì các cửa sổ lặp nhau.

\subsection{Cùng một ngân sách tinh chỉnh}

Siêu tham số của M2, M4, M5, M6 và M7 được chọn từ một lưới cho trước, bằng cùng một quy tắc và cùng ngân sách tinh chỉnh, rồi cố định khi đăng ký trước.

Siêu tham số là “núm vặn” phải đặt trước khi mô hình học dữ liệu, ví dụ mức co rút. Vặn kỹ cho mô hình yêu thích và qua loa với mô hình khác thì cuộc đua không công bằng; cùng ngân sách tinh chỉnh cho mọi mô hình cơ hội như nhau, và khóa vào đăng ký trước chặn việc vặn lại sau khi thấy kết quả. Suy đoán của bản đọc: M1 và M3 không nằm trong danh sách vì M1 chia đều, M3 dùng ước lượng mẫu nguyên xi.

Để đo mức hối tiếc, cần một mô hình chọn trước: tại đầu mỗi cửa sổ đánh giá, đó là mô hình có lợi suất tương đương chắc chắn ngoài mẫu cao nhất trong 36 tháng liền trước; mô hình được chọn tái cân bằng hằng tháng suốt cửa sổ (mục 4.3.6).

Đây là quy tắc mô phỏng của đề tài, chọn theo thành tích gần nhất; đề cương không nói các tổ chức thật chọn như vậy.

\section{Đăng ký trước}

Đăng ký trước là viết dự đoán và cách kiểm tra vào một “phong bì” có dấu thời gian trước khi xem dữ liệu mới. Trước trận đấu, bạn đưa trọng tài phong bì ghi đội nào thắng, tỷ số bao nhiêu; sau trận ai cũng mở được và biết bạn không sửa theo kết quả.

“Phong bì” ở đây là kế hoạch phân tích trên OSF Registries, nơi lưu các đăng ký trước; “khóa” nghĩa là không sửa được nội dung, dấu thời gian là bằng chứng về thời điểm.

\subsection{Điều đề cương thừa nhận}

Tác giả đã dùng dữ liệu Việt Nam 2019--2025, gồm hai giai đoạn căng thẳng, trong Nguyen và Nguyen (2026a), nên không thể khẳng định chưa từng thấy kết quả ở chế độ căng thẳng.

Đã thấy đề thi năm ngoái rồi mới viết “dự đoán” cho năm nay: nói là dự đoán trước khi xem đề thì sai; trung thực là khai “tôi đã thấy đề năm ngoái, đây là những phần tôi đã thấy”.

Đề cương làm đúng vậy: đăng ký trước và khai rõ những kết quả đã xem. Giấu điều này thì người đọc tin kết quả hơn mức cho phép; khai ra làm đề tài kém hoàn hảo nhưng cho người đọc thước đo chính xác về độ tin cậy.

Đề cương không nói tác giả đã xem kết quả cụ thể nào; đừng suy đoán.

\subsection{Tệp đăng ký cố định những gì}

Tệp đăng ký cố định:

\begin{itemize}
\item
  độ dài các cửa sổ;
\item
  siêu tham số của bảy mô hình;
\item
  số chế độ và ngưỡng xác suất lọc;
\item
  mức ý nghĩa;
\item
  nhóm kết quả chính;
\item
  quy tắc hiệu chỉnh so sánh bội;
\item
  quy tắc gọi kết quả là không kết luận.
\end{itemize}

Mỗi mục là một “núm vặn” mà người phân tích có thể bị cám dỗ vặn lại sau khi thấy dữ liệu; khóa trước là tự cột tay mình. Tiêu chí ủng hộ của từng giả thuyết (mục 4.3.7--4.3.8) cũng là phần của kế hoạch.

\subsection{Mẫu kiểm tra riêng và “lệch khỏi kế hoạch”}

Dữ liệu phát sinh sau ngày đăng ký là mẫu kiểm tra riêng, và mọi thay đổi sau đó được ghi là lệch khỏi kế hoạch.

Sau khi niêm phong, những trận chưa từng thấy là “sân thật”; mọi chỉnh sửa sau đó, dù chính đáng, phải ghi đúng tên là chỉnh sửa.

\subsection{Đăng ký trước giảm chứ không loại thiên lệch}

Đề cương không tuyên bố đăng ký trước xóa sạch thiên lệch. Diễn giải của bản đọc: tác giả đã thấy dữ liệu 2019--2025, nên lựa chọn thiết kế có thể đã vô tình chịu ảnh hưởng; đăng ký trước hạn chế việc điều chỉnh về sau, đi kèm việc thừa nhận giới hạn đó.

\subsection{Lịch trình đăng ký}

Theo Bảng 5: quý 2 năm thứ nhất ước lượng thử bảy mô hình trên 2019--2025 để đo thời gian tính và số quan sát ở chế độ căng thẳng, quyết định phương pháp ước lượng cho M5; quý 4 đăng ký trước.

Các quyết định từ ước lượng thử (phương pháp cho M5, biện minh độ dài cửa sổ) được chốt trước quý 4, rồi mới khóa vào đăng ký. Đề cương không nói ước lượng thử còn ảnh hưởng đến lựa chọn nào khác.

\section{Những gì mục 4.3.5 không nói}

\begin{itemize}
\item
  Mục này không nói bảy mô hình bao quát mọi cách chia tiền có thể có; đây là tập do đề tài chọn.
\item
  Mục này không khẳng định mô hình nào sẽ thắng. Đây là đề cương, không có kết quả.
\item
  Mục này không khẳng định tác giả chưa xem kết quả. Chính mục này nói ngược lại.
\item
  Mục này không nói đăng ký trước loại hết thiên lệch.
\end{itemize}

\section{Ghi nhớ khi đi thi}

\begin{tcolorbox}[breakable]
Bảy mô hình: M1 1/N (DeMiguel và cộng sự, 2009); M2 phương sai nhỏ nhất với hiệp phương sai co rút (Ledoit và Wolf, 2004); M3 trung bình-phương sai với ước lượng mẫu (Markowitz, 1952); M4 Black-Litterman với quan điểm từ mô hình dự báo (Black và Litterman, 1992); M5 Black-Litterman có chuyển chế độ (mở rộng từ Nguyen và Nguyen, 2026a); M6 danh mục vững chắc với nhiều niềm tin tiên nghiệm (Garlappi và cộng sự, 2007); M7 giảm tỷ trọng rủi ro theo xác suất căng thẳng (Kritzman và cộng sự, 2012).

M5 do tác giả phát triển, xử lý như các mô hình khác, báo cáo cả khi có lẫn khi không có M5; H1 lặp lại trên tập mô hình tĩnh M1, M2, M3, M4, M6.

Ước lượng 104 tuần, đánh giá 6 tháng không chồng lấn, chọn trước theo 36 tháng liền trước; các độ dài là dự kiến.

Đã dùng dữ liệu 2019--2025 (hai giai đoạn căng thẳng) trong bài trước, nên không thể khẳng định chưa xem kết quả ở chế độ căng thẳng.

Đăng ký trước trên OSF Registries quý 4 năm thứ nhất, khóa có dấu thời gian, khai rõ kết quả đã xem; dữ liệu sau ngày đăng ký là mẫu kiểm tra riêng; thay đổi ghi là lệch khỏi kế hoạch.
\end{tcolorbox}

\section{Tự kiểm tra}

{\itshape 1. Kể tên bảy mô hình M1 đến M7.\par}

Đáp án: M1 1/N; M2 phương sai nhỏ nhất với hiệp phương sai co rút; M3 trung bình-phương sai với ước lượng mẫu; M4 Black-Litterman với quan điểm từ mô hình dự báo; M5 Black-Litterman có chuyển chế độ; M6 danh mục vững chắc với nhiều niềm tin tiên nghiệm; M7 giảm tỷ trọng rủi ro theo xác suất chế độ căng thẳng.

{\itshape 2. Vì sao tác giả không thể khẳng định chưa xem kết quả ở chế độ căng thẳng?\par}

Đáp án: vì đã dùng dữ liệu Việt Nam 2019--2025, gồm hai giai đoạn căng thẳng, trong Nguyen và Nguyen (2026a).

{\itshape 3. Tệp đăng ký trước cố định những gì?\par}

Đáp án: độ dài các cửa sổ, siêu tham số của bảy mô hình, số chế độ và ngưỡng xác suất lọc, mức ý nghĩa, nhóm kết quả chính, quy tắc hiệu chỉnh so sánh bội, quy tắc gọi kết quả là không kết luận.

{\itshape 4. Dữ liệu phát sinh sau ngày đăng ký được dùng thế nào, và thay đổi sau đó được ghi ra sao?\par}

Đáp án: là mẫu kiểm tra riêng; mọi thay đổi sau đó ghi là lệch khỏi kế hoạch.

{\itshape 5. M5 có gì đặc biệt, và đề cương xử lý ra sao?\par}

Đáp án: M5 do tác giả phát triển (mở rộng từ Nguyen và Nguyen, 2026a); được xử lý như các mô hình khác, báo cáo cả khi có lẫn khi không có M5.

\chapter{Mục 4.3.6 của đề cương: Chọn thước chấm điểm, vì sao lợi suất tương đương chắc chắn là thước chính và Sharpe chỉ là thước phụ}

Trước mọi phép đo rủi ro phải trả lời: lấy gì làm điểm số cho mỗi mô hình? Đề cương chọn lợi suất tương đương chắc chắn sau chi phí giao dịch, đo ngoài mẫu. Tỷ số Sharpe chỉ là thước phụ, vì khi lợi suất vượt trội trung bình âm (đúng tình huống của chế độ căng thẳng), danh mục biến động hơn lại có Sharpe ít âm hơn. Trên điểm số này, đề cương dựng ba thước đo: độ phân tán mô hình, mức hối tiếc và kích thước tập tin cậy mô hình (Hansen và cộng sự, 2011).

Hãy đọc chậm: mọi công thức về sau đứng trên điểm số này. Hiểu nó thì hiểu vì sao mức hối tiếc được tính như vậy, vì sao phải chia cho sai số chuẩn bootstrap, và vì sao phương trình (5) có dạng như vậy.

\section{Bảy cách chia tiền cần một cách chấm điểm chung}

Tập M gồm bảy mô hình ở Bảng 3, mỗi mô hình là một cách chia tiền vào cổ phiếu.

Bảy đội thi chạy cần một cách chấm chung; chấm mỗi đội một kiểu thì mọi so sánh vô nghĩa.

Mục 4.3.6 làm ba việc: chọn điểm số cho từng mô hình trên một cửa sổ đánh giá; từ bảy điểm số định nghĩa ba thước đo rủi ro lựa chọn mô hình; quy định các cửa sổ thời gian. Phần này và hai phần kế tiếp đi qua cả ba.

Rủi ro lựa chọn mô hình là tổn thất kỳ vọng do phải chọn trước khi biết mô hình nào tốt nhất, đo bằng mức hối tiếc. Muốn nói “tốt nhất” hay “thiệt bao nhiêu” thì phải có điểm số.

\section{Lợi suất tương đương chắc chắn là gì}

Bạn được chọn: nhận chắc một khoản tiền, hoặc tham gia một khoản đầu tư lúc lãi nhiều lúc lỗ. Khoản tiền chắc chắn tối thiểu khiến bạn bỏ qua khoản đầu tư rủi ro chính là giá trị tương đương chắc chắn của nó; tính theo tỷ lệ thì là lợi suất tương đương chắc chắn.

Đó là giá bán lại một tấm vé số rủi ro, tính bằng tiền mặt chắc chắn. Người rất ngại rủi ro chấp nhận ít tiền mặt hơn để đổi tấm vé; cùng một tấm vé, con số khác nhau tùy mức ngại rủi ro.

Càng cao càng tốt: sau khi trừ sự khó chịu vì dao động, mô hình vẫn để lại nhiều giá trị. Đề cương tính trên lợi suất sau chi phí giao dịch, nên mô hình mua bán lại quá nhiều bị trừ điểm vì phí; và đo ngoài mẫu, tức trên khoảng thời gian mô hình chưa thấy khi được ước lượng.

Đề cương tính từ trung bình và phương sai lợi suất trên cửa sổ đánh giá kết thúc tại t, không nói quy ra năm. Các ví dụ dưới chỉ dùng đơn vị minh họa.

\section{Phương trình (1), từng ký hiệu một}

Chỉ số m chỉ mô hình, t chỉ thời điểm kết thúc cửa sổ đánh giá.

\(U_{m,t}\  = \ \mu_{m,t}\  - \ \frac{\rho}{2}\sigma_{m,t}^{2}\) (1)

\(U_{m,t}\) là điểm của mô hình m trên cửa sổ kết thúc tại t (chữ U gợi “mức thỏa dụng”). Cùng một mô hình có điểm khác nhau ở các cửa sổ khác nhau.

\(M\) là tập bảy mô hình ở Bảng 3; chữ M in hoa chỉ tập, chữ m thường chỉ một mô hình.

\(\mu_{m,t}\) là trung bình lợi suất danh mục của mô hình m trên cửa sổ, sau chi phí giao dịch: phần thưởng bình quân. Chữ $\mu$ (đọc là “muy”) là ký hiệu quen thuộc của trung bình.

\(\sigma_{m,t}^{2}\) là phương sai lợi suất trên cùng cửa sổ, sau chi phí: thước đo cảm giác “xóc” trên đường đi.

\(\rho\) là hệ số ngại rủi ro, ấn định trước: giá trị cơ sở 5, kiểm tra độ nhạy ở 3 và 10. Số \(\rho\) càng lớn, bạn càng ghét dao động.

Bằng lời: lấy phần thưởng bình quân, trừ một khoản phạt bằng nửa hệ số ngại rủi ro nhân phương sai. Khoản phạt là cái giá của dao động; phần còn lại là lợi suất tương đương chắc chắn. Ba bước: nhân phương sai với \(\rho\) rồi chia đôi; lấy trung bình trừ khoản đó; ghi kết quả làm điểm của mô hình m ở cửa sổ t.

Hai tuyến xe đến cùng một nơi: tuyến A nhanh hơn nhưng xóc dữ, tuyến B chậm hơn chút mà êm. Điểm tổng hợp lấy tốc độ trung bình trừ phần mệt do xóc, mức trừ tùy độ say xe. Người say nặng (\(\rho\) lớn) chọn B; người chịu xóc giỏi (\(\rho\) nhỏ) có thể chọn A.

\begin{tcolorbox}[breakable]
Ví dụ giả định (số tự đặt để minh họa): ba mô hình trên một cửa sổ, lãi suất phi rủi ro coi bằng 0, \(\rho\) bằng 5 nên khoản phạt là 2,5 lần phương sai.

A: trung bình 8\%, độ lệch chuẩn 20\%, phương sai 0,04; phạt 0,10; điểm 0,08 − 0,10 = −2\%.

B: trung bình 6\%, độ lệch chuẩn 10\%, phương sai 0,01; phạt 0,025; điểm 3,5\%.

C: trung bình 2\%, độ lệch chuẩn 5\%, phương sai 0,0025; phạt 0,00625; điểm khoảng 1,4\%.

Xếp hạng là B, C, A. A có lợi suất trung bình cao nhất nhưng đứng cuối vì dao động quá lớn so với mức ngại rủi ro đã chọn.
\end{tcolorbox}

Điểm số này không phải lợi suất, và hạng theo nó khác hạng theo lợi suất. Một mô hình có thể “thắng” về lợi suất chỉ nhờ chịu rủi ro cao hơn; với nhà đầu tư ngại rủi ro, chiến thắng đó không đủ giá trị.

\section{Hệ số ngại rủi ro: cơ sở 5, thử thêm 3 và 10}

Điểm số phụ thuộc \(\rho\), nên xếp hạng có thể đổi khi \(\rho\) đổi. Đề cương ấn định \(\rho\) bằng 5 làm cơ sở rồi thử độ nhạy ở 3 và 10, để xem kết luận có đứng vững không. Đề cương không nêu lý do riêng của con số 5.

Nhớ ba con số và vai trò: 5 là cơ sở; 3 và 10 là hai cực để thử độ nhạy.

Ví dụ dưới chỉ chứng minh điểm số có thể đổi hạng theo \(\rho\).

\begin{tcolorbox}[breakable]
Ví dụ giả định (số tự đặt để minh họa): A có trung bình 10\%, phương sai 0,04; B có trung bình 5\%, phương sai 0,01.

\(\rho\) bằng 3 (phạt 1,5 lần phương sai): A = 0,10 − 0,06 = 0,04; B = 0,05 − 0,015 = 0,035. A trên.

\(\rho\) bằng 5 (phạt 2,5 lần): A = 0,10 − 0,10 = 0; B = 0,05 − 0,025 = 0,025. B trên.

\(\rho\) bằng 10 (phạt 5 lần): A = 0,10 − 0,20 = −0,10; B = 0,05 − 0,05 = 0. B trên, cách xa hơn.
\end{tcolorbox}

Thứ hạng đổi giữa \(\rho\) bằng 3 và \(\rho\) bằng 5. Nếu kết luận chỉ đúng ở một mức \(\rho\) và đổi chiều ở mức khác, người đọc cần biết; đó là ích lợi của việc thử độ nhạy theo \(\rho\).

\section{Vì sao tỷ số Sharpe chỉ là thước đo phụ}

Tỷ số Sharpe lấy lợi suất vượt trội trung bình (phần vượt lãi suất phi rủi ro) chia cho độ biến động: mỗi đơn vị rủi ro mang lại bao nhiêu phần thưởng. Thường thì càng cao càng tốt.

Vấn đề nằm ở chữ “càng cao càng tốt”. Khi lợi suất vượt trội trung bình âm, đúng tình huống của chế độ căng thẳng, quan hệ “danh mục tốt hơn thì điểm cao hơn” không còn.

Lý do: tử số âm, chia cho độ biến động lớn hơn thì kết quả gần 0 hơn, tức “đẹp hơn”. Danh mục xóc hơn được chấm cao hơn dù lợi suất vượt trội như nhau.

Hai học sinh cùng bị trừ 6 điểm. Cách chấm chia số điểm bị trừ cho độ ồn của lớp mỗi em ngồi; em ở lớp ồn hơn bị chia cho số lớn hơn nên trông đỡ tệ hơn. Chính cái ồn giúp em trông tốt hơn.

\begin{tcolorbox}[breakable]
Ví dụ giả định (số tự đặt để minh họa): lãi suất phi rủi ro coi bằng 0, \(\rho\) bằng 5. Trong một cửa sổ căng thẳng, hai mô hình đều có trung bình −6\%.

A: độ lệch chuẩn 30\%, phương sai 0,09; Sharpe = −0,06 / 0,30 = −0,2.

B: độ lệch chuẩn 10\%, phương sai 0,01; Sharpe = −0,06 / 0,10 = −0,6.

Theo Sharpe, A (−0,2) có vẻ tốt hơn B (−0,6), dù A dao động gấp ba mà lỗ như nhau.

Theo điểm tương đương chắc chắn với \(\rho\) bằng 5: A = −0,06 − 2,5 × 0,09 = −0,285; B = −0,06 − 2,5 × 0,01 = −0,085. B tốt hơn, đúng trực giác.
\end{tcolorbox}

Khi lợi suất như nhau, thước đo này trừ nhiều hơn cho mô hình dao động mạnh, nên mô hình êm hơn luôn điểm cao hơn; Sharpe làm ngược lại khi tử số âm.

Vì vậy đề cương chỉ báo cáo Sharpe như thước đo phụ, ở các cửa sổ có lợi suất vượt trội trung bình dương.

Sharpe không bị bỏ: mục 4.3.7 vẫn so sánh Sharpe giữa các mô hình bằng khoảng tin cậy bootstrap và dùng tỷ số Sharpe hiệu chỉnh để hiệu chỉnh thiên lệch chọn lọc.

\section{Thước đo này trong dòng suy luận chung}

Theo bản đọc, có hai điểm nối. Thứ nhất, DeMiguel và cộng sự (2009) cũng xét lợi suất tương đương chắc chắn bên cạnh Sharpe và vòng quay, nên đây là thước đo có sẵn trong dòng nghiên cứu mà đề cương dẫn.

Thứ hai, thước đo này đặt trên kết quả của quyết định phân bổ, không trên dự báo rủi ro, nên các thước đo dựng trên nó nối được với hậu quả của quyết định, đúng khoảng trống (1).

Nhớ cách dùng thuật ngữ: độ phân tán mô hình và kích thước tập tin cậy chỉ báo bất định; chỉ mức hối tiếc đo rủi ro lựa chọn.

\section{Phương trình (2): điểm trung bình của bảy mô hình}

Công thức thứ hai ngắn nhưng làm nền cho công thức thứ ba.

\({\overline{U}}_{t}\  = \ \frac{1}{|M|}\sum_{m \in M}^{}U_{m,t}\) (2)

\({\overline{U}}_{t}\) là trung bình điểm của bảy mô hình trên cùng cửa sổ; gạch ngang trên đầu là ký hiệu “trung bình”. Cộng điểm của mọi mô hình m trong \(M\) rồi chia cho 7.

Bảy đội có bảy điểm, điểm trung bình lớp cho biết mức chung, chưa cho biết các đội chênh nhau bao nhiêu. Phương trình (3) đo điều đó.

\begin{tcolorbox}[breakable]
Ví dụ giả định (số tự đặt để minh họa): điểm của bảy mô hình trên một cửa sổ là 2,0; 3,0; −2,0; 1,0; 2,6; 3,5 và 1,4 (phần trăm).

Tổng 11,5, chia 7 được khoảng 1,64\%.
\end{tcolorbox}

Mọi mô hình có trọng số ngang nhau khi tính trung bình; đề cương không nêu lý do, nhưng tính đều là tự nhiên khi mục đích là xem các mô hình khác nhau bao nhiêu.

\section{Ba thước đo rủi ro, nhìn trước}

Bảng dưới là bản đồ ba thước đo; phần kế tiếp giải thích đầy đủ.

{\small\begin{longtable}[]{@{}
  >{\raggedright\arraybackslash}p{(\columnwidth - 6\tabcolsep) * \real{0.1577}}
  >{\raggedright\arraybackslash}p{(\columnwidth - 6\tabcolsep) * \real{0.1777}}
  >{\raggedright\arraybackslash}p{(\columnwidth - 6\tabcolsep) * \real{0.4933}}
  >{\raggedright\arraybackslash}p{(\columnwidth - 6\tabcolsep) * \real{0.1713}}@{}}
\toprule\noalign{}
\begin{minipage}[b]{\linewidth}\raggedright
Thước đo
\end{minipage} & \begin{minipage}[b]{\linewidth}\raggedright
Ký hiệu
\end{minipage} & \begin{minipage}[b]{\linewidth}\raggedright
Đo cái gì
\end{minipage} & \begin{minipage}[b]{\linewidth}\raggedright
Vai trò trong H1
\end{minipage} \\
\midrule\noalign{}
\endhead
\bottomrule\noalign{}
\endlastfoot
Mức hối tiếc & \(R_{t}\) & Khoản thiệt vì mô hình chọn trước không phải mô hình tốt nhất sau thực tế & Kết quả chính \\
Độ phân tán mô hình & \(D_{t}\) & Các mô hình chênh nhau nhiều hay ít & Kết quả phụ \\
Kích thước tập tin cậy mô hình & không có ký hiệu riêng & Số mô hình chưa thể loại sau vòng sơ loại thống kê & Kết quả phụ \\
\end{longtable}}

Độ phân tán và mức hối tiếc sinh ra từ cùng tập điểm số, nên liên quan nhưng không thay nhau. Mức hối tiếc trả lời “tôi thiệt bao nhiêu vì đã chọn mô hình ấy”; độ phân tán trả lời “các mô hình khác nhau bao nhiêu”; tập tin cậy trả lời “còn bao nhiêu mô hình chưa thể loại”.

\begin{tcolorbox}[breakable]
\textbf{Ghi nhớ khi đi thi}

• Điểm số mỗi mô hình: lợi suất tương đương chắc chắn sau chi phí, đo ngoài mẫu: \(U_{m,t}\) = \(\mu_{m,t}\) − (\(\rho\)/2) × \(\sigma_{m,t}^{2}\).

• Hệ số ngại rủi ro \(\rho\): cơ sở 5, độ nhạy 3 và 10; tập M có 7 mô hình.

• Khi lợi suất vượt trội trung bình âm, danh mục biến động hơn có Sharpe ít âm hơn; Sharpe chỉ là thước phụ, báo cáo ở cửa sổ có lợi suất vượt trội trung bình dương.

• Ba thước đo: độ phân tán mô hình, mức hối tiếc, kích thước tập tin cậy mô hình (mức ý nghĩa 10\% ấn định trước).

• Độ phân tán và tập tin cậy chỉ báo bất định; mức hối tiếc đo rủi ro lựa chọn.
\end{tcolorbox}

\section{Tự kiểm tra}

{\itshape 1. Lợi suất tương đương chắc chắn tính từ những đại lượng nào, và hệ số ngại rủi ro cơ sở là bao nhiêu?\par}

Đáp án: trung bình lợi suất sau chi phí trừ nửa hệ số ngại rủi ro nhân phương sai lợi suất, trên cửa sổ đánh giá. Cơ sở là 5, độ nhạy ở 3 và 10.

{\itshape 2. Vì sao tỷ số Sharpe chỉ là thước đo phụ trong đề cương?\par}

Đáp án: khi lợi suất vượt trội trung bình âm, đúng ở chế độ căng thẳng, danh mục biến động hơn có Sharpe ít âm hơn. Vì vậy Sharpe chỉ báo cáo ở cửa sổ có lợi suất vượt trội trung bình dương.

{\itshape 3. Ba thước đo rủi ro lựa chọn mô hình là gì, thước đo nào là kết quả chính của H1?\par}

Đáp án: độ phân tán mô hình, mức hối tiếc, kích thước tập tin cậy mô hình. Kết quả chính là mức hối tiếc; hai thước đo kia là kết quả phụ.

{\itshape 4. Độ phân tán mô hình và kích thước tập tin cậy đo rủi ro hay bất định?\par}

Đáp án: bất định, không phải rủi ro.

{\itshape 5. Vì sao tính sau chi phí giao dịch?\par}

Đáp án: đề cương đo hiệu quả ngoài mẫu sau chi phí giao dịch, nên trung bình và phương sai trong công thức là của lợi suất sau chi phí; mô hình mua bán nhiều bị trừ điểm vì phí.

\chapter{Mục 4.3.6 của đề cương (tiếp): Độ phân tán mô hình, mức hối tiếc và tập tin cậy mô hình}

Từ bảy điểm số trên một cửa sổ đánh giá, đề cương dựng ba thước đo. Độ phân tán mô hình \(D_{t}\) đo bảy mô hình chênh nhau bao nhiêu. Mức hối tiếc \(R_{t}\) đo khoản thiệt vì mô hình chọn trước kém mô hình tốt nhất sau thực tế. Kích thước tập tin cậy mô hình đếm số mô hình chưa thể loại. Mức hối tiếc là thước đo rủi ro chính; hai thước kia chỉ báo bất định. Điểm tinh tế nhất: mức hối tiếc không bao giờ âm, và kỳ vọng của nó vẫn dương khi các mô hình ngang chất lượng.

Phần này giải thích từng công thức, dành nhiều chỗ cho câu hỏi khó nhất (vì sao kỳ vọng dương khi không mô hình nào hơn), rồi nói vì sao tập tin cậy lớn ở chế độ căng thẳng chưa chắc là tin xấu.

\section{Phương trình (3): độ phân tán mô hình}

Độ phân tán mô hình \(D_{t}\) là độ lệch chuẩn mẫu của điểm số theo các mô hình trong cùng một cửa sổ (tính qua các mô hình, không qua thời gian):

\(D_{t}\  = \ \sqrt{\frac{1}{|M|\  - \ 1}\sum_{m \in M}^{}\left( U_{m,t}\  - \ {\overline{U}}_{t} \right)^{2}}\) (3)

\(D_{t}\) là độ phân tán tại cửa sổ kết thúc ở t; \(U_{m,t}\) là điểm của mô hình m (phương trình (1)); \({\overline{U}}_{t}\) là trung bình điểm (phương trình (2)). Tổng chạy qua 7 mô hình, nên mẫu số trong dấu căn là 7 − 1 = 6.

Bốn bước: lấy điểm của từng mô hình trừ trung bình; bình phương để mọi số dương và độ lệch lớn bị nhấn mạnh; cộng lại rồi chia cho 6; lấy căn bậc hai để về cùng đơn vị với điểm số.

Chia cho số quan sát trừ 1 là cách hiệu chỉnh thường dùng của độ lệch chuẩn mẫu, vì trung bình đã tính từ chính các điểm đó.

Bảy đội tiếp sức nộp thành tích: gần bằng nhau thì độ phân tán nhỏ; có đội rất nhanh, đội rất chậm thì lớn. Độ phân tán không nói đội nào nhanh nhất, chỉ nói các đội khác nhau nhiều hay ít.

\begin{tcolorbox}[breakable]
Ví dụ giả định (số tự đặt để minh họa): bảy điểm số ở ví dụ trước: 2,0; 3,0; −2,0; 1,0; 2,6; 3,5; 1,4 (phần trăm); trung bình khoảng 1,64.

Độ lệch so với trung bình: 0,36; 1,36; −3,64; −0,64; 0,96; 1,86; −0,24.

Bình phương: 0,13; 1,85; 13,27; 0,41; 0,92; 3,45; 0,06; tổng khoảng 20,08.

20,08 / 6 ≈ 3,35; căn bậc hai ≈ 1,83. Độ phân tán khoảng 1,83 điểm phần trăm.

Riêng mô hình có điểm −2,0 đóng góp 13,27 trong 20,08: một mô hình lệch xa kéo độ phân tán lên rất mạnh vì bước bình phương.
\end{tcolorbox}

Đề cương coi độ phân tán là chỉ báo bất định: các mô hình cho kết quả giống nhau thì chọn mô hình nào cũng như nhau; khác nhau xa thì việc chọn có thể quan trọng.

Vì sao cần cả độ phân tán lẫn mức hối tiếc? Theo bản đọc: độ phân tán lớn chưa chắc gây thiệt nếu mô hình được chọn lại là mô hình tốt. Một thước nói về độ rộng khoảng cách giữa các mô hình, thước kia nói về cú trượt chân của một lần chọn cụ thể.

\section{Mức hối tiếc: khoản thiệt vì không chọn mô hình tốt nhất sau thực tế}

Mức hối tiếc (tiếng Anh: regret) là thước đo trung tâm. Trước mùa giải bạn chọn một cách chia tiền; cuối mùa mới biết cách nào tốt nhất; chênh lệch điểm giữa cách tốt nhất và cách đã chọn là mức hối tiếc.

Phải chọn trước vì nhà đầu tư chỉ biết kết quả sau khi đã bỏ tiền. Quy tắc mô phỏng của đề cương: tại đầu mỗi cửa sổ đánh giá, chọn mô hình có lợi suất tương đương chắc chắn ngoài mẫu cao nhất trong 36 tháng liền trước; mô hình đó tái cân bằng hằng tháng suốt cửa sổ.

\section{Phương trình (4), từng ký hiệu một}

\(R_{t}\  = \ U_{o,t}\  - \ U_{c,t}\) (4)

\(R_{t}\) là mức hối tiếc tại cửa sổ kết thúc ở t. \(U_{c,t}\) là điểm của mô hình chọn trước (c: chọn). \(U_{o,t}\) là điểm cao nhất trong bảy mô hình, nhìn lại sau thực tế (o: tốt nhất).

Ba bước: đầu cửa sổ, nhìn lại 36 tháng, chọn mô hình dẫn đầu; hết cửa sổ, tính điểm của cả bảy và lấy giá trị lớn nhất làm \(U_{o,t}\); lấy điểm của mô hình đã chọn làm \(U_{c,t}\) và tính \(U_{o,t}\) trừ \(U_{c,t}\).

Chọn một trong bảy đội đua theo thành tích mùa trước; hết mùa, so thu nhập của đội đã chọn với đội tốt nhất mùa đó. Khoảng chênh là khoản đáng lẽ có thêm; chọn đúng đội tốt nhất thì bằng 0.

\begin{tcolorbox}[breakable]
Ví dụ giả định (số tự đặt để minh họa): điểm trong cửa sổ: M1 2,0; M2 3,0; M3 −2,0; M4 1,0; M5 2,6; M6 3,5; M7 1,4 (phần trăm).

Tốt nhất sau thực tế là M6, nên \(U_{o,t}\) = 3,5.

Nếu M3 dẫn đầu 36 tháng trước nên được chọn: mức hối tiếc = 3,5 − (−2,0) = 5,5 điểm phần trăm.

Nếu M6 được chọn: 3,5 − 3,5 = 0.

Nếu M2 được chọn: 3,5 − 3,0 = 0,5.
\end{tcolorbox}

Mức hối tiếc phản ánh cả chất lượng quy tắc chọn lẫn may rủi. Đề cương tính nó ở mỗi cửa sổ và dùng chuỗi \(R_{t}\) theo thời gian để so giữa các chế độ ở phương trình (5).

Hai điểm nên nhớ: mức hối tiếc so với mô hình tốt nhất trong tập, không so với một chuẩn bên ngoài; và quy tắc chọn theo 36 tháng là quy tắc dùng trong thước đo, không phải khuyến nghị cách nhà đầu tư nên chọn.

\section{Vì sao mức hối tiếc không bao giờ âm}

Đề cương nêu: vì mô hình chọn trước thuộc tập M, mức hối tiếc không âm. Mô hình chọn trước nằm trong bảy mô hình; không phần tử nào của tập lớn hơn phần tử lớn nhất của chính tập đó. Vậy \(U_{o,t}\) luôn lớn hơn hoặc bằng \(U_{c,t}\).

Điểm của một học sinh không thể cao hơn điểm cao nhất lớp, vì điểm cao nhất được tính cả bạn đó.

Ngoại lệ (mục 4.3.8): khi kiểm định H3, cơ chế (ii) và (iii) được loại khỏi tập M, nên mức hối tiếc của chúng có thể âm, tức cơ chế có thể vượt cả mô hình tốt nhất trong tập.

\section{Vì sao kỳ vọng của mức hối tiếc vẫn dương khi các mô hình ngang chất lượng}

Đây là điểm khó nhất. Đề cương viết: kỳ vọng của mức hối tiếc vẫn dương khi các mô hình ngang chất lượng, vì giá trị lớn nhất của các đại lượng có nhiễu tăng theo phương sai của nhiễu. Ta gỡ thành ba ý.

\subsection{Ý thứ nhất: điểm số đo được luôn có nhiễu}

Điểm số trên cửa sổ 6 tháng chỉ dựa vào khoảng 6 tháng dữ liệu, nên trung bình và phương sai ước lượng có sai số. Hai mô hình giỏi ngang nhau vẫn cho hai điểm khác nhau chỉ vì may rủi; phần chênh đó là nhiễu ước lượng.

Bảy cầu thủ ngang tài cùng sút mười quả phạt đền: người vào bảy, người chín, người năm. Chênh lệch chỉ do may rủi, không phản ánh tài năng.

\subsection{Ý thứ hai: giá trị lớn nhất của các số nhiễu thường dương}

Nếu bảy mô hình ngang nhau, điểm mỗi mô hình bằng mức chung cộng một số nhiễu trung bình 0. Nhưng mức hối tiếc không dùng trung bình bảy số nhiễu; nó dùng điểm lớn nhất, tức mức chung cộng số nhiễu lớn nhất.

Số lớn nhất trong bảy số ngẫu nhiên trung bình 0 gần như chắc chắn dương; chỉ âm khi cả bảy cùng âm, khả năng rất nhỏ. Vì vậy phần nhiễu mà \(U_{o,t}\) nhặt được có kỳ vọng dương: mô hình tốt nhất sau thực tế trông tốt hơn mức thật của nó.

Trong bảy cầu thủ ngang tài, thế nào cũng có người gặp may nhất; lấy người đó làm chuẩn thì chuẩn luôn cao hơn mức trung bình thật.

\subsection{Ý thứ ba: mô hình chọn trước không hưởng cái may đó}

Diễn giải của bản đọc: mô hình chọn trước được chọn bằng 36 tháng quá khứ, không bằng chính may rủi của cửa sổ đánh giá, nên phần nhiễu của nó có trung bình 0, không bị “chọn lọc” cho cao lên.

Vậy \(U_{o,t}\) mang phần nhiễu lớn nhất, kỳ vọng dương; \(U_{c,t}\) mang phần nhiễu không chọn lọc, kỳ vọng 0. Hiệu của chúng có kỳ vọng dương dù không mô hình nào thật sự hơn.

Nhiễu càng lớn thì số lớn nhất trong bảy số nhiễu càng lớn. Theo bản đọc: nếu ở chế độ căng thẳng biến động cao làm điểm số kém ổn định hơn, mức hối tiếc thô ở đó có thể lớn hơn dù bản thân việc chọn không nguy hiểm hơn.

\begin{tcolorbox}[breakable]
Ví dụ giả định (số tự đặt để minh họa): bảy mô hình hoàn toàn ngang chất lượng; nhiễu đo của mỗi mô hình độc lập, phân phối chuẩn, trung bình 0, độ lệch chuẩn 1 điểm phần trăm.

Mức hối tiếc trung bình (so với một mô hình chọn ngẫu nhiên) bằng kỳ vọng của số lớn nhất trong bảy nhiễu, khoảng 1,35 điểm phần trăm.

Nhiễu có độ lệch chuẩn gấp đôi thì mức hối tiếc trung bình cũng gấp đôi, khoảng 2,7, dù các mô hình vẫn ngang nhau.

Số mô hình cũng ảnh hưởng: với độ lệch chuẩn 1, kỳ vọng của số lớn nhất khoảng 0,56 với 2 mô hình, 0,85 với 3, 1,16 với 5, 1,35 với 7 và 1,54 với 10.
\end{tcolorbox}

Hai nhận xét: mức hối tiếc thô tăng theo độ lớn của nhiễu dù chất lượng thật không đổi; và tăng theo số mô hình, nên khi so tập có và không có M5 phải để ý số mô hình (nhận xét thứ hai là của bản đọc, đúng dưới giả định nhiễu chuẩn độc lập).

Hệ quả cho H1: thấy mức hối tiếc thô cao hơn ở chế độ căng thẳng thì chưa biết do rủi ro thật hay do nhiễu lớn hơn. Đề cương xử lý ở hai tầng: chia cho sai số chuẩn bootstrap (dưới đây) và hiệu chuẩn bằng phân phối rỗng (mục 4.3.7).

\section{Chia cho sai số chuẩn bootstrap: đo bằng đơn vị của nhiễu}

Trong phân tích chính, độ phân tán và mức hối tiếc được chia cho sai số chuẩn bootstrap của chênh lệch lợi suất tương đương chắc chắn giữa các mô hình trong cùng cửa sổ, để loại tác động cơ học của nhiễu.

Sai số chuẩn cho biết một con số ước lượng dao động bao nhiêu nếu đo lại bằng mẫu khác: thước đo của nhiễu. Bootstrap ước lượng nó bằng cách bốc lại nhiều lần từ chính dữ liệu (chi tiết ở mục 4.3.7).

Chia một số đo cho thước đo nhiễu tương ứng là chuẩn hóa: kết quả tính bằng đơn vị nhiễu.

Cân hàng trên đường xóc thì kim nhảy nhiều hơn. Muốn so hai hộp trên đường xóc với hai hộp trên đường êm, phải chia số chênh cho mức rung của cân ở từng đoạn đường; thương số mới là phép so công bằng.

\begin{tcolorbox}[breakable]
Ví dụ giả định (số tự đặt để minh họa): cửa sổ bình thường, mức hối tiếc thô 1,0, sai số chuẩn 0,8; sau khi chia: 1,25.

Cửa sổ căng thẳng, mức hối tiếc thô 3,0 (gấp ba), sai số chuẩn 2,5 vì dữ liệu ồn hơn; sau khi chia: 1,2.

Số thô tăng gấp ba, số đã chia gần như bằng nhau: phần lớn mức tăng đến từ nhiễu.

Nếu mức hối tiếc thô ở cửa sổ căng thẳng là 6,0 với sai số chuẩn 2,5, số đã chia là 2,4, gần gấp đôi 1,25: lúc đó mới có dấu hiệu mức tăng không chỉ do nhiễu.
\end{tcolorbox}

Đề cương không nói phép chia xóa sạch tác động cơ học; còn bước hiệu chuẩn bằng phân phối rỗng, tạo một thế giới giả trong đó các mô hình ngang nhau để xem mức hối tiếc vẫn tăng bao nhiêu chỉ vì nhiễu (mục 4.3.7).

Đề cương chưa nói số lần lặp bootstrap cho sai số chuẩn trong từng cửa sổ (con số 1.000 chỉ gắn với phân phối rỗng), cũng chưa nói cách gộp sai số chuẩn của nhiều cặp mô hình cho \(D_{t}\) hay \(R_{t}\). Đừng bổ sung chi tiết đề cương chưa quy định.

\section{Kích thước tập tin cậy mô hình}

Thước đo thứ ba là kích thước tập tin cậy mô hình (Hansen và cộng sự, 2011) ở mức ý nghĩa 10\% ấn định trước: số mô hình chưa thể bị loại là kém hơn.

Trong vòng sơ loại, ban giám khảo chỉ loại người chắc chắn kém hơn người giỏi nhất; ai còn có thể ngang ngửa trong giới hạn may rủi thì ở lại.

Đề cương không mô tả thuật toán, chỉ nêu ý nghĩa và mức 10\% quyết định trước khi xem dữ liệu. Tập tin cậy không chỉ ra một người thắng mà liệt kê những người chưa thể loại.

Dữ liệu kém thông tin cho tập lớn, dữ liệu nhiều thông tin cho tập nhỏ: dữ liệu ít hoặc ồn thì kiểm định không đủ sức loại ai.

Vì vậy tập lớn hơn ở chế độ căng thẳng có thể chỉ phản ánh dữ liệu kém thông tin hơn; đề cương báo cáo kích thước tập như thước đo mô tả, kèm số quan sát và công suất kiểm định. Công suất là khả năng phát hiện khác biệt khi khác biệt có thật; công suất thấp thì danh sách chưa thể loại dài ra.

\begin{tcolorbox}[breakable]
Ví dụ giả định (số tự đặt để minh họa): cửa sổ bình thường, nhiều quan sát, dữ liệu ít ồn: tập còn 3 mô hình.

Cửa sổ căng thẳng, dữ liệu ồn hơn, ít quan sát độc lập hơn: tập còn 6 mô hình.

Đọc sai: “căng thẳng làm mọi mô hình ngang nhau, chọn gì cũng được”. Đọc đúng hơn: “kiểm định ở chế độ căng thẳng ít sức hơn, chưa loại được nhiều”. Cùng con số 6, khác ý nghĩa; vì thế phải ghi kèm số quan sát và công suất.
\end{tcolorbox}

Theo bản đọc, tập tin cậy là kết quả phụ vì kích thước của nó phụ thuộc cả chất lượng mô hình lẫn lượng thông tin của dữ liệu, nên dễ đọc nhầm hơn mức hối tiếc.

\section{Ghép ba thước đo: bất định và rủi ro}

Nhắc lại mục 4.1: rủi ro mô hình là hậu quả bất lợi từ mô hình sai hoặc dùng sai; bất định mô hình là nhiều mô hình khả dĩ cho kết quả khác nhau; rủi ro lựa chọn mô hình là tổn thất kỳ vọng đo bằng mức hối tiếc. Độ phân tán và kích thước tập tin cậy chỉ báo bất định.

{\small\begin{longtable}[]{@{}
  >{\raggedright\arraybackslash}p{(\columnwidth - 6\tabcolsep) * \real{0.1381}}
  >{\raggedright\arraybackslash}p{(\columnwidth - 6\tabcolsep) * \real{0.3574}}
  >{\raggedright\arraybackslash}p{(\columnwidth - 6\tabcolsep) * \real{0.1623}}
  >{\raggedright\arraybackslash}p{(\columnwidth - 6\tabcolsep) * \real{0.3422}}@{}}
\toprule\noalign{}
\begin{minipage}[b]{\linewidth}\raggedright
Thước đo
\end{minipage} & \begin{minipage}[b]{\linewidth}\raggedright
Công thức hoặc cách tính
\end{minipage} & \begin{minipage}[b]{\linewidth}\raggedright
Thuộc nhóm
\end{minipage} & \begin{minipage}[b]{\linewidth}\raggedright
Câu hỏi nó trả lời
\end{minipage} \\
\midrule\noalign{}
\endhead
\bottomrule\noalign{}
\endlastfoot
Độ phân tán mô hình \(D_{t}\) & độ lệch chuẩn mẫu của \(U_{m,t}\) theo m, phương trình (3) & Chỉ báo của bất định & Các mô hình khác nhau nhiều hay ít? \\
Mức hối tiếc \(R_{t}\) & \(U_{o,t}\) trừ \(U_{c,t}\), phương trình (4) & Thước đo rủi ro lựa chọn & Chọn trước thì thiệt bao nhiêu so với mô hình tốt nhất sau thực tế? \\
Kích thước tập tin cậy mô hình & số mô hình chưa thể loại ở mức ý nghĩa 10\% & Chỉ báo của bất định & Còn bao nhiêu mô hình chưa thể loại? \\
\end{longtable}}

Đề cương dùng hai phiên bản: bản thô cho mức chênh nguyên trạng; bản đã chia cho sai số chuẩn bootstrap cho mức chênh so với nhiễu. Phân tích chính dùng bản đã chia.

\begin{tcolorbox}[breakable]
\textbf{Ghi nhớ khi đi thi}

• \(D_{t}\): độ lệch chuẩn mẫu của điểm số giữa các mô hình, chia cho M − 1 trong dấu căn (phương trình (3)).

• \(R_{t}\): điểm cao nhất sau thực tế trừ điểm của mô hình chọn trước (phương trình (4)); mô hình chọn trước có lợi suất tương đương chắc chắn ngoài mẫu cao nhất trong 36 tháng liền trước, chọn tại đầu mỗi cửa sổ đánh giá.

• Mức hối tiếc không âm vì mô hình chọn trước thuộc M; kỳ vọng vẫn dương khi các mô hình ngang nhau, vì giá trị lớn nhất của các đại lượng có nhiễu tăng theo phương sai của nhiễu.

• Cơ chế (ii) và (iii) bị loại khỏi M khi kiểm định H3, nên mức hối tiếc của chúng có thể âm.

• Phân tích chính chia \(D_{t}\) và \(R_{t}\) cho sai số chuẩn bootstrap của chênh lệch điểm số giữa các mô hình trong cùng cửa sổ.

• Tập tin cậy ở mức 10\%; tập lớn hơn ở chế độ căng thẳng có thể chỉ do dữ liệu kém thông tin, nên báo cáo mô tả kèm số quan sát và công suất.
\end{tcolorbox}

\section{Tự kiểm tra}

{\itshape 1. Mô tả bốn bước tính \(D_{t}\) và mẫu số trong dấu căn.\par}

Đáp án: lấy điểm từng mô hình trừ trung bình, bình phương, cộng lại, chia cho M − 1 = 6, lấy căn bậc hai.

{\itshape 2. Mức hối tiếc tính từ những đại lượng nào, mô hình chọn trước là mô hình nào?\par}

Đáp án: \(U_{o,t}\) (điểm cao nhất trong M sau thực tế) trừ \(U_{c,t}\) (điểm của mô hình chọn trước). Mô hình chọn trước có lợi suất tương đương chắc chắn ngoài mẫu cao nhất trong 36 tháng liền trước.

{\itshape 3. Vì sao mức hối tiếc không âm, nhưng kỳ vọng vẫn dương khi các mô hình ngang nhau?\par}

Đáp án: không âm vì mô hình chọn trước thuộc M. Kỳ vọng dương vì mô hình tốt nhất sau thực tế nhặt lấy giá trị lớn nhất của các đại lượng có nhiễu, và giá trị đó tăng theo phương sai của nhiễu.

{\itshape 4. Vì sao chia cho sai số chuẩn bootstrap?\par}

Đáp án: để loại tác động cơ học của nhiễu ước lượng.

{\itshape 5. Vì sao tập tin cậy lớn hơn ở chế độ căng thẳng chưa chắc nghĩa là rủi ro cao hơn?\par}

Đáp án: có thể chỉ do dữ liệu kém thông tin; đề cương báo cáo kích thước tập dạng mô tả, kèm số quan sát và công suất.

{\itshape 6. Ba thuật ngữ khác nhau thế nào, độ phân tán thuộc nhóm nào?\par}

Đáp án: rủi ro mô hình là hậu quả bất lợi từ mô hình sai hoặc dùng sai; bất định mô hình là nhiều mô hình khả dĩ cho kết quả khác nhau; rủi ro lựa chọn mô hình là tổn thất kỳ vọng đo bằng mức hối tiếc. Độ phân tán chỉ báo bất định.

\chapter{Mục 4.3.6 của đề cương (tiếp): Ba cửa sổ thời gian, siêu tham số, M3 và M5}

Để phép so sánh bảy mô hình công bằng và lặp lại được, đề cương quy định trước bốn thứ: độ dài các cửa sổ (104 tuần ước lượng, 6 tháng đánh giá, 36 tháng chọn mô hình); cách tinh chỉnh siêu tham số (cùng lưới, cùng quy tắc, cùng ngân sách); cách xử lý M3 khi dữ liệu ít; và cách báo cáo M5 (cả khi có lẫn khi không có M5). Các độ dài cửa sổ là dự kiến, biện minh sau ước lượng thử ở năm thứ nhất.

Mục tiêu: trả lời được những câu như “vì sao ba cửa sổ dài khác nhau” hay “vì sao báo cáo cả khi không có M5”.

\section{Ba cửa sổ thời gian, ba việc khác nhau}

Hình 2 của đề cương vẽ ba cửa sổ và thời điểm chọn mô hình. Bảng tóm tắt:

{\small\begin{longtable}[]{@{}
  >{\raggedright\arraybackslash}p{(\columnwidth - 6\tabcolsep) * \real{0.1954}}
  >{\raggedright\arraybackslash}p{(\columnwidth - 6\tabcolsep) * \real{0.3563}}
  >{\raggedright\arraybackslash}p{(\columnwidth - 6\tabcolsep) * \real{0.3793}}
  >{\raggedright\arraybackslash}p{(\columnwidth - 6\tabcolsep) * \real{0.0691}}@{}}
\toprule\noalign{}
\begin{minipage}[b]{\linewidth}\raggedright
Cửa sổ
\end{minipage} & \begin{minipage}[b]{\linewidth}\raggedright
Độ dài (dự kiến)
\end{minipage} & \begin{minipage}[b]{\linewidth}\raggedright
Dùng để làm gì
\end{minipage} & \begin{minipage}[b]{\linewidth}\raggedright
Mục
\end{minipage} \\
\midrule\noalign{}
\endhead
\bottomrule\noalign{}
\endlastfoot
Cửa sổ ước lượng & 104 tuần & Ước lượng tham số của mỗi mô hình & 4.3.6 \\
Cửa sổ đánh giá & 6 tháng, không chồng lấn & Đo lợi suất tương đương chắc chắn ngoài mẫu và mức hối tiếc & 4.3.6, Hình 2 \\
Cửa sổ nhìn lại để chọn mô hình & 36 tháng liền trước đầu cửa sổ đánh giá & Chọn trước mô hình có lợi suất tương đương chắc chắn ngoài mẫu cao nhất & 4.3.6 \\
\end{longtable}}

Huấn luyện viên có ba cuốn sổ: sổ tập luyện trước mùa giải (ước lượng), sổ thi đấu trong mùa (đánh giá), sổ thành tích ba năm gần đây để quyết định tin ai (chọn mô hình). Ba việc khác nhau nên độ dài khác nhau.

\subsection{Cửa sổ 104 tuần để ước lượng}

104 tuần dữ liệu tuần là hai năm. Ước lượng là dùng 104 tuần lợi suất để tính đầu vào (trung bình, hiệp phương sai) rồi ra trọng số danh mục.

Nhắc lại ba tần suất: ngày cho chế độ, tuần cho hiệp phương sai, tháng cho tái cân bằng.

Vì sao 104? Đề cương nói đây là giá trị dự kiến, biện minh sau ước lượng thử; con số là điểm xuất phát, và được cố định khi đăng ký trước.

\subsection{Cửa sổ đánh giá 6 tháng, không chồng lấn}

Cửa sổ đánh giá là đoạn ngoài mẫu mà mô hình chưa thấy khi ước lượng; phương trình (1) tính trên đúng cửa sổ này. Mỗi cửa sổ cho bảy điểm số, từ đó một giá trị \(D_{t}\) và một giá trị \(R_{t}\).

“Không chồng lấn” nghĩa là các cửa sổ đánh giá kế tiếp không đè lên nhau; cửa sổ trượt hằng tháng chỉ dùng cho phân tích phụ.

Vì sao quan trọng (giải thích chung)? Nếu mỗi lần đo nhiệt độ lấy trung bình bảy ngày gần nhất, hai lần đo liền nhau dùng chung sáu ngày, gần như giống nhau: nhiều con số hơn mà không thêm thông tin. Cửa sổ không chồng lấn giúp mỗi giá trị \(R_{t}\) mang thông tin mới.

Đổi lại, số điểm dữ liệu ít đi. Với mẫu 2007--2026, đề tài dự kiến khoảng 28 cửa sổ đánh giá độc lập, và dưới 10 cửa sổ độc lập ở chế độ căng thẳng thì kết quả gọi là không kết luận (mục 4.3.10).

\begin{tcolorbox}[breakable]
Ví dụ giả định (số tự đặt để minh họa): một giai đoạn căng thẳng dài 14 tháng chứa tối đa 2 cửa sổ 6 tháng không chồng lấn.

Cửa sổ trượt hằng tháng cho 9 cửa sổ trọn vẹn (bắt đầu từ tháng 1 đến tháng 9), nhưng chúng đè lên nhau nhiều, nên chỉ có khoảng 2 cửa sổ độc lập.

Đếm cửa sổ độc lập, không đếm cửa sổ trượt, mới trả lời đúng câu “đã đủ bằng chứng chưa”.
\end{tcolorbox}

\subsection{Cửa sổ 36 tháng để chọn mô hình}

Tại đầu mỗi cửa sổ đánh giá, nhà đầu tư chọn trước mô hình có lợi suất tương đương chắc chắn ngoài mẫu cao nhất trong 36 tháng liền trước. Dữ liệu để chọn chỉ gồm những gì đã xảy ra trước thời điểm quyết định, nên không nhìn trước.

Cùng tinh thần với rổ cổ phiếu xác định bằng dữ liệu có tại ngày đó và mô hình chế độ ước lượng lại bằng dữ liệu đến ngày đó.

Vì cần 104 tuần để ước lượng và 36 tháng kết quả ngoài mẫu để chọn, mức hối tiếc chỉ tính được từ khoảng năm thứ năm của mẫu.

\section{Siêu tham số: cùng lưới, cùng quy tắc, cùng ngân sách}

Siêu tham số là “núm vặn” phải đặt trước khi mô hình học dữ liệu. Ví dụ chung: mô hình co rút có núm cho biết co rút nhiều hay ít; Black-Litterman có núm cho biết tin quan điểm đến đâu. M1 chia đều nên không có núm; M3 dùng ước lượng mẫu nguyên xi.

Siêu tham số của M2, M4, M5, M6 và M7 được chọn từ một lưới cho trước (danh sách giá trị thử), bằng cùng một quy tắc và cùng ngân sách tinh chỉnh (mức công sức thử), rồi cố định khi đăng ký trước.

Cho một vận động viên tập thử hai mươi lần, người khác hai lần, thì người tập nhiều thắng là chuyện thường. So công bằng đòi cùng số buổi tập thử và cùng cách chọn buổi tốt nhất.

Theo bản đọc, đây là cùng loại vấn đề với giá trị lớn nhất của các đại lượng có nhiễu: thử nhiều tổ hợp rồi báo cáo tổ hợp thắng thì nhặt cả phần may. Mục 4.3.7 còn hiệu chỉnh thiên lệch chọn lọc bằng tỷ số Sharpe hiệu chỉnh.

\begin{tcolorbox}[breakable]
Ví dụ giả định (số tự đặt để minh họa): hai mô hình chất lượng như nhau; P thử 5 tổ hợp, Q thử 50, cả hai báo cáo tổ hợp tốt nhất.

Số lớn nhất của 50 số nhiễu thường lớn hơn số lớn nhất của 5 số nhiễu, nên Q thường trông tốt hơn.

Chiến thắng đến từ số lần thử, không từ chất lượng; cùng ngân sách tránh khoảng chênh giả này.
\end{tcolorbox}

Thêm hai chi tiết: quan điểm của M4 lấy từ một mô hình dự báo đơn giản cố định khi đăng ký trước, nên không được đổi sau khi thấy dữ liệu; và đăng ký trước chỉ giảm chứ không loại hẳn thiên lệch, vì tác giả đã dùng dữ liệu 2019--2025 ở bài trước.

\section{M3 và nghịch đảo giả Moore-Penrose}

M3 là trung bình-phương sai với ước lượng mẫu (Markowitz, 1952): dùng nguyên trung bình và hiệp phương sai tính từ cửa sổ, không co rút.

Hai nhận định của nhóm (1) áp dụng thẳng vào M3: tối ưu hóa dồn trọng số vào tài sản có sai số ước lượng lớn (Michaud, 1989); thay tham số tổng thể bằng ước lượng mẫu làm hiệu quả ngoài mẫu giảm mạnh (Kan và Zhou, 2007).

Điều kiện kỹ thuật (mục 4.3.5): chỉ ước lượng trực tiếp M3 khi tỷ lệ số cổ phiếu trên số quan sát tuần không vượt 0,5; vượt ngưỡng thì dùng nghịch đảo giả Moore-Penrose và ghi rõ.

Công thức trung bình-phương sai cần nghịch đảo ma trận hiệp phương sai, giống phép chia trong số thường. Số cổ phiếu quá lớn so với số quan sát thì ma trận thiếu thông tin, phép “chia” không làm được hoặc rất bất ổn (kiến thức toán chung; đề cương chỉ nêu ngưỡng 0,5).

Có 30 ẩn số mà chỉ 20 phương trình thì có vô số nghiệm. Nghịch đảo giả chọn một nghiệm đặc biệt (nghiệm “ngắn nhất”) để vẫn tính tiếp được: lối thoát khi phép chia thường không dùng được.

Đề cương không giấu việc này: kết quả nào của M3 dùng nghịch đảo giả sẽ được ghi rõ, để người đọc tự đánh giá độ tin cậy.

\begin{tcolorbox}[breakable]
Ví dụ giả định (số tự đặt để minh họa): 20 cổ phiếu, 100 quan sát tuần, tỷ lệ 0,2, không vượt 0,5: ước lượng bình thường.

60 cổ phiếu trên 100 quan sát tuần, tỷ lệ 0,6: dùng nghịch đảo giả và ghi rõ.

Hai con số chỉ để minh họa. Với rổ 30 mã và 104 quan sát tuần, tỷ lệ khoảng 0,29; đề cương không tự nêu phép đối chiếu này.
\end{tcolorbox}

\section{M5: mô hình do chính tác giả phát triển}

M5 là Black-Litterman có chuyển chế độ, mở rộng từ Nguyen và Nguyen (2026a), mô hình duy nhất trong tập do tác giả phát triển. Theo bản đọc, điều này đặt ra nguy cơ mô hình của chính tác giả chi phối kết luận.

Đề cương xử lý bằng hai biện pháp: M5 dùng cùng lưới, cùng quy tắc, cùng ngân sách tinh chỉnh như các mô hình khác; và kết quả báo cáo cả khi có lẫn khi không có M5 (khi đó tập còn sáu mô hình).

Một thí sinh đồng thời là người soạn đề: dù công bằng đến đâu, khán giả vẫn muốn biết kết quả khi người đó vắng mặt. Báo cáo cả hai cho người đọc biết kết luận có phụ thuộc mô hình của tác giả hay không.

Hệ quả kỹ thuật (suy ra từ cấu trúc giá trị lớn nhất): bỏ M5 thì mọi phép tính chạy trên sáu mô hình, và mức hối tiếc thô có xu hướng tăng theo số mô hình, nên khi so phải nhớ số mô hình đã đổi.

Biện pháp thứ ba (mục 4.3.7): lặp lại kiểm định H1 trên tập chỉ gồm mô hình tĩnh M1, M2, M3, M4, M6, để loại tác động của M5 và M7.

\section{Vì sao ấn định trước nhiều thứ như vậy}

Theo bản đọc, các quy định về cửa sổ, siêu tham số và M5 cùng một nguyên tắc: quyết định trước, ghi lại, không đổi sau khi thấy dữ liệu, tức đăng ký trước (mục 4.3.5).

Nếu được tự do đổi cửa sổ hay siêu tham số sau khi thấy kết quả, người nghiên cứu có thể vô tình hay cố ý chọn cấu hình cho kết quả đẹp; đăng ký trước cắt con đường đó. Giới hạn đã nói: tác giả đã dùng dữ liệu 2019--2025, nên kế hoạch khai rõ những kết quả đã xem.

Dữ liệu phát sinh sau ngày đăng ký là mẫu kiểm tra riêng; mọi thay đổi sau đó ghi là lệch khỏi kế hoạch.

\begin{tcolorbox}[breakable]
\textbf{Ghi nhớ khi đi thi}

• Ước lượng 104 tuần; đánh giá 6 tháng không chồng lấn; chọn trước tại đầu mỗi cửa sổ đánh giá theo 36 tháng liền trước.

• Các độ dài là dự kiến, biện minh sau ước lượng thử; cửa sổ trượt hằng tháng chỉ là phân tích phụ; khoảng 28 cửa sổ đánh giá với mẫu 2007--2026.

• Siêu tham số của M2, M4, M5, M6, M7: cùng lưới, cùng quy tắc, cùng ngân sách, cố định khi đăng ký trước; quan điểm của M4 từ một mô hình dự báo đơn giản cố định khi đăng ký trước.

• M3 chỉ ước lượng trực tiếp khi số cổ phiếu trên số quan sát tuần không vượt 0,5; vượt thì dùng nghịch đảo giả Moore-Penrose và ghi rõ.

• M5 do tác giả phát triển, xử lý như mọi mô hình khác, báo cáo cả khi có lẫn khi không có M5.
\end{tcolorbox}

\section{Tự kiểm tra}

{\itshape 1. Nêu độ dài ba cửa sổ và việc của từng cửa sổ.\par}

Đáp án: ước lượng 104 tuần; đánh giá 6 tháng không chồng lấn; chọn mô hình theo 36 tháng liền trước đầu cửa sổ đánh giá. Cả ba là giá trị dự kiến.

{\itshape 2. Siêu tham số của những mô hình nào được chọn từ một lưới cho trước, và quy tắc chung là gì?\par}

Đáp án: M2, M4, M5, M6, M7; cùng quy tắc, cùng ngân sách tinh chỉnh, cố định khi đăng ký trước.

{\itshape 3. Khi nào M3 dùng nghịch đảo giả Moore-Penrose?\par}

Đáp án: khi tỷ lệ số cổ phiếu trên số quan sát tuần vượt 0,5; việc dùng nghịch đảo giả được ghi rõ trong kết quả.

{\itshape 4. Vì sao đề cương báo cáo kết quả có và không có M5?\par}

Đáp án: vì M5 do tác giả phát triển; báo cáo cả hai cho biết kết luận có phụ thuộc mô hình của tác giả không. Ngoài ra H1 được lặp lại trên tập mô hình tĩnh.

{\itshape 5. Đề cương thừa nhận giới hạn gì về đăng ký trước?\par}

Đáp án: tác giả đã dùng dữ liệu 2019--2025, gồm hai giai đoạn căng thẳng, trong Nguyen và Nguyen (2026a), nên không thể khẳng định chưa xem kết quả ở chế độ căng thẳng; kế hoạch khai rõ những kết quả đã xem.

\chapter{Mục 4.3.7 của đề cương: Kiểm định giả thuyết bằng bootstrap, tỷ số Sharpe hiệu chỉnh, sai số chuẩn bền vững và phương trình (5)}

Mục 4.3.7 trả lời: làm sao biết một chênh lệch trong dữ liệu là thật chứ không phải nhiễu? Nửa đầu mục giới thiệu ba công cụ: bootstrap (bốc lại nhiều lần từ chính dữ liệu để xem kết quả dao động bao nhiêu); tỷ số Sharpe hiệu chỉnh (trừ phần thành tích đẹp do chọn lọc); và hồi quy phương trình (5) với sai số chuẩn bền vững, để đo rủi ro lựa chọn mô hình ở chế độ căng thẳng so với bình thường sau khi tính đến các biến kênh.

Phần này đi qua nửa đầu; nửa sau (phân phối rỗng, ba điều kiện ủng hộ H1, tập mô hình tĩnh, phần của \texorpdfstring{$\beta$}{beta} do biến kênh giải thích, hệ số phóng đại phương sai) ở phần tiếp theo.

\section{Dữ liệu tài chính ít và ồn}

Chuỗi lợi suất có ba đặc điểm gây khó: mẫu không lớn, nhất là số giai đoạn căng thẳng; các tuần kề nhau có thể liên quan nhau nên không độc lập; và độ biến động thay đổi theo thời gian.

Công cụ cổ điển thường giả định mẫu độc lập và biến động đều; giả định sai thì kết luận tự tin hơn mức nên có. Theo bản đọc, mục 4.3.7 chọn công cụ chịu được các đặc điểm ấy: bootstrap giữ thứ tự thời gian, và sai số chuẩn bền vững trước phương sai thay đổi và tự tương quan.

\section{Bootstrap theo khối}

Nhắc lại (nhóm (3) của mục 4.2): bootstrap coi dữ liệu như một túi quan sát, bốc có hoàn lại cho đủ một mẫu mới cùng cỡ, tính lại con số cần quan tâm, lặp nhiều lần (ví dụ 1.000). Độ rộng của các con số thu được cho biết con số gốc dao động bao nhiêu.

Với chuỗi thời gian, bốc từng tuần lẻ sẽ xáo trộn thứ tự và phá vỡ các cụm tuần êm hay dữ, nên mẫu giả kém giống thật.

Vì vậy đề cương dùng bootstrap chuỗi thời gian lấy mẫu lặp theo khối liên tiếp (Ledoit và Wolf, 2008): bốc từng nắm tuần liền nhau rồi ghép lại, giữ nguyên thứ tự bên trong mỗi nắm.

\begin{tcolorbox}[breakable]
Ví dụ giả định (số tự đặt để minh họa): 12 tuần chia thành bốn khối ba tuần: A (1--3), B (4--6), C (7--9), D (10--12).

Một lần bốc có hoàn lại bốn khối, ví dụ C, A, C, D: mẫu giả là tuần 7--9, 1--3, 7--9, 10--12. Khối C xuất hiện hai lần, B không xuất hiện.

Thứ tự trong mỗi khối được giữ, nên các cụm êm hay dữ vẫn nằm cạnh nhau. Từ mẫu giả tính lại, ví dụ, chênh lệch điểm số giữa hai mô hình; lặp 1.000 lần được 1.000 chênh lệch.
\end{tcolorbox}

Đề cương chưa nêu độ dài khối; hãy nói thẳng nếu được hỏi. Nguyên tắc chung: khối đủ dài để giữ liên hệ theo thời gian, nhưng không dài đến mức còn quá ít khối để bốc.

Hai biến thể: khối liên tiếp của Ledoit và Wolf (2008) cho khoảng tin cậy của lợi suất tương đương chắc chắn và tỷ số Sharpe; bootstrap khối dừng của Politis và Romano (1994) cho các hệ số hồi quy. Cùng nguyên tắc: bốc từng nắm liền nhau.

Chữ “dừng” là thuật ngữ của Politis và Romano; hiểu khái quát (ngoài đề cương), biến thể này cho độ dài khối thay đổi ngẫu nhiên.

Với mỗi mẫu giả, chạy lại hồi quy và ghi hệ số; khoảng tin cậy là khoảng chứa phần lớn các hệ số thu được. Không chứa 0 thì có cơ sở nói hệ số khác 0 một cách có hệ thống; chứa 0 thì chênh lệch có thể chỉ là nhiễu. Đề cương dùng chính điều kiện “khoảng tin cậy loại trừ 0” trong tiêu chí ủng hộ giả thuyết.

\begin{tcolorbox}[breakable]
Ví dụ giả định (số tự đặt để minh họa): hệ số ước lượng 0,9; 1.000 lần bootstrap khối, lấy 90\% ở giữa được khoảng 0,3 đến 1,5.

Khoảng không chứa 0: hệ số được coi là khác 0.

Nếu khoảng là −0,2 đến 2,0 thì chứa 0, không loại trừ được khả năng hệ số thật bằng 0. Mức 90\% chỉ để minh họa; đề cương không nêu mức tin cậy ở đây.
\end{tcolorbox}

\section{Tỷ số Sharpe hiệu chỉnh}

Đề cương hiệu chỉnh thiên lệch chọn lọc bằng tỷ số Sharpe hiệu chỉnh (Bailey và López de Prado, 2014), công cụ cho thiên lệch chọn lọc và quá khớp khi kiểm định ngược.

Thử nhiều thứ rồi giữ thứ trông tốt nhất thì thứ đó thường tốt hơn mức thật, vì nó nhặt cả phần may: cùng cơ chế với giá trị lớn nhất của các đại lượng có nhiễu ở phần mức hối tiếc.

Một trăm người đoán tung đồng xu mười lần; có người đoán đúng chín, mười lần chỉ nhờ may. Phỏng vấn người giỏi nhất để hỏi bí quyết thì nghe chuyện thuyết phục mà chẳng có bí quyết nào.

Theo bản đọc, đề tài có nhiều chỗ sinh thiên lệch chọn lọc (bảy mô hình, nhiều siêu tham số, nhiều cửa sổ, nhiều biến kênh), nên đề cương dùng cả quy trình (cùng ngân sách tinh chỉnh, đăng ký trước, hiệu chỉnh so sánh bội Holm) lẫn tỷ số Sharpe hiệu chỉnh.

Đề cương không viết công thức và không nói áp dụng cụ thể cho đại lượng nào; chỉ nên nói đúng như vậy.

\section{So sánh giữa các mô hình bằng khoảng tin cậy}

Với hai mô hình A và B, mỗi mẫu giả bốc theo khối cho một hiệu điểm số; sau nhiều lần lặp được một đám hiệu, và khoảng tin cậy cho biết hiệu thật có thể nằm ở đâu.

Không chứa 0 thì A và B khác nhau vượt quá nhiễu; chứa 0 thì chưa đủ cơ sở. Cùng cách làm cho tỷ số Sharpe, vốn chỉ là thước phụ báo cáo ở cửa sổ có lợi suất vượt trội trung bình dương.

Cùng kỹ thuật xuất hiện ở hai chỗ: dựng khoảng tin cậy cho chênh lệch giữa các mô hình, và lấy sai số chuẩn của chính chênh lệch đó làm mẫu số khi chuẩn hóa mức hối tiếc và độ phân tán (mục 4.3.6).

\section{Phương trình (5): đường xấp xỉ liên hệ giữa rủi ro và chế độ}

Có thước đo rủi ro (\(D_{t}\) hay \(R_{t}\)) và công cụ đo nhiễu rồi, việc còn lại là trả lời H1: rủi ro có cao hơn ở chế độ căng thẳng không? Đề cương dùng hồi quy, tức đường xấp xỉ cho biết khi một đại lượng thay đổi thì đại lượng kia thay đổi trung bình bao nhiêu.

Để kiểm định H1, H2 và H4, đề tài hồi quy các thước đo rủi ro mô hình theo biến chế độ và các biến kênh. H3 không dùng hồi quy này mà so cơ chế giảm thiểu với chọn một mô hình (mục 4.3.8).

\(Y_{t}\  = \ \alpha\  + \ \beta \cdot S_{t}\  + \ \gamma'X_{t}\  + \ \varepsilon_{t}\) (5)

\(Y_{t}\) là biến phụ thuộc. Đề cương nói \(Y_{t}\) là \(D_{t}\) hoặc \(R_{t}\): hai hồi quy riêng, một cho độ phân tán, một cho mức hối tiếc. Mỗi cửa sổ đánh giá là một quan sát; trong phân tích chính, hai đại lượng là bản đã chia cho sai số chuẩn bootstrap.

\(\alpha\) là hệ số chặn: mức nền của \(Y_{t}\), giá trị mà \(Y_{t}\) nhận khi mọi biến giải thích bằng 0, chỗ đường xấp xỉ cắt trục tung.

\(S_{t}\) là biến chế độ: tỷ lệ số ngày trong cửa sổ có xác suất lọc của chế độ căng thẳng vượt ngưỡng đăng ký trước. Nó chạy từ 0 (không ngày nào vượt ngưỡng) đến 1 (mọi ngày đều vượt).

\(\beta\) là hệ số của \(S_{t}\), đo chênh lệch của \(Y_{t}\) giữa chế độ căng thẳng và chế độ bình thường. Vì \(S_{t}\) chạy từ 0 đến 1, \(\beta\) là mức thay đổi trung bình của \(Y_{t}\) khi đi từ cửa sổ hoàn toàn bình thường tới cửa sổ hoàn toàn căng thẳng, các thứ khác giữ nguyên.

\(X_{t}\) là véc-tơ biến kênh (một danh sách đại lượng đi cùng nhau), \(\gamma\) là véc-tơ hệ số; mỗi biến kênh có hệ số riêng, và \(\gamma\) gom các hệ số đó.

\(\varepsilon_{t}\) là sai số: phần của \(Y_{t}\) các biến không giải thích được, khoảng cách từ từng điểm đến đường xấp xỉ.

Bằng lời: rủi ro lựa chọn trong một cửa sổ bằng mức nền, cộng phần do cửa sổ nằm ở chế độ căng thẳng nhiều hay ít, cộng phần do các biến kênh, cộng phần ngẫu nhiên. Bốn phần cộng lại cho \(Y_{t}\).

\begin{tcolorbox}[breakable]
Ví dụ giả định (số tự đặt để minh họa): cửa sổ 120 ngày, 30 ngày có xác suất lọc vượt ngưỡng, nên \(S_{t}\) = 0,25.

Giả sử \(\alpha\) = 0,5; \(\beta\) = 1,2; hệ số của độ phi thanh khoản 0,4; hệ số của tương quan đuôi −0,3; cửa sổ có độ phi thanh khoản 1,5 (đơn vị tùy ý) và tương quan đuôi 0,6.

Phần nền 0,5; phần chế độ 1,2 × 0,25 = 0,3; phần phi thanh khoản 0,4 × 1,5 = 0,6; phần tương quan đuôi −0,3 × 0,6 = −0,18.

Giá trị dự đoán của \(Y_{t}\) là 0,5 + 0,3 + 0,6 − 0,18 = 1,22; giá trị thật lệch khỏi đó một lượng bằng \(\varepsilon_{t}\).

Nếu cửa sổ hoàn toàn căng thẳng, \(S_{t}\) = 1, phần chế độ là 1,2 thay vì 0,3, tức tăng 0,9, bằng \(\beta\) nhân phần chênh của \(S_{t}\) (1,2 × 0,75).
\end{tcolorbox}

H1 chỉ cần \(S_{t}\), nhưng Câu hỏi 2 muốn biết biến kênh nào có liên hệ thống kê với rủi ro. Đưa biến kênh vào cùng phương trình tách hai câu: sau khi tính biến kênh, chế độ căng thẳng còn liên hệ thế nào với rủi ro; và biến kênh chiếm bao nhiêu phần của hiệu ứng chế độ (hiệu \texorpdfstring{$\beta$}{beta}0 − \texorpdfstring{$\beta$}{beta}1, phần tiếp theo).

Biến động thực hiện luôn có mặt trong véc-tơ biến kênh với vai trò biến kiểm soát; bầy đàn chỉ ở mức thăm dò. Như vậy phương trình (5) phục vụ cả H1, H2 và H4.

\section{Sai số chuẩn bền vững: khi dao động không đều và có “trí nhớ”}

Đề cương hồi quy với sai số chuẩn bền vững trước phương sai thay đổi (độ lớn sai số không đều giữa các quan sát) và tự tương quan (sai số của quan sát này liên quan sai số của quan sát kề bên). Sai số chuẩn này vẫn đáng tin khi hai hiện tượng đó có mặt.

Đo chiều cao đứa trẻ mỗi tháng: tháng mệt thì đo ẩu, sai số lúc to lúc nhỏ; thước lệch thì lệch lặp lại mỗi tháng. Tin rằng sai số đều và độc lập là tự tin quá mức.

Theo bản đọc, dữ liệu của đề tài có thể có cả hai: cửa sổ căng thẳng biến động hơn nên sai số lớn hơn; cửa sổ kề nhau có sai số tương tự vì giai đoạn thị trường kéo dài.

\begin{tcolorbox}[breakable]
Ví dụ giả định (số tự đặt để minh họa): \texorpdfstring{$\beta$}{beta} ước lượng 0,9; coi sai số độc lập và đều thì sai số chuẩn 0,30, tỷ số 3,0, trông rất có ý nghĩa.

Với sai số chuẩn bền vững là 0,60, tỷ số chỉ còn 1,5. Cùng ước lượng, độ tin cậy giảm hẳn khi khai đúng tính chất của sai số.

Sai số chuẩn bền vững không đổi \texorpdfstring{$\beta$}{beta}; nó chỉ đổi mức ta dám tin vào \texorpdfstring{$\beta$}{beta}.
\end{tcolorbox}

Đừng nhầm hai công cụ: sai số chuẩn bền vững tính sai số chuẩn cho hồi quy; bootstrap khối dừng dựng khoảng tin cậy dùng trong tiêu chí ủng hộ giả thuyết. Đề cương dùng cả hai cho phương trình (5).

\section{Các biến kênh: đường dẫn khả dĩ từ căng thẳng đến rủi ro lựa chọn}

Biến kênh gồm độ phi thanh khoản Amihud, tương quan đuôi và tương quan hiệu chỉnh theo biến động, độ phân tán chéo của lợi suất; biến động thực hiện luôn là biến kiểm soát. Đề cương nêu dấu dự đoán cho độ phi thanh khoản, tương quan đuôi và độ phân tán chéo, không nêu cho tương quan hiệu chỉnh theo biến động.

\subsection{Độ phi thanh khoản: khó mua bán đến mức nào}

Độ phi thanh khoản Amihud (Amihud, 2002) đo mức một lệnh mua bán vừa phải cũng làm giá nhích mạnh. Đề cương không viết công thức.

H2 dự đoán dấu dương: thanh khoản cạn làm các mô hình có vòng quay khác nhau cách xa nhau. Giải thích của bản đọc: mô hình mua bán nhiều chịu chi phí lớn hơn khi khó mua bán, nên tách khỏi mô hình mua bán ít; các mô hình xa nhau hơn thì độ phân tán và mức hối tiếc có thể tăng.

Phố ngập thì xe chạy nhiều chặng thiệt hơn xe chạy một chặng; khoảng cách giữa hai loại xe giãn ra.

\subsection{Tương quan đuôi: các cổ phiếu cùng sụp một lúc}

Tương quan đuôi đo mức các cổ phiếu cùng rơi mạnh một lúc. Longin và Solnik (2001) bác bỏ phân phối chuẩn đa biến ở đuôi âm: khi thị trường xấu, các mã có thể cùng đi xuống theo cách mô hình chuẩn không dự đoán được.

H2 dự đoán dấu ÂM, dễ gây ngạc nhiên: tương quan cao làm các danh mục chỉ mua hội tụ. Giải thích của bản đọc: danh mục chỉ mua không bán khống, nên khi mọi mã cùng chuyển động, các cách chia tiền cho kết quả giống nhau hơn; chọn mô hình nào cũng gần như nhau, nên độ phân tán và mức hối tiếc có thể giảm.

Bảy chiếc thuyền gặp cơn bão làm mọi thứ chìm cùng lúc: bảy cách chèo khác nhau, kết quả như nhau.

Đây chỉ là dự đoán sẽ được kiểm định; nếu dữ liệu cho dấu khác, H2 không được ủng hộ.

\subsection{Tương quan hiệu chỉnh theo biến động}

Forbes và Rigobon (2002): hệ số tương quan phụ thuộc vào biến động; sau hiệu chỉnh, gần như không còn dấu hiệu lây nhiễm trong ba cuộc khủng hoảng họ khảo sát.

Thấy tương quan tăng ở chế độ căng thẳng thì phải hỏi tăng thật hay chỉ vì biến động tăng; đó là khả năng đối lập của H2. Theo bản đọc, vì vậy danh sách biến kênh có cả tương quan đuôi lẫn tương quan hiệu chỉnh theo biến động.

Trong phòng ồn, hai người cùng nói to vì tiếng ồn, bạn tưởng họ bàn bạc với nhau. Hiệu chỉnh theo biến động trừ phần cùng phản ứng với tiếng ồn.

\subsection{Độ phân tán chéo của lợi suất: thước đo bầy đàn ở cấp thị trường}

Độ phân tán chéo của lợi suất là thước đo bầy đàn ở cấp thị trường (Christie và Huang, 1995; Chang và cộng sự, 2000). Đề cương không giải thích cơ chế vì sao H4 dự đoán dấu dương.

H4 (thăm dò): sau khi kiểm soát biến động thực hiện, độ phân tán chéo cao đi cùng rủi ro lựa chọn mô hình cao. Đề tài không quan sát hành vi chọn mô hình của tổ chức, nên nếu có liên hệ, đó không phải bằng chứng tổ chức bầy đàn khi chọn mô hình.

\subsection{Bảng tóm tắt biến kênh}

{\small\begin{longtable}[]{@{}
  >{\raggedright\arraybackslash}p{(\columnwidth - 8\tabcolsep) * \real{0.1713}}
  >{\raggedright\arraybackslash}p{(\columnwidth - 8\tabcolsep) * \real{0.2028}}
  >{\raggedright\arraybackslash}p{(\columnwidth - 8\tabcolsep) * \real{0.1940}}
  >{\raggedright\arraybackslash}p{(\columnwidth - 8\tabcolsep) * \real{0.0797}}
  >{\raggedright\arraybackslash}p{(\columnwidth - 8\tabcolsep) * \real{0.3523}}@{}}
\toprule\noalign{}
\begin{minipage}[b]{\linewidth}\raggedright
Biến kênh
\end{minipage} & \begin{minipage}[b]{\linewidth}\raggedright
Nguồn
\end{minipage} & \begin{minipage}[b]{\linewidth}\raggedright
Dự đoán dấu hệ số
\end{minipage} & \begin{minipage}[b]{\linewidth}\raggedright
Giả thuyết
\end{minipage} & \begin{minipage}[b]{\linewidth}\raggedright
Lý do (theo đề cương)
\end{minipage} \\
\midrule\noalign{}
\endhead
\bottomrule\noalign{}
\endlastfoot
Độ phi thanh khoản Amihud & Amihud (2002) & Dương & H2 & Thanh khoản cạn làm các mô hình có vòng quay khác nhau cách xa nhau \\
Tương quan đuôi và tương quan hiệu chỉnh theo biến động & Longin và Solnik (2001); Forbes và Rigobon (2002) & Âm (tương quan đuôi); đề cương không nêu dấu cho tương quan hiệu chỉnh & H2 & Tương quan cao làm các danh mục chỉ mua hội tụ \\
Độ phân tán chéo của lợi suất & Christie và Huang (1995); Chang và cộng sự (2000) & Dương & H4 (thăm dò) & Đề cương không nêu lý do riêng; thước đo bầy đàn ở cấp thị trường \\
\end{longtable}}

\begin{tcolorbox}[breakable]
\textbf{Ghi nhớ khi đi thi}

• Bootstrap khối liên tiếp (Ledoit và Wolf, 2008) cho khoảng tin cậy của lợi suất tương đương chắc chắn và Sharpe; bootstrap khối dừng (Politis và Romano, 1994) cho hệ số hồi quy.

• Thiên lệch chọn lọc được hiệu chỉnh bằng tỷ số Sharpe hiệu chỉnh (Bailey và López de Prado, 2014).

• H1, H2, H4 kiểm định bằng hồi quy theo biến chế độ và biến kênh, với sai số chuẩn bền vững trước phương sai thay đổi và tự tương quan.

• Phương trình (5): \(Y_{t}\) = \(\alpha\) + \(\beta\) \(S_{t}\) + \(\gamma\) nhân \(X_{t}\) + \(\varepsilon_{t}\); \(Y_{t}\) là \(D_{t}\) hoặc \(R_{t}\); \(\beta\) đo chênh lệch giữa chế độ căng thẳng và bình thường.

• \(S_{t}\) là tỷ lệ số ngày trong cửa sổ có xác suất lọc của chế độ căng thẳng vượt ngưỡng đăng ký trước.

• Biến kênh: độ phi thanh khoản (dự đoán dương), tương quan đuôi (âm), tương quan hiệu chỉnh theo biến động, độ phân tán chéo (dương, H4 thăm dò); biến động thực hiện là biến kiểm soát.
\end{tcolorbox}

\section{Tự kiểm tra}

{\itshape 1. Đề cương dùng hai biến thể bootstrap nào, cho đại lượng nào?\par}

Đáp án: khối liên tiếp (Ledoit và Wolf, 2008) cho khoảng tin cậy của lợi suất tương đương chắc chắn và Sharpe; khối dừng (Politis và Romano, 1994) cho hệ số hồi quy.

{\itshape 2. Tỷ số Sharpe hiệu chỉnh dùng để làm gì?\par}

Đáp án: hiệu chỉnh thiên lệch chọn lọc (Bailey và López de Prado, 2014).

{\itshape 3. Hãy đọc từng thành phần của phương trình (5).\par}

Đáp án: \(Y_{t}\) là \(D_{t}\) hoặc \(R_{t}\); \(\alpha\) là hệ số chặn; \(\beta\) đo chênh lệch của \(Y_{t}\) giữa chế độ căng thẳng và bình thường; \(S_{t}\) là tỷ lệ ngày có xác suất lọc vượt ngưỡng đăng ký trước; \(X_{t}\) là véc-tơ biến kênh với hệ số \(\gamma\); \(\varepsilon_{t}\) là sai số.

{\itshape 4. Giả thuyết nào dùng hồi quy phương trình (5)? Giả thuyết nào không?\par}

Đáp án: H1, H2, H4 dùng phương trình (5); H3 được đánh giá bằng so sánh cơ chế giảm thiểu (mục 4.3.8).

{\itshape 5. Dấu dự đoán của độ phi thanh khoản và tương quan đuôi trong H2, và lý do?\par}

Đáp án: độ phi thanh khoản dương, vì thanh khoản cạn làm các mô hình có vòng quay khác nhau cách xa nhau; tương quan đuôi âm, vì tương quan cao làm các danh mục chỉ mua hội tụ.

\chapter{Mục 4.3.7 của đề cương (tiếp): Phân phối rỗng, ba điều kiện ủng hộ H1, mô hình tĩnh, phần của beta do biến kênh và hệ số phóng đại phương sai}

Nửa sau của mục 4.3.7 trả lời: làm sao tách rủi ro lựa chọn thật khỏi nhiễu ước lượng và biến động? Đề cương dựng một thế giới giả trong đó các mô hình ngang nhau (phân phối rỗng, 1.000 lần bootstrap khối), và chỉ coi H1 được ủng hộ khi thỏa đồng thời ba điều kiện. Kết quả được lặp lại trên năm mô hình tĩnh; phần của \texorpdfstring{$\beta$}{beta} do biến kênh giải thích đo bằng hiệu \texorpdfstring{$\beta$}{beta}0 − \texorpdfstring{$\beta$}{beta}1; hệ số phóng đại phương sai được kiểm tra với ngưỡng 5; mọi kết quả là liên hệ có điều kiện, không phải nhân quả.

\section{Vì sao cần phân phối rỗng}

Mục 4.3.6 cho thấy mức hối tiếc có kỳ vọng dương ngay cả khi các mô hình ngang nhau, và nhiễu càng lớn thì mức hối tiếc thô càng lớn. Nếu chế độ căng thẳng nhiễu hơn, một hệ số \texorpdfstring{$\beta$}{beta} dương có thể chỉ phản ánh nhiễu lớn hơn, chưa chắc rủi ro lựa chọn thật sự cao hơn.

Bước thứ nhất đã làm: chia \(D_{t}\) và \(R_{t}\) cho sai số chuẩn bootstrap. Bước thứ hai: hiệu chuẩn \texorpdfstring{$\beta$}{beta} bằng phân phối rỗng, để tách rủi ro lựa chọn khỏi nhiễu ước lượng và biến động. Hiệu chuẩn là chỉnh thước đo cho khớp một chuẩn; chuẩn ở đây là kết quả sẽ thấy nếu các mô hình ngang nhau.

Phân phối rỗng là “kết quả trông thế nào nếu thực ra không có gì xảy ra”: không mô hình nào hơn mô hình nào, mọi chênh lệch chỉ là nhiễu. \texorpdfstring{$\beta$}{beta} thật vượt rõ mức mà thế giới giả tạo ra thì có cơ sở nói còn gì đó ngoài nhiễu.

Nghi đồng xu thiên vị vì tung 100 lần được 58 mặt ngửa? Hãy tung một đồng xu công bằng 100 lần, lặp hàng nghìn lần, xem 58 nằm đâu trong đám kết quả. Đám đó là phân phối rỗng; 58 nằm ở rìa xa mới đáng ngờ.

Khác với đồng xu, phân phối rỗng ở đây không nhất thiết có tâm ở 0: ngay trong thế giới giả, mức hối tiếc vẫn dương và có thể đổi theo mức nhiễu của từng chế độ. Vì vậy đề cương so với phân vị 95\% của phân phối rỗng, không chỉ so với 0.

\section{Dựng phân phối rỗng từng bước}

Đề cương mô tả: phân phối rỗng gồm 1.000 lần tái lấy mẫu bootstrap khối chuỗi lợi suất của bảy mô hình; trước khi tái lấy mẫu, loại chênh lệch trung bình giữa các mô hình trong từng chế độ nhưng giữ nguyên tương quan chéo và biến động. Gỡ thành bốn bước:

\textbf{Bước một: lấy chuỗi lợi suất danh mục của bảy mô hình} trên toàn mẫu, sau chi phí: bảy chuỗi song song.

\textbf{Bước hai: loại chênh lệch trung bình trong từng chế độ.} Trong mỗi chế độ, dịch chuỗi của từng mô hình sao cho bảy mô hình có cùng lợi suất trung bình (cách hiểu tự nhiên của bản đọc). Đó là thế giới “các mô hình ngang chất lượng”.

Làm riêng từng chế độ vì mô hình có thể hơn ở chế độ này, kém ở chế độ kia, và đề tài quan tâm đúng sự đổi thứ hạng đó; thế giới giả phải xóa cả lợi thế theo chế độ.

\textbf{Bước ba: giữ nguyên tương quan chéo và biến động}, tức mức các mô hình cùng lên xuống và độ lớn dao động. Nhờ vậy nhiễu trong thế giới giả có đúng mức và đúng cấu trúc như nhiễu thật.

\textbf{Bước bốn: tái lấy mẫu bootstrap khối 1.000 lần.} Mỗi lần cho một thế giới giả với bảy chuỗi; theo bản đọc, với mỗi thế giới ta tính lại mức hối tiếc hay độ phân tán, hồi quy theo biến chế độ và ghi \texorpdfstring{$\beta$}{beta}. Sau 1.000 lần có 1.000 hệ số \texorpdfstring{$\beta$}{beta} rỗng.

Đề cương chưa nói cách gắn nhãn chế độ vào chuỗi đã tái lấy mẫu; chi tiết này sẽ nằm trong tệp đăng ký hoặc luận án, đừng đoán.

Muốn biết lớp buổi sáng có thật giỏi hơn lớp buổi chiều: san cho hai lớp cùng điểm trung bình nhưng giữ cách điểm dao động của từng em, rồi mô phỏng hàng nghìn lần. Chênh lệch nào xuất hiện trong mô phỏng chỉ có thể do may; đó là thước để đánh giá chênh lệch thật.

\begin{tcolorbox}[breakable]
Ví dụ giả định (số tự đặt để minh họa): 1.000 hệ số \texorpdfstring{$\beta$}{beta} rỗng của mức hối tiếc nằm từ −0,2 đến 0,7; giá trị thứ 950 khi xếp tăng dần là 0,40: phân vị 95\%.

Tức là khi các mô hình ngang nhau, nhiễu một mình chỉ có 5\% khả năng tạo ra \texorpdfstring{$\beta$}{beta} lớn hơn 0,40.

Trường hợp A: \texorpdfstring{$\beta$}{beta} thật là 0,90, vượt hẳn 0,40; nhiễu khó giải thích.

Trường hợp B: \texorpdfstring{$\beta$}{beta} thật là 0,30, dương nhưng dưới 0,40; không đủ ủng hộ H1.

Trong trường hợp B, chỉ so với 0 thì dễ tưởng có hiệu ứng; phân phối rỗng chặn cách đọc đó.
\end{tcolorbox}

Vì sao cần cả phép chia lẫn phân phối rỗng? Theo bản đọc: phép chia đưa mọi cửa sổ về cùng thang nhiễu; phân phối rỗng kiểm tra phần cơ học còn sót sau phép chia.

\section{Ba điều kiện ủng hộ H1}

Đề cương coi H1 được ủng hộ khi thỏa đồng thời ba điều kiện: \texorpdfstring{$\beta$}{beta} của \(R_{t}\) dương và vượt phân vị 95\% của phân phối rỗng; khoảng tin cậy bootstrap khối của \texorpdfstring{$\beta$}{beta} loại trừ 0; và \texorpdfstring{$\beta$}{beta} của \(D_{t}\) cùng dấu.

{\small\begin{longtable}[]{@{}
  >{\raggedright\arraybackslash}p{(\columnwidth - 4\tabcolsep) * \real{0.0849}}
  >{\raggedright\arraybackslash}p{(\columnwidth - 4\tabcolsep) * \real{0.4390}}
  >{\raggedright\arraybackslash}p{(\columnwidth - 4\tabcolsep) * \real{0.4761}}@{}}
\toprule\noalign{}
\begin{minipage}[b]{\linewidth}\raggedright
Điều kiện
\end{minipage} & \begin{minipage}[b]{\linewidth}\raggedright
Nội dung
\end{minipage} & \begin{minipage}[b]{\linewidth}\raggedright
Ý nghĩa đơn giản
\end{minipage} \\
\midrule\noalign{}
\endhead
\bottomrule\noalign{}
\endlastfoot
1 & \texorpdfstring{$\beta$}{beta} của \(R_{t}\) dương và vượt phân vị 95\% của phân phối rỗng & Mức hối tiếc cao hơn ở chế độ căng thẳng, vượt cái nhiễu một mình tạo ra \\
2 & Khoảng tin cậy bootstrap khối của \texorpdfstring{$\beta$}{beta} loại trừ 0 (\texorpdfstring{$\beta$}{beta} của \(R_{t}\)) & Ước lượng đủ chắc, không chứa khả năng bằng 0 \\
3 & \texorpdfstring{$\beta$}{beta} của độ phân tán mô hình cùng dấu & Thước đo phụ đi cùng chiều thước đo chính \\
\end{longtable}}

Điều kiện 1 nói về dấu và về nhiễu; điều kiện 2 về độ chắc chắn của ước lượng; điều kiện 3 về sự nhất quán giữa thước đo chính và phụ. Thiếu điều kiện nào cũng mất một lớp bảo vệ.

Theo bản đọc, yêu cầu cùng dấu loại những bức tranh lệch pha: khoảng cách giữa các mô hình không rộng ra mà thiệt hại lại lớn hơn, hoặc ngược lại.

Kích thước tập tin cậy mô hình là kết quả phụ, báo cáo mô tả, không nằm trong điều kiện ủng hộ. Các tiêu chí này thuộc kế hoạch đăng ký trước, tức được viết trước khi xem dữ liệu mới.

\begin{tcolorbox}[breakable]
Ví dụ giả định (số tự đặt để minh họa): phân vị 95\% của phân phối rỗng là 0,40.

Kịch bản 1: \texorpdfstring{$\beta$}{beta} của mức hối tiếc 0,90, khoảng tin cậy 0,30 đến 1,50, \texorpdfstring{$\beta$}{beta} của độ phân tán 0,60. Ba điều kiện cùng đúng: H1 được ủng hộ.

Kịch bản 2: \texorpdfstring{$\beta$}{beta} 0,90 nhưng khoảng tin cậy −0,20 đến 2,00, chứa 0: điều kiện 2 sai, H1 không được ủng hộ dù ước lượng điểm lớn.

Kịch bản 3: \texorpdfstring{$\beta$}{beta} 0,30, dưới 0,40, dù khoảng tin cậy 0,10 đến 0,50: điều kiện 1 sai.

Kịch bản 4: \texorpdfstring{$\beta$}{beta} 0,90, đủ hai điều kiện đầu, nhưng \texorpdfstring{$\beta$}{beta} của độ phân tán −0,10: điều kiện 3 sai.
\end{tcolorbox}

Mục 4.3.10 thêm: khi số cửa sổ độc lập ở chế độ căng thẳng dưới 10, kết quả gọi là không kết luận, không gọi là ủng hộ hay bác bỏ. Tiêu chí trên áp dụng trong khuôn khổ quy tắc đó.

Nếu dữ liệu bác bỏ giả thuyết, đề tài vẫn trả lời được câu hỏi: rủi ro không đổi theo chế độ, hoặc 1/N đủ dùng ở Việt Nam.

\section{Lặp lại trên tập mô hình tĩnh}

Đề tài lặp lại kiểm định trên tập chỉ gồm M1, M2, M3, M4, M6, để loại tác động của M5 và M7, vốn chỉ đổi hành vi khi xác suất căng thẳng tăng.

Năm đội luôn chạy một chiến thuật, hai đội đổi chiến thuật ngay khi trời có dấu hiệu mưa. Thấy các đội chênh nhau hơn khi mưa thì khó biết do thời tiết hay do hai đội chủ động đổi. Chỉ giữ năm đội không đổi thì phần chênh còn lại không đến từ việc đổi chiến thuật.

Đề cương chưa quy định cách gọi kết luận nếu tập đầy đủ và tập tĩnh cho kết quả khác nhau; hãy nói đúng như vậy.

Ba phép kiểm tra liên quan M5 và M7:

{\small\begin{longtable}[]{@{}
  >{\raggedright\arraybackslash}p{(\columnwidth - 6\tabcolsep) * \real{0.1865}}
  >{\raggedright\arraybackslash}p{(\columnwidth - 6\tabcolsep) * \real{0.1813}}
  >{\raggedright\arraybackslash}p{(\columnwidth - 6\tabcolsep) * \real{0.0797}}
  >{\raggedright\arraybackslash}p{(\columnwidth - 6\tabcolsep) * \real{0.5526}}@{}}
\toprule\noalign{}
\begin{minipage}[b]{\linewidth}\raggedright
Phép kiểm tra
\end{minipage} & \begin{minipage}[b]{\linewidth}\raggedright
Tập mô hình
\end{minipage} & \begin{minipage}[b]{\linewidth}\raggedright
Mục
\end{minipage} & \begin{minipage}[b]{\linewidth}\raggedright
Mục đích
\end{minipage} \\
\midrule\noalign{}
\endhead
\bottomrule\noalign{}
\endlastfoot
Phân tích chính & Cả bảy mô hình (M = 7) & 4.3.7 & Kiểm định H1 \\
Báo cáo không có M5 & Sáu mô hình (bỏ M5) & 4.3.5 & Xem kết quả có phụ thuộc mô hình do tác giả phát triển không \\
Lặp lại trên mô hình tĩnh & M1, M2, M3, M4, M6 & 4.3.7 & Loại tác động của M5 và M7, vốn đổi hành vi theo xác suất căng thẳng \\
\end{longtable}}

\section{Phần của \texorpdfstring{$\beta$}{beta} do biến kênh giải thích: hiệu \texorpdfstring{$\beta$}{beta}0 trừ \texorpdfstring{$\beta$}{beta}1}

Câu hỏi 2 hỏi các biến kênh giải thích bao nhiêu phần chênh lệch giữa hai chế độ, tức bao nhiêu phần của \texorpdfstring{$\beta$}{beta} ở phương trình (5). Đề cương dùng một phép trừ.

Phần đó ước lượng bằng hiệu \texorpdfstring{$\beta$}{beta}0 − \texorpdfstring{$\beta$}{beta}1: \texorpdfstring{$\beta$}{beta}0 là hệ số của \(X_{t}\) khi phương trình không có \(X_{t}\); \texorpdfstring{$\beta$}{beta}1 là hệ số của \(S_{t}\) khi có \(S_{t}\); kèm khoảng tin cậy bootstrap khối. Hiệu cho biết phần liên hệ giữa chế độ căng thẳng và rủi ro mà biến kênh “nuốt” mất. (Ký hiệu \(S_{t}\) là biến chế độ.)

Ba bước: chạy hồi quy chỉ có \(S_{t}\), ghi \texorpdfstring{$\beta$}{beta}0; chạy hồi quy có cả \(S_{t}\) và \(X_{t}\), ghi \texorpdfstring{$\beta$}{beta}1; lấy hiệu và dựng khoảng tin cậy bằng bootstrap khối.

Khu phố nhiều ô tô thì nhiều tai nạn. Thêm biến tốc độ trung bình vào phân tích mà liên hệ giữa số ô tô và tai nạn co lại nhiều, phần co lại là phần tốc độ giải thích. \texorpdfstring{$\beta$}{beta}0 − \texorpdfstring{$\beta$}{beta}1 là phần co lại ấy.

\begin{tcolorbox}[breakable]
Ví dụ giả định (số tự đặt để minh họa): hồi quy chỉ có \(S_{t}\) cho \texorpdfstring{$\beta$}{beta}0 = 1,20; có cả biến kênh cho \texorpdfstring{$\beta$}{beta}1 = 0,75.

Hiệu 0,45: trong mức tăng 1,20, biến kênh chiếm 0,45, phần còn lại 0,75 vẫn gắn với chế độ căng thẳng.

Theo tỷ lệ là 37,5\%; đề cương chỉ nêu hiệu, cách đổi sang tỷ lệ là của bản đọc.

Khoảng tin cậy của hiệu là 0,10 đến 0,80 thì loại trừ 0, có cơ sở nói biến kênh chiếm một phần dương; −0,10 đến 0,90 thì chưa.
\end{tcolorbox}

Nếu \texorpdfstring{$\beta$}{beta}1 lớn hơn \texorpdfstring{$\beta$}{beta}0 thì hiệu âm; đề cương không bàn trường hợp này, đừng diễn giải như một kết luận.

Phép tính này thuộc Câu hỏi 2 (H2, H4), không thuộc tiêu chí ủng hộ H1. H2 được ủng hộ khi hệ số độ phi thanh khoản dương, hệ số tương quan đuôi âm, và khoảng tin cậy bootstrap khối của cả hai loại trừ 0; H4 chỉ báo cáo như thăm dò.

\section{Hệ số phóng đại phương sai, ngưỡng 5}

Vì \(S_{t}\) phụ thuộc biến động và lợi suất, còn \(X_{t}\) gồm những biến cùng phụ thuộc biến động, đề tài báo cáo hệ số phóng đại phương sai với ngưỡng 5 ấn định trước. Vượt ngưỡng thì báo cáo hai phương trình tách riêng và không diễn giải riêng từng hệ số.

Vấn đề là đa cộng tuyến: các biến giải thích đi cùng nhau quá chặt; biến động tăng thì tất cả cùng tăng, nên hồi quy khó tách phần đóng góp riêng của từng biến.

Hệ số phóng đại phương sai cho biết phương sai của một hệ số bị phình lên bao nhiêu lần vì biến đó bị các biến khác giải thích. Kiến thức thống kê chung (đề cương không viết công thức): hệ số = 1 / (1 − R²) của hồi quy biến đó theo các biến còn lại; bằng 5 ứng với R² = 0,8.

Hai người cùng đẩy một chiếc xe, cùng lúc, cùng lực: xe đi tới nhưng không thể nói mỗi người đẩy bao nhiêu phần. Hồi quy có thể cho một hệ số rất lớn và một hệ số rất nhỏ mà tổng vẫn đúng.

\begin{tcolorbox}[breakable]
Ví dụ giả định (số tự đặt để minh họa): R² = 0,9 thì hệ số = 10, vượt ngưỡng.

R² = 0,5 thì hệ số = 2, dưới ngưỡng.
\end{tcolorbox}

Vượt ngưỡng thì đề cương không cố cứu bằng phép tính: báo cáo hai phương trình tách riêng. Đề cương không nói rõ mỗi phương trình gồm biến nào.

Theo bản đọc: các biến luôn đi cùng nhau thì dữ liệu không đủ thông tin để tách hiệu ứng riêng, nên không diễn giải riêng; tinh thần gần với quy tắc không kết luận.

\section{Liên hệ có điều kiện, không phải nhân quả}

Đề tài gọi các kết quả là liên hệ có điều kiện, không gọi là quan hệ nhân quả. Liên hệ có điều kiện: sau khi giữ các biến khác (biến động thực hiện, biến kênh), hai đại lượng vẫn đi cùng nhau. Nhân quả: đại lượng này gây ra thay đổi ở đại lượng kia.

Ngày bán nhiều kem cũng bán nhiều đồ bơi, kể cả khi đã tính ngày cuối tuần; nhưng kem không làm tăng bán đồ bơi, cả hai cùng theo thời tiết nóng.

Nếu hệ số độ phi thanh khoản dương, chỉ được nói thanh khoản cạn đi cùng rủi ro lựa chọn cao hơn sau khi kiểm soát biến động, không nói gây ra. Tương tự với \texorpdfstring{$\beta$}{beta} của chế độ căng thẳng.

Theo bản đọc: biến chế độ và biến kênh cùng phụ thuộc biến động, giai đoạn căng thẳng độc lập ít, và đề tài chỉ quan sát dữ liệu lịch sử chứ không thể chủ động thay đổi một biến rồi quan sát.

\section{Ba điều nhắc cuối về giới hạn}

Thứ nhất: mục 4.3.7 là phương pháp dự kiến, chưa có kết quả; mọi con số trong ví dụ đều tự đặt.

Thứ hai (mục 4.3.10): cạnh mỗi hệ số, đề tài báo cáo số cửa sổ độc lập và số giai đoạn căng thẳng độc lập; dưới 10 cửa sổ độc lập ở chế độ căng thẳng thì kết quả là không kết luận, kèm kích thước hiệu ứng và khoảng tin cậy theo từng giai đoạn.

Thứ ba: hiệu chỉnh so sánh bội theo Holm (1979) ở mức 10\% cho bốn giả thuyết, để tránh việc kiểm nhiều giả thuyết thì có một cái trông đúng chỉ do may. Đề cương không nói Holm kết hợp với phân vị 95\% của phân phối rỗng ra sao.

\begin{tcolorbox}[breakable]
\textbf{Ghi nhớ khi đi thi}

• Phân phối rỗng: 1.000 lần bootstrap khối chuỗi lợi suất của bảy mô hình, sau khi loại chênh lệch trung bình trong từng chế độ, giữ tương quan chéo và biến động.

• H1 được ủng hộ khi đồng thời: \texorpdfstring{$\beta$}{beta} của mức hối tiếc dương và vượt phân vị 95\% của phân phối rỗng; khoảng tin cậy bootstrap khối của \texorpdfstring{$\beta$}{beta} loại trừ 0; \texorpdfstring{$\beta$}{beta} của độ phân tán cùng dấu.

• Lặp lại trên tập mô hình tĩnh M1, M2, M3, M4, M6 để loại tác động của M5 và M7.

• Phần của \texorpdfstring{$\beta$}{beta} do biến kênh: \texorpdfstring{$\beta$}{beta}0 − \texorpdfstring{$\beta$}{beta}1 (không có và có biến kênh), kèm khoảng tin cậy bootstrap khối.

• Hệ số phóng đại phương sai, ngưỡng 5; vượt thì hai phương trình tách riêng, không diễn giải riêng từng hệ số.

• Liên hệ có điều kiện, không phải nhân quả; dưới 10 cửa sổ độc lập ở chế độ căng thẳng thì không kết luận.
\end{tcolorbox}

\section{Tự kiểm tra}

{\itshape 1. Phân phối rỗng được dựng bằng những thao tác nào và dùng để làm gì?\par}

Đáp án: 1.000 lần tái lấy mẫu bootstrap khối chuỗi lợi suất của bảy mô hình, sau khi loại chênh lệch trung bình giữa các mô hình trong từng chế độ, giữ nguyên tương quan chéo và biến động; dùng để hiệu chuẩn \texorpdfstring{$\beta$}{beta}, tách rủi ro lựa chọn khỏi nhiễu và biến động.

{\itshape 2. Nêu ba điều kiện để H1 được ủng hộ.\par}

Đáp án: \texorpdfstring{$\beta$}{beta} của mức hối tiếc dương và vượt phân vị 95\% của phân phối rỗng; khoảng tin cậy bootstrap khối của \texorpdfstring{$\beta$}{beta} loại trừ 0; \texorpdfstring{$\beta$}{beta} của độ phân tán mô hình cùng dấu.

{\itshape 3. Vì sao đề cương lặp lại phân tích trên M1, M2, M3, M4, M6?\par}

Đáp án: để loại tác động của M5 và M7, vốn chỉ đổi hành vi khi xác suất căng thẳng tăng.

{\itshape 4. Phần của \texorpdfstring{$\beta$}{beta} được giải thích bởi biến kênh ước lượng thế nào?\par}

Đáp án: bằng hiệu \texorpdfstring{$\beta$}{beta}0 − \texorpdfstring{$\beta$}{beta}1 giữa hệ số \texorpdfstring{$\beta$}{beta} của phương trình (5) khi không có và khi có biến kênh, kèm khoảng tin cậy bootstrap khối.

{\itshape 5. Hệ số phóng đại phương sai dùng để làm gì, ngưỡng bao nhiêu, vượt ngưỡng thì sao?\par}

Đáp án: vì biến chế độ và biến kênh cùng phụ thuộc biến động; ngưỡng 5 ấn định trước; vượt thì báo cáo hai phương trình tách riêng, không diễn giải riêng từng hệ số.

{\itshape 6. Kết quả của mục 4.3.7 được gọi là gì, không được gọi là gì, và khi nào kết quả là không kết luận?\par}

Đáp án: liên hệ có điều kiện, không phải nhân quả; dưới 10 cửa sổ độc lập ở chế độ căng thẳng thì không kết luận.

\chapter{Mục 4.3.8 của đề cương: Ba cơ chế giảm thiểu rủi ro lựa chọn mô hình}

Mục 4.3.8 trả lời câu hỏi thứ ba của đề cương: cơ chế nào giảm được rủi ro lựa chọn mô hình. Tác giả đề cương đánh giá ba cơ chế, mỗi cơ chế kèm một đối chứng hoặc một mốc so sánh. Cả ba được chấm bằng sáu tiêu chí, tính sau chi phí giao dịch. Mọi nội dung dưới đây là kế hoạch kiểm tra, chưa phải kết quả.

\section{Bối cảnh: vì sao cần một “thuốc giảm đau”}

Trước khi nói về cách giảm, bạn cần nhớ rủi ro đang cần giảm là gì. Hình dung như sau: trước mùa giải, một tổ chức đầu tư phải chọn một cách chia tiền vào các cổ phiếu. Đề cương gọi mỗi cách chia là một mô hình danh mục, và dùng bảy mô hình, từ M1 đến M7.

Theo quy tắc của đề cương, mô hình được chọn trước là mô hình có mức thỏa dụng tương đương chắc chắn (lợi suất tương đương chắc chắn, certainty equivalent) cao nhất trong 36 tháng trước ngày tái cân bằng. Sau thực tế mới biết cách nào hóa ra tốt nhất. Chênh lệch giữa cách tốt nhất sau thực tế và cách đã chọn là mức hối tiếc, tức khoản thiệt vì đã không chọn mô hình tốt nhất sau thực tế.

Đề cương nêu giả thuyết H1: mức hối tiếc ở chế độ căng thẳng (stress regime) cao hơn ở chế độ bình thường. Nếu H1 được ủng hộ, câu hỏi tiếp theo là làm gì với nó. Mục 4.3.8 đưa ra ba hướng trả lời, và giả thuyết H3 là phát biểu có thể kiểm định của hướng trả lời đó.

Ba hướng này khác nhau ở cách nghĩ. Hướng thứ nhất là đừng chọn một mô hình, hãy kết hợp nhiều mô hình. Hướng thứ hai là chọn một danh mục chịu được tình huống xấu nhất trong nhiều niềm tin khác nhau. Hướng thứ ba là đổi hành vi danh mục khi xác suất thị trường rơi vào chế độ căng thẳng tăng lên. Cách tóm tắt ba hướng này là cách diễn giải của người viết bài đọc này, còn tên và nguồn của từng cơ chế là của đề cương.

\section{Cơ chế (i): kết hợp mô hình thay vì chọn một mô hình}

Cơ chế (i) là kết hợp mô hình theo trọng số Bayes, hoặc chỉ kết hợp các mô hình nằm trong tập tin cậy mô hình. Đối chứng của nó là kết hợp quy tắc 1/N với quy tắc trung bình-phương sai.

Hình dung như sau: thay vì giao toàn bộ tiền cho một cố vấn, bạn lập một hội đồng cố vấn và để mỗi người cầm một phần tiền. Nếu cố vấn được chọn hóa ra dở trong mùa bão, hội đồng vẫn còn những người khác đỡ lại. Đó là trực giác đứng sau việc kết hợp mô hình.

\subsection{Trọng số Bayes: chia phần theo mức được dữ liệu ủng hộ}

Đề cương dẫn Hoeting và cộng sự (1999), người trình bày phương pháp trung bình hóa mô hình Bayes (trung bình hóa mô hình Bayes, Bayesian model averaging). Đề cương cũng dẫn Avramov (2002), người áp dụng trung bình hóa mô hình Bayes vào dự báo lợi suất cổ phiếu, với hiệu quả ngoài mẫu vượt các tiêu chí chọn mô hình.

Giải thích chung: ở mức trực giác, trọng số Bayes cho mỗi mô hình một phần nhất định thay vì chỉ giữ một mô hình duy nhất và bỏ phần còn lại.

Bài gốc không viết công thức tính trọng số cho cơ chế này, nên đề cương chưa quy định cách tính cụ thể ở mục 4.3.8.

\subsection{Chỉ kết hợp trong tập tin cậy mô hình: bỏ những ứng viên đã bị loại}

Phương án thứ hai của cơ chế (i) là chỉ kết hợp các mô hình thuộc tập tin cậy mô hình. tập tin cậy mô hình (tập tin cậy mô hình, Model Confidence Set) là danh sách những ứng viên chưa thể loại sau vòng sơ loại thống kê, theo Hansen và cộng sự (2011). Mức ý nghĩa của tập tin cậy mô hình trong đề cương được ấn định trước ở 10\%.

Nếu hội đồng cố vấn của bạn gồm cả những người đã bị chứng minh là kém, việc kết hợp có thể kéo kết quả xuống. Chỉ giữ những người chưa bị loại là cách tránh điều đó. Đây là hình dung minh họa, không phải tuyên bố của bài gốc.

\subsection{Đối chứng của cơ chế (i): kết hợp 1/N với trung bình-phương sai}

Đối chứng là kết hợp quy tắc 1/N với quy tắc trung bình-phương sai, theo Tu và Zhou (2011). Đề cương nêu rằng nghiên cứu này kết hợp 1/N với trung bình-phương sai và vượt 1/N trong hầu hết kịch bản mô phỏng.

Vai trò chung của một đối chứng là cho bạn một thước để so. Nếu cách kết hợp theo trọng số Bayes không hơn nổi một cách kết hợp đơn giản, việc dùng cách phức tạp khó biện minh. Đề cương không nói điều này bằng những lời đó, đây là cách giải thích vai trò của đối chứng.

\section{Cơ chế (ii): danh mục vững chắc}

Cơ chế (ii) là danh mục vững chắc, tối ưu hóa cho tình huống xấu nhất trong nhiều niềm tin ban đầu, theo Garlappi và cộng sự (2007).

Hình dung như sau: bạn chuẩn bị cho chuyến đi biển khi có nhiều bản dự báo thời tiết khác nhau. Một người chọn tin bản lạc quan nhất. Người thứ hai kiểm tra bản xấu nhất rồi chọn kế hoạch ít tệ nhất trong đó. Cách của người thứ hai gần với tinh thần của danh mục vững chắc.

Trong ngôn ngữ của đề cương, mỗi bản dự báo là một niềm tin ban đầu. Ở phần cơ sở lý thuyết, đề cương mô tả Garlappi và cộng sự là mô hình hóa nhà đầu tư né tránh mơ hồ với nhiều niềm tin ban đầu, và được danh mục ổn định hơn. Đề cương nêu Gilboa và Schmeidler (1989) với tiêu chí maxmin (chọn phương án tốt nhất trong tình huống xấu nhất) và nhiều niềm tin ban đầu.

Đề cương không nêu đối chứng riêng cho cơ chế (ii). Ở các phần khác, đề cương nêu rằng không hướng khắc phục nào thắng ổn định theo DeMiguel và cộng sự (2009) trên bảy bộ dữ liệu.

Cách hiểu của bản đọc này: vì không hướng khắc phục nào thắng ổn định, cơ chế (ii) cũng cần được kiểm tra, không thể coi là chắc chắn tốt.

\section{Cơ chế (iii): quy tắc chuyển chế độ}

Cơ chế (iii) là quy tắc chuyển chế độ theo Kritzman và cộng sự (2012). Đối chứng của nó, cùng với M5, là giá trị của chuyển chế độ trong quyết định danh mục, theo Tu (2010).

Hình dung như sau: người đi biển nghe dự báo bão đang đến gần. Anh ta không đổi kế hoạch ngay lập tức mà thu hẹp chuyến đi theo xác suất bão. Quy tắc chuyển chế độ làm điều tương tự với danh mục: giảm rủi ro khi xác suất thị trường ở chế độ căng thẳng tăng. M7 trong Bảng 3 mô tả đúng tinh thần này.

Về nguồn gốc, đề cương nói Kritzman và cộng sự (2012) dự báo chế độ bằng mô hình Markov, và chiến lược động vượt chiến lược tĩnh trong kiểm tra ngược (kiểm định ngược). Còn Tu (2010) nêu tổn thất lợi suất tương đương chắc chắn do bỏ qua chuyển chế độ thường trên 2\% mỗi năm và có thể tới 10\%. Đây là bằng chứng từ tài liệu, không phải kết quả của đề tài.

Đối chứng M5 hoạt động theo một logic khác. M5 là Black-Litterman có chuyển chế độ, mở rộng từ bài báo của chính tác giả đề cương. Khi đặt M5 cạnh các mô hình khác, đề cương hỏi chuyển chế độ có đáng giá trong quyết định danh mục hay không, chứ chưa khẳng định điều đó.

\section{Vì sao cơ chế (ii) và (iii) bị loại khỏi tập mô hình khi kiểm định H3}

Ở mục 4.3.6, đề cương nói các cơ chế (ii) và (iii) được loại khỏi tập M khi kiểm định H3, và mức hối tiếc của chúng có thể âm. Giải thích dưới đây là cách hiểu của bản đọc này, suy ra từ định nghĩa mức hối tiếc.

Theo định nghĩa, mức hối tiếc là chênh lệch giữa lợi suất tương đương chắc chắn cao nhất trong M và lợi suất tương đương chắc chắn của mô hình chọn trước. Vì mô hình chọn trước thuộc M, mức hối tiếc không âm. Nhưng khi một cơ chế nằm ngoài M, lợi suất tương đương chắc chắn của nó có thể cao hơn lợi suất tương đương chắc chắn cao nhất trong M. Khi đó mức hối tiếc tính theo mốc của M sẽ âm. Bài gốc chỉ nói mức hối tiếc của chúng có thể âm, phần lý giải này là của bản đọc này.

\begin{tcolorbox}[breakable]
Ví dụ giả định (số tự đặt để minh họa, không phải số liệu của đề cương): trong một cửa sổ, lợi suất tương đương chắc chắn cao nhất trong tập M là 2,0 đơn vị.

Mô hình chọn trước có lợi suất tương đương chắc chắn là 1,5 đơn vị, nên mức hối tiếc của việc chọn trước là 0,5 đơn vị.

Một cơ chế nằm ngoài M có lợi suất tương đương chắc chắn là 2,4 đơn vị. Nếu so cơ chế đó với lợi suất tương đương chắc chắn cao nhất của M, chênh lệch là 2,0 trừ 2,4, tức âm 0,4 đơn vị.
\end{tcolorbox}

Phép tính trên cho thấy mức hối tiếc âm là dấu hiệu cơ chế đã vượt qua mốc tốt nhất của cả tập mô hình. Cách đọc này là suy ra từ định nghĩa mức hối tiếc ở mục 4.3.6, chứ bài gốc không đưa ví dụ số.

Bài gốc không nói rõ M6 và M7 có bị loại khỏi M hay không, dù Bảng 3 có M6 theo Garlappi và M7 theo Kritzman. Bài gốc chỉ nói các cơ chế (ii) và (iii) được loại khỏi M khi kiểm định H3.

\section{Sáu tiêu chí đánh giá}

Tiêu chí của mục 4.3.8 gồm sáu thứ: lợi suất tương đương chắc chắn, mức hối tiếc, độ sụt giảm tối đa, ES mức 95\%, vòng quay, và hiệu quả sau chi phí giao dịch. Mỗi tiêu chí nhìn một góc khác nhau của cùng một danh mục.

\subsection{lợi suất tương đương chắc chắn: số tiền chắc chắn tương đương}

lợi suất tương đương chắc chắn là số tiền chắc chắn mà bạn sẵn lòng nhận để đổi lấy một khoản đầu tư có rủi ro. lợi suất tương đương chắc chắn càng cao càng tốt. Trong đề cương, lợi suất tương đương chắc chắn của mô hình m là \(U_{m,t}\  = \ \mu_{m,t}\  - \ (\rho/2)\sigma_{m,t}^{2}\), với \(\mu_{m,t}\) là trung bình lợi suất sau chi phí, \(\sigma_{m,t}^{2}\) là phương sai và \(\rho\) là hệ số ngại rủi ro, cơ sở bằng 5.

Nói bằng lời: lợi suất tương đương chắc chắn lấy lợi suất trung bình rồi trừ đi một khoản phạt tỷ lệ với phương sai. Người ngại rủi ro nhiều thì khoản phạt lớn. Đó là lý do đề cương dùng lợi suất tương đương chắc chắn làm thước chính thay vì chỉ nhìn lợi suất.

\subsection{Mức hối tiếc: khoản thiệt vì chọn sai}

Mức hối tiếc là tiêu chí trung tâm của cả đề cương. Ở mục 4.3.8, mức hối tiếc cho biết cơ chế giảm thiểu có làm khoản thiệt do chọn mô hình nhỏ đi hay không. H3 kỳ vọng mức hối tiếc thấp hơn so với chọn một mô hình ở chế độ căng thẳng.

\subsection{Độ sụt giảm tối đa: cú rơi sâu nhất}

Độ sụt giảm tối đa (maximum drawdown) là mức giảm lớn nhất từ đỉnh. Đây là thước đo cảm giác đau của nhà đầu tư: khi đã lên tới đỉnh, danh mục có thể rơi sâu đến đâu trước khi hồi lại.

\begin{tcolorbox}[breakable]
Ví dụ giả định (số tự đặt để minh họa, không phải số liệu của đề cương): giá trị danh mục đạt đỉnh 100 đơn vị, sau đó giảm xuống thấp nhất là 80 đơn vị rồi hồi lên.

Mức giảm lớn nhất từ đỉnh trong trường hợp này là 20 đơn vị, tức 20\% so với đỉnh.
\end{tcolorbox}

\subsection{ES mức 95\%: trung bình của những ngày tệ nhất}

ES (giá trị thiệt hại kỳ vọng, expected shortfall) mức 95\% là tổn thất trung bình trong 5\% ngày xấu nhất. Độ sụt giảm tối đa chỉ nhìn một cú rơi duy nhất. ES nhìn cả nhóm những ngày tệ, nên bổ sung cho nó.

\begin{tcolorbox}[breakable]
Ví dụ giả định (số tự đặt để minh họa, không phải số liệu của đề cương): trong 100 ngày, năm ngày xấu nhất lỗ lần lượt 1, 2, 3, 4 và 5 phần trăm.

Tổn thất trung bình của năm ngày đó là 3 phần trăm, và đó là ES mức 95\% của chuỗi ngày giả định này.
\end{tcolorbox}

\subsection{Vòng quay: tần suất mua bán lại}

Vòng quay (turnover) là tỷ lệ mua bán lại danh mục. Một danh mục đổi tay nhiều làm chi phí giao dịch tăng. Nhờ có tiêu chí này, bạn thấy được cơ chế nào thắng chỉ vì mua bán rất nhiều và cơ chế nào thắng mà không tốn nhiều chi phí.

Đề cương gắn vòng quay với H2: thanh khoản cạn làm tăng chênh lệch giữa các mô hình có vòng quay khác nhau. Vòng quay vì vậy xuất hiện cả ở phần giải thích rủi ro lẫn phần giảm thiểu.

\subsection{Hiệu quả sau chi phí giao dịch}

Tiêu chí cuối là hiệu quả sau chi phí giao dịch. Chi phí giao dịch được tham số hóa và được kiểm tra độ nhạy. Điều đó nghĩa là đề cương cho chi phí chạy qua nhiều mức khác nhau để xem kết luận có đổi hay không.

Đề cương chưa quy định mức chi phí cụ thể ở mục 4.3.8. Ở mục dữ liệu, đề cương nói biên độ dao động giá, thuế và phí giao dịch theo từng thời kỳ được ghi kèm nguồn ở chương 3 của luận án.

\section{Tiêu chí nào quyết định H3}

Không phải sáu tiêu chí có vai trò bằng nhau trong kết luận về H3. Khi tổng kết kết quả dự kiến, đề cương nói H3 được ủng hộ khi lợi suất tương đương chắc chắn sau chi phí cao hơn, mức hối tiếc thấp hơn và độ sụt giảm tối đa không cao hơn so với chọn một mô hình ở chế độ căng thẳng, với khoảng tin cậy bootstrap khối loại trừ 0.

Bài gốc không nói ES và vòng quay là điều kiện bắt buộc để ủng hộ H3. Chúng nằm trong danh sách thước đo của H3 ở Bảng 2 và trong tiêu chí của mục 4.3.8, nhưng điều kiện ủng hộ nêu ba tiêu chí còn lại.

Khả năng đối lập của H3 là kết hợp mô hình không vượt việc chọn một mô hình sau chi phí giao dịch và sau hiệu chỉnh thiên lệch chọn lọc, theo Bailey và López de Prado (2014). Nói cách khác, nếu cơ chế chỉ thắng khi bỏ qua chi phí hoặc khi quên mất rằng nhiều mô hình đã được thử, H3 sẽ không đứng vững.

\section{Điều mục 4.3.8 không nói}

Mục 4.3.8 không tuyên bố cơ chế nào sẽ thắng. Nó chỉ nêu danh sách ba cơ chế, các đối chứng và sáu tiêu chí. Mọi kết luận về hiệu quả của cơ chế chờ dữ liệu.

Kết quả cũng có thể không ủng hộ H3. Theo mục 4.3.10, khi đó đề tài vẫn trả lời được câu hỏi: rủi ro không đổi theo chế độ, hoặc quy tắc 1/N đủ dùng ở Việt Nam.

\begin{tcolorbox}[breakable]
\textbf{Ghi nhớ khi đi thi}

• Có ba cơ chế: (i) kết hợp mô hình theo trọng số Bayes hoặc chỉ trong tập tin cậy mô hình, đối chứng là kết hợp 1/N với trung bình-phương sai; (ii) danh mục vững chắc với nhiều niềm tin tiên nghiệm; (iii) quy tắc chuyển chế độ, đối chứng cùng M5 là giá trị của chuyển chế độ.

• Sáu tiêu chí: lợi suất tương đương chắc chắn, mức hối tiếc, độ sụt giảm tối đa (mức giảm lớn nhất từ đỉnh), giá trị thiệt hại kỳ vọng mức 95\% (tổn thất trung bình trong 5\% ngày xấu nhất), vòng quay (tỷ lệ mua bán lại danh mục), hiệu quả sau chi phí.

• Chi phí giao dịch được tham số hóa và kiểm tra độ nhạy.

• Khi kiểm định H3, cơ chế (ii) và (iii) bị loại khỏi tập M, nên mức hối tiếc của chúng có thể âm.

• H3 được ủng hộ khi lợi suất tương đương chắc chắn sau chi phí cao hơn, mức hối tiếc thấp hơn, độ sụt giảm tối đa không cao hơn, với khoảng tin cậy bootstrap khối loại trừ 0.
\end{tcolorbox}

\section{Tự kiểm tra}

{\itshape 1. Hãy kể tên ba cơ chế giảm thiểu của mục 4.3.8 và đối chứng của cơ chế (i).\par}

Đáp án: Cơ chế (i) kết hợp mô hình theo trọng số Bayes hoặc chỉ trong tập tin cậy mô hình, đối chứng là kết hợp 1/N với trung bình-phương sai. Cơ chế (ii) danh mục vững chắc với nhiều niềm tin tiên nghiệm. Cơ chế (iii) quy tắc chuyển chế độ theo Kritzman và cộng sự (2012).

{\itshape 2. Độ sụt giảm tối đa và ES mức 95\% khác nhau ở điểm nào?\par}

Đáp án: Độ sụt giảm tối đa là mức giảm lớn nhất từ đỉnh. ES mức 95\% là tổn thất trung bình trong 5\% ngày xấu nhất.

{\itshape 3. Vì sao mức hối tiếc của cơ chế (ii) và (iii) có thể âm?\par}

Đáp án: Vì khi kiểm định H3, chúng được loại khỏi tập M, nên mức hối tiếc của chúng có thể âm. Cách hiểu của bản đọc này: lợi suất tương đương chắc chắn của chúng có thể cao hơn lợi suất tương đương chắc chắn cao nhất trong M.

{\itshape 4. H3 được ủng hộ khi nào?\par}

Đáp án: Khi lợi suất tương đương chắc chắn sau chi phí cao hơn, mức hối tiếc thấp hơn và độ sụt giảm tối đa không cao hơn so với chọn một mô hình ở chế độ căng thẳng, với khoảng tin cậy bootstrap khối loại trừ 0.

{\itshape 5. Điều gì xảy ra với chi phí giao dịch trong mục 4.3.8?\par}

Đáp án: Chi phí giao dịch được tham số hóa và kiểm tra độ nhạy.

\chapter{Mục 4.3.10 của đề cương: Độ tin cậy và giới hạn, hay khi nào đề tài được phép nói không biết}

Mục 4.3.10 nói thẳng một điều bất lợi: trong dữ liệu có ít giai đoạn căng thẳng độc lập, nên các kiểm định có sức mạnh thấp. Vì vậy tác giả đề cương đặt một quy tắc cứng. Khi số cửa sổ độc lập ở chế độ căng thẳng dưới 10, kết quả được gọi là không kết luận, không phải ủng hộ và cũng không phải bác bỏ. Mục này cũng nêu cách hiệu chỉnh so sánh bội, phạm vi kết luận, và cam kết rằng kết quả có thể bác bỏ giả thuyết.

\section{Vấn đề gốc: bão hiếm, nên khó biết chắc}

Chế độ căng thẳng (stress regime) là trạng thái thống kê do mô hình chuyển chế độ ước lượng, không phải tên một cuộc khủng hoảng. Đề cương cố ý tách hai khái niệm này. Khủng hoảng là sự kiện lịch sử, còn chế độ căng thẳng là chế độ có phương sai điều kiện cao nhất và lợi suất kỳ vọng thấp hơn chế độ bình thường.

Hình dung như sau: bạn muốn biết một loại áo mưa có thực sự tốt hơn loại khác trong mùa bão. Nếu bạn chỉ trải qua hai cơn bão trong đời, mọi so sánh đều mỏng manh. Dù bạn có hàng nghìn ngày nắng để đo, những ngày nắng không trả lời được câu hỏi về mùa bão.

Đề cương nói rằng Việt Nam là bối cảnh kiểm chứng vì thị trường có ít cuộc căng thẳng và thanh khoản hạn chế. Đề cương cũng nói đó không phải lý do của khoảng trống nghiên cứu. Cách hiểu của bản đọc này: chính đặc điểm đó nối với giới hạn mà mục 4.3.10 đang xử lý.

Thêm vào đó, mỗi mô hình được ước lượng trên cửa sổ 104 tuần và đánh giá trên cửa sổ 6 tháng không chồng lấn. Các giá trị này là dự kiến, và tác giả đề cương biện minh chúng sau khi ước lượng thử ở năm 1. Giải thích chung: cửa sổ đánh giá càng dài thì trong cùng một mẫu càng ít cửa sổ, nên con số 10 trong quy tắc càng đáng chú ý.

\section{Cửa sổ độc lập và giai đoạn căng thẳng độc lập}

Đề cương yêu cầu báo cáo hai con số cạnh mỗi hệ số: số cửa sổ độc lập và số giai đoạn căng thẳng độc lập. Cửa sổ độc lập là cửa sổ đánh giá không chồng lấn với cửa sổ khác. Cách hiểu này là của bản đọc này: bài gốc không định nghĩa chi tiết thuật ngữ, ngoài việc nêu cửa sổ đánh giá 6 tháng không chồng lấn.

Vì sao cần cả hai con số? Hình dung như sau: một cơn bão kéo dài nhiều ngày vẫn chỉ là một cơn bão. Nếu bạn đếm mỗi ngày mưa như một bằng chứng riêng, bạn sẽ tưởng mình có rất nhiều bằng chứng, trong khi thật ra chỉ có một. Số giai đoạn độc lập là cách đếm cơn bão, không đếm ngày mưa.

Đề cương không quy định thêm ngưỡng cho số giai đoạn căng thẳng độc lập ngoài việc báo cáo nó. Ngưỡng cứng duy nhất được nêu là 10 cửa sổ độc lập ở chế độ căng thẳng.

\section{Ba kết cục thay vì hai}

Trong thống kê thông thường, người ta hay chia kết quả thành hai nhóm: ủng hộ giả thuyết hoặc không. Đề cương chia thành ba nhóm. Nhóm thứ ba là không kết luận.

{\small\begin{longtable}[]{@{}
  >{\raggedright\arraybackslash}p{(\columnwidth - 2\tabcolsep) * \real{0.6830}}
  >{\raggedright\arraybackslash}p{(\columnwidth - 2\tabcolsep) * \real{0.3170}}@{}}
\toprule\noalign{}
\begin{minipage}[b]{\linewidth}\raggedright
Tình huống
\end{minipage} & \begin{minipage}[b]{\linewidth}\raggedright
Cách gọi trong đề cương
\end{minipage} \\
\midrule\noalign{}
\endhead
\bottomrule\noalign{}
\endlastfoot
Đủ cửa sổ độc lập và dữ liệu ủng hộ theo tiêu chí đã đăng ký & Ủng hộ giả thuyết \\
Đủ cửa sổ độc lập và dữ liệu không ủng hộ & Bác bỏ giả thuyết \\
Dưới 10 cửa sổ độc lập ở chế độ căng thẳng & Không kết luận \\
\end{longtable}}

Bảng trên là cách người viết tóm tắt quy tắc của mục 4.3.10. Hai dòng đầu là suy ra từ cách đề cương đặt điều kiện ủng hộ ở các mục 4.3.7 và 4, còn dòng thứ ba là quy tắc nguyên văn.

Khi kết quả là không kết luận, đề cương không bỏ trống. Tác giả đề cương báo cáo kích thước hiệu ứng (effect size, độ lớn của chênh lệch đo được) cùng khoảng tin cậy theo từng giai đoạn. Người đọc vẫn thấy dữ liệu, chỉ là không có một câu kết luận chắc chắn.

Giải thích chung: quy tắc này chặn trước việc cố diễn giải dữ liệu mỏng, vì nếu số cửa sổ độc lập quá ít, một chênh lệch đẹp mắt có thể chỉ là ngẫu nhiên.

Đề cương nêu quy tắc gọi kết quả là không kết luận trong tệp đăng ký cố định trước.

\begin{tcolorbox}[breakable]
Ví dụ giả định (số tự đặt để minh họa, không phải số liệu của đề cương): một nghiên cứu chỉ có 6 cửa sổ độc lập ở chế độ căng thẳng và thấy chênh lệch mức hối tiếc có dấu dương.

Theo quy tắc dưới 10 cửa sổ, nghiên cứu đó phải gọi kết quả là không kết luận, rồi báo cáo kích thước hiệu ứng và khoảng tin cậy theo từng giai đoạn.
\end{tcolorbox}

Con số 6 trong ví dụ là số tự đặt. Con số 10 là ngưỡng thật của đề cương.

\section{Hiệu chỉnh so sánh bội Holm: vì sao phải làm}

Đề cương hiệu chỉnh so sánh bội theo Holm (1979) ở mức 10\% cho bốn giả thuyết.

Giải thích chung: lý do nằm ở tính chất của việc kiểm định nhiều giả thuyết cùng lúc.

Hình dung như sau: nếu bạn mua một tấm vé số, cơ hội trúng rất nhỏ. Nếu bạn mua bốn tấm, cơ hội có ít nhất một tấm trúng cao hơn. Kiểm định bốn giả thuyết cũng giống vậy: càng nhiều giả thuyết thì càng dễ có một giả thuyết trông như được ủng hộ chỉ vì may. Hiệu chỉnh so sánh bội là cách siết ngưỡng để bù lại.

Phương pháp Holm làm việc theo thứ tự. Bạn xếp các giá trị p từ nhỏ đến lớn, rồi so giá trị nhỏ nhất với ngưỡng chặt nhất, giá trị nhỏ thứ hai với ngưỡng nới hơn một chút, và cứ thế. Dừng ngay ở giả thuyết đầu tiên không vượt ngưỡng. Phần mô tả quy trình này là kiến thức thống kê chung, không phải chi tiết có trong bài gốc.

\begin{tcolorbox}[breakable]
Ví dụ giả định (số tự đặt để minh họa, không phải số liệu của đề cương): bốn giá trị p được xếp thành 0,004; 0,02; 0,03 và 0,2, với mức 10\%.

Theo thủ tục Holm, giá trị nhỏ nhất so với 0,10 chia 4, tức 0,025; giá trị nhỏ thứ hai so với 0,10 chia 3, tức khoảng 0,033; giá trị thứ ba so với 0,10 chia 2, tức 0,05; giá trị lớn nhất so với 0,10.

Trong ví dụ này, 0,004 và 0,02 qua ngưỡng, 0,03 cũng qua ngưỡng 0,05, còn 0,2 không qua 0,10.
\end{tcolorbox}

Vì ví dụ chỉ để minh họa cách tính, bạn không nên rút ra điều gì về kết quả thật của đề tài. Đề cương chỉ nêu tên phương pháp, mức 10\% và số giả thuyết là bốn.

Hiệu chỉnh Holm áp dụng cho bốn giả thuyết H1 đến H4. Bài gốc nói Holm áp dụng cho bốn giả thuyết, nhưng không nói rõ H4, vốn chỉ được báo cáo là kết quả thăm dò, được xử lý thế nào trong thủ tục này; đây là điểm người đọc nên hỏi lại tác giả đề cương nếu cần.

\section{Giới hạn phạm vi: chỉ nói về Việt Nam}

Đề cương giới hạn kết luận trong phạm vi Việt Nam. Đề tài dùng dữ liệu HOSE. Cách hiểu của bản đọc này: kết quả vì vậy không tự động áp dụng cho thị trường khác. Đề cương cũng dùng chữ phạm vi thực nghiệm ngay từ phần lĩnh vực nghiên cứu.

\section{Nhận diện chế độ cũng là một dạng rủi ro mô hình}

Đề cương nêu rằng việc nhận diện chế độ là một dạng rủi ro mô hình.

Giải thích chung: nếu mô hình chuyển chế độ gọi sai chế độ, các phân tích dùng biến chế độ đó phía sau có thể lệch.

Cách xử lý của đề cương là đưa các định nghĩa chế độ thay thế vào kiểm tra độ bền (robustness check). Các định nghĩa thay thế đó dựa vào ngưỡng sụt giảm và biến động thực hiện. Giải thích chung: nếu kết luận đổi khi định nghĩa chế độ đổi, bạn biết kết luận phụ thuộc vào cách đặt tên cho cơn bão.

Kiểm định điểm gãy của Bai và Perron (1998) chỉ dùng toàn mẫu để mô tả và đối chiếu ở chương 3 của luận án. Nó không làm biến chế độ ở các chương 4 đến 6. Biến chế độ của các chương này lấy từ xác suất lọc, tính từ tham số chuyển chế độ được ước lượng lại theo cửa sổ mở rộng chỉ bằng dữ liệu đến ngày đó.

\section{Khả năng thực thi: thị trường mỏng thì khó mua bán thật}

Danh mục tính trên giấy khác với danh mục mua bán thật. Đề cương nói khả năng thực thi trên HOSE phụ thuộc thanh khoản, nên dùng rổ cổ phiếu thanh khoản cao và kiểm tra độ nhạy theo chi phí.

Rổ được dùng gồm 30 cổ phiếu có giá trị giao dịch bình quân ba tháng và vốn hóa tự do chuyển nhượng lớn nhất tại mỗi ngày tái cân bằng. Rổ được xác định chỉ bằng dữ liệu có tại ngày đó nhằm tránh thiên lệch sống sót và nhìn trước. Mục 4.3.8 và 3.6 cùng cho chi phí giao dịch chạy qua các mức khác nhau để xem kết luận có đổi hay không.

\section{Kết quả có thể bác bỏ giả thuyết}

Một câu đáng chú ý của mục 4.3.10 là: kết quả có thể bác bỏ giả thuyết. Khi đó đề tài vẫn trả lời câu hỏi nghiên cứu: rủi ro lựa chọn mô hình không đổi theo chế độ, hoặc quy tắc 1/N đủ dùng ở Việt Nam.

Điều này khớp với cách đề cương thiết kế giả thuyết. Mỗi giả thuyết đi kèm một khả năng đối lập. H1 có khả năng đối lập là rủi ro không đổi theo chế độ, hoặc 1/N thắng ổn định. H2 có khả năng đối lập là tương quan tăng chỉ vì biến động tăng. H3 có khả năng đối lập là kết hợp mô hình không vượt việc chọn một mô hình sau chi phí giao dịch. H4 có khả năng đối lập là bầy đàn không xuất hiện ở tần suất dữ liệu của đề tài.

Hình dung như sau: một nhà khoa học bày sẵn trên bàn cách để mình sai. Đề cương làm điều đó bằng việc ghi trước cả tiêu chí ủng hộ lẫn quy tắc gọi kết quả là không kết luận vào tệp đăng ký trước.

\section{Những giới hạn khác rải rác trong đề cương}

Bên cạnh mục 4.3.10, đề cương còn nêu vài giới hạn ở các mục khác mà bạn nên đọc chung. Chúng không thay thế mục 4.3.10, nhưng cùng với mục 4.3.10 tạo thành bức tranh đầy đủ hơn về mức chắc chắn của đề tài.

\begin{itemize}
\item
  Quan hệ được báo cáo chỉ là liên hệ có điều kiện, không phải quan hệ nhân quả.
\item
  H4 chỉ báo cáo là kết quả thăm dò.
\item
  Tác giả đề cương đã dùng dữ liệu Việt Nam 2019 đến 2025, gồm hai giai đoạn căng thẳng, trong bài báo trước. Vì vậy không thể khẳng định chưa xem kết quả ở chế độ căng thẳng, và việc đăng ký trước chỉ giảm chứ không loại thiên lệch.
\item
  Tuyên bố tính mới chỉ nói theo hiểu biết của tác giả đề cương, trên một phạm vi tìm kiếm có giới hạn, không gồm Scopus, Web of Science và nghiên cứu tiếng Việt trên tạp chí trong nước.
\item
  Quy trình giám sát mô hình ở mục 4.4 là đề xuất chưa kiểm nghiệm với người dùng.
\item
  Đề tài không quan sát hành vi chọn mô hình của tổ chức đầu tư, và không đánh giá mức đầy đủ của khung quy định.
\item
  Các giá trị cửa sổ 104 tuần, 6 tháng, 36 tháng, hệ số ngại rủi ro 5 và số mô hình 7 là giá trị dự kiến hoặc cơ sở, được biện minh sau ước lượng thử.
\end{itemize}

\section{Mối liên hệ với phần còn lại của đề cương}

Mục 4.3.10 là chỗ tập hợp các giới hạn và cam kết của đề cương. Quy tắc dưới 10 cửa sổ gắn với mục 4.3.7 và 4, vì tiêu chí ủng hộ giả thuyết được đặt cùng nơi. Quy tắc Holm gắn với kế hoạch đăng ký trước ở mục 4.3.5, vì tệp đăng ký cố định quy tắc hiệu chỉnh so sánh bội và quy tắc gọi kết quả là không kết luận.

Nhờ vậy, khi đọc mục 4.4, bạn cần luôn nhớ rằng mọi chữ ủng hộ nằm sau hai cánh cửa: tiêu chí thống kê của từng giả thuyết, và quy tắc số cửa sổ độc lập. Chỉ khi qua cả hai, câu kết luận mới được viết là ủng hộ.

\begin{tcolorbox}[breakable]
\textbf{Ghi nhớ khi đi thi}

• Số giai đoạn căng thẳng độc lập ít nên kiểm định có sức mạnh thấp; đề cương báo cáo số cửa sổ độc lập và số giai đoạn căng thẳng độc lập cạnh mỗi hệ số.

• Dưới 10 cửa sổ độc lập ở chế độ căng thẳng thì gọi là không kết luận, không ủng hộ, không bác bỏ; khi đó báo cáo kích thước hiệu ứng và khoảng tin cậy theo từng giai đoạn.

• Hiệu chỉnh so sánh bội Holm (1979) ở mức 10\% cho bốn giả thuyết; kết luận giới hạn trong phạm vi Việt Nam.

• Nhận diện chế độ là một dạng rủi ro mô hình, nên định nghĩa thay thế nằm trong kiểm tra độ bền.

• Kết quả có thể bác bỏ giả thuyết; khi đó đề tài vẫn trả lời: rủi ro không đổi theo chế độ, hoặc 1/N đủ dùng ở Việt Nam.
\end{tcolorbox}

\section{Tự kiểm tra}

{\itshape 1. Quy tắc nào khiến một kết quả được gọi là không kết luận?\par}

Đáp án: Khi số cửa sổ độc lập ở chế độ căng thẳng dưới 10. Khi đó kết quả không được gọi là ủng hộ hay bác bỏ.

{\itshape 2. Khi kết quả là không kết luận, đề tài báo cáo gì?\par}

Đáp án: Kích thước hiệu ứng cùng khoảng tin cậy theo từng giai đoạn.

{\itshape 3. Hiệu chỉnh Holm áp dụng ở mức nào và cho mấy giả thuyết?\par}

Đáp án: Mức 10\% cho bốn giả thuyết.

{\itshape 4. Vì sao đề cương xem nhận diện chế độ là một dạng rủi ro mô hình, và xử lý thế nào?\par}

Đáp án: Vì xác định chế độ cũng là dùng một mô hình. Các định nghĩa chế độ thay thế, theo ngưỡng sụt giảm và biến động thực hiện, nằm trong kiểm tra độ bền.

{\itshape 5. Nếu H1 bị bác bỏ, đề tài có thất bại không?\par}

Đáp án: Không. Đề tài vẫn trả lời câu hỏi nghiên cứu: rủi ro lựa chọn mô hình không đổi theo chế độ, hoặc quy tắc 1/N đủ dùng ở Việt Nam.

\chapter{Mục 4.3.10 của đề cương: Khả năng thực hiện và minh bạch}

Mục 4.3.10 trả lời câu hỏi người phản biện nào cũng hỏi: đề tài này làm được không, và người khác kiểm tra được không. Tác giả đề cương nêu ba nhóm bảo đảm. Về dữ liệu, nguồn và độ dài được xác nhận trong quý đầu của năm 1, và nếu thiếu thì thu hẹp. Về công cụ, các mô hình chạy được bằng phần mềm thống kê thông dụng, và quyết định về phương pháp phương pháp tính toán Bayes xấp xỉ cho M5 được hoãn tới ước lượng thử ở quý 2 năm 1. Về minh bạch, mã nguồn và tệp đăng ký trước được lưu lại.

\section{Dữ liệu: xác nhận sớm, thu hẹp nếu cần}

Đề tài dùng dữ liệu công khai của thị trường chứng khoán Việt Nam. Dữ liệu gồm giá đóng cửa điều chỉnh theo ngày của cổ phiếu niêm yết trên HOSE và của VN-Index, lợi suất trái phiếu chính phủ kỳ hạn một năm làm lãi suất phi rủi ro, và khối lượng giao dịch. Nguồn là HOSE và nhà cung cấp dữ liệu thương mại.

Mục 4.3.10 nói tác giả đề cương xác nhận nguồn và độ dài dữ liệu trong quý đầu của năm 1. Việc này cũng khớp với Bảng 5, nơi quý 1 năm 1 mở đầu bằng việc xác nhận nguồn, độ dài và giấy phép dữ liệu, cùng ngày HOSE công bố VN30 và VN100.

Hình dung như sau: trước khi xây nhà, bạn kiểm tra xem kho có đủ gạch không. Nếu thiếu, bạn không bỏ dở mà vẽ lại căn nhà nhỏ hơn. Mục 4.3.10 nói đúng điều đó: nếu dữ liệu không đủ, đề tài thu hẹp rổ cổ phiếu hoặc giai đoạn mẫu.

Cần phân biệt hai chuyện. Giai đoạn mẫu dự kiến từ 2007 đến 2026, tùy dữ liệu sẵn có, và gồm khủng hoảng tài chính toàn cầu, đại dịch COVID-19 năm 2020, cùng căng thẳng trái phiếu doanh nghiệp từ quý IV/2022. Đó là dự kiến, còn độ dài thực tế phụ thuộc kết quả xác nhận ở quý 1 năm 1.

\subsection{Những điều dữ liệu cần ghi kèm}

Đề cương cho biết biên độ dao động giá, mức thuế và phí giao dịch theo từng thời kỳ và điều khoản giấy phép của nhà cung cấp dữ liệu được ghi kèm nguồn trong chương 3 của luận án. Dữ liệu là dữ liệu công khai và không chứa thông tin cá nhân.

Bài gốc không nói tên nhà cung cấp dữ liệu thương mại, không nói chi phí mua dữ liệu và không nói điều khoản giấy phép cụ thể. Theo Bảng 5, giấy phép dữ liệu được xác nhận ở quý 1 năm 1.

\subsection{Vì sao rổ chỉ có 30 cổ phiếu}

Rổ cổ phiếu tại mỗi ngày tái cân bằng là 30 cổ phiếu có giá trị giao dịch bình quân ba tháng và vốn hóa tự do chuyển nhượng lớn nhất. Rổ được xác định chỉ bằng dữ liệu có tại ngày đó. Cách chọn này nhằm tránh hai thiên lệch: thiên lệch sống sót và thiên lệch nhìn trước.

Thiên lệch sống sót xảy ra khi bạn chỉ xét những cổ phiếu còn tồn tại đến hôm nay và quên những cổ phiếu đã biến mất. Thiên lệch nhìn trước xảy ra khi bạn dùng thông tin mà thật ra chưa có tại thời điểm quyết định. Hai khái niệm này được giải thích theo nghĩa thông dụng, không có định nghĩa riêng trong bài gốc.

\begin{tcolorbox}[breakable]
Ví dụ giả định (số tự đặt để minh họa, không phải số liệu của đề cương): giả sử bạn chấm điểm các đội bóng của giải đấu chỉ dựa trên những đội còn thi đấu ở mùa cuối cùng.

Các đội đã bị loại trong những mùa trước không có mặt trong bảng điểm, nên kết quả trông đẹp hơn thực tế. Đó là cách thiên lệch sống sót làm sai lệch bức tranh.
\end{tcolorbox}

VN30 và VN100 chỉ dùng để đối chiếu và làm kiểm tra độ bền, từ thời điểm HOSE công bố chúng. Thời điểm này tác giả đề cương xác nhận trong quý 1 năm 1.

\section{Công cụ: dùng phần mềm thông dụng, một điểm cần thử là phương pháp tính toán Bayes xấp xỉ}

Mục 4.3.10 nói các mô hình ở Bảng 3 ước lượng được bằng phần mềm thống kê thông dụng. Bài gốc không nêu tên phần mềm nào, nên đề cương chưa quy định điểm này.

Điểm cần thử nghiệm là M5. M5 là Black-Litterman có chuyển chế độ, do tác giả đề cương phát triển, mở rộng từ bài báo Nguyen và Nguyen (2026a). Trong bài báo đó, tác giả đã dùng mô hình định giá tài sản vốn chuyển chế độ và phương pháp phương pháp tính toán Bayes xấp xỉ.

\subsection{phương pháp tính toán Bayes xấp xỉ là gì, giải thích từ đầu}

phương pháp tính toán Bayes xấp xỉ là tên viết tắt của hai ý ghép lại. Đề cương gọi nó là phương pháp tính toán Bayes xấp xỉ kết hợp Monte Carlo chuỗi Markov. Phần giải thích dưới đây là giải thích chung của người viết, vì bài gốc không mô tả cách chạy của phương pháp.

Ý thứ nhất là tính toán Bayes xấp xỉ (ABC, Approximate Bayesian Computation). Hình dung như sau: bạn muốn biết một công thức nấu ăn bí mật, nhưng chỉ có món ăn thật để nếm. Bạn thử nhiều công thức khác nhau, nấu từng món, rồi giữ lại những công thức cho ra món nếm giống món thật. Bạn không cần tính chính xác khả năng của từng công thức, chỉ cần so món nấu thử với món thật.

Ý thứ hai là Monte Carlo chuỗi Markov (MCMC, Markov chain Monte Carlo). Hình dung như sau: bạn thả một người đi bộ ngẫu nhiên trong một vùng đồi núi, với quy tắc là người đó ưa ở những chỗ cao hơn. Sau rất nhiều bước, nơi người đó ghé thăm nhiều nhất cho bạn biết nơi nào đáng tin nhất. Mỗi bước chỉ phụ thuộc vào bước ngay trước, nên đó là chuỗi Markov.

Ghép hai ý lại, phương pháp tính toán Bayes xấp xỉ dùng cách đi bộ ngẫu nhiên để lần theo những bộ tham số cho dữ liệu mô phỏng gần với dữ liệu thật. Đó là cách nói ngắn gọn để người đọc nắm tinh thần, không thay thế mô tả kĩ thuật.

\subsection{Vì sao đề cương hoãn quyết định đến quý 2 năm 1}

Giải thích chung: những phương pháp mô phỏng như trên thường tốn nhiều thời gian máy.

Đề cương nói ước lượng thử ở quý 2 năm 1 quyết định giữ hay thay phương pháp tính toán Bayes xấp xỉ bằng ước lượng nhanh hơn cho M5. Bảng 5 nói ước lượng thử này chạy bảy mô hình trên giai đoạn 2019 đến 2025, để đo thời gian tính và số quan sát ở chế độ căng thẳng.

Điểm đáng chú ý là đề cương chưa chốt trước một phương pháp cho M5: ước lượng thử quyết định giữ hay thay. Bảng 5 nói ước lượng thử đo thời gian tính. Bài gốc không nói ước lượng nhanh hơn đó là phương pháp nào.

Lựa chọn này cũng khớp với điều đề cương nói ở mục 4.3.6: siêu tham số của M2, M4, M5, M6 và M7 được chọn từ một lưới cho trước bằng cùng một quy tắc và cùng ngân sách tinh chỉnh, cố định khi đăng ký trước. Cách hiểu của bản đọc này: vì siêu tham số cố định khi đăng ký trước ở quý 4 năm 1, các lựa chọn của M5 cần xong trước thời điểm đó.

\subsection{Bai-Perron trong một bản thảo}

Mục 4.3.10 nói tác giả đề cương dùng phương pháp Bai-Perron trong một bản thảo nêu ở mục 8. Bai và Perron (1998) là nguồn của phương pháp tìm nhiều điểm gãy cấu trúc.

Giải thích chung: điểm gãy cấu trúc là thời điểm mà quy luật của dữ liệu đổi hẳn.

Trong đề tài này, kiểm định Bai-Perron chỉ dùng toàn mẫu để mô tả và đối chiếu ở chương 3 của luận án. Nó không làm biến chế độ ở chương 4 đến 6.

\section{Minh bạch: để người khác kiểm lại được}

Mục 4.3.10 cam kết hai việc và một mô tả. Tác giả đề cương lưu mã nguồn và tệp đăng ký trước, và mô tả đầy đủ cách xử lý dữ liệu trong luận án.

Tệp đăng ký trước liên hệ chặt với cam kết này. Theo mục 4.3.5, kế hoạch phân tích được đăng ký trên OSF Registries vào quý 4 năm 1, ở chế độ khóa có dấu thời gian. Hình dung như sau: bạn viết dự đoán và cách kiểm tra vào một phong bì có dấu thời gian trước khi xem dữ liệu mới. Dấu thời gian cho thấy bạn đã viết gì và viết khi nào.

Tệp đăng ký cố định nhiều thứ: độ dài các cửa sổ, siêu tham số của bảy mô hình, số chế độ và ngưỡng xác suất lọc, mức ý nghĩa, nhóm kết quả chính, quy tắc hiệu chỉnh so sánh bội và quy tắc gọi kết quả là không kết luận. Tiêu chí ủng hộ từng giả thuyết cũng ghi trong tệp.

Dữ liệu phát sinh sau ngày đăng ký được giữ làm mẫu kiểm tra riêng. Mọi thay đổi sau đó ghi là lệch khỏi kế hoạch. Nghĩa là nếu tác giả đề cương đổi ý, việc đổi ý không được giấu đi mà phải ghi nhận.

\subsection{Một giới hạn của đăng ký trước}

Đề cương nói rõ tác giả đã dùng dữ liệu Việt Nam 2019 đến 2025, gồm hai giai đoạn căng thẳng, trong Nguyen và Nguyen (2026a). Vì vậy tác giả không thể khẳng định chưa xem kết quả ở chế độ căng thẳng. Đăng ký trước chỉ giảm chứ không loại thiên lệch. Để minh bạch, tệp đăng ký khai rõ những kết quả đã xem.

Giải thích chung, để nhớ khi đi thi: minh bạch ở đây không có nghĩa là hoàn hảo, mà là nêu rõ điều đã biết và điều chưa chắc.

\section{Rủi ro thực hiện và cách xử lý}

Mục 7.4 tóm tắt hai nhóm rủi ro. Rủi ro về dữ liệu và về số cuộc khủng hoảng trong mẫu ít được xử lý theo mục 4.3.10 và 3.7. Rủi ro tạp chí từ chối bản thảo được xử lý bằng cách chuyển sang tạp chí khác trong cùng năm.

Bảng dưới đây ghép từng rủi ro với cách ứng phó được nêu ở các mục 4.3.5, 3.6, 3.7 và 6.

{\small\begin{longtable}[]{@{}
  >{\raggedright\arraybackslash}p{(\columnwidth - 2\tabcolsep) * \real{0.3385}}
  >{\raggedright\arraybackslash}p{(\columnwidth - 2\tabcolsep) * \real{0.6615}}@{}}
\toprule\noalign{}
\begin{minipage}[b]{\linewidth}\raggedright
Rủi ro
\end{minipage} & \begin{minipage}[b]{\linewidth}\raggedright
Cách ứng phó theo đề cương
\end{minipage} \\
\midrule\noalign{}
\endhead
\bottomrule\noalign{}
\endlastfoot
Dữ liệu không đủ nguồn hoặc độ dài & Xác nhận trong quý đầu năm 1; thu hẹp rổ hoặc giai đoạn mẫu \\
phương pháp tính toán Bayes xấp xỉ quá chậm cho M5 & Ước lượng thử quý 2 năm 1; quyết định giữ hay thay bằng ước lượng nhanh hơn \\
Có thể bị nghi ngờ về việc điều chỉnh sau khi xem kết quả & Lưu tệp đăng ký trước và mã nguồn; ghi mọi thay đổi là lệch khỏi kế hoạch \\
Ít cuộc căng thẳng trong mẫu & Quy tắc dưới 10 cửa sổ độc lập là không kết luận (mục 4.3.10) \\
\end{longtable}}

Bảng trên gom lại các điều ở mục 4.3.10, 3.7 và 6. Cách chia thành bốn dòng là cách trình bày của người viết bài đọc này, và dòng thứ ba là cách ghép của bản đọc này: bài gốc không gọi đó là một rủi ro thực hiện.

\section{Tại sao mục này quan trọng với hội đồng}

Giải thích chung: một đề cương có thể có ý tưởng tốt nhưng không làm được vì thiếu dữ liệu hoặc công cụ.

Mục 4.3.10 trả lời bằng ba việc cụ thể và có mốc thời gian: quý 1 năm 1 cho dữ liệu, quý 2 năm 1 cho công cụ, quý 4 năm 1 cho đăng ký trước. Bảng 5 xác nhận các mốc đó.

Mục 4.3.10 cũng thừa nhận phần chưa chắc: nguồn dữ liệu thương mại, độ dài mẫu, phương pháp tính toán Bayes xấp xỉ. Với những điểm này, đề cương đặt điều kiện để đổi thay vì cam kết một con đường duy nhất.

\begin{tcolorbox}[breakable]
\textbf{Ghi nhớ khi đi thi}

• Dữ liệu: xác nhận nguồn và độ dài trong quý đầu năm 1; thu hẹp rổ cổ phiếu hoặc giai đoạn mẫu nếu dữ liệu không đủ.

• Công cụ: các mô hình ở Bảng 3 ước lượng được bằng phần mềm thống kê thông dụng; đã dùng mô hình định giá tài sản vốn chuyển chế độ và phương pháp tính toán Bayes xấp xỉ trong Nguyen và Nguyen (2026a), dùng Bai-Perron trong một bản thảo ở mục 8.

• Ước lượng thử quý 2 năm 1 quyết định giữ hay thay phương pháp tính toán Bayes xấp xỉ bằng ước lượng nhanh hơn cho M5.

• Minh bạch: lưu mã nguồn và tệp đăng ký trước; mô tả đầy đủ cách xử lý dữ liệu trong luận án.

• Đăng ký trước trên OSF Registries ở quý 4 năm 1 chỉ giảm chứ không loại thiên lệch, vì tác giả đã xem dữ liệu 2019 đến 2025.
\end{tcolorbox}

\section{Tự kiểm tra}

{\itshape 1. Khi nào tác giả đề cương xác nhận nguồn và độ dài dữ liệu, và nếu dữ liệu không đủ thì làm gì?\par}

Đáp án: Trong quý đầu của năm 1; nếu không đủ thì thu hẹp rổ cổ phiếu hoặc giai đoạn mẫu.

{\itshape 2. phương pháp tính toán Bayes xấp xỉ viết tắt của cụm gì trong đề cương, và vì sao nó gắn với M5?\par}

Đáp án: Phương pháp tính toán Bayes xấp xỉ kết hợp Monte Carlo chuỗi Markov. Tác giả đã dùng nó cùng mô hình định giá tài sản vốn chuyển chế độ trong Nguyen và Nguyen (2026a), và M5 là mở rộng từ bài báo đó.

{\itshape 3. Điều gì quyết định việc giữ hay thay phương pháp tính toán Bayes xấp xỉ?\par}

Đáp án: Ước lượng thử ở quý 2 năm 1, đo thời gian tính và số quan sát ở chế độ căng thẳng.

{\itshape 4. Mục 4.3.10 cam kết lưu những gì để bảo đảm minh bạch?\par}

Đáp án: Mã nguồn và tệp đăng ký trước; và mô tả đầy đủ cách xử lý dữ liệu trong luận án.

{\itshape 5. Vì sao đề cương nói đăng ký trước chỉ giảm chứ không loại thiên lệch?\par}

Đáp án: Vì tác giả đề cương đã dùng dữ liệu Việt Nam 2019 đến 2025, gồm hai giai đoạn căng thẳng, trong bài báo trước, nên không thể khẳng định chưa xem kết quả ở chế độ căng thẳng.

\chapter{Mục 4.4 của đề cương: Dự kiến kết quả đạt được}

Mục 4.4 nói đề tài hi vọng thu được gì, với điều kiện dữ liệu ủng hộ các giả thuyết. Có ba nhóm kết quả dự kiến: lý thuyết, thực nghiệm, và thực tiễn cùng chính sách. Mục này cũng nêu tiêu chí để gọi một giả thuyết là được ủng hộ, và danh sách sản phẩm công bố. Điều bạn cần nhớ trước hết là chữ nếu: đề cương chưa có kết quả nào, mọi thứ ở đây là mục tiêu và kịch bản.

\section{Đọc đúng chữ nếu}

Đề cương mở đầu mục 4.4 bằng câu: nếu dữ liệu ủng hộ các giả thuyết, đề tài dự kiến có ba nhóm kết quả. Mọi câu sau đó chỉ đúng trong kịch bản đó.

Hình dung như sau: một kiến trúc sư đưa bản vẽ và nói nếu nền đất đủ chắc, căn nhà sẽ có ba tầng. Bản vẽ không phải là căn nhà đã xây. Mục 4.4 là bản vẽ của đề tài, và nền đất là dữ liệu thị trường chứng khoán Việt Nam.

Điều này khớp với mục 4.3.10, nơi đề cương nói kết quả có thể bác bỏ giả thuyết và khi đó đề tài vẫn trả lời câu hỏi nghiên cứu. Vì vậy ba nhóm kết quả dưới đây chỉ là một nhánh của các khả năng.

\section{Nhóm kết quả 1: lý thuyết}

Nhóm thứ nhất là một khung khái niệm và thước đo rủi ro lựa chọn mô hình theo chế độ, cho quyết định phân bổ danh mục.

Nói đơn giản, đề tài muốn để lại một cách nghĩ và một cây thước. Cách nghĩ là: khi chọn một cách chia tiền trong nhiều cách, có một khoản rủi ro riêng, và khoản này có thể đổi theo thời tiết thị trường. Cây thước là bộ ba độ phân tán mô hình \(D_{t}\), mức hối tiếc \(R_{t}\) và kích thước tập tin cậy mô hình, với lợi suất tương đương chắc chắn làm thang đo hiệu quả.

Đề cương nhấn mạnh độ phân tán mô hình và kích thước tập tin cậy mô hình là chỉ báo của bất định, không phải của rủi ro. Chỉ mức hối tiếc đo rủi ro lựa chọn theo đúng định nghĩa. Vì vậy ba thuật ngữ bất định mô hình, rủi ro mô hình và rủi ro lựa chọn mô hình được dùng đúng nghĩa của chúng.

Đề cương cho biết khoảng trống nằm ở hai chỗ: thước đo rủi ro mô hình hiện có đo chênh lệch của dự báo rủi ro, chưa đo hậu quả của việc chọn sai mô hình lên trọng số và kết quả của danh mục; còn nghiên cứu danh mục so sánh mô hình bằng thứ hạng trung bình trên toàn mẫu. Cách hiểu của bản đọc này: khung khái niệm và thước đo ứng với khoảng trống thứ nhất.

\section{Nhóm kết quả 2: thực nghiệm}

Nhóm thứ hai là bằng chứng cho thị trường Việt Nam về mức độ, thời điểm và kênh của rủi ro lựa chọn mô hình.

Ba chữ này tương ứng ba câu hỏi nhỏ. Mức độ là rủi ro lớn đến đâu, tương ứng với kích thước của hệ số \(\beta\) ở phương trình (5). Thời điểm là rủi ro lớn lên vào chế độ nào, tương ứng với so sánh chế độ căng thẳng và bình thường. Kênh là liên hệ giữa rủi ro với biến động, thanh khoản, tương quan cực trị và bầy đàn.

Cách đối ứng ba chữ này với phần phương pháp là cách giải thích của người viết bài đọc này. Đề cương nêu ba chữ ấy mà không phân tích từng chữ.

Đề cương gọi kết quả là liên hệ có điều kiện, không gọi là quan hệ nhân quả. Nghĩa là khi bạn nói thanh khoản cạn đi cùng rủi ro lựa chọn cao hơn, bạn chưa được nói thanh khoản cạn gây ra rủi ro đó.

\section{Nhóm kết quả 3: thực tiễn và chính sách}

Nhóm thứ ba là một quy trình giám sát mô hình gồm ba bước. Đề cương gọi nó là đề xuất suy ra từ H3, chưa được kiểm nghiệm với người dùng, và làm tư liệu tham khảo cho tổ chức đầu tư và cơ quan giám sát.

Ba bước như sau, nêu đúng thứ tự trong đề cương.

{\itshape 1. Khai báo tập mô hình.

2. Đo độ phân tán và tập tin cậy mô hình định kỳ.

3. Kích hoạt kết hợp mô hình khi thị trường ở chế độ căng thẳng.\par}

Hình dung như sau: một đội y tế duy trì một danh sách các phương án điều trị chính thức, rồi khám lại định kỳ xem các phương án còn khác nhau nhiều không, và chỉ khi gặp ca nguy kịch mới lập hội chẩn nhiều bác sĩ. Ba bước của quy trình có hình dạng tương tự. Đây là phép ví von, không phải phát biểu của bài gốc.

Cách hiểu của bản đọc này: bước 1 khai báo tập mô hình có nghĩa là nêu rõ những cách chia tiền nào đang được xét. Bước 2 đo độ phân tán mô hình và tập tin cậy mô hình để biết các cách chia tiền đang gần nhau hay xa nhau. Bước 3 kích hoạt kết hợp mô hình, gần với cơ chế (i), khi thị trường ở chế độ căng thẳng.

Bài gốc nói cả quy trình là đề xuất suy ra từ H3, tức từ việc kết hợp mô hình giảm mức hối tiếc, nếu H3 được ủng hộ.

Có ba điều đề cương không nói. Đề tài không quan sát hành vi chọn mô hình của tổ chức đầu tư. Đề tài không đánh giá mức đầy đủ của khung quy định. Và quy trình chưa được thử với người dùng thực.

Ngoài ra, các văn bản giám sát nước ngoài mà đề cương nhắc tới, SR 11-7 và SR 26-2 của Cục Dự trữ Liên bang Mỹ, không ràng buộc tổ chức Việt Nam. Đề tài dùng chúng làm tham chiếu quốc tế. Thông tư 83/2025/TT-NHNN áp dụng cho ngân hàng, trong khi đối tượng dùng kết quả của đề tài là công ty quản lý quỹ và tổ chức đầu tư cổ phiếu.

\section{Khi nào gọi một giả thuyết là được ủng hộ}

Đề cương nêu tiêu chí ủng hộ cho từng giả thuyết, và ghi các tiêu chí vào tệp đăng ký trước. Nhờ đó, tiêu chí được cố định trước khi tác giả xem dữ liệu mới.

\subsection{Tiêu chí của H1}

Kết quả ủng hộ H1 theo tiêu chí ở mục 4.3.7. Nói cụ thể, H1 được ủng hộ chỉ khi hệ số \(\beta\) của mức hối tiếc \(R_{t}\) dương, vượt phân vị 95\% của phân phối rỗng và có khoảng tin cậy bootstrap khối loại trừ 0, và hệ số \(\beta\) của độ phân tán mô hình \(D_{t}\) cùng dấu.

Phân phối rỗng là kết quả sẽ trông thế nào nếu thực ra không có gì xảy ra. Đề cương tạo nó bằng 1.000 lần tái lấy mẫu bootstrap khối chuỗi lợi suất của bảy mô hình, sau khi loại chênh lệch trung bình giữa các mô hình trong từng chế độ.

Đề cương cũng lặp lại kiểm tra trên tập chỉ gồm các mô hình tĩnh M1, M2, M3, M4 và M6, để loại tác động của M5 và M7, vốn chỉ đổi hành vi khi xác suất căng thẳng tăng.

\subsection{Tiêu chí của H2}

H2 được ủng hộ khi hệ số của độ phi thanh khoản dương, hệ số của tương quan đuôi âm, và khoảng tin cậy bootstrap khối của cả hai loại trừ 0.

Cần nhớ hai dấu ngược nhau. độ phi thanh khoản dương nghĩa là thanh khoản cạn làm tăng chênh lệch giữa các mô hình có vòng quay khác nhau. Tương quan đuôi âm nghĩa là tương quan cao làm các danh mục chỉ mua hội tụ, tức khiến các cách chia tiền gần giống nhau hơn và rủi ro lựa chọn giảm. Cách diễn đạt hai dấu lấy từ Bảng 2 của đề cương; vế rủi ro lựa chọn giảm là suy ra của bản đọc này từ dấu âm.

\subsection{Tiêu chí của H3}

H3 được ủng hộ khi lợi suất tương đương chắc chắn sau chi phí cao hơn, mức hối tiếc thấp hơn và độ sụt giảm tối đa không cao hơn so với chọn một mô hình ở chế độ căng thẳng, với khoảng tin cậy bootstrap khối loại trừ 0.

\subsection{H4: chỉ thăm dò}

H4 chỉ được báo cáo là kết quả thăm dò. Đề cương không đặt một bộ tiêu chí ủng hộ cho H4 như ba giả thuyết kia. Ở phần cơ sở lý thuyết, đề cương cũng nêu rằng bằng chứng về bầy đàn trái chiều giữa các thị trường: Christie và Huang (1995) không thấy bầy đàn khi giá biến động mạnh trên dữ liệu Mỹ, còn Chang và cộng sự (2000) thấy bằng chứng có ý nghĩa ở Hàn Quốc và Đài Loan, không thấy ở Mỹ và Hồng Kông.

Hình dung như sau: bầy đàn là đàn cừu chạy theo con đầu đàn thay vì tự nhìn đường. Đề cương chỉ dùng độ phân tán chéo của lợi suất như một thước đo bầy đàn ở cấp thị trường, và không quan sát hành vi chọn mô hình của tổ chức đầu tư.

\section{Ba lớp lọc trước khi viết chữ ủng hộ}

Cách đọc của bản này: một giả thuyết chỉ được viết là ủng hộ nếu vượt qua nhiều lớp. Gom các điều ở mục 4.3.7, 3.6 và 4 lại, bạn có thể thấy ba lớp. Lớp đầu là tiêu chí thống kê riêng của từng giả thuyết. Lớp hai là số cửa sổ độc lập ở chế độ căng thẳng phải từ 10 trở lên. Lớp ba là hiệu chỉnh so sánh bội Holm ở mức 10\% cho bốn giả thuyết.

Cách đếm ba lớp là cách tổng hợp của người viết. Đề cương không dùng cụm ba lớp lọc, và bài gốc không nói rõ hiệu chỉnh Holm được ghép với tiêu chí ủng hộ từng giả thuyết theo thứ tự nào.

Theo cách đọc đó, nếu H1 qua cả ba lớp, kết quả là ủng hộ. Nếu số cửa sổ độc lập dưới 10, kết quả là không kết luận. Nếu dữ liệu đủ nhưng tiêu chí không đạt, đề tài ghi nhận rằng giả thuyết không được ủng hộ, và theo mục 4.3.10, vẫn trả lời được câu hỏi nghiên cứu.

\section{Sản phẩm công bố dự kiến}

Sản phẩm công bố dự kiến gồm hai đến ba bài báo trên tạp chí khoa học, một đến hai báo cáo hội thảo và luận án, theo quy định công bố của chương trình đào tạo tiến sĩ. Tác giả đề cương coi đây là mục tiêu, và kết quả phụ thuộc vào dữ liệu.

Bảng 5 cho biết thời điểm: bài báo 1 và báo cáo hội thảo ở năm 2, bài báo 2 và 3 ở năm 3, luận án ở năm 4. Bảng 5 cũng cho biết bài báo 1 gắn với việc đo rủi ro lựa chọn mô hình và kiểm định H1. Các bài báo 2 và 3 được viết sau khi phân tích biến kênh và đánh giá cơ chế giảm thiểu.

Bài gốc không nói tên tạp chí hay hội thảo dự kiến cho các sản phẩm này, và không nói bài 2 và bài 3 mỗi bài viết về nội dung nào. Bài gốc chỉ nêu hai việc năm 3, phân tích biến kênh (H2, H4) và đánh giá cơ chế giảm thiểu (H3), kèm việc viết bài báo 2 và 3.

\section{Nối với phần còn lại}

Ba nhóm kết quả đều nối với mục tiêu cụ thể ở mục 4.1. Mục tiêu (1) là xây thước đo và xác định chế độ, ứng với nhóm lý thuyết. Mục tiêu (2) và (3) là kiểm định rủi ro theo chế độ và xác định biến kênh, ứng với nhóm thực nghiệm. Mục tiêu (4) là đánh giá cơ chế giảm thiểu và rút ra hàm ý quản trị, ứng với nhóm thực tiễn và chính sách. Việc ghép mục tiêu với nhóm kết quả là cách đối chiếu của người viết.

\begin{tcolorbox}[breakable]
\textbf{Ghi nhớ khi đi thi}

• Mục 4.4 là kịch bản nếu dữ liệu ủng hộ giả thuyết; đề cương chưa có kết quả thực nghiệm nào.

• Ba nhóm kết quả: (1) lý thuyết, khung khái niệm và thước đo; (2) thực nghiệm, bằng chứng Việt Nam về mức độ, thời điểm, kênh; (3) thực tiễn và chính sách, quy trình giám sát mô hình ba bước.

• Quy trình ba bước: khai báo tập mô hình; đo phân tán và tập tin cậy mô hình định kỳ; kích hoạt kết hợp mô hình khi thị trường ở chế độ căng thẳng. Đây là đề xuất suy ra từ H3, chưa kiểm nghiệm với người dùng.

• Tiêu chí ủng hộ: H1 theo mục 4.3.7; H2 khi độ phi thanh khoản dương, tương quan đuôi âm, khoảng tin cậy bootstrap khối của cả hai loại trừ 0; H3 khi lợi suất tương đương chắc chắn sau chi phí cao hơn, mức hối tiếc thấp hơn, độ sụt giảm tối đa không cao hơn; H4 chỉ thăm dò.

• Sản phẩm dự kiến: hai đến ba bài báo tạp chí, một đến hai báo cáo hội thảo và luận án; đây là mục tiêu, kết quả phụ thuộc dữ liệu.
\end{tcolorbox}

\section{Tự kiểm tra}

{\itshape 1. Ba nhóm kết quả dự kiến của đề tài là gì?\par}

Đáp án: Lý thuyết (khung khái niệm và thước đo rủi ro lựa chọn mô hình theo chế độ); thực nghiệm (bằng chứng Việt Nam về mức độ, thời điểm, kênh); thực tiễn và chính sách (quy trình giám sát mô hình ba bước).

{\itshape 2. Quy trình giám sát gồm những bước nào, và trạng thái của nó là gì?\par}

Đáp án: Khai báo tập mô hình; đo phân tán và tập tin cậy mô hình định kỳ; kích hoạt kết hợp mô hình khi thị trường ở chế độ căng thẳng. Đó là đề xuất suy ra từ H3, chưa được kiểm nghiệm với người dùng.

{\itshape 3. Dấu mong đợi của hệ số độ phi thanh khoản và hệ số tương quan đuôi trong tiêu chí ủng hộ H2 là gì?\par}

Đáp án: Hệ số độ phi thanh khoản dương, hệ số tương quan đuôi âm, và khoảng tin cậy bootstrap khối của cả hai loại trừ 0.

{\itshape 4. H1 chỉ được ủng hộ khi nào?\par}

Đáp án: Khi hệ số \texorpdfstring{$\beta$}{beta} của R\_t dương, vượt phân vị 95\% của phân phối rỗng và có khoảng tin cậy bootstrap khối loại trừ 0, và hệ số \texorpdfstring{$\beta$}{beta} của D\_t cùng dấu.

{\itshape 5. Sản phẩm công bố dự kiến là gì và đề cương coi nó thế nào?\par}

Đáp án: Hai đến ba bài báo tạp chí, một đến hai báo cáo hội thảo và luận án; đề cương coi đây là mục tiêu và kết quả phụ thuộc vào dữ liệu.

\chapter{Mục 5 của đề cương: Lý do lựa chọn đơn vị đào tạo}

Tác giả đề cương chọn Trường Quốc tế, Đại học Quốc gia Hà Nội vì ba lý do. Thứ nhất, định hướng của chương trình Tiến sĩ Kinh tế và Quản lý khớp với tính giao thoa của đề tài. Thứ hai, ngôn ngữ đào tạo thuận lợi cho công bố quốc tế, và tác giả có chứng chỉ tiếng Anh Bậc 4. Thứ ba, ngành tốt nghiệp của tác giả thuộc nhóm cần học bổ sung 9 tín chỉ.

\section{Lý do thứ nhất: định hướng chương trình khớp với đề tài}

Theo Thông báo tuyển sinh năm 2026, chương trình Tiến sĩ Kinh tế và Quản lý có tính liên ngành và tập trung vào kinh tế hành vi và quản lý trong lĩnh vực kinh tế. Tác giả đề cương nói đề tài giao thoa giữa tài chính định lượng, quản lý và hành vi, nên khớp với định hướng đó.

Để hiểu lý do này, bạn cần nhớ lại ba gốc của đề tài; cách gọi ba gốc là cách sắp xếp của bản đọc này, dựa trên mục lĩnh vực nghiên cứu. Gốc thứ nhất là tài chính định lượng và quản trị rủi ro, nơi nằm toàn bộ phần đo lường: bảy mô hình, mức hối tiếc, lợi suất tương đương chắc chắn. Gốc thứ hai là quản lý, ở phần thiết kế cơ chế giám sát mô hình, tức quy trình ba bước của mục 4.4. Gốc thứ ba là kinh tế hành vi, nhưng ở một phép thử giới hạn.

Phép thử giới hạn đó là: độ phân tán chéo của lợi suất, một thước đo bầy đàn ở cấp thị trường, có đi cùng rủi ro lựa chọn mô hình sau khi kiểm soát biến động hay không. Đề tài không quan sát hành vi chọn mô hình của tổ chức đầu tư. Vì vậy phần hành vi của đề tài nhỏ và có ranh giới rõ.

Hình dung như sau: đề tài giống một ngôi nhà có phần móng là tài chính định lượng, một căn phòng nhỏ dành cho kinh tế hành vi và một cầu thang dẫn sang quản lý. Chương trình đào tạo được chọn có cả ba thứ đó trong định hướng. Phép ví von này chỉ để giúp bạn nhớ.

Bài gốc không mô tả thêm nội dung chương trình đào tạo, không nêu tên giảng viên và không nêu thứ hạng của đơn vị. Những điều này đề cương chưa quy định ở mục 5. Mục 8 của đề cương, về người hướng dẫn, chưa điền tên ở bản này.

\section{Lý do thứ hai: ngôn ngữ đào tạo và chứng chỉ tiếng Anh}

Ngôn ngữ đào tạo là tiếng Anh hoặc tiếng Việt. Tác giả đề cương nói điều này thuận lợi cho công bố quốc tế. Tác giả có chứng chỉ tiếng Anh Bậc 4 theo Khung năng lực ngoại ngữ 6 bậc dùng cho Việt Nam.

Chứng chỉ cụ thể ở mục 8 là VSTEP, điểm tổng 6,5, do Đại học Phenikaa cấp ngày 26/01/2026. Bảng 5 ghi chứng chỉ này là Bậc 4 (VSTEP). Bài gốc không nói thêm về điểm của từng kĩ năng.

Bạn nên thấy mối nối giữa lý do này và mục 4.4. Sản phẩm dự kiến gồm hai đến ba bài báo tạp chí và một đến hai báo cáo hội thảo, và Bảng 7 liệt kê chín bản thảo đã gửi tạp chí, trong đó có tạp chí xếp hạng Scopus Câu hỏi 2 và Câu hỏi 3. Mối nối này là do người viết bài đọc này đối chiếu, không phải điều đề cương tự nói.

\section{Lý do thứ ba: ngành tốt nghiệp và 9 tín chỉ bổ sung}

Tác giả đề cương tốt nghiệp ngành Kinh tế, chuyên ngành Đầu tư. Theo đối chiếu của tác giả, ngành này thuộc Nhóm 1 (Kinh tế học) trong Phụ lục 1 của Thông báo, nên cần học bổ sung 9 tín chỉ. Trường sẽ xác định điều kiện cuối cùng.

Có hai điểm cần chú ý. Thứ nhất, đây là kết quả đối chiếu của chính tác giả, chưa phải quyết định của Trường. Thứ hai, 9 tín chỉ là điều kiện cho người học thuộc nhóm ngành này, và Bảng 5 xếp việc học bổ sung vào quý 1 và quý 2 của năm 1.

Học phần bổ sung được nêu trong Bảng 5 là Hành vi tổ chức và Lãnh đạo, kèm hai học phần tự chọn. Học phần Hành vi tổ chức và Lãnh đạo gần với phần quản lý và hành vi của đề tài, khớp với lý do thứ nhất. Bài gốc không nêu tên hai học phần tự chọn. Cách nối này là suy ra của người viết, vì đề cương không giải thích vì sao chọn các học phần đó.

\begin{tcolorbox}[breakable]
\textbf{Ghi nhớ khi đi thi}

• Ba lý do: định hướng chương trình khớp với đề tài; ngôn ngữ đào tạo và chứng chỉ tiếng Anh Bậc 4; ngành tốt nghiệp thuộc Nhóm 1 cần học bổ sung 9 tín chỉ.

• Chương trình là Tiến sĩ Kinh tế và Quản lý, liên ngành, tập trung kinh tế hành vi và quản lý trong lĩnh vực kinh tế (Thông báo tuyển sinh năm 2026).

• Điều kiện cuối cùng về học bổ sung do Trường xác định.
\end{tcolorbox}

\section{Tự kiểm tra}

{\itshape 1. Nêu ba lý do chọn đơn vị đào tạo.\par}

Đáp án: Định hướng liên ngành về kinh tế hành vi và quản lý khớp với đề tài; ngôn ngữ đào tạo tiếng Anh hoặc tiếng Việt và chứng chỉ tiếng Anh Bậc 4; ngành Kinh tế chuyên ngành Đầu tư thuộc Nhóm 1 nên học bổ sung 9 tín chỉ.

{\itshape 2. Ai xác định điều kiện cuối cùng về học bổ sung?\par}

Đáp án: Trường.

{\itshape 3. Đề tài giao thoa với những lĩnh vực nào?\par}

Đáp án: Tài chính định lượng, quản lý và hành vi.

\chapter{Mục 7.4 của đề cương: Kế hoạch thực hiện trong 4 năm}

Mục 7.4 xếp toàn bộ công việc vào bốn năm đào tạo, với năm 1 dành cho chuẩn bị, năm 2 cho đo rủi ro và kiểm định H1, năm 3 cho biến kênh và cơ chế giảm thiểu, năm 4 cho luận án và bảo vệ. Thời gian đào tạo là 4 năm đối với người có bằng đại học.

Cách hiểu của bản đọc này: trật tự này có logic, vì mỗi bước dựa vào bước trước, và các thứ cần cố định được đặt trước phần phân tích chính.

\section{Bảng 5 trình bày lại theo năm và theo quý}

Bảng 5 gốc ghi theo năm. Dưới đây là cách xếp lại theo từng quý để bạn nhìn ra dòng chảy công việc. Cách xếp lại này chỉ đổi hình thức, không thêm việc nào.

{\small\begin{longtable}[]{@{}
  >{\raggedright\arraybackslash}p{(\columnwidth - 6\tabcolsep) * \real{0.0797}}
  >{\raggedright\arraybackslash}p{(\columnwidth - 6\tabcolsep) * \real{0.1007}}
  >{\raggedright\arraybackslash}p{(\columnwidth - 6\tabcolsep) * \real{0.5669}}
  >{\raggedright\arraybackslash}p{(\columnwidth - 6\tabcolsep) * \real{0.2527}}@{}}
\toprule\noalign{}
\begin{minipage}[b]{\linewidth}\raggedright
Năm
\end{minipage} & \begin{minipage}[b]{\linewidth}\raggedright
Quý
\end{minipage} & \begin{minipage}[b]{\linewidth}\raggedright
Công việc
\end{minipage} & \begin{minipage}[b]{\linewidth}\raggedright
Sản phẩm
\end{minipage} \\
\midrule\noalign{}
\endhead
\bottomrule\noalign{}
\endlastfoot
1 & Quý 1 & Xác nhận nguồn, độ dài và giấy phép dữ liệu, ngày HOSE công bố VN30 và VN100; xác định khung quy định và mức dùng mô hình định lượng của công ty quản lý quỹ & Chưa ghi riêng cho quý \\
1 & Quý 1 đến quý 2 & Học bổ sung 9 tín chỉ (Hành vi tổ chức và Lãnh đạo, hai học phần tự chọn) và các học phần theo chương trình & Chưa ghi riêng cho quý \\
1 & Quý 2 & Ước lượng thử bảy mô hình trên giai đoạn 2019-2025 để đo thời gian tính và số quan sát ở chế độ căng thẳng; quyết định giữ hay thay phương pháp tính toán Bayes xấp xỉ cho M5 & Chưa ghi riêng cho quý \\
1 & Quý 3 & Làm sạch dữ liệu; đọc toàn văn các nguồn ở mục 4.2 & Chưa ghi riêng cho quý \\
1 & Quý 4 & Đăng ký trước kế hoạch phân tích trên OSF Registries & Sản phẩm cả năm: tổng quan nghiên cứu; kế hoạch phân tích đăng ký trước; tập dữ liệu sạch \\
2 & Quý 1 đến quý 2 & Xác định chế độ (mục 4.3.4) và ước lượng bảy mô hình & Chưa ghi riêng cho quý \\
2 & Quý 3 đến quý 4 & Đo rủi ro lựa chọn mô hình, kiểm định H1, viết bài báo 1 & Sản phẩm cả năm: bài báo 1; báo cáo hội thảo \\
3 & Quý 1 đến quý 2 & Phân tích các biến kênh (H2, H4) & Chưa ghi riêng cho quý \\
3 & Quý 3 đến quý 4 & Đánh giá cơ chế giảm thiểu (H3), viết bài báo 2 và 3 & Sản phẩm cả năm: bài báo 2 và 3 \\
4 & Quý 1 đến quý 2 & Hoàn thiện các chương luận án, phần Kết luận và hàm ý quản trị & Chưa ghi riêng cho quý \\
4 & Quý 3 đến quý 4 & Chỉnh sửa theo góp ý và bảo vệ & Sản phẩm cả năm: luận án \\
\end{longtable}}

\section{Năm 1: chuẩn bị nền và cố định kế hoạch}

Theo Bảng 5, sản phẩm của năm 1 không gồm bài báo. Cách hiểu của bản đọc này: mọi việc ở đây nhằm chuẩn bị cho các năm sau. Sản phẩm là tổng quan nghiên cứu, kế hoạch phân tích đăng ký trước và tập dữ liệu sạch.

\subsection{Quý 1: biết mình có gì trong tay}

Việc đầu tiên là xác nhận nguồn, độ dài và giấy phép dữ liệu, cùng ngày HOSE công bố VN30 và VN100. Cách hiểu của bản đọc này, dựa trên mục 4.3.10: nếu dữ liệu không đủ, đề tài phải thu hẹp rổ cổ phiếu hoặc giai đoạn mẫu, và quyết định đó cần biết sớm.

Việc thứ hai của quý 1 là xác định khung quy định và mức dùng mô hình định lượng của công ty quản lý quỹ. Mục 4.1 của đề cương nói khung Thông tư 83/2025/TT-NHNN áp dụng cho ngân hàng, còn đối tượng dùng kết quả đề tài là công ty quản lý quỹ và tổ chức đầu tư cổ phiếu. Vì vậy năm 1 cần xác định khung áp dụng cho họ. Đề tài không đánh giá mức đầy đủ của khung quy định.

\subsection{Quý 1 đến quý 2: học phần bổ sung và học phần chương trình}

Tác giả đề cương học bổ sung 9 tín chỉ, gồm học phần Hành vi tổ chức và Lãnh đạo và hai học phần tự chọn, cùng các học phần theo chương trình. Mốc này gắn trực tiếp với lý do thứ ba ở mục 5.

\subsection{Quý 2: chạy thử}

Quý 2 ước lượng thử bảy mô hình trên giai đoạn 2019 đến 2025. Hai mục đích được nêu rõ: đo thời gian tính và số quan sát ở chế độ căng thẳng. Từ đó quyết định giữ hay thay phương pháp tính toán Bayes xấp xỉ cho M5.

Hình dung như sau: trước khi chạy marathon thật, bạn chạy thử một quãng để biết giày có vừa không và mình thở thế nào. Chạy thử không phải cuộc đua. Tuy nhiên, vì dữ liệu 2019 đến 2025 đã được dùng trong bài báo trước, đề cương nói rõ không thể khẳng định chưa xem kết quả ở chế độ căng thẳng.

Ước lượng thử cũng là điều kiện để biện minh các giá trị dự kiến. Đề cương nói độ dài cửa sổ 104 tuần và 6 tháng là dự kiến, và tác giả biện minh chúng sau khi ước lượng thử ở năm 1.

\subsection{Quý 3: làm sạch dữ liệu và đọc toàn văn}

Quý 3 làm sạch dữ liệu và đọc toàn văn các nguồn ở mục 4.2. Với dữ liệu, rổ cổ phiếu phải xác định chỉ bằng thông tin có tại từng ngày tái cân bằng. Với tài liệu, đề cương thừa nhận tác giả mới đọc tóm tắt và thông tin thư mục của các nguồn; đọc toàn văn và mở rộng phạm vi tìm kiếm nằm trong kế hoạch năm 1.

\subsection{Quý 4: đăng ký trước}

Quý 4 đăng ký trước kế hoạch phân tích trên OSF Registries. Việc này đặt sau ước lượng thử và đọc toàn văn. Tệp đăng ký cố định độ dài cửa sổ, siêu tham số của bảy mô hình, số chế độ và ngưỡng xác suất lọc, mức ý nghĩa, nhóm kết quả chính, quy tắc hiệu chỉnh so sánh bội và quy tắc không kết luận.

Giải thích chung: những thứ này chỉ cố định hợp lý khi ước lượng thử đã cho thông tin.

Sau ngày đăng ký, dữ liệu phát sinh được giữ làm mẫu kiểm tra riêng, và mọi thay đổi ghi là lệch khỏi kế hoạch.

\section{Năm 2: chế độ, bảy mô hình và bài báo đầu tiên}

Quý 1 và quý 2 xác định chế độ theo mục 4.3.4 và ước lượng bảy mô hình. Xác định chế độ nghĩa là chạy mô hình chuyển chế độ Markov cho lợi suất VN-Index, ước lượng lại theo cửa sổ mở rộng tại mỗi ngày tái cân bằng, rồi lấy xác suất lọc của chế độ căng thẳng. Ước lượng bảy mô hình nghĩa là chạy từ M1 đến M7 trên cửa sổ 104 tuần và đánh giá trên cửa sổ 6 tháng không chồng lấn.

Quý 3 và quý 4 đo rủi ro lựa chọn mô hình, kiểm định H1 và viết bài báo 1. Sản phẩm của năm 2 là bài báo 1 và báo cáo hội thảo.

Cách hiểu của bản đọc này về thứ tự: chỉ khi có chế độ và kết quả của bảy mô hình mới tính được mức hối tiếc, độ phân tán mô hình và tập tin cậy mô hình để đưa vào phương trình (5).

Bài gốc không nói rõ bài báo 1 và báo cáo hội thảo viết về chính xác điều gì ngoài việc chúng nằm ở năm 2, sau việc đo rủi ro và kiểm định H1.

\section{Năm 3: kênh và cơ chế giảm thiểu}

Quý 1 và quý 2 phân tích các biến kênh, tức H2 và H4. Các biến kênh gồm độ phi thanh khoản, tương quan đuôi và tương quan hiệu chỉnh theo biến động, và độ phân tán chéo của lợi suất. Phần của hệ số \(\beta\) được giải thích bởi biến kênh bằng hiệu \(\beta_{0}\  - \ \beta_{1}\) giữa hệ số của phương trình (5) không có và có biến kênh.

Quý 3 và quý 4 đánh giá cơ chế giảm thiểu (H3) và viết bài báo 2 và 3. Thứ tự ở đây tuân theo mạch logic của đề cương: phải biết rủi ro nằm ở đâu và đi qua kênh nào thì mới đánh giá cơ chế giảm. Đây là cách giải thích của người viết, mạch năm bước nêu ở mục 4.3.9 là căn cứ.

\section{Năm 4: luận án và bảo vệ}

Quý 1 và quý 2 hoàn thiện các chương luận án, phần Kết luận và hàm ý quản trị. Phần Kết luận trả lời Câu hỏi 1 đến Câu hỏi 3, nêu hàm ý quản trị, hạn chế (số giai đoạn căng thẳng, phạm vi Việt Nam) và hướng nghiên cứu tiếp.

Quý 3 và quý 4 chỉnh sửa theo góp ý và bảo vệ. Sản phẩm là luận án.

\section{Rủi ro kế hoạch và cách xử lý}

Đề cương nêu hai nhóm rủi ro. Rủi ro về dữ liệu và về số cuộc khủng hoảng trong mẫu ít được xử lý theo mục 4.3.10 và 3.7. Rủi ro tạp chí từ chối bản thảo được xử lý bằng cách chuyển sang tạp chí khác trong cùng năm.

Bạn nên để ý rằng phần rủi ro của mục 7.4 không có dòng riêng về thời gian tính. Quý 2 năm 1 có việc đo thời gian tính và quyết định giữ hay thay phương pháp tính toán Bayes xấp xỉ cho M5. Cách hiểu của bản đọc này: đó là chỗ đề cương xử lý mối lo về thời gian tính, dù bài gốc không gọi nó là rủi ro.

\begin{tcolorbox}[breakable]
\textbf{Ghi nhớ khi đi thi}

• Thời gian đào tạo là 4 năm đối với người có bằng đại học.

• Năm 1: xác nhận dữ liệu và khung quy định (quý 1); học bổ sung 9 tín chỉ (quý 1 đến 2); ước lượng thử và quyết định phương pháp tính toán Bayes xấp xỉ (quý 2); làm sạch dữ liệu và đọc toàn văn (quý 3); đăng ký trước trên OSF (quý 4).

• Năm 2: xác định chế độ và ước lượng bảy mô hình (quý 1 đến 2); đo rủi ro, kiểm định H1, viết bài báo 1 (quý 3 đến 4); sản phẩm là bài báo 1 và báo cáo hội thảo.

• Năm 3: biến kênh H2 và H4 (quý 1 đến 2); cơ chế giảm thiểu H3 và bài báo 2, 3 (quý 3 đến 4).

• Năm 4: hoàn thiện luận án và hàm ý quản trị (quý 1 đến 2); chỉnh sửa và bảo vệ (quý 3 đến 4).

• Rủi ro tạp chí từ chối xử lý bằng cách chuyển tạp chí khác trong cùng năm.
\end{tcolorbox}

\section{Tự kiểm tra}

{\itshape 1. Việc nào nằm ở quý 4 của năm 1, và sản phẩm của năm 1 là gì?\par}

Đáp án: Đăng ký trước kế hoạch phân tích trên OSF Registries. Sản phẩm của năm 1 là tổng quan nghiên cứu, kế hoạch phân tích đăng ký trước và tập dữ liệu sạch.

{\itshape 2. Quý 2 năm 1 làm gì và để quyết định điều gì?\par}

Đáp án: Ước lượng thử bảy mô hình trên giai đoạn 2019-2025 để đo thời gian tính và số quan sát ở chế độ căng thẳng, rồi quyết định giữ hay thay phương pháp tính toán Bayes xấp xỉ cho M5.

{\itshape 3. Năm 3 kiểm định giả thuyết nào, theo thứ tự nào?\par}

Đáp án: Quý 1 đến 2 phân tích biến kênh (H2, H4); quý 3 đến 4 đánh giá cơ chế giảm thiểu (H3) và viết bài báo 2 và 3.

{\itshape 4. Rủi ro tạp chí từ chối bản thảo được xử lý thế nào?\par}

Đáp án: Chuyển sang tạp chí khác trong cùng năm.

\chapter{Mục 8 của đề cương: Kinh nghiệm, kiến thức và sự chuẩn bị của tác giả đề cương}

Mục 8 cho hội đồng thấy tác giả đề cương đã có gì trong tay trước khi bắt đầu. Mục này gồm ba bảng: Bảng 5 về học vấn, ngoại ngữ và nghiên cứu khoa học sinh viên; Bảng 6 về ba công trình đã công bố hoặc đã được chấp nhận đăng; Bảng 7 về chín bản thảo đã gửi tạp chí nhưng chưa công bố tại ngày 30/09/2026. Trong chín bản thảo, ba bản liên quan trực tiếp tới đề tài. Cần đọc kĩ trạng thái từng bài, vì một bài đã công bố, một bài đã được chấp nhận và một bài đang chờ xét là ba mức chắc chắn khác nhau.

\section{Bảng 5: học vấn, ngoại ngữ và nghiên cứu khoa học sinh viên}

{\small\begin{longtable}[]{@{}
  >{\raggedright\arraybackslash}p{(\columnwidth - 2\tabcolsep) * \real{0.1802}}
  >{\raggedright\arraybackslash}p{(\columnwidth - 2\tabcolsep) * \real{0.8198}}@{}}
\toprule\noalign{}
\begin{minipage}[b]{\linewidth}\raggedright
Nội dung
\end{minipage} & \begin{minipage}[b]{\linewidth}\raggedright
Chi tiết
\end{minipage} \\
\midrule\noalign{}
\endhead
\bottomrule\noalign{}
\endlastfoot
Học vấn & Cử nhân Kinh tế, chuyên ngành Đầu tư, Học viện Chính sách và Phát triển, hệ chính quy khóa 2022-2026, xếp loại Giỏi (điểm trung bình 3,41/4 hay 8,06/10). Các học phần liên quan: Quản lý danh mục đầu tư (9,9/10), Xác suất và thống kê (10/10), Kinh tế lượng (8,1/10), Phân tích báo cáo tài chính (8,6/10) \\
Đang theo học & Ngành Kế toán, hệ vừa làm vừa học, Đại học Kinh tế Quốc dân (năm thứ 3); ngành Công nghệ thông tin, hệ đào tạo từ xa, Học viện Công nghệ Bưu chính Viễn thông (năm thứ 3); chương trình thạc sĩ trực tuyến tại WorldQuant University từ 07/2026 \\
Ngoại ngữ & Tiếng Anh Bậc 4 (VSTEP), điểm tổng 6,5, do Đại học Phenikaa cấp ngày 26/01/2026 \\
Nghiên cứu khoa học sinh viên & Chủ nhiệm đề tài về mô hình Black-Litterman cho ngành ngân hàng trên thị trường chứng khoán Việt Nam, đạt Giải Nhất cấp Khoa và Giải Nhì cấp Học viện năm 2026; đề tài đang dự Giải thưởng Khoa học và Công nghệ dành cho sinh viên năm 2026 của Bộ Giáo dục và Đào tạo; Giải Nhì nghiên cứu khoa học sinh viên cấp Khoa năm 2025; Giải Khuyến khích cuộc thi Data Explorer \\
\end{longtable}}

\subsection{Học vấn: mỗi môn gắn với phần nào của đề tài}

Tác giả đề cương tốt nghiệp cử nhân Kinh tế, chuyên ngành Đầu tư, hệ chính quy khóa 2022 đến 2026, xếp loại Giỏi. Điểm trung bình là 3,41 trên thang 4, tương đương 8,06 trên thang 10.

Bốn học phần được liệt kê là các học phần liên quan, với điểm số trên thang 10. Quản lý danh mục đầu tư 9,9. Xác suất và thống kê 10. Kinh tế lượng 8,1. Phân tích báo cáo tài chính 8,6.

Bạn có thể đối chiếu bốn học phần này với các việc của đề tài. Quản lý danh mục đầu tư gần với việc chia tiền vào cổ phiếu, tức bảy mô hình danh mục. Xác suất và thống kê gần với bootstrap, phân phối rỗng và khoảng tin cậy. Kinh tế lượng gần với hồi quy ở phương trình (5) và sai số chuẩn bền sai số chuẩn bền vững. Phân tích báo cáo tài chính gần với dữ liệu về doanh nghiệp niêm yết. Việc đối chiếu này là cách giải thích của người viết, đề cương chỉ nêu tên và điểm của học phần.

\subsection{Đang theo học: hai ngành cùng lúc và một chương trình thạc sĩ}

Bảng 5 ghi ba chương trình đang theo học. Ngành Kế toán hệ vừa làm vừa học tại Đại học Kinh tế Quốc dân ở năm thứ 3. Ngành Công nghệ thông tin hệ đào tạo từ xa tại Học viện Công nghệ Bưu chính Viễn thông ở năm thứ 3. Chương trình thạc sĩ trực tuyến tại WorldQuant University từ 07/2026.

Bài gốc chỉ liệt kê. Đề cương không nói các chương trình này đóng góp gì cho đề tài, và không nêu kết quả học tập ở các chương trình đó.

\subsection{Ngoại ngữ}

Tiếng Anh Bậc 4 (VSTEP), điểm tổng 6,5, do Đại học Phenikaa cấp ngày 26/01/2026. Chứng chỉ này là căn cứ cho lý do thứ hai ở mục 5.

\subsection{Nghiên cứu khoa học sinh viên}

Tác giả đề cương là chủ nhiệm một đề tài về nâng cao hiệu quả danh mục đầu tư bằng mô hình Black-Litterman, áp dụng cho ngành ngân hàng trên thị trường chứng khoán Việt Nam. Đề tài đạt Giải Nhất cấp Khoa và Giải Nhì cấp Học viện năm 2026, và đang dự Giải thưởng Khoa học và Công nghệ dành cho sinh viên năm 2026 do Bộ Giáo dục và Đào tạo tổ chức.

Hai điều cần phân biệt. Giải Nhất cấp Khoa và Giải Nhì cấp Học viện đã đạt. Giải thưởng của Bộ Giáo dục và Đào tạo là đang dự, chưa có kết quả trong đề cương. Ngoài ra tác giả còn đạt Giải Nhì nghiên cứu khoa học sinh viên cấp Khoa năm 2025 và Giải Khuyến khích cuộc thi Data Explorer.

Black-Litterman cũng là tên của hai mô hình M4 và M5 trong Bảng 3. Do vậy đề tài sinh viên này có chung tên mô hình với hai trong bảy mô hình của đề cương. Bài gốc không nói thêm mối quan hệ nào giữa đề tài sinh viên và đề tài tiến sĩ.

\section{Bảng 6: công trình đã công bố hoặc đã được chấp nhận đăng}

{\small\begin{longtable}[]{@{}
  >{\raggedright\arraybackslash}p{(\columnwidth - 6\tabcolsep) * \real{0.0797}}
  >{\raggedright\arraybackslash}p{(\columnwidth - 6\tabcolsep) * \real{0.4019}}
  >{\raggedright\arraybackslash}p{(\columnwidth - 6\tabcolsep) * \real{0.3969}}
  >{\raggedright\arraybackslash}p{(\columnwidth - 6\tabcolsep) * \real{0.1215}}@{}}
\toprule\noalign{}
\begin{minipage}[b]{\linewidth}\raggedright
TT
\end{minipage} & \begin{minipage}[b]{\linewidth}\raggedright
Công trình
\end{minipage} & \begin{minipage}[b]{\linewidth}\raggedright
Tạp chí; ISSN
\end{minipage} & \begin{minipage}[b]{\linewidth}\raggedright
Trạng thái
\end{minipage} \\
\midrule\noalign{}
\endhead
\bottomrule\noalign{}
\endlastfoot
1 & Nguyễn Thanh Bình, Nguyễn Văn Trung*, Tiền gửi không kỳ hạn, hiệu quả hoạt động và ổn định tài chính ngân hàng: Bằng chứng từ mô hình ngưỡng tại Việt Nam & Tạp chí Kinh tế - Luật và Ngân hàng; ISSN 3030-4199; Tập 28, Số 9, tr. 52-64 & Đã công bố \\
2 & Nguyễn Thanh Bình, Nguyễn Văn Trung, Black-Litterman portfolio optimization using regime switching mô hình định giá tài sản vốn and phương pháp tính toán Bayes xấp xỉ: Empirical evidence from the Vietnamese stock market period 2019 - 2025 & Tạp chí Nghiên cứu Chính sách và Phát triển; ISSN 3030-4091; Tập 3, Số 1, tr. 81-99 & Đã công bố \\
3 & Nguyễn Thanh Bình, Nguyễn Văn Trung*, Nguyễn Anh Tuấn, GRI adoption and corporate brownwashing: Board governance evidence from ASEAN-5 & International Journal of Management and Sustainability; Scopus Câu hỏi 3 & Đã được chấp nhận đăng năm 2026, chưa xuất bản tại ngày 30/09/2026 \\
\end{longtable}}

(*) Tác giả liên hệ.

\subsection{Ba mức trạng thái trong Bảng 6}

Bài 1 và bài 2 ở trạng thái đã công bố, có số tập, số kỳ và trang trong bảng. Bài 3 ở trạng thái đã được chấp nhận đăng, và đề cương ghi rõ chưa xuất bản tại ngày 30/09/2026.

Giải thích chung: đã công bố nghĩa là bài đã in ở tạp chí; đã được chấp nhận đăng nghĩa là tạp chí đã đồng ý đăng nhưng bài chưa xuất bản.

Đây là điểm bạn cần giữ đúng khi trả lời hội đồng. Bài 3 không phải là bài đã công bố. Bảng 6 gồm cả hai loại, theo đúng tiêu đề của bảng là đã công bố hoặc đã được chấp nhận đăng.

\subsection{Bài 1: tiền gửi không kỳ hạn và ổn định ngân hàng}

Tiêu đề của bài 1 là Tiền gửi không kỳ hạn, hiệu quả hoạt động và ổn định tài chính ngân hàng: Bằng chứng từ mô hình ngưỡng tại Việt Nam. Bài đăng ở Tạp chí Kinh tế - Luật và Ngân hàng, ISSN 3030-4199, tập 28, số 9, trang 52 đến 64. Bài này ở trạng thái đã công bố. Tác giả đề cương là tác giả liên hệ, đánh dấu bằng dấu sao.

Một mô hình ngưỡng là mô hình cho phép quan hệ giữa các biến đổi dạng khi một biến vượt qua một mức nhất định. Đây là cách giải thích chung về cụm từ trong tiêu đề. Bài gốc không nói bài này kết luận điều gì, và không nêu mối liên hệ trực tiếp với đề tài.

Đề cương dẫn bài này là Nguyễn và Nguyễn (2026b) trong danh mục tài liệu tiếng Việt.

\subsection{Bài 2: Black-Litterman với mô hình định giá tài sản vốn chuyển chế độ và phương pháp tính toán Bayes xấp xỉ}

Bài 2 là bài mà đề cương gọi là Nguyen và Nguyen (2026a). Tiêu đề: Black-Litterman portfolio optimization using regime switching mô hình định giá tài sản vốn and phương pháp tính toán Bayes xấp xỉ: Empirical evidence from the Vietnamese stock market period 2019 - 2025. Bài đăng ở Tạp chí Nghiên cứu Chính sách và Phát triển, ISSN 3030-4091, tập 3, số 1, trang 81 đến 99, và ở trạng thái đã công bố.

Cách hiểu của bản đọc này: bài này gắn chặt với đề tài, vì đề cương dựa vào nó ở nhiều chỗ. M5 là Black-Litterman có chuyển chế độ, mở rộng từ bài này. Tác giả đã dùng mô hình định giá tài sản vốn chuyển chế độ và phương pháp tính toán Bayes xấp xỉ trong bài này. Dữ liệu Việt Nam 2019 đến 2025, gồm hai giai đoạn căng thẳng, cũng đã được dùng trong bài này.

Đề cương cũng đặt hai ràng buộc về bài này. Một là kết quả của bài chỉ làm đối chứng và không được báo cáo lại như kết quả mới của luận án. Hai là vì đã dùng dữ liệu 2019 đến 2025 ở đó, tác giả không thể khẳng định chưa xem kết quả ở chế độ căng thẳng.

So với bài này, đề tài làm ba điều khác. Đề tài đặt M5 cạnh sáu mô hình khác. Đề tài mở mẫu về trước 2019 nếu dữ liệu cho phép. Và đề tài đo chênh lệch giữa các mô hình thay vì đánh giá riêng M5.

\subsection{Bài 3: GRI và brownwashing}

Bài 3 có tiêu đề GRI adoption and corporate brownwashing: Board governance evidence from ASEAN-5, đăng ở International Journal of Management and Sustainability, Scopus Câu hỏi 3. Trạng thái: đã được chấp nhận đăng năm 2026, chưa xuất bản tại ngày 30/09/2026. Tác giả đề cương là tác giả liên hệ, cùng hai đồng tác giả.

Scopus Câu hỏi 3 là một cách xếp hạng tạp chí trong cơ sở dữ liệu Scopus, theo nhóm tư phân vị. Câu hỏi 3 là nhóm thứ ba. Đây là giải thích chung về cách ghi, bài gốc không giải thích.

Bài gốc không nêu mức liên quan trực tiếp của bài 3 với đề tài. Cột liên quan chỉ xuất hiện ở Bảng 7, không có ở Bảng 6.

Dấu sao tác giả liên hệ trong Bảng 6 có ở bài 1 và bài 3, không có ở bài 2 theo cách ghi của bảng. Bài gốc không giải thích sự khác nhau này.

\section{Bảng 7: chín bản thảo đã gửi, chưa công bố tại ngày 30/09/2026}

{\small\begin{longtable}[]{@{}
  >{\raggedright\arraybackslash}p{(\columnwidth - 10\tabcolsep) * \real{0.0787}}
  >{\raggedright\arraybackslash}p{(\columnwidth - 10\tabcolsep) * \real{0.4382}}
  >{\raggedright\arraybackslash}p{(\columnwidth - 10\tabcolsep) * \real{0.0787}}
  >{\raggedright\arraybackslash}p{(\columnwidth - 10\tabcolsep) * \real{0.2229}}
  >{\raggedright\arraybackslash}p{(\columnwidth - 10\tabcolsep) * \real{0.1609}}
  >{\raggedright\arraybackslash}p{(\columnwidth - 10\tabcolsep) * \real{0.0205}}@{}}
\toprule\noalign{}
\begin{minipage}[b]{\linewidth}\raggedright
TT
\end{minipage} & \begin{minipage}[b]{\linewidth}\raggedright
Tên bài báo
\end{minipage} & \begin{minipage}[b]{\linewidth}\raggedright
Số tác giả
\end{minipage} & \begin{minipage}[b]{\linewidth}\raggedright
Tạp chí
\end{minipage} & \begin{minipage}[b]{\linewidth}\raggedright
Trạng thái
\end{minipage} & \begin{minipage}[b]{\linewidth}\raggedright
Liên quan đến đề tài
\end{minipage} \\
\midrule\noalign{}
\endhead
\bottomrule\noalign{}
\endlastfoot
1 & Portfolio optimization with the inverse Black-Litterman framework and clustering machine-learning models: Evidence from Vietnamese bank stocks & 5 & Multidisciplinary Science Journal; Scopus Câu hỏi 3 & Đã sửa xong vòng 2 & Có \\
2 & Tác động của cấu trúc sở hữu lên ổn định tài chính của các ngân hàng thương mại Việt Nam & 2 & Tạp chí Kinh tế - Luật và Ngân hàng & Qua 2 vòng phản biện, đang chỉnh sửa & Không \\
3 & Limits to arbitrage in the Vietnamese physical gold market: The failure of short-term momentum & 3 & Journal of Economic and Banking Studies & Đang sửa vòng 1 & Không \\
4 & Corporate Brownwashing and Firm Valuation in ASEAN-5 Disclosure Regimes & 2 & Cogent Business \& Management; Scopus Câu hỏi 2 & Đang phản biện vòng 1 & Không \\
5 & Structural breaks in Vietnamese stock market liquidity: Evidence from the Amihud illiquidity measure and Bai-Perron regime-shift detection & 6 & Asian Journal of Economics and Banking & Chờ phân công phản biện & Có (kênh thanh khoản, điểm gãy) \\
6 & Nested equity index correlations overstate true co-movement: Evidence from Vietnam & 4 & Asia-Pacific Financial Markets; Scopus Câu hỏi 2 & Ban biên tập đang xem xét & Có (kênh tương quan) \\
7 & Who gains from a market upgrade? Stock liquidity and prices around Vietnam\textquotesingle s FTSE Russell reclassification & 4 & Finance Research Open & Ban biên tập đang xem xét & Không \\
8 & Rated at the peak? Firm valuation around the first LSEG ESG score in five Southeast Asian markets & 4 & Sustainable Finance Review & Chờ xử lý tại tòa soạn & Không \\
9 & Sales-based real earnings management and the cost of debt: Evidence from five ASEAN markets & 5 & Accounting Open & Ban biên tập đang xem xét & Không \\
\end{longtable}}

Tác giả đề cương là tác giả liên hệ ở cả chín bản thảo.

\subsection{Đọc đúng chữ trạng thái}

Chín bản thảo đều chưa công bố tại ngày 30/09/2026. Mọi điều đề cương nói về chúng đều ở dạng trạng thái, không phải kết quả đã được kiểm chứng bởi quá trình phản biện.

Cách hiểu của bản đọc này: vì vậy bạn không nên dùng nội dung các bản thảo như bằng chứng cho đề tài.

Các trạng thái trong Bảng 7 gồm: đã sửa xong vòng 2 (số 1); qua 2 vòng phản biện, đang chỉnh sửa (số 2); đang sửa vòng 1 (số 3); đang phản biện vòng 1 (số 4); chờ phân công phản biện (số 5); ban biên tập đang xem xét (số 6, 7 và 9); chờ xử lý tại tòa soạn (số 8). Chú giải từng trạng thái ở đoạn dưới là giải thích chung về quy trình xét duyệt, không phải thông tin của bài gốc.

Hình dung như sau: một bài báo đi qua nhiều trạm. Trạm đầu là tòa soạn nhận bài. Sau đó biên tập viên cân nhắc có gửi đi phản biện hay không. Phản biện viên đọc và góp ý. Tác giả sửa theo góp ý, có thể qua nhiều vòng. Mỗi trạm chưa phải đích. Bài gốc chỉ nói bài đang ở trạm nào.

\subsection{Từng bản thảo, giải thích nhanh}

Phần dưới đây giải thích tiêu đề bằng lời thường và nhắc lại trạng thái, không thêm thông tin về nội dung bài.

Bản thảo 1 nói về tối ưu hóa danh mục với khung Black-Litterman đảo ngược và các mô hình học máy phân cụm, cho cổ phiếu ngân hàng Việt Nam. Bài có 5 tác giả, gửi Multidisciplinary Science Journal (Scopus Câu hỏi 3), đã sửa xong vòng 2. Cột liên quan ghi Có.

Bản thảo 2 nói về tác động của cấu trúc sở hữu lên ổn định tài chính của các ngân hàng thương mại Việt Nam. Bài có 2 tác giả, gửi Tạp chí Kinh tế - Luật và Ngân hàng, đã qua 2 vòng phản biện và đang chỉnh sửa. Cột liên quan ghi Không.

Bản thảo 3 nói về giới hạn của kinh doanh chênh lệch trên thị trường vàng vật chất Việt Nam, với phụ đề là sự thất bại của đà tăng ngắn hạn. Bài có 3 tác giả, gửi Journal of Economic and Banking Studies, đang sửa vòng 1. Cột liên quan ghi Không.

Bản thảo 4 nói về brownwashing của doanh nghiệp và định giá doanh nghiệp trong các chế độ công bố thông tin ở ASEAN-5. Bài có 2 tác giả, gửi Cogent Business \& Management (Scopus Câu hỏi 2), đang phản biện vòng 1. Cột liên quan ghi Không.

Bản thảo 5 nói về các điểm gãy cấu trúc trong thanh khoản thị trường chứng khoán Việt Nam, đo bằng thước đo bất thanh khoản Amihud và phát hiện chuyển chế độ bằng Bai-Perron. Bài có 6 tác giả, gửi Asian Journal of Economics and Banking, chờ phân công phản biện. Cột liên quan ghi Có, với ghi chú kênh thanh khoản và điểm gãy.

Bản thảo 6 nói về việc các chỉ số cổ phiếu lồng nhau làm tương quan trông lớn hơn mức đồng biến động thật, với bằng chứng từ Việt Nam; đây là cách tiêu đề phát biểu, không phải kết quả đã được kiểm chứng. Bài có 4 tác giả, gửi Asia-Pacific Financial Markets (Scopus Câu hỏi 2), ban biên tập đang xem xét. Cột liên quan ghi Có, với ghi chú kênh tương quan.

Bản thảo 7 hỏi ai được lợi khi thị trường được nâng hạng, xét thanh khoản và giá cổ phiếu quanh lần tái phân loại của FTSE Russell với Việt Nam. Bài có 4 tác giả, gửi Finance Research Open, ban biên tập đang xem xét. Cột liên quan ghi Không.

Bản thảo 8 hỏi liệu doanh nghiệp có được định giá ở đỉnh hay không, quanh lần đầu LSEG công bố điểm ESG ở năm thị trường Đông Nam Á. Bài có 4 tác giả, gửi Sustainable Finance Review, chờ xử lý tại tòa soạn. Cột liên quan ghi Không.

Bản thảo 9 nói về quản trị lợi nhuận thực dựa trên doanh thu và chi phí nợ, với bằng chứng từ năm thị trường ASEAN. Bài có 5 tác giả, gửi Accounting Open, ban biên tập đang xem xét. Cột liên quan ghi Không.

Phần mô tả bằng lời thường ở trên bám sát tiêu đề. Bài gốc không nêu phát hiện của bất kỳ bản thảo nào.

\section{Vì sao ba bản thảo liên quan trực tiếp}

Đề cương nói ba bản thảo liên quan trực tiếp chuẩn bị cho phần xác định chế độ, kênh thanh khoản và kênh tương quan của đề tài. Ba bản thảo đó là các số 1, 5 và 6 trong Bảng 7.

Bản thảo 5 gắn chặt với kênh thanh khoản. Thước đo Amihud trong tiêu đề chính là độ phi thanh khoản của đề cương, biến kênh thanh khoản của phương trình (5), và đề cương trích Amihud (2002). Phần điểm gãy trong tiêu đề gắn với Bai-Perron, vốn xuất hiện ở mục 4.3.4 như công cụ mô tả và đối chiếu, và ở mục 4.3.10 như công cụ tác giả đã dùng trong một bản thảo.

Bản thảo 6 gắn với kênh tương quan. Tương quan đuôi và tương quan hiệu chỉnh theo biến động là biến kênh trong phương trình (5), với cơ sở Longin và Solnik (2001) và Forbes và Rigobon (2002). Tiêu đề bản thảo 6 nói về chỉ số lồng nhau làm tương quan trông lớn hơn, nên cùng chủ đề đo tương quan, theo cột liên quan ghi kênh tương quan.

Bản thảo 1 ghi Có ở cột liên quan nhưng không có ghi chú nêu phần nào của đề tài. Tiêu đề của nó nhắc Black-Litterman, mô hình xuất hiện ở M4 và M5. Bài gốc không nói rõ bản thảo 1 chuẩn bị cho phần nào trong ba phần xác định chế độ, kênh thanh khoản và kênh tương quan. Bản thảo 5 và 6 thì có ghi chú kênh thanh khoản và kênh tương quan.

Điều cần giữ khi đi thi: ba bản thảo này chưa được công bố. Cách hiểu của bản đọc này: chúng là sự chuẩn bị, chưa phải kết quả cho đề tài. Đề cương chỉ nói chúng chuẩn bị cho các phần xác định chế độ, thanh khoản và tương quan, và không trình bày chúng như bằng chứng cho H1 đến H4.

\begin{tcolorbox}[breakable]
\textbf{Ghi nhớ khi đi thi}

• Bảng 5: cử nhân Kinh tế (Đầu tư), xếp loại Giỏi (3,41/4 hay 8,06/10); Quản lý danh mục đầu tư 9,9, Xác suất và thống kê 10, Kinh tế lượng 8,1, Phân tích báo cáo tài chính 8,6; tiếng Anh Bậc 4 (VSTEP) điểm tổng 6,5.

• Bảng 6 có ba công trình: hai bài đã công bố, một bài đã được chấp nhận đăng năm 2026 và chưa xuất bản tại ngày 30/09/2026.

• Bài 2 trong Bảng 6 là Nguyen và Nguyen (2026a), nguồn của mô hình định giá tài sản vốn chuyển chế độ, phương pháp tính toán Bayes xấp xỉ và M5; kết quả của bài chỉ làm đối chứng.

• Bảng 7 có chín bản thảo chưa công bố tại ngày 30/09/2026; tác giả đề cương là tác giả liên hệ ở cả chín.

• Ba bản thảo liên quan trực tiếp (số 1, 5, 6) chuẩn bị cho phần xác định chế độ, kênh thanh khoản và kênh tương quan; số 1 đã sửa xong vòng 2, số 5 chờ phân công phản biện, số 6 ban biên tập đang xem xét.
\end{tcolorbox}

\section{Tự kiểm tra}

{\itshape 1. Bài nào trong Bảng 6 chưa xuất bản tại ngày 30/09/2026, và trạng thái chính xác là gì?\par}

Đáp án: Bài 3, GRI adoption and corporate brownwashing, ở trạng thái đã được chấp nhận đăng năm 2026, chưa xuất bản tại ngày 30/09/2026.

{\itshape 2. Bản thảo nào trong Bảng 7 đã sửa xong vòng 2 và nó có liên quan đến đề tài không?\par}

Đáp án: Bản thảo số 1, về Black-Litterman đảo ngược và các mô hình học máy phân cụm cho cổ phiếu ngân hàng Việt Nam; cột liên quan ghi Có.

{\itshape 3. Ba bản thảo liên quan trực tiếp chuẩn bị cho những phần nào của đề tài?\par}

Đáp án: Phần xác định chế độ, kênh thanh khoản và kênh tương quan. Bản thảo số 5 ghi kênh thanh khoản và điểm gãy; số 6 ghi kênh tương quan; bản thảo số 1 ghi Có.

{\itshape 4. Tác giả đề cương là tác giả liên hệ ở những bản thảo nào của Bảng 7?\par}

Đáp án: Cả chín bản thảo.

{\itshape 5. Vì sao bài Nguyen và Nguyen (2026a) quan trọng với đề tài, và đề cương hạn chế việc dùng kết quả của nó thế nào?\par}

Đáp án: M5 mở rộng từ bài này, tác giả đã dùng mô hình định giá tài sản vốn chuyển chế độ và phương pháp tính toán Bayes xấp xỉ ở đó; kết quả bài chỉ làm đối chứng và không được báo cáo lại như kết quả mới.

\chapter{Mười hai điều cần nhớ về toàn đề cương}

Mười hai điều dưới đây gom toàn bộ đề cương thành những ý bạn có thể nói lại trong vài phút. Mỗi điều có chỉ dẫn mục gốc để bạn mở lại phần đầy đủ. Hãy nhớ rằng đây là đề cương, chưa có kết quả thực nghiệm nào.

\subsection{1. Câu hỏi trung tâm: chọn một mô hình cũng là một rủi ro}

Tổ chức đầu tư phân bổ danh mục bằng mô hình thống kê, và DeMiguel và cộng sự (2009) không thấy mô hình nào nhất quán vượt quy tắc 1/N trên bảy bộ dữ liệu. Vì vậy việc chọn mô hình là một quyết định có rủi ro riêng. Rủi ro lựa chọn mô hình là hậu quả bất lợi kỳ vọng, đo bằng mức hối tiếc, do phải chọn một mô hình trước khi biết mô hình nào tốt nhất. Đề tài chuyển từ đánh giá một mô hình sang câu hỏi chọn giữa nhiều mô hình. (Mục 4.1, 2.1)

\subsection{2. Đây là đề cương, chưa có kết quả}

Mọi kết quả trong đề cương là dự kiến hoặc kịch bản nếu dữ liệu ủng hộ. Đề cương dùng các chữ dự kiến, nếu, sẽ kiểm định, và nêu cả khả năng giả thuyết bị bác bỏ. Bạn không được trình bày bất kỳ phát hiện nào như đã có. (Mục 4.3.10, 4)

Lời khuyên của bản đọc này: khi đi thi, câu an toàn là đề tài sẽ kiểm định, không phải đề tài đã chứng minh.

\subsection{3. Bảy cách chia tiền và quy tắc chọn trước}

Tập mô hình cạnh tranh gồm bảy mô hình: M1 phân bổ đều 1/N, M2 phương sai nhỏ nhất với hiệp phương sai co rút, M3 trung bình-phương sai, M4 Black-Litterman, M5 Black-Litterman có chuyển chế độ, M6 danh mục vững chắc với nhiều niềm tin tiên nghiệm, M7 giảm rủi ro theo xác suất chế độ căng thẳng. Mô hình chọn trước là mô hình có lợi suất tương đương chắc chắn cao nhất trong 36 tháng trước ngày tái cân bằng. Mỗi mô hình ước lượng trên cửa sổ 104 tuần và đánh giá trên cửa sổ 6 tháng không chồng lấn, những giá trị dự kiến. Cấu hình cơ sở là danh mục chỉ mua (long-only) trên cổ phiếu niêm yết trên HOSE. (Mục 4.3.5, 3.3)

\subsection{4. Chế độ thị trường là trạng thái thống kê, không phải tên cuộc khủng hoảng}

Chế độ căng thẳng do mô hình chuyển chế độ Markov (Hamilton, 1989) ước lượng cho lợi suất VN-Index, là chế độ có phương sai điều kiện cao nhất và lợi suất kỳ vọng thấp hơn chế độ bình thường. Tham số được ước lượng lại theo cửa sổ mở rộng tại mỗi ngày tái cân bằng, chỉ bằng dữ liệu đến ngày đó, để tránh thiên lệch nhìn trước. Khủng hoảng là sự kiện lịch sử và đề tài đối chiếu hai khái niệm này, không dùng thay nhau. (Mục 4.1, 3.1.2)

\subsection{5. Ba thước đo rủi ro và một thang đo hiệu quả}

Thang đo hiệu quả là lợi suất tương đương chắc chắn sau chi phí giao dịch, \(U_{m,t}\  = \ \mu_{m,t}\  - \ (\rho/2)\sigma_{m,t}^{2}\) với \(\rho\) cơ sở bằng 5. Ba thước đo rủi ro là độ phân tán mô hình \(D_{t}\), mức hối tiếc \(R_{t}\) và kích thước tập tin cậy mô hình ở mức ý nghĩa 10\%. Chỉ mức hối tiếc đo rủi ro, còn độ phân tán và tập tin cậy mô hình là chỉ báo của bất định. Tỷ số Sharpe chỉ là thước đo phụ, vì nó không còn đơn điệu theo chất lượng danh mục khi lợi suất vượt trội trung bình âm. (Mục 4.1, 3.3)

\subsection{6. Bốn giả thuyết, mỗi giả thuyết có một khả năng đối lập}

H1: rủi ro lựa chọn mô hình ở chế độ căng thẳng cao hơn bình thường. H2: sau khi kiểm soát biến động thực hiện, hệ số độ phi thanh khoản dương và hệ số tương quan đuôi âm. H3: kết hợp mô hình và danh mục vững chắc giảm mức hối tiếc và độ sụt giảm tối đa, tăng lợi suất tương đương chắc chắn, so với chọn một mô hình ở chế độ căng thẳng, sau chi phí giao dịch. H4 là thăm dò: độ phân tán chéo của lợi suất cao đi cùng rủi ro cao. Mỗi giả thuyết đi kèm một khả năng đối lập riêng; chẳng hạn khả năng đối lập của H1 là rủi ro không đổi theo chế độ, hoặc 1/N thắng ổn định. (Mục 4.3.2)

\subsection{7. Kiểm định bằng hồi quy, bootstrap khối và phân phối rỗng}

Phương trình (5) hồi quy \(D_{t}\) hoặc \(R_{t}\) theo biến chế độ \(S_{t}\) và các biến kênh, với sai số chuẩn bền sai số chuẩn bền vững. Hệ số \(\beta\) đo chênh lệch ở chế độ căng thẳng so với bình thường. H1 chỉ được ủng hộ khi \(\beta\) của \(R_{t}\) dương, vượt phân vị 95\% của phân phối rỗng gồm 1.000 lần tái lấy mẫu bootstrap khối, có khoảng tin cậy loại trừ 0, và \(\beta\) của \(D_{t}\) cùng dấu. Kết quả gọi là liên hệ có điều kiện, không gọi là quan hệ nhân quả. (Mục 4.3.7)

\subsection{8. Ba cơ chế giảm thiểu và sáu tiêu chí}

Cơ chế (i) kết hợp mô hình theo trọng số Bayes hoặc chỉ trong tập tin cậy mô hình, với đối chứng là kết hợp 1/N và trung bình-phương sai. Cơ chế (ii) danh mục vững chắc tối ưu cho tình huống xấu nhất trong nhiều niềm tin ban đầu. Cơ chế (iii) quy tắc chuyển chế độ, với đối chứng cùng M5 là giá trị của chuyển chế độ. Sáu tiêu chí gồm lợi suất tương đương chắc chắn, mức hối tiếc, độ sụt giảm tối đa, giá trị thiệt hại kỳ vọng mức 95\%, vòng quay và hiệu quả sau chi phí. (Mục 4.3.8)

\subsection{9. Quy tắc dưới 10 cửa sổ và quyền được nói không kết luận}

Số giai đoạn căng thẳng độc lập ít nên kiểm định có sức mạnh thấp. Khi số cửa sổ độc lập ở chế độ căng thẳng dưới 10, kết quả được gọi là không kết luận, không ủng hộ cũng không bác bỏ. Hiệu chỉnh so sánh bội Holm ở mức 10\% áp dụng cho bốn giả thuyết, và kết luận giới hạn trong phạm vi Việt Nam. Kết quả có thể bác bỏ giả thuyết, khi đó đề tài vẫn trả lời: rủi ro không đổi theo chế độ, hoặc 1/N đủ dùng. (Mục 4.3.10)

\subsection{10. Đăng ký trước và minh bạch, kèm một giới hạn}

Kế hoạch phân tích đăng ký trên OSF Registries vào quý 4 năm 1, khóa có dấu thời gian, cố định cửa sổ, siêu tham số, số chế độ, mức ý nghĩa và quy tắc không kết luận. Mã nguồn và tệp đăng ký được lưu, cách xử lý dữ liệu được mô tả đầy đủ. Nhưng tác giả đã dùng dữ liệu 2019 đến 2025 trong bài báo trước, nên không thể khẳng định chưa xem kết quả ở chế độ căng thẳng. Đăng ký trước chỉ giảm, không loại thiên lệch. (Mục 4.3.5, 3.7)

\subsection{11. Kết quả dự kiến chia làm ba nhóm, và quy trình giám sát chưa được kiểm nghiệm}

Nếu dữ liệu ủng hộ, kết quả gồm khung khái niệm và thước đo, bằng chứng Việt Nam về mức độ, thời điểm và kênh, và một quy trình giám sát mô hình ba bước. Ba bước là khai báo tập mô hình, đo phân tán và tập tin cậy mô hình định kỳ, kích hoạt kết hợp mô hình khi thị trường ở chế độ căng thẳng. Quy trình là đề xuất suy ra từ H3, chưa kiểm nghiệm với người dùng. Sản phẩm công bố dự kiến là hai đến ba bài báo, một đến hai báo cáo hội thảo và luận án, tùy dữ liệu. (Mục 4.4)

\subsection{12. Kế hoạch 4 năm, sự chuẩn bị của tác giả và phạm vi của tuyên bố tính mới}

Năm 1 xác nhận dữ liệu, học bổ sung 9 tín chỉ, ước lượng thử và đăng ký trước; năm 2 xác định chế độ, ước lượng bảy mô hình và kiểm định H1; năm 3 phân tích biến kênh và cơ chế giảm thiểu; năm 4 hoàn thiện luận án và bảo vệ. Tác giả có hai bài đã công bố, một bài đã được chấp nhận đăng và chín bản thảo đã gửi, trong đó ba bản liên quan trực tiếp. Tuyên bố tính mới chỉ nói theo hiểu biết của tác giả, trên phạm vi tìm kiếm không gồm Scopus, Web of Science và nghiên cứu tiếng Việt trên tạp chí trong nước. (Mục 4.1, 2.3, 6, 7)

\chapter{Thuật ngữ cần nhớ}

Bảng dưới đây gom các thuật ngữ xuất hiện trong toàn đề cương, xếp theo nhóm để bạn dễ ôn. Mỗi dòng có tên tiếng Việt, tên tiếng Anh để nhận ra trong tài liệu gốc, phần giải nghĩa bằng lời thường, và mục của đề cương nơi thuật ngữ xuất hiện. Cột mục gốc ghi mục có thuật ngữ, không phải nơi duy nhất nhắc đến nó.

\section{Nhóm 1: ba loại rủi ro và các khái niệm cốt lõi}

{\small\begin{longtable}[]{@{}
  >{\raggedright\arraybackslash}p{(\columnwidth - 6\tabcolsep) * \real{0.1643}}
  >{\raggedright\arraybackslash}p{(\columnwidth - 6\tabcolsep) * \real{0.1508}}
  >{\raggedright\arraybackslash}p{(\columnwidth - 6\tabcolsep) * \real{0.6224}}
  >{\raggedright\arraybackslash}p{(\columnwidth - 6\tabcolsep) * \real{0.0625}}@{}}
\toprule\noalign{}
\begin{minipage}[b]{\linewidth}\raggedright
Thuật ngữ
\end{minipage} & \begin{minipage}[b]{\linewidth}\raggedright
Tiếng Anh
\end{minipage} & \begin{minipage}[b]{\linewidth}\raggedright
Giải nghĩa
\end{minipage} & \begin{minipage}[b]{\linewidth}\raggedright
Mục gốc
\end{minipage} \\
\midrule\noalign{}
\endhead
\bottomrule\noalign{}
\endlastfoot
Rủi ro mô hình & model risk & Khả năng phát sinh hậu quả bất lợi từ quyết định dựa trên mô hình sai hoặc dùng sai. Khái niệm này theo thư giám sát SR 11-7 của Cục Dự trữ Liên bang Mỹ. & 2.1 \\
Bất định mô hình & model uncertainty & Tình huống có nhiều mô hình khả dĩ cho kết quả khác nhau, và bạn chưa biết mô hình nào đúng hơn. & 2.1 \\
Rủi ro lựa chọn mô hình & model selection risk & Hậu quả bất lợi kỳ vọng, đo bằng regret, do phải chọn một mô hình trước khi biết mô hình nào tốt nhất. Đây là đối tượng trung tâm của đề tài. & 2.1 \\
Regret & regret & Khoản thiệt vì đã không chọn mô hình tốt nhất sau thực tế. Bằng CEQ cao nhất trong bảy mô hình trừ CEQ của mô hình chọn trước. & 3.3 \\
Độ phân tán mô hình & model dispersion, D\_t & Độ lệch chuẩn chéo giữa các mô hình của CEQ. Chỉ báo mức bất định, cho biết bảy cách chia tiền cho kết quả khác nhau đến đâu. & 3.3 \\
Kích thước MCS & MCS size & Số mô hình còn lại trong tập tin cậy mô hình. Cũng là chỉ báo của bất định, không phải của rủi ro. & 3.3 \\
Tập tin cậy mô hình & Model Confidence Set, MCS & Danh sách những ứng viên chưa thể loại sau vòng sơ loại thống kê, theo Hansen và cộng sự (2011). Mức ý nghĩa của đề tài ấn định trước ở 10\%. & 2.2.3, 3.3 \\
Tập mô hình cạnh tranh & competing model set & Bảy mô hình M1 đến M7 được đặt cạnh nhau để so, thay vì đánh giá riêng từng mô hình. & 3.2 \\
Mô hình danh mục & portfolio model & Một cách chia tiền vào các cổ phiếu. Bảy mô hình là bảy cách chia tiền khác nhau. & 3.2 \\
Quản trị danh mục & portfolio management & Việc quyết định chia vốn vào các tài sản và điều chỉnh theo thời gian. Phạm vi của đề tài là danh mục cổ phiếu niêm yết trên HOSE. & 1 \\
Khoảng trống nghiên cứu & research gap & Điều các công trình trước chưa làm. Đề cương nêu bốn khoảng trống G1 đến G4 ở Bảng 1. & 2.3 \\
Tính mới & novelty & Điều đề tài làm khác các công trình trước. Đề cương chỉ tuyên bố theo hiểu biết của tác giả, trong một phạm vi tìm kiếm có giới hạn. & 2.3 \\
\end{longtable}}

\section{Nhóm 2: chế độ thị trường và nhận diện chế độ}

{\small\begin{longtable}[]{@{}
  >{\raggedright\arraybackslash}p{(\columnwidth - 6\tabcolsep) * \real{0.1652}}
  >{\raggedright\arraybackslash}p{(\columnwidth - 6\tabcolsep) * \real{0.1551}}
  >{\raggedright\arraybackslash}p{(\columnwidth - 6\tabcolsep) * \real{0.6050}}
  >{\raggedright\arraybackslash}p{(\columnwidth - 6\tabcolsep) * \real{0.0746}}@{}}
\toprule\noalign{}
\begin{minipage}[b]{\linewidth}\raggedright
Thuật ngữ
\end{minipage} & \begin{minipage}[b]{\linewidth}\raggedright
Tiếng Anh
\end{minipage} & \begin{minipage}[b]{\linewidth}\raggedright
Giải nghĩa
\end{minipage} & \begin{minipage}[b]{\linewidth}\raggedright
Mục gốc
\end{minipage} \\
\midrule\noalign{}
\endhead
\bottomrule\noalign{}
\endlastfoot
Chế độ thị trường & market regime & Trạng thái của thị trường, giống thời tiết. Chế độ bình thường như trời quang, chế độ căng thẳng như mùa bão. & 3.1.2 \\
Chế độ bình thường & normal regime & Chế độ có phương sai điều kiện thấp hơn và lợi suất kỳ vọng cao hơn chế độ căng thẳng. & 3.1.2 \\
Chế độ căng thẳng & stress regime & Chế độ có phương sai điều kiện cao nhất và lợi suất kỳ vọng thấp hơn chế độ bình thường. Là trạng thái thống kê, không phải tên một cuộc khủng hoảng. & 3.1.2 \\
Khủng hoảng & crisis & Sự kiện lịch sử như khủng hoảng tài chính toàn cầu hay đại dịch COVID-19. Khác với chế độ căng thẳng, và hai khái niệm không dùng thay nhau. & 1.1, 2.1 \\
Mô hình chuyển chế độ Markov & Markov regime-switching model & Mô hình cho rằng thị trường ở một trong vài trạng thái ẩn và chuyển giữa chúng với những xác suất nhất định, theo Hamilton (1989). & 2.2.4, 3.1.2 \\
Chuỗi Markov & Markov chain & Quá trình mà bước kế tiếp chỉ phụ thuộc vào trạng thái hiện tại. Có mặt trong cả mô hình chuyển chế độ lẫn tên ABC-MCMC. & 3.1.2, 3.7 \\
Xác suất lọc & filtered probability & Xác suất tại ngày t rằng thị trường đang ở chế độ căng thẳng, tính chỉ bằng dữ liệu đến ngày đó. & 3.1.2 \\
Cửa sổ mở rộng & expanding window & Cách ước lượng lại tham số với dữ liệu ngày càng dài, luôn bắt đầu từ đầu mẫu đến ngày hiện tại. & 3.1.2 \\
Điểm gãy cấu trúc & structural break & Thời điểm quy luật của dữ liệu đổi hẳn. Đề tài dùng kiểm định Bai-Perron (1998) chỉ để mô tả và đối chiếu. & 2.2.4, 3.1.2 \\
Tiêu chí thông tin & information criterion & Chỉ số dùng để chọn số chế độ, hai hay ba, trên giai đoạn ước lượng ban đầu. & 3.1.2 \\
Thiên lệch nhìn trước & look-ahead bias & Sai lệch do dùng thông tin chưa có tại thời điểm ra quyết định. Đề tài tránh bằng cách chỉ dùng dữ liệu đến ngày đó. & 3.1.1, 3.1.2 \\
Thiên lệch sống sót & survivorship bias & Sai lệch do chỉ xét các cổ phiếu còn tồn tại và bỏ qua những cổ phiếu đã biến mất. & 3.1.1 \\
\end{longtable}}

\section{Nhóm 3: dữ liệu và thị trường Việt Nam}

{\small\begin{longtable}[]{@{}
  >{\raggedright\arraybackslash}p{(\columnwidth - 6\tabcolsep) * \real{0.1690}}
  >{\raggedright\arraybackslash}p{(\columnwidth - 6\tabcolsep) * \real{0.1748}}
  >{\raggedright\arraybackslash}p{(\columnwidth - 6\tabcolsep) * \real{0.5965}}
  >{\raggedright\arraybackslash}p{(\columnwidth - 6\tabcolsep) * \real{0.0597}}@{}}
\toprule\noalign{}
\begin{minipage}[b]{\linewidth}\raggedright
Thuật ngữ
\end{minipage} & \begin{minipage}[b]{\linewidth}\raggedright
Tiếng Anh
\end{minipage} & \begin{minipage}[b]{\linewidth}\raggedright
Giải nghĩa
\end{minipage} & \begin{minipage}[b]{\linewidth}\raggedright
Mục gốc
\end{minipage} \\
\midrule\noalign{}
\endhead
\bottomrule\noalign{}
\endlastfoot
HOSE & Ho Chi Minh Stock Exchange & Sở Giao dịch Chứng khoán Thành phố Hồ Chí Minh, nơi niêm yết các cổ phiếu của đề tài. & 3.1.1 \\
VN-Index & VN-Index & Chỉ số toàn thị trường HOSE. Lợi suất của nó được dùng để xác định chế độ. & 3.1.1 \\
VN30 và VN100 & VN30, VN100 & Hai rổ gồm 30 và 100 cổ phiếu có vốn hóa và thanh khoản lớn nhất. Đề tài chỉ dùng để đối chiếu và kiểm tra độ bền. & 3.1.1 \\
Vốn hóa tự do chuyển nhượng & free-float market capitalization & Giá trị thị trường của phần cổ phiếu được tự do giao dịch. Dùng cùng giá trị giao dịch để chọn 30 cổ phiếu của rổ. & 3.1.1 \\
Giá đóng cửa điều chỉnh & adjusted closing price & Giá cuối ngày đã điều chỉnh theo các sự kiện của doanh nghiệp. Là một trong các dữ liệu của đề tài. & 3.1.1 \\
Lãi suất phi rủi ro & risk-free rate & Lãi suất của khoản đầu tư coi như không rủi ro. Đề tài lấy lợi suất trái phiếu chính phủ kỳ hạn một năm. & 3.1.1 \\
Biên độ dao động giá & price limit & Mức giá được phép tăng giảm tối đa trong một phiên. Đề cương ghi kèm nguồn ở chương 3 của luận án. & 3.1.1 \\
Công ty quản lý quỹ & fund management company & Đối tượng dùng kết quả của đề tài cùng các tổ chức đầu tư cổ phiếu, khác với ngân hàng thương mại. & 1.1 \\
Thông tư 83/2025/TT-NHNN & Circular 83/2025/TT-NHNN & Văn bản của Ngân hàng Nhà nước về hệ thống kiểm soát nội bộ của ngân hàng thương mại. Đề tài không đánh giá khung quy định. & 1.1 \\
SR 11-7 và SR 26-2 & SR 11-7, SR 26-2 & Thư giám sát về quản lý rủi ro mô hình của Cục Dự trữ Liên bang Mỹ. Không ràng buộc tổ chức Việt Nam, chỉ là tham chiếu quốc tế. & 1.1 \\
Ngân hàng thương mại & commercial bank & Loại tổ chức mà Thông tư 83/2025 áp dụng. Khác với công ty quản lý quỹ, đối tượng dùng kết quả của đề tài. & 1.1 \\
\end{longtable}}

\section{Nhóm 4: các mô hình danh mục}

{\small\begin{longtable}[]{@{}
  >{\raggedright\arraybackslash}p{(\columnwidth - 6\tabcolsep) * \real{0.1788}}
  >{\raggedright\arraybackslash}p{(\columnwidth - 6\tabcolsep) * \real{0.1728}}
  >{\raggedright\arraybackslash}p{(\columnwidth - 6\tabcolsep) * \real{0.5893}}
  >{\raggedright\arraybackslash}p{(\columnwidth - 6\tabcolsep) * \real{0.0591}}@{}}
\toprule\noalign{}
\begin{minipage}[b]{\linewidth}\raggedright
Thuật ngữ
\end{minipage} & \begin{minipage}[b]{\linewidth}\raggedright
Tiếng Anh
\end{minipage} & \begin{minipage}[b]{\linewidth}\raggedright
Giải nghĩa
\end{minipage} & \begin{minipage}[b]{\linewidth}\raggedright
Mục gốc
\end{minipage} \\
\midrule\noalign{}
\endhead
\bottomrule\noalign{}
\endlastfoot
Phân bổ đều 1/N & 1/N rule, naive diversification & Chia đều vốn cho N tài sản, như chia đều trứng vào N giỏ. Là M1 của đề tài. & 3.2 \\
Trung bình-phương sai & mean-variance & Khung của Markowitz (1952): chọn danh mục có lợi suất kỳ vọng cao nhất với phương sai cho trước. Là M3. & 1.1, 3.2 \\
Phương sai nhỏ nhất & minimum variance & Chọn danh mục có độ biến động thấp nhất. M2 làm điều này với hiệp phương sai co rút. & 3.2 \\
Hiệp phương sai & covariance & Thước đo hai tài sản cùng lên xuống mạnh đến đâu. Sai số khi ước lượng nó truyền sang trọng số danh mục. & 1.1 \\
Co rút hiệp phương sai & covariance shrinkage & Kéo ước lượng hiệp phương sai về một mốc ổn định hơn, theo Ledoit và Wolf (2004), để giảm sai số ước lượng. & 2.2.1 \\
Black-Litterman & Black-Litterman & Mô hình kết hợp cân bằng toàn cầu của CAPM với quan điểm nhà đầu tư, theo Black và Litterman (1992). Có mặt ở M4 và M5. & 2.2.1, 3.2 \\
CAPM & Capital Asset Pricing Model & Mô hình định giá tài sản vốn, làm nền cho cân bằng toàn cầu trong Black-Litterman. Bài Nguyen và Nguyen (2026a), nguồn của M5, dùng CAPM chuyển chế độ. & 2.2.1, 3.7 \\
Niềm tin ban đầu & prior & Giả định về tham số mà bạn có trước khi thấy dữ liệu. Danh mục vững chắc dùng nhiều niềm tin ban đầu cùng lúc. & 2.2.1 \\
Danh mục vững chắc đa prior & robust multi-prior portfolio & Danh mục tối ưu cho tình huống xấu nhất trong nhiều niềm tin ban đầu, theo Garlappi và cộng sự (2007). Là M6. & 3.2, 3.5 \\
Né tránh mơ hồ & ambiguity aversion & Thái độ của nhà đầu tư không chắc chắn về chính xác phân phối của lợi suất và ưa phương án ít bị tổn hại trong tình huống xấu. & 2.2.1 \\
Maxmin & maxmin & Tiêu chí chọn phương án tốt nhất trong tình huống xấu nhất, theo Gilboa và Schmeidler (1989). & 2.2.2 \\
Giảm rủi ro theo xác suất chế độ căng thẳng & regime-probability de-risking & Quy tắc giảm rủi ro khi xác suất thị trường ở chế độ căng thẳng tăng, theo Kritzman và cộng sự (2012). Là M7. & 3.2, 3.5 \\
Nghịch đảo giả Moore-Penrose & Moore-Penrose pseudo-inverse & Phép nghịch đảo tổng quát dùng khi ma trận không khả nghịch thông thường. M3 dùng khi số cổ phiếu chia số quan sát tuần vượt 0,5. & 3.3 \\
Long-only (chỉ mua) & long-only & Danh mục chỉ được mua, không được bán khống. Là cấu hình cơ sở của đề tài. & 3.1.1 \\
\end{longtable}}

\section{Nhóm 5: thang đo hiệu quả và rủi ro}

{\small\begin{longtable}[]{@{}
  >{\raggedright\arraybackslash}p{(\columnwidth - 6\tabcolsep) * \real{0.1899}}
  >{\raggedright\arraybackslash}p{(\columnwidth - 6\tabcolsep) * \real{0.1805}}
  >{\raggedright\arraybackslash}p{(\columnwidth - 6\tabcolsep) * \real{0.5723}}
  >{\raggedright\arraybackslash}p{(\columnwidth - 6\tabcolsep) * \real{0.0574}}@{}}
\toprule\noalign{}
\begin{minipage}[b]{\linewidth}\raggedright
Thuật ngữ
\end{minipage} & \begin{minipage}[b]{\linewidth}\raggedright
Tiếng Anh
\end{minipage} & \begin{minipage}[b]{\linewidth}\raggedright
Giải nghĩa
\end{minipage} & \begin{minipage}[b]{\linewidth}\raggedright
Mục gốc
\end{minipage} \\
\midrule\noalign{}
\endhead
\bottomrule\noalign{}
\endlastfoot
Mức thỏa dụng tương đương chắc chắn & certainty equivalent, CEQ & Số tiền chắc chắn bạn sẵn lòng nhận để đổi lấy một khoản đầu tư có rủi ro. Bằng lợi suất trung bình trừ một khoản phạt theo phương sai. Càng cao càng tốt. & 3.3 \\
Hệ số ngại rủi ro & risk aversion coefficient, $\rho$ & Mức bạn ngại biến động. Cơ sở bằng 5, kiểm tra độ nhạy ở 3 và 10. & 3.3 \\
Tỷ số Sharpe & Sharpe ratio & Lợi suất vượt lãi suất phi rủi ro trên mỗi đơn vị độ lệch chuẩn. Chỉ là thước đo phụ của đề tài. & 1.1, 3.3 \\
Tỷ số Sharpe hiệu chỉnh & deflated Sharpe ratio & Tỷ số Sharpe được hiệu chỉnh cho thiên lệch chọn lọc và quá khớp backtest, theo Bailey và López de Prado (2014). & 2.2.3, 3.4 \\
Độ sụt giảm tối đa & maximum drawdown & Mức giảm lớn nhất từ đỉnh của giá trị danh mục. & 3.5 \\
Giá trị thiệt hại kỳ vọng 95\% & expected shortfall, ES 95\% & Tổn thất trung bình trong 5\% ngày xấu nhất. & 3.5 \\
Vòng quay danh mục & turnover & Tỷ lệ mua bán lại danh mục. Vòng quay cao kéo theo chi phí giao dịch cao. & 3.5 \\
Chi phí giao dịch & transaction cost & Khoản phí và chênh lệch mua bán khi đổi danh mục. Đề tài tham số hóa và kiểm tra độ nhạy. & 3.5 \\
Tái cân bằng & rebalancing & Việc điều chỉnh lại trọng số danh mục theo chu kỳ. Danh mục cơ sở tái cân bằng hằng tháng. & 3.1.1 \\
Ngoài mẫu và trong mẫu & out-of-sample, in-sample & Ngoài mẫu là kết quả trên dữ liệu chưa dùng để ước lượng. Đề tài đo hiệu quả ngoài mẫu của mỗi mô hình. & 3.3 \\
Cửa sổ ước lượng và cửa sổ đánh giá & estimation window, evaluation window & Cửa sổ ước lượng dài 104 tuần để ước lượng mô hình. Cửa sổ đánh giá dài 6 tháng, không chồng lấn, để chấm kết quả. Các giá trị này là dự kiến. & 3.3 \\
Siêu tham số & hyperparameter & Tham số đặt trước khi ước lượng, chọn từ một lưới cho trước và cố định khi đăng ký trước. & 3.3 \\
\end{longtable}}

\section{Nhóm 6: kiểm định và thống kê}

{\small\begin{longtable}[]{@{}
  >{\raggedright\arraybackslash}p{(\columnwidth - 6\tabcolsep) * \real{0.1396}}
  >{\raggedright\arraybackslash}p{(\columnwidth - 6\tabcolsep) * \real{0.1992}}
  >{\raggedright\arraybackslash}p{(\columnwidth - 6\tabcolsep) * \real{0.6008}}
  >{\raggedright\arraybackslash}p{(\columnwidth - 6\tabcolsep) * \real{0.0603}}@{}}
\toprule\noalign{}
\begin{minipage}[b]{\linewidth}\raggedright
Thuật ngữ
\end{minipage} & \begin{minipage}[b]{\linewidth}\raggedright
Tiếng Anh
\end{minipage} & \begin{minipage}[b]{\linewidth}\raggedright
Giải nghĩa
\end{minipage} & \begin{minipage}[b]{\linewidth}\raggedright
Mục gốc
\end{minipage} \\
\midrule\noalign{}
\endhead
\bottomrule\noalign{}
\endlastfoot
Hồi quy & regression & Đường xấp xỉ mối liên hệ giữa hai đại lượng. Phương trình (5) là hồi quy của D\_t hoặc R\_t theo chế độ và biến kênh. & 3.4 \\
Hệ số \texorpdfstring{$\beta$}{beta} & beta coefficient & Hệ số đo chênh lệch của thước đo rủi ro ở chế độ căng thẳng so với chế độ bình thường. & 3.4 \\
Biến kiểm soát & control variable & Biến đưa vào hồi quy để tách ảnh hưởng của nó. Đề tài kiểm soát biến động thực hiện. & 2.4 \\
Biến kênh & channel variable & Biến có thể nối chế độ căng thẳng với rủi ro: thanh khoản, tương quan cực trị, bầy đàn. & 1.2, 3.4 \\
Biến động thực hiện & realized volatility & Mức biến động đo từ dữ liệu giá đã xảy ra. Đề tài kiểm soát nó trước khi đọc biến kênh. & 2.4 \\
Sai số chuẩn bền & heteroskedasticity and autocorrelation consistent standard error, HAC & Cách tính sai số chuẩn không bị lệch khi phương sai thay đổi và sai số tự tương quan. & 3.4 \\
Bootstrap & bootstrap & Bốc lại nhiều lần từ chính túi dữ liệu để xem kết quả dao động bao nhiêu. & 3.4 \\
Bootstrap khối & block bootstrap & Bootstrap bốc từng nắm liền nhau, giữ thứ tự thời gian. Dùng cho chuỗi thời gian. & 3.4 \\
Bootstrap khối dừng & stationary bootstrap & Biến thể bootstrap khối, theo Politis và Romano (1994), dùng cho hệ số hồi quy. & 3.4 \\
Phân phối rỗng & null distribution & Kết quả sẽ trông thế nào nếu thực ra không có gì xảy ra. Dùng làm thước để so. Đề tài lập nó bằng 1.000 lần tái lấy mẫu. & 3.4 \\
Phân vị 95\% & 95th percentile & Giá trị mà 95\% kết quả của phân phối rỗng không vượt. H1 đòi hỏi hệ số \texorpdfstring{$\beta$}{beta} vượt mức này. & 3.4 \\
Khoảng tin cậy & confidence interval & Khoảng giá trị hợp lý cho một hệ số. Khi khoảng loại trừ 0, hệ số được coi là khác 0. & 3.4 \\
Hệ số nhân tử lạm phát phương sai & variance inflation factor, VIF & Chỉ số cho biết các biến giải thích chồng lấn nhau đến đâu. Ngưỡng 5, ấn định trước. & 3.4 \\
Liên hệ có điều kiện & conditional association & Mối liên hệ thống kê sau khi kiểm soát các biến khác, chưa chứng minh nguyên nhân. & 3.4 \\
Nhân quả & causality & Quan hệ trong đó một biến gây ra biến kia. Đề tài không gọi kết quả là nhân quả. & 3.4 \\
ILLIQ & Amihud illiquidity measure & Thước đo mức khó bán của cổ phiếu, theo Amihud (2002). Giá trị cao nghĩa là thanh khoản cạn. & 2.4, 3.4 \\
Tương quan đuôi & tail correlation & Tương quan giữa các tài sản trong những ngày cực xấu, theo Longin và Solnik (2001). & 2.4, 3.4 \\
Độ phân tán chéo của lợi suất & cross-sectional dispersion of returns & Mức khác nhau giữa lợi suất các cổ phiếu cùng ngày. Đề tài dùng làm thước đo bầy đàn cấp thị trường. & 1, 3.4 \\
Bầy đàn & herding & Đàn cừu chạy theo con đầu đàn thay vì tự nhìn đường. Làm theo số đông thay vì thông tin riêng. & 2.2.5 \\
\end{longtable}}

\section{Nhóm 7: độ tin cậy, minh bạch và phương pháp tính}

{\small\begin{longtable}[]{@{}
  >{\raggedright\arraybackslash}p{(\columnwidth - 6\tabcolsep) * \real{0.1576}}
  >{\raggedright\arraybackslash}p{(\columnwidth - 6\tabcolsep) * \real{0.1843}}
  >{\raggedright\arraybackslash}p{(\columnwidth - 6\tabcolsep) * \real{0.5981}}
  >{\raggedright\arraybackslash}p{(\columnwidth - 6\tabcolsep) * \real{0.0600}}@{}}
\toprule\noalign{}
\begin{minipage}[b]{\linewidth}\raggedright
Thuật ngữ
\end{minipage} & \begin{minipage}[b]{\linewidth}\raggedright
Tiếng Anh
\end{minipage} & \begin{minipage}[b]{\linewidth}\raggedright
Giải nghĩa
\end{minipage} & \begin{minipage}[b]{\linewidth}\raggedright
Mục gốc
\end{minipage} \\
\midrule\noalign{}
\endhead
\bottomrule\noalign{}
\endlastfoot
Hiệu chỉnh so sánh bội Holm & Holm correction & Cách siết ngưỡng khi kiểm định nhiều giả thuyết cùng lúc, theo Holm (1979). Đề tài dùng mức 10\% cho bốn giả thuyết. & 3.6 \\
Không kết luận & inconclusive & Cách gọi kết quả khi số cửa sổ độc lập ở chế độ căng thẳng dưới 10. Không ủng hộ, không bác bỏ. & 3.6 \\
Cửa sổ độc lập & independent window & Theo cách hiểu của bản đọc này, là cửa sổ đánh giá không chồng lấn với cửa sổ khác; bài gốc không định nghĩa chi tiết. Đề tài báo cáo số này cạnh mỗi hệ số. & 3.6 \\
Sức mạnh kiểm định & statistical power & Khả năng phát hiện một hiệu ứng có thật. Thấp khi có ít giai đoạn căng thẳng độc lập. & 3.6 \\
Kích thước hiệu ứng & effect size & Độ lớn của chênh lệch đo được. Báo cáo theo từng giai đoạn khi kết quả không kết luận. & 3.6 \\
Kiểm tra độ bền & robustness check & Lặp lại phân tích với định nghĩa hoặc thiết lập thay thế để xem kết luận có đổi không. & 3.1.2, 3.6 \\
Thiên lệch chọn lọc & selection bias & Sai lệch do chọn kết quả đẹp nhất trong nhiều lần thử. Hiệu chỉnh bằng tỷ số Sharpe hiệu chỉnh. & 2.2.3, 3.4 \\
Quá khớp kiểm định ngược & backtest overfitting & Khi mô hình khớp quá sát dữ liệu quá khứ nên kém hiệu quả ở tương lai. & 2.2.3 \\
Kiểm định ngược & backtest & Thử một chiến lược trên dữ liệu quá khứ để xem nó đã hoạt động ra sao. & 2.2.4 \\
Đăng ký trước & pre-registration & Viết dự đoán và cách kiểm tra vào phong bì có dấu thời gian trước khi xem dữ liệu mới. & 3.2 \\
OSF Registries & Open Science Framework Registries & Nơi đăng ký kế hoạch phân tích ở chế độ khóa có dấu thời gian, vào quý 4 năm 1. & 3.2 \\
Lệch khỏi kế hoạch & deviation from plan & Mọi thay đổi sau ngày đăng ký, phải ghi rõ là lệch khỏi kế hoạch. & 3.2 \\
Mẫu kiểm tra riêng & hold-out sample & Dữ liệu phát sinh sau ngày đăng ký, được giữ riêng để kiểm tra. & 3.2 \\
ABC-MCMC & Approximate Bayesian Computation with Markov chain Monte Carlo & Tính toán Bayes xấp xỉ kết hợp Monte Carlo chuỗi Markov. Tác giả đã dùng trong Nguyen và Nguyen (2026a); có thể thay bằng ước lượng nhanh hơn cho M5. & 3.7 \\
Mã nguồn & source code & Mã chương trình tác giả cam kết lưu để người khác có thể kiểm tra. & 3.7 \\
\end{longtable}}

\section{Nhóm 8: kế hoạch, chuẩn bị và trình bày}

{\small\begin{longtable}[]{@{}
  >{\raggedright\arraybackslash}p{(\columnwidth - 6\tabcolsep) * \real{0.1673}}
  >{\raggedright\arraybackslash}p{(\columnwidth - 6\tabcolsep) * \real{0.1923}}
  >{\raggedright\arraybackslash}p{(\columnwidth - 6\tabcolsep) * \real{0.5820}}
  >{\raggedright\arraybackslash}p{(\columnwidth - 6\tabcolsep) * \real{0.0584}}@{}}
\toprule\noalign{}
\begin{minipage}[b]{\linewidth}\raggedright
Thuật ngữ
\end{minipage} & \begin{minipage}[b]{\linewidth}\raggedright
Tiếng Anh
\end{minipage} & \begin{minipage}[b]{\linewidth}\raggedright
Giải nghĩa
\end{minipage} & \begin{minipage}[b]{\linewidth}\raggedright
Mục gốc
\end{minipage} \\
\midrule\noalign{}
\endhead
\bottomrule\noalign{}
\endlastfoot
Kinh tế hành vi & behavioral economics & Ngành xem xét ảnh hưởng của tâm lý lên quyết định kinh tế. Đề tài giao thoa ở một phép thử giới hạn. & 1, 5 \\
Tài chính định lượng & quantitative finance & Ngành dùng mô hình toán và thống kê để phân tích tài chính. Là lĩnh vực chính của đề tài. & 1 \\
Quy trình giám sát mô hình ba bước & three-step model oversight procedure & Khai báo tập mô hình; đo phân tán và MCS định kỳ; kích hoạt kết hợp mô hình khi thị trường ở chế độ căng thẳng. Là đề xuất chưa kiểm nghiệm. & 4 \\
Kết quả thăm dò & exploratory result & Kết quả chỉ báo cáo ở mức thăm dò. H4 thuộc loại này. & 2.4, 4 \\
Khả năng đối lập & rival explanation & Cách giải thích khác cho cùng hiện tượng. Mỗi giả thuyết có một khả năng đối lập riêng. & 2.4 \\
Nhóm 1 (Kinh tế học) & Group 1 (Economics) & Nhóm ngành trong Phụ lục 1 của Thông báo tuyển sinh. Ngành của tác giả thuộc nhóm này và cần học bổ sung 9 tín chỉ. & 5 \\
VSTEP Bậc 4 & VSTEP level 4 & Chứng chỉ tiếng Anh Bậc 4 theo Khung năng lực ngoại ngữ 6 bậc dùng cho Việt Nam, điểm tổng 6,5. & 5, 7 \\
Tác giả liên hệ & corresponding author & Tác giả chịu trách nhiệm trao đổi với tạp chí. Trong Bảng 6, dấu sao có ở bài 1 và bài 3, không có ở bài 2. Bảng 7 ghi tác giả đề cương là tác giả liên hệ ở cả chín bản thảo. & 7 \\
Scopus Q2 và Q3 & Scopus quartile Q2, Q3 & Cách xếp hạng tạp chí trong Scopus theo nhóm tư phân vị. Q2 là nhóm thứ hai, Q3 là nhóm thứ ba. & 7 \\
Bản thảo đã gửi & submitted manuscript & Bài đã gửi tạp chí nhưng chưa công bố. Bảng 7 có chín bản thảo. & 7 \\
Đã được chấp nhận đăng & accepted for publication & Bài được tạp chí đồng ý đăng nhưng chưa xuất bản. Bài 3 của Bảng 6 ở trạng thái này. & 7 \\
\end{longtable}}
